% =====================================================================
%  main.tex  --  GENERATED by generate_paper_draft.py.  Do not edit by
%  hand: edit the section file named in the \input and re-run.
%
%  Target venue formatting is not fixed here.  The document class below
%  is a neutral placeholder that compiles; switching to IEEEtran or
%  elsarticle requires changing only this preamble, because no section
%  file uses a class-specific command except \CIRCLE/\LEFTcircle/\Circle
%  in the related-work table (wasysym) and table* in two places.
%
%  ------------------------------------------------------------------
%  SWITCHING VENUE CLASS.  Three lines change; no section file changes.
%
%    IEEE Transactions / IEEE Access
%      1. \documentclass[journal]{IEEEtran}
%      2. \bibliographystyle{IEEEtran}
%      3. delete \usepackage{lmodern} and the geometry line -- IEEEtran
%         sets its own margins and font, and geometry fights it.
%      Note: IEEEtran wants \IEEEauthorblockN/\IEEEauthorblockA inside
%      \author, so the author block below is rewritten, not just filled.
%
%    Elsevier (Applied Energy, Energy)
%      1. \documentclass[preprint,12pt]{elsarticle}
%      2. \bibliographystyle{elsarticle-num}
%      3. delete the geometry line; wrap the author block in
%         \begin{frontmatter}...\end{frontmatter} with \author and
%         \affiliation, and move % =====================================================================
%  abstract.tex  --  PHASE VII, Task 14E
%
%  Rewritten from the Phase I draft.  Four things changed and each is a
%  consequence of a measured result rather than a stylistic choice:
%
%   1. The forecasting component is a gradient-boosted tree, not the
%      Seq2Seq LSTM.  Phase V: the sequence model is beaten by an
%      untrained persistence reference by 7.95 s of per-window MAE.
%   2. The decision-layer saving is 2.9 +/- 0.6 USD/yr, not 7.2, and it
%      is NOT claimed as a saving -- its interval spans zero.  The claim
%      is the paired contrast against the static rule.
%   3. The calendar share of the duration forecaster is 82% (Shapley),
%      not 94.5% (split gain).
%   4. The physics accept/reject loop is named as a contribution.
%
%  NO placeholder survives.  Every number below is traceable to a
%  stored output; see the section named beside it in the source.
% =====================================================================

\begin{abstract}
Industrial machines consume a substantial fraction of their electrical
energy while energised but not producing, yet eliminating this standby
consumption in an existing facility is obstructed by three coupled
difficulties: per-timestamp operating-state annotations do not exist and
are prohibitively expensive to obtain; the duration of a future
non-productive interval is unknown at the moment a shutdown decision must
be made; and shutting a machine down is economical only when that
duration exceeds a break-even point set by restart cost and standby
power. Existing work addresses these in isolation---non-intrusive load
monitoring infers states but neither forecasts nor acts, while
energy-aware scheduling decides optimally but assumes the idle duration
is given. This paper proposes a unified, label-free framework that closes
the loop from raw electrical measurements to an economically justified
shutdown decision, and evaluates it end to end on IMDELD, a public
benchmark of $1$\,Hz measurements from eight heavy machines in a
full-scale plant. The framework comprises (i)~\emph{physics-gated
unsupervised state inference}, in which a Gaussian hidden Markov model is
fitted without labels and the resulting partition is then \emph{accepted
or rejected} against circuit-level checks external to the model;
(ii)~\emph{non-productive duration forecasting} benchmarked against an
untrained reference; and (iii)~\emph{optimisation-based decision
support}, in which the break-even duration is derived as the solution of
a constrained cost minimisation rather than supplied as a threshold. The
accept/reject step is load-bearing: fitted at ten seeds the labeller
places the standby cluster on a physically different state in three, and
conditioning on a passing score collapses the standby-hour spread from
$\pm33.7\%$ to $\pm3.0\%$. Temporal decoding reduces state flicker from
$66.9\%$ to $44.5\%$ and an explicit minimum-dwell constraint removes it
entirely, at a cost of $0.19\%$ of rows relabelled. Two results are
reported against expectation: of seven forecasters, five neural sequence
models are beaten by an untrained persistence reference, and only a
gradient-boosted tree clears it ($11.50$ against $12.95$\,s of
per-window MAE, $p_{\text{Holm}}=7\times10^{-4}$); and the proposed
policy cannot be shown to save money against the plant's own behaviour on
an 18-day held-out sample ($+2.9$\,USD/yr, $95\%$ CI $[-56,+46]$), though
it beats the fixed-threshold rule it replaces by $+68.9$\,USD/yr
($p=0.001$). Across the plant the state definition is physically valid on
the four motor-driven machines and correctly rejected on the four
auxiliary ones, and the facility's total attainable standby saving is
under $90$\,USD/yr. The framework's demonstrated value on this facility
is therefore decision quality and validated state inference rather than
recovered energy, and it requires only the voltage and current channels a
standard retrofit meter already provides.
\end{abstract}
 inside it.
%
%    Springer (SCIS, LNCS)
%      1. \documentclass{sn-jnl}  (or {llncs})
%      2. \bibliographystyle{sn-mathphys}
%      3. delete geometry and lmodern.
%
%  In every case the figures are already sized for a 3.5in single
%  column and a 7.16in double column, so nothing needs re-rendering.
%  ------------------------------------------------------------------
% =====================================================================
\documentclass[10pt,twocolumn]{article}

\usepackage[T1]{fontenc}
\usepackage[utf8]{inputenc}
\usepackage{lmodern}
\usepackage[margin=0.75in]{geometry}
\usepackage{amsmath,amssymb}
\usepackage{booktabs}
\usepackage{threeparttable}
\usepackage{multirow}
\usepackage[nointegrals]{wasysym}   % \CIRCLE \LEFTcircle \Circle
\usepackage{graphicx}
\usepackage{amsthm}
\usepackage[hidelinks]{hyperref}

\newtheorem{proposition}{Proposition}

\graphicspath{%
  {../outputs/phase7/figures/}%
  {../outputs/phase6/figures/}%
  {../outputs/phase5/figures/}%
  {../outputs/phase4/figures/}%
  {../outputs/phase3/figures/pelletizer-I/}%
  {../outputs/phase2/figures/pelletizer-I/}%
}

\title{A Physics-Gated, Label-Free Framework for Industrial Standby
Energy Elimination: State Inference, Duration Forecasting and
Optimisation-Based Shutdown Decisions}

% ---------------------------------------------------------------------
%  AUTHOR BLOCK.  Every FILL_IN below must be replaced before
%  submission.  `generate_paper_draft.py` counts them and prints the
%  count on every run; it does not refuse to assemble, because an
%  unfilled author list is an incomplete submission rather than an
%  inconsistent document.
% ---------------------------------------------------------------------
\author{%
  FILL\_IN Author Name$^{1}$\thanks{Corresponding author:
  \texttt{FILL\_IN author@institution.edu}}\\[2pt]
  \normalsize $^{1}$FILL\_IN Department, FILL\_IN Institution,
  FILL\_IN City, FILL\_IN Country%
}
\date{}

\begin{document}
\maketitle

% =====================================================================
%  abstract.tex  --  PHASE VII, Task 14E
%
%  Rewritten from the Phase I draft.  Four things changed and each is a
%  consequence of a measured result rather than a stylistic choice:
%
%   1. The forecasting component is a gradient-boosted tree, not the
%      Seq2Seq LSTM.  Phase V: the sequence model is beaten by an
%      untrained persistence reference by 7.95 s of per-window MAE.
%   2. The decision-layer saving is 2.9 +/- 0.6 USD/yr, not 7.2, and it
%      is NOT claimed as a saving -- its interval spans zero.  The claim
%      is the paired contrast against the static rule.
%   3. The calendar share of the duration forecaster is 82% (Shapley),
%      not 94.5% (split gain).
%   4. The physics accept/reject loop is named as a contribution.
%
%  NO placeholder survives.  Every number below is traceable to a
%  stored output; see the section named beside it in the source.
% =====================================================================

\begin{abstract}
Industrial machines consume a substantial fraction of their electrical
energy while energised but not producing, yet eliminating this standby
consumption in an existing facility is obstructed by three coupled
difficulties: per-timestamp operating-state annotations do not exist and
are prohibitively expensive to obtain; the duration of a future
non-productive interval is unknown at the moment a shutdown decision must
be made; and shutting a machine down is economical only when that
duration exceeds a break-even point set by restart cost and standby
power. Existing work addresses these in isolation---non-intrusive load
monitoring infers states but neither forecasts nor acts, while
energy-aware scheduling decides optimally but assumes the idle duration
is given. This paper proposes a unified, label-free framework that closes
the loop from raw electrical measurements to an economically justified
shutdown decision, and evaluates it end to end on IMDELD, a public
benchmark of $1$\,Hz measurements from eight heavy machines in a
full-scale plant. The framework comprises (i)~\emph{physics-gated
unsupervised state inference}, in which a Gaussian hidden Markov model is
fitted without labels and the resulting partition is then \emph{accepted
or rejected} against circuit-level checks external to the model;
(ii)~\emph{non-productive duration forecasting} benchmarked against an
untrained reference; and (iii)~\emph{optimisation-based decision
support}, in which the break-even duration is derived as the solution of
a constrained cost minimisation rather than supplied as a threshold. The
accept/reject step is load-bearing: fitted at ten seeds the labeller
places the standby cluster on a physically different state in three, and
conditioning on a passing score collapses the standby-hour spread from
$\pm33.7\%$ to $\pm3.0\%$. Temporal decoding reduces state flicker from
$66.9\%$ to $44.5\%$ and an explicit minimum-dwell constraint removes it
entirely, at a cost of $0.19\%$ of rows relabelled. Two results are
reported against expectation: of seven forecasters, five neural sequence
models are beaten by an untrained persistence reference, and only a
gradient-boosted tree clears it ($11.50$ against $12.95$\,s of
per-window MAE, $p_{\text{Holm}}=7\times10^{-4}$); and the proposed
policy cannot be shown to save money against the plant's own behaviour on
an 18-day held-out sample ($+2.9$\,USD/yr, $95\%$ CI $[-56,+46]$), though
it beats the fixed-threshold rule it replaces by $+68.9$\,USD/yr
($p=0.001$). Across the plant the state definition is physically valid on
the four motor-driven machines and correctly rejected on the four
auxiliary ones, and the facility's total attainable standby saving is
under $90$\,USD/yr. The framework's demonstrated value on this facility
is therefore decision quality and validated state inference rather than
recovered energy, and it requires only the voltage and current channels a
standard retrofit meter already provides.
\end{abstract}


% --- Introduction, positioning, contributions
% =====================================================================
%  introduction.tex  --  PHASE VII, Task 14E
%
%  Section 1.  Context and problem, the positioning statement, the
%  contribution list, and the organisation paragraph.
%
%  Supersedes the corresponding blocks of novelty_and_contributions.tex
%  (Phase I), which recorded the design rationale for the reframing and
%  is retained for that purpose only.  Changes from the Phase I text:
%
%   * component (ii) is a duration forecaster benchmarked against an
%     untrained reference, not "a Seq2Seq LSTM" -- Phase V;
%   * the accept/reject loop is promoted to a named contribution
%     -- Phase V Task 12B;
%   * the decision contribution claims the paired contrast against the
%     static rule, not a saving against the status quo -- Phase V 11B;
%   * the sustainability/deployment forward references are removed;
%     those analyses are not in this paper.
%
%  Drop-in usage: % =====================================================================
%  introduction.tex  --  PHASE VII, Task 14E
%
%  Section 1.  Context and problem, the positioning statement, the
%  contribution list, and the organisation paragraph.
%
%  Supersedes the corresponding blocks of novelty_and_contributions.tex
%  (Phase I), which recorded the design rationale for the reframing and
%  is retained for that purpose only.  Changes from the Phase I text:
%
%   * component (ii) is a duration forecaster benchmarked against an
%     untrained reference, not "a Seq2Seq LSTM" -- Phase V;
%   * the accept/reject loop is promoted to a named contribution
%     -- Phase V Task 12B;
%   * the decision contribution claims the paired contrast against the
%     static rule, not a saving against the status quo -- Phase V 11B;
%   * the sustainability/deployment forward references are removed;
%     those analyses are not in this paper.
%
%  Drop-in usage: % =====================================================================
%  introduction.tex  --  PHASE VII, Task 14E
%
%  Section 1.  Context and problem, the positioning statement, the
%  contribution list, and the organisation paragraph.
%
%  Supersedes the corresponding blocks of novelty_and_contributions.tex
%  (Phase I), which recorded the design rationale for the reframing and
%  is retained for that purpose only.  Changes from the Phase I text:
%
%   * component (ii) is a duration forecaster benchmarked against an
%     untrained reference, not "a Seq2Seq LSTM" -- Phase V;
%   * the accept/reject loop is promoted to a named contribution
%     -- Phase V Task 12B;
%   * the decision contribution claims the paired contrast against the
%     static rule, not a saving against the status quo -- Phase V 11B;
%   * the sustainability/deployment forward references are removed;
%     those analyses are not in this paper.
%
%  Drop-in usage: \input{introduction}
%  Bibliography: paper/references.bib
% =====================================================================

\section{Introduction}
\label{sec:intro}

Electric motor systems account for the largest single share of
industrial electricity use, and a substantial part of that share is
drawn while the machine is energised but producing nothing. An induction
motor left running between production runs continues to draw magnetising
current: real power falls to a fraction of its loaded value, but it does
not fall to zero, and the machine's power factor collapses. ISO~14955-1
\cite{iso14955} names this operating mode --- ready for production,
performing no value-adding function --- and studies of discrete-part
manufacturing find that such load-independent consumption frequently
exceeds the energy actually delivered to the workpiece
\cite{gutowski2006electrical,dahmus2004environmental}. Eliminating it
requires no capital equipment and no process change. It requires only
knowing when a machine is idle, and for how long it will remain so.

Neither is known in an existing facility. Three difficulties compound.

\textbf{There are no labels.} A per-timestamp record of which state each
machine occupied would reduce state identification to a standard
classification problem, but obtaining one means instrumenting every
machine's contactor and control system, or having an operator observe
each machine for months. Neither is affordable at facility scale and
neither transfers to a new machine without repeating the effort. Public
industrial benchmarks reflect this: IMDELD
\cite{martins2018imdeld,imdeld2018dataset}, the benchmark used
throughout this paper, publishes $1$\,Hz electrical measurements from
eight heavy machines and no state annotation at all. Any method that
must be validated against ground truth is therefore inapplicable, and
so is any evaluation metric that assumes it.

\textbf{The duration is unknown when the decision must be made.} The
shutdown trade-off is well understood: de-energising saves standby power
but incurs a restart cost, so it pays only when the idle interval
exceeds a break-even duration \cite{mouzon2007operational}. The
operations-research literature solves the resulting scheduling problem
in considerable depth --- and universally assumes the idle interval is
given, either by a deterministic production schedule or by an
analyst-calibrated arrival process. In a retrofit installation the only
available signal is the machine's own measured electrical behaviour, and
nothing in that literature converts one into the other.

\textbf{The two capabilities have matured separately and remain
unconnected.} Non-intrusive load monitoring can infer what a machine is
doing from its electrical signature \cite{hart1992nonintrusive}, but its
objective is retrospective attribution: it does not forecast and it does
not act. Energy-aware scheduling can act optimally given a duration it
cannot obtain. \emph{The former produces labels no controller consumes;
the latter consumes durations no inference procedure supplies.}
Section~\ref{sec:related} develops this gap across four bodies of
literature and states it precisely.

\medskip
\noindent
To close it we propose a \emph{unified data-driven framework for
industrial standby energy elimination}. Rather than treating state
inference, forecasting and control as separate problems, the framework
treats them as one pipeline whose stages are designed against each
other's requirements: the inference stage is constrained to produce a
temporally coherent sequence because duration forecasting is otherwise
ill-posed, and the forecasting stage is constrained to predict a
duration rather than a class because the decision stage consumes a
duration. The framework has three components.

\begin{enumerate}
  \item \textbf{Physics-gated unsupervised state inference.} A Gaussian
  hidden Markov model \cite{rabiner1989tutorial} estimates the emission
  distributions and transition dynamics of the machine's operating
  states directly from electrical measurements, requiring no
  annotation, and an explicit minimum-dwell constraint enforces
  physically admissible state durations at decode time. Because no
  ground-truth labels exist against which the partition could be
  scored, the fit is then \emph{accepted or rejected} against checks
  drawn from sources external to the model: induction-motor
  power-factor and no-load current physics, variance and power
  ordering, agreement with independently derived threshold classifiers,
  and the plant's published shift calendar.

  \item \textbf{Non-productive duration forecasting.} The recent
  inferred-state history is mapped to the duration of the upcoming
  non-productive interval. Predicting duration rather than the next
  state label is what makes the output actionable, since the break-even
  criterion is a statement about how long a machine will remain idle.
  Six trained forecasters --- recurrent, convolutional, transformer and
  tree-based --- are compared under a common protocol \emph{and against
  an untrained persistence reference}, which is the comparison that
  determines whether the stage earns its place at all.

  \item \textbf{Optimisation-based shutdown decision support.} The
  decision is posed as minimisation of total cost --- standby energy,
  restart, and production delay --- subject to thermal cool-down and
  restart-frequency constraints, in the spirit of
  \cite{mouzon2007operational} but with the idle duration supplied by
  component~(ii) instead of by a known schedule. The break-even
  duration is therefore a \emph{derived} quantity of the optimisation
  rather than an operator-supplied threshold, and because the duration
  is unknown at decision time the policy is benchmarked against the
  $2$-competitive forecast-free rent-or-buy rule
  \cite{karlin1994competitive} rather than against trivial policies.
\end{enumerate}

\medskip
\noindent
\textbf{Scope of the claim.} We claim no novelty in the Gaussian hidden
Markov model, in any of the forecasting architectures, or in the
break-even criterion, each of which is established. The contribution is
the framework that connects them --- specifically, the demonstration
that label-free state inference can be made trustworthy enough, and
temporally coherent enough, to serve as the input to an economic control
decision on real industrial data --- together with an evaluation
honest enough to identify which of its own components do not earn their
place.

\subsection*{Contributions}

\begin{itemize}
  \item \textbf{A physics accept/reject loop, and the measurement that
  shows it is load-bearing.} Fitting the labeller at ten seeds on
  identical preprocessed data moves the headline standby total by a
  coefficient of variation of $33.7\%$, and in three seeds of ten the
  cluster named \textsc{standby} is physically a different state
  --- once a loaded running state at $13$\,kW, twice the machine
  switched off at under $1$\,W. The physics battery scores exactly those
  three below threshold, and conditioning on a passing score collapses
  the spread to $3.0\%$. \emph{The pipeline is reproducible because of
  the physics validation, not despite the clustering}
  (Sections~\ref{sec:method:accept_reject},
  \ref{sec:significance:seeds}).

  \item \textbf{An account of what temporal coherence actually
  requires.} Viterbi decoding over a learned transition matrix reduces
  flicker substantially but --- contrary to the usual assumption ---
  cannot eliminate it, and neither can strengthening the transition
  prior. The penalty a transition matrix can impose on a state change is
  bounded by $\log(p_{\text{self}}/p_{\text{switch}})$, measured here at
  $9.23$ nats, whereas full-covariance Gaussian emissions fitted to
  low-noise electrical data differ by a median of $27.8$ nats and a
  $90$th percentile of $41\,701$ nats between adjacent states. At
  $55.7\%$ of timesteps the emission term is arithmetically unable to be
  overruled. We show the remedy is an explicit minimum-dwell constraint,
  that it removes flicker entirely at a cost of $0.19\%$ of rows, and
  that it does so on all eight machines
  (Sections~\ref{sec:flicker}, \ref{sec:multimachine:general}).

  \item \textbf{A validation methodology for state partitions in the
  absence of ground truth.} Circuit-level physics, model-level
  diagnostics, and independently derived Otsu-thresholded classifiers
  \cite{otsu1979threshold} are combined by consensus
  \cite{strehl2002cluster}, none of them using the model's own internal
  statistics. Applied across eight machines the battery accepts the four
  motor-driven ones and rejects the four auxiliary ones, on which there
  is no energised-idle state to find --- which is the behaviour that
  makes it a validation procedure rather than a formality
  (Sections~\ref{sec:validation}, \ref{sec:multimachine:states}).

  \item \textbf{A forecasting comparison anchored on not forecasting,
  and a negative result reported as such.} Under one protocol with
  seeds, an untrained persistence reference beats all five neural
  sequence models, the originally proposed sequence-to-sequence LSTM by
  the largest margin of the seven ($7.95$\,s of per-window MAE,
  $p=5.7\times10^{-44}$). A gradient-boosted tree is the only trained
  forecaster that clears the reference, and it is the forecasting
  component of the framework as reported here
  (Sections~\ref{sec:forecasting}, \ref{sec:significance:forecast}).

  \item \textbf{A decision layer in which the break-even duration is
  derived, with an invariance result and a competitive benchmark.} We
  replace the fixed-threshold rule with a constrained cost
  minimisation, show that the reduced problem admits an exact offline
  optimum, and measure the restart coefficients from the record instead
  of assuming them --- finding the studied machine has no measurable
  electrical restart penalty, so its restart cost is a production-delay
  cost. We show analytically and numerically that the break-even
  duration is \emph{independent of electricity tariff} whenever the
  restart cost is purely electrical, and report the decision surface
  over the axes that do carry information. The policy beats the fixed
  threshold it replaces by $+68.9$\,USD/yr ($p=0.001$); it is not shown
  to beat the plant's own behaviour, and we say so
  (Sections~\ref{sec:decision}, \ref{sec:significance:decision}).

  \item \textbf{End-to-end evaluation with ablation, significance
  testing and transfer.} Every stage is evaluated across eight machines
  and two sites, each component's marginal contribution is priced by
  ablation, every headline difference carries a dependence-aware
  interval and a pre-registered practical threshold, and the two designs
  that cannot detect anything are identified as such rather than
  reported as null results (Sections~\ref{sec:multimachine}--\ref{sec:ablation}).

  \item \textbf{Reproducibility.} The full pipeline --- preprocessing,
  inference, validation harness, forecasting, decision layer and every
  experiment reported here --- is released as open source, and all
  results are obtained from public benchmark data.
\end{itemize}

\subsection*{Organisation}

Section~\ref{sec:related} reviews related work across machine state
identification, non-intrusive load monitoring, industrial energy
optimisation and predictive shutdown, and states the resulting gap.
Section~\ref{sec:method} presents the framework.
Section~\ref{sec:states} reports the state layer, including the
comparison against clustering baselines, the flicker analysis, and the
ground-truth-free validation. Section~\ref{sec:forecasting} reports the
forecasting comparison and Section~\ref{sec:significance} the
statistical analysis on which every subsequent claim depends.
Section~\ref{sec:decision} develops the decision layer.
Sections~\ref{sec:multimachine} and~\ref{sec:crossmachine} extend the
evaluation to all eight machines, to transfer between them, and to a
second site. Section~\ref{sec:explain} attributes what each model uses,
Section~\ref{sec:ablation} prices each component, and
Section~\ref{sec:limitations} states the boundaries of the evidence
before Section~\ref{sec:conclusion} concludes.

%  Bibliography: paper/references.bib
% =====================================================================

\section{Introduction}
\label{sec:intro}

Electric motor systems account for the largest single share of
industrial electricity use, and a substantial part of that share is
drawn while the machine is energised but producing nothing. An induction
motor left running between production runs continues to draw magnetising
current: real power falls to a fraction of its loaded value, but it does
not fall to zero, and the machine's power factor collapses. ISO~14955-1
\cite{iso14955} names this operating mode --- ready for production,
performing no value-adding function --- and studies of discrete-part
manufacturing find that such load-independent consumption frequently
exceeds the energy actually delivered to the workpiece
\cite{gutowski2006electrical,dahmus2004environmental}. Eliminating it
requires no capital equipment and no process change. It requires only
knowing when a machine is idle, and for how long it will remain so.

Neither is known in an existing facility. Three difficulties compound.

\textbf{There are no labels.} A per-timestamp record of which state each
machine occupied would reduce state identification to a standard
classification problem, but obtaining one means instrumenting every
machine's contactor and control system, or having an operator observe
each machine for months. Neither is affordable at facility scale and
neither transfers to a new machine without repeating the effort. Public
industrial benchmarks reflect this: IMDELD
\cite{martins2018imdeld,imdeld2018dataset}, the benchmark used
throughout this paper, publishes $1$\,Hz electrical measurements from
eight heavy machines and no state annotation at all. Any method that
must be validated against ground truth is therefore inapplicable, and
so is any evaluation metric that assumes it.

\textbf{The duration is unknown when the decision must be made.} The
shutdown trade-off is well understood: de-energising saves standby power
but incurs a restart cost, so it pays only when the idle interval
exceeds a break-even duration \cite{mouzon2007operational}. The
operations-research literature solves the resulting scheduling problem
in considerable depth --- and universally assumes the idle interval is
given, either by a deterministic production schedule or by an
analyst-calibrated arrival process. In a retrofit installation the only
available signal is the machine's own measured electrical behaviour, and
nothing in that literature converts one into the other.

\textbf{The two capabilities have matured separately and remain
unconnected.} Non-intrusive load monitoring can infer what a machine is
doing from its electrical signature \cite{hart1992nonintrusive}, but its
objective is retrospective attribution: it does not forecast and it does
not act. Energy-aware scheduling can act optimally given a duration it
cannot obtain. \emph{The former produces labels no controller consumes;
the latter consumes durations no inference procedure supplies.}
Section~\ref{sec:related} develops this gap across four bodies of
literature and states it precisely.

\medskip
\noindent
To close it we propose a \emph{unified data-driven framework for
industrial standby energy elimination}. Rather than treating state
inference, forecasting and control as separate problems, the framework
treats them as one pipeline whose stages are designed against each
other's requirements: the inference stage is constrained to produce a
temporally coherent sequence because duration forecasting is otherwise
ill-posed, and the forecasting stage is constrained to predict a
duration rather than a class because the decision stage consumes a
duration. The framework has three components.

\begin{enumerate}
  \item \textbf{Physics-gated unsupervised state inference.} A Gaussian
  hidden Markov model \cite{rabiner1989tutorial} estimates the emission
  distributions and transition dynamics of the machine's operating
  states directly from electrical measurements, requiring no
  annotation, and an explicit minimum-dwell constraint enforces
  physically admissible state durations at decode time. Because no
  ground-truth labels exist against which the partition could be
  scored, the fit is then \emph{accepted or rejected} against checks
  drawn from sources external to the model: induction-motor
  power-factor and no-load current physics, variance and power
  ordering, agreement with independently derived threshold classifiers,
  and the plant's published shift calendar.

  \item \textbf{Non-productive duration forecasting.} The recent
  inferred-state history is mapped to the duration of the upcoming
  non-productive interval. Predicting duration rather than the next
  state label is what makes the output actionable, since the break-even
  criterion is a statement about how long a machine will remain idle.
  Six trained forecasters --- recurrent, convolutional, transformer and
  tree-based --- are compared under a common protocol \emph{and against
  an untrained persistence reference}, which is the comparison that
  determines whether the stage earns its place at all.

  \item \textbf{Optimisation-based shutdown decision support.} The
  decision is posed as minimisation of total cost --- standby energy,
  restart, and production delay --- subject to thermal cool-down and
  restart-frequency constraints, in the spirit of
  \cite{mouzon2007operational} but with the idle duration supplied by
  component~(ii) instead of by a known schedule. The break-even
  duration is therefore a \emph{derived} quantity of the optimisation
  rather than an operator-supplied threshold, and because the duration
  is unknown at decision time the policy is benchmarked against the
  $2$-competitive forecast-free rent-or-buy rule
  \cite{karlin1994competitive} rather than against trivial policies.
\end{enumerate}

\medskip
\noindent
\textbf{Scope of the claim.} We claim no novelty in the Gaussian hidden
Markov model, in any of the forecasting architectures, or in the
break-even criterion, each of which is established. The contribution is
the framework that connects them --- specifically, the demonstration
that label-free state inference can be made trustworthy enough, and
temporally coherent enough, to serve as the input to an economic control
decision on real industrial data --- together with an evaluation
honest enough to identify which of its own components do not earn their
place.

\subsection*{Contributions}

\begin{itemize}
  \item \textbf{A physics accept/reject loop, and the measurement that
  shows it is load-bearing.} Fitting the labeller at ten seeds on
  identical preprocessed data moves the headline standby total by a
  coefficient of variation of $33.7\%$, and in three seeds of ten the
  cluster named \textsc{standby} is physically a different state
  --- once a loaded running state at $13$\,kW, twice the machine
  switched off at under $1$\,W. The physics battery scores exactly those
  three below threshold, and conditioning on a passing score collapses
  the spread to $3.0\%$. \emph{The pipeline is reproducible because of
  the physics validation, not despite the clustering}
  (Sections~\ref{sec:method:accept_reject},
  \ref{sec:significance:seeds}).

  \item \textbf{An account of what temporal coherence actually
  requires.} Viterbi decoding over a learned transition matrix reduces
  flicker substantially but --- contrary to the usual assumption ---
  cannot eliminate it, and neither can strengthening the transition
  prior. The penalty a transition matrix can impose on a state change is
  bounded by $\log(p_{\text{self}}/p_{\text{switch}})$, measured here at
  $9.23$ nats, whereas full-covariance Gaussian emissions fitted to
  low-noise electrical data differ by a median of $27.8$ nats and a
  $90$th percentile of $41\,701$ nats between adjacent states. At
  $55.7\%$ of timesteps the emission term is arithmetically unable to be
  overruled. We show the remedy is an explicit minimum-dwell constraint,
  that it removes flicker entirely at a cost of $0.19\%$ of rows, and
  that it does so on all eight machines
  (Sections~\ref{sec:flicker}, \ref{sec:multimachine:general}).

  \item \textbf{A validation methodology for state partitions in the
  absence of ground truth.} Circuit-level physics, model-level
  diagnostics, and independently derived Otsu-thresholded classifiers
  \cite{otsu1979threshold} are combined by consensus
  \cite{strehl2002cluster}, none of them using the model's own internal
  statistics. Applied across eight machines the battery accepts the four
  motor-driven ones and rejects the four auxiliary ones, on which there
  is no energised-idle state to find --- which is the behaviour that
  makes it a validation procedure rather than a formality
  (Sections~\ref{sec:validation}, \ref{sec:multimachine:states}).

  \item \textbf{A forecasting comparison anchored on not forecasting,
  and a negative result reported as such.} Under one protocol with
  seeds, an untrained persistence reference beats all five neural
  sequence models, the originally proposed sequence-to-sequence LSTM by
  the largest margin of the seven ($7.95$\,s of per-window MAE,
  $p=5.7\times10^{-44}$). A gradient-boosted tree is the only trained
  forecaster that clears the reference, and it is the forecasting
  component of the framework as reported here
  (Sections~\ref{sec:forecasting}, \ref{sec:significance:forecast}).

  \item \textbf{A decision layer in which the break-even duration is
  derived, with an invariance result and a competitive benchmark.} We
  replace the fixed-threshold rule with a constrained cost
  minimisation, show that the reduced problem admits an exact offline
  optimum, and measure the restart coefficients from the record instead
  of assuming them --- finding the studied machine has no measurable
  electrical restart penalty, so its restart cost is a production-delay
  cost. We show analytically and numerically that the break-even
  duration is \emph{independent of electricity tariff} whenever the
  restart cost is purely electrical, and report the decision surface
  over the axes that do carry information. The policy beats the fixed
  threshold it replaces by $+68.9$\,USD/yr ($p=0.001$); it is not shown
  to beat the plant's own behaviour, and we say so
  (Sections~\ref{sec:decision}, \ref{sec:significance:decision}).

  \item \textbf{End-to-end evaluation with ablation, significance
  testing and transfer.} Every stage is evaluated across eight machines
  and two sites, each component's marginal contribution is priced by
  ablation, every headline difference carries a dependence-aware
  interval and a pre-registered practical threshold, and the two designs
  that cannot detect anything are identified as such rather than
  reported as null results (Sections~\ref{sec:multimachine}--\ref{sec:ablation}).

  \item \textbf{Reproducibility.} The full pipeline --- preprocessing,
  inference, validation harness, forecasting, decision layer and every
  experiment reported here --- is released as open source, and all
  results are obtained from public benchmark data.
\end{itemize}

\subsection*{Organisation}

Section~\ref{sec:related} reviews related work across machine state
identification, non-intrusive load monitoring, industrial energy
optimisation and predictive shutdown, and states the resulting gap.
Section~\ref{sec:method} presents the framework.
Section~\ref{sec:states} reports the state layer, including the
comparison against clustering baselines, the flicker analysis, and the
ground-truth-free validation. Section~\ref{sec:forecasting} reports the
forecasting comparison and Section~\ref{sec:significance} the
statistical analysis on which every subsequent claim depends.
Section~\ref{sec:decision} develops the decision layer.
Sections~\ref{sec:multimachine} and~\ref{sec:crossmachine} extend the
evaluation to all eight machines, to transfer between them, and to a
second site. Section~\ref{sec:explain} attributes what each model uses,
Section~\ref{sec:ablation} prices each component, and
Section~\ref{sec:limitations} states the boundaries of the evidence
before Section~\ref{sec:conclusion} concludes.

%  Bibliography: paper/references.bib
% =====================================================================

\section{Introduction}
\label{sec:intro}

Electric motor systems account for the largest single share of
industrial electricity use, and a substantial part of that share is
drawn while the machine is energised but producing nothing. An induction
motor left running between production runs continues to draw magnetising
current: real power falls to a fraction of its loaded value, but it does
not fall to zero, and the machine's power factor collapses. ISO~14955-1
\cite{iso14955} names this operating mode --- ready for production,
performing no value-adding function --- and studies of discrete-part
manufacturing find that such load-independent consumption frequently
exceeds the energy actually delivered to the workpiece
\cite{gutowski2006electrical,dahmus2004environmental}. Eliminating it
requires no capital equipment and no process change. It requires only
knowing when a machine is idle, and for how long it will remain so.

Neither is known in an existing facility. Three difficulties compound.

\textbf{There are no labels.} A per-timestamp record of which state each
machine occupied would reduce state identification to a standard
classification problem, but obtaining one means instrumenting every
machine's contactor and control system, or having an operator observe
each machine for months. Neither is affordable at facility scale and
neither transfers to a new machine without repeating the effort. Public
industrial benchmarks reflect this: IMDELD
\cite{martins2018imdeld,imdeld2018dataset}, the benchmark used
throughout this paper, publishes $1$\,Hz electrical measurements from
eight heavy machines and no state annotation at all. Any method that
must be validated against ground truth is therefore inapplicable, and
so is any evaluation metric that assumes it.

\textbf{The duration is unknown when the decision must be made.} The
shutdown trade-off is well understood: de-energising saves standby power
but incurs a restart cost, so it pays only when the idle interval
exceeds a break-even duration \cite{mouzon2007operational}. The
operations-research literature solves the resulting scheduling problem
in considerable depth --- and universally assumes the idle interval is
given, either by a deterministic production schedule or by an
analyst-calibrated arrival process. In a retrofit installation the only
available signal is the machine's own measured electrical behaviour, and
nothing in that literature converts one into the other.

\textbf{The two capabilities have matured separately and remain
unconnected.} Non-intrusive load monitoring can infer what a machine is
doing from its electrical signature \cite{hart1992nonintrusive}, but its
objective is retrospective attribution: it does not forecast and it does
not act. Energy-aware scheduling can act optimally given a duration it
cannot obtain. \emph{The former produces labels no controller consumes;
the latter consumes durations no inference procedure supplies.}
Section~\ref{sec:related} develops this gap across four bodies of
literature and states it precisely.

\medskip
\noindent
To close it we propose a \emph{unified data-driven framework for
industrial standby energy elimination}. Rather than treating state
inference, forecasting and control as separate problems, the framework
treats them as one pipeline whose stages are designed against each
other's requirements: the inference stage is constrained to produce a
temporally coherent sequence because duration forecasting is otherwise
ill-posed, and the forecasting stage is constrained to predict a
duration rather than a class because the decision stage consumes a
duration. The framework has three components.

\begin{enumerate}
  \item \textbf{Physics-gated unsupervised state inference.} A Gaussian
  hidden Markov model \cite{rabiner1989tutorial} estimates the emission
  distributions and transition dynamics of the machine's operating
  states directly from electrical measurements, requiring no
  annotation, and an explicit minimum-dwell constraint enforces
  physically admissible state durations at decode time. Because no
  ground-truth labels exist against which the partition could be
  scored, the fit is then \emph{accepted or rejected} against checks
  drawn from sources external to the model: induction-motor
  power-factor and no-load current physics, variance and power
  ordering, agreement with independently derived threshold classifiers,
  and the plant's published shift calendar.

  \item \textbf{Non-productive duration forecasting.} The recent
  inferred-state history is mapped to the duration of the upcoming
  non-productive interval. Predicting duration rather than the next
  state label is what makes the output actionable, since the break-even
  criterion is a statement about how long a machine will remain idle.
  Six trained forecasters --- recurrent, convolutional, transformer and
  tree-based --- are compared under a common protocol \emph{and against
  an untrained persistence reference}, which is the comparison that
  determines whether the stage earns its place at all.

  \item \textbf{Optimisation-based shutdown decision support.} The
  decision is posed as minimisation of total cost --- standby energy,
  restart, and production delay --- subject to thermal cool-down and
  restart-frequency constraints, in the spirit of
  \cite{mouzon2007operational} but with the idle duration supplied by
  component~(ii) instead of by a known schedule. The break-even
  duration is therefore a \emph{derived} quantity of the optimisation
  rather than an operator-supplied threshold, and because the duration
  is unknown at decision time the policy is benchmarked against the
  $2$-competitive forecast-free rent-or-buy rule
  \cite{karlin1994competitive} rather than against trivial policies.
\end{enumerate}

\medskip
\noindent
\textbf{Scope of the claim.} We claim no novelty in the Gaussian hidden
Markov model, in any of the forecasting architectures, or in the
break-even criterion, each of which is established. The contribution is
the framework that connects them --- specifically, the demonstration
that label-free state inference can be made trustworthy enough, and
temporally coherent enough, to serve as the input to an economic control
decision on real industrial data --- together with an evaluation
honest enough to identify which of its own components do not earn their
place.

\subsection*{Contributions}

\begin{itemize}
  \item \textbf{A physics accept/reject loop, and the measurement that
  shows it is load-bearing.} Fitting the labeller at ten seeds on
  identical preprocessed data moves the headline standby total by a
  coefficient of variation of $33.7\%$, and in three seeds of ten the
  cluster named \textsc{standby} is physically a different state
  --- once a loaded running state at $13$\,kW, twice the machine
  switched off at under $1$\,W. The physics battery scores exactly those
  three below threshold, and conditioning on a passing score collapses
  the spread to $3.0\%$. \emph{The pipeline is reproducible because of
  the physics validation, not despite the clustering}
  (Sections~\ref{sec:method:accept_reject},
  \ref{sec:significance:seeds}).

  \item \textbf{An account of what temporal coherence actually
  requires.} Viterbi decoding over a learned transition matrix reduces
  flicker substantially but --- contrary to the usual assumption ---
  cannot eliminate it, and neither can strengthening the transition
  prior. The penalty a transition matrix can impose on a state change is
  bounded by $\log(p_{\text{self}}/p_{\text{switch}})$, measured here at
  $9.23$ nats, whereas full-covariance Gaussian emissions fitted to
  low-noise electrical data differ by a median of $27.8$ nats and a
  $90$th percentile of $41\,701$ nats between adjacent states. At
  $55.7\%$ of timesteps the emission term is arithmetically unable to be
  overruled. We show the remedy is an explicit minimum-dwell constraint,
  that it removes flicker entirely at a cost of $0.19\%$ of rows, and
  that it does so on all eight machines
  (Sections~\ref{sec:flicker}, \ref{sec:multimachine:general}).

  \item \textbf{A validation methodology for state partitions in the
  absence of ground truth.} Circuit-level physics, model-level
  diagnostics, and independently derived Otsu-thresholded classifiers
  \cite{otsu1979threshold} are combined by consensus
  \cite{strehl2002cluster}, none of them using the model's own internal
  statistics. Applied across eight machines the battery accepts the four
  motor-driven ones and rejects the four auxiliary ones, on which there
  is no energised-idle state to find --- which is the behaviour that
  makes it a validation procedure rather than a formality
  (Sections~\ref{sec:validation}, \ref{sec:multimachine:states}).

  \item \textbf{A forecasting comparison anchored on not forecasting,
  and a negative result reported as such.} Under one protocol with
  seeds, an untrained persistence reference beats all five neural
  sequence models, the originally proposed sequence-to-sequence LSTM by
  the largest margin of the seven ($7.95$\,s of per-window MAE,
  $p=5.7\times10^{-44}$). A gradient-boosted tree is the only trained
  forecaster that clears the reference, and it is the forecasting
  component of the framework as reported here
  (Sections~\ref{sec:forecasting}, \ref{sec:significance:forecast}).

  \item \textbf{A decision layer in which the break-even duration is
  derived, with an invariance result and a competitive benchmark.} We
  replace the fixed-threshold rule with a constrained cost
  minimisation, show that the reduced problem admits an exact offline
  optimum, and measure the restart coefficients from the record instead
  of assuming them --- finding the studied machine has no measurable
  electrical restart penalty, so its restart cost is a production-delay
  cost. We show analytically and numerically that the break-even
  duration is \emph{independent of electricity tariff} whenever the
  restart cost is purely electrical, and report the decision surface
  over the axes that do carry information. The policy beats the fixed
  threshold it replaces by $+68.9$\,USD/yr ($p=0.001$); it is not shown
  to beat the plant's own behaviour, and we say so
  (Sections~\ref{sec:decision}, \ref{sec:significance:decision}).

  \item \textbf{End-to-end evaluation with ablation, significance
  testing and transfer.} Every stage is evaluated across eight machines
  and two sites, each component's marginal contribution is priced by
  ablation, every headline difference carries a dependence-aware
  interval and a pre-registered practical threshold, and the two designs
  that cannot detect anything are identified as such rather than
  reported as null results (Sections~\ref{sec:multimachine}--\ref{sec:ablation}).

  \item \textbf{Reproducibility.} The full pipeline --- preprocessing,
  inference, validation harness, forecasting, decision layer and every
  experiment reported here --- is released as open source, and all
  results are obtained from public benchmark data.
\end{itemize}

\subsection*{Organisation}

Section~\ref{sec:related} reviews related work across machine state
identification, non-intrusive load monitoring, industrial energy
optimisation and predictive shutdown, and states the resulting gap.
Section~\ref{sec:method} presents the framework.
Section~\ref{sec:states} reports the state layer, including the
comparison against clustering baselines, the flicker analysis, and the
ground-truth-free validation. Section~\ref{sec:forecasting} reports the
forecasting comparison and Section~\ref{sec:significance} the
statistical analysis on which every subsequent claim depends.
Section~\ref{sec:decision} develops the decision layer.
Sections~\ref{sec:multimachine} and~\ref{sec:crossmachine} extend the
evaluation to all eight machines, to transfer between them, and to a
second site. Section~\ref{sec:explain} attributes what each model uses,
Section~\ref{sec:ablation} prices each component, and
Section~\ref{sec:limitations} states the boundaries of the evidence
before Section~\ref{sec:conclusion} concludes.


% --- Four-theme review and the research gap
% =====================================================================
%  related_work.tex  --  PHASE I, Task 1
%
%  Structured four-theme literature review + cross-comparison table
%  + explicit research-gap statement.
%
%  Drop-in usage (IEEEtran or elsarticle):
%      % =====================================================================
%  related_work.tex  --  PHASE I, Task 1
%
%  Structured four-theme literature review + cross-comparison table
%  + explicit research-gap statement.
%
%  Drop-in usage (IEEEtran or elsarticle):
%      % =====================================================================
%  related_work.tex  --  PHASE I, Task 1
%
%  Structured four-theme literature review + cross-comparison table
%  + explicit research-gap statement.
%
%  Drop-in usage (IEEEtran or elsarticle):
%      \input{related_work}
%  Requires in the preamble:
%      \usepackage{booktabs}
%      \usepackage{threeparttable}   % for the table notes
%      \usepackage{multirow}
%      \usepackage{wasysym}          % \CIRCLE \LEFTcircle \Circle
%  Note: wasysym redefines \Box/\diamond; if that clashes with your
%  template, load it as \usepackage[nointegrals]{wasysym} or swap the
%  harvey-ball glyphs for {\ding{108}}/{\ding{109}} from pifont.
%  Bibliography: paper/references.bib
% =====================================================================

\section{Related Work}
\label{sec:related}

The problem addressed in this paper sits at the intersection of four
research communities that have so far developed largely in isolation:
(i)~unsupervised identification of machine operating states from
electrical measurements, (ii)~non-intrusive load monitoring (NILM),
(iii)~industrial energy optimisation at the process and system level,
and (iv)~predictive shutdown and decision support for manufacturing
equipment. We review each in turn, then summarise the resulting gap.


% ---------------------------------------------------------------------
\subsection{Machine State Identification}
\label{sec:rw:states}

The recognition that an electrically driven machine occupies a small
number of discrete operating regimes---typically \textsc{off},
\textsc{standby} (energised but not producing), and \textsc{working}---is
the foundation of essentially all energy-state analysis. ISO~14955-1
\cite{iso14955} formalises this decomposition as a design-time
requirement for energy-efficient machine tools, defining standby as the
operating mode in which a machine is ready for production but performs
no value-adding function. This definition is the one adopted throughout
the present work.

\subsubsection{Supervised approaches}
Where per-timestamp annotations are available, state identification
reduces to a standard classification problem, and support-vector
machines, random forests, and convolutional classifiers all perform
well. The practical obstacle in industrial deployments is that such
annotations essentially never exist: obtaining them requires either
instrumenting the machine's contactor and PLC directly or having an
operator observe each machine over an extended period. Both are
prohibitively expensive at facility scale, and neither transfers to a
new machine without repeating the effort. Supervised methods are
therefore poorly matched to the retrofit scenario that motivates this
paper.

\subsubsection{Unsupervised approaches}
Unsupervised state inference avoids the annotation cost entirely.
$k$-means \cite{lloyd1982least} is the simplest option but imposes
isotropic, equal-variance clusters, which is a poor model for
electrical states whose dispersion grows with load. Gaussian mixture
models (GMM) \cite{mclachlan2000finite}, fitted by
expectation--maximisation \cite{dempster1977em}, relax both assumptions:
each state receives its own mean and full covariance, and assignments
are soft, which matters precisely at the \textsc{standby}/\textsc{working}
boundary where the two regimes overlap. Density-based methods such as
DBSCAN \cite{ester1996dbscan} and HDBSCAN \cite{campello2015hdbscan}
make no parametric shape assumption and identify noise points
explicitly, while spectral clustering \cite{ng2001spectral} operates on
the affinity graph and can recover non-convex structure. Model order is
conventionally selected by the Bayesian information criterion
\cite{schwarz1978bic} or by silhouette analysis
\cite{rousseeuw1987silhouettes}.

Neither criterion selects a state count on records of the size
considered here, and this is worth stating because it determines the
design of the method proposed below. The BIC penalty grows as
$\tfrac{1}{2}p\log n$ while the likelihood gained by splitting a
non-Gaussian component grows linearly in $n$, so at $n>10^{6}$ samples
per machine the penalty is negligible and a mixture model keeps buying
components in order to approximate a distribution that is not a mixture
of $k$ Gaussians for any $k$. Scanning $k\in[2,8]$ over the eight
machines of Section~\ref{sec:multimachine}, BIC is still falling at the
boundary on seven of them, and the improvement from $k=4$ to $k=8$ is
$3.1$--$33.1\%$ of the criterion's own value. Silhouette behaves
similarly. Statistical order selection must therefore be supplemented
by an external, physical admissibility criterion rather than merely
confirmed by one --- which is the role of the validation battery of
Section~\ref{sec:method:accept_reject}, and the reason it is part of
the method here rather than an appendix to it.

\subsubsection{Temporal structure}
A limitation common to all of the above is that they treat samples as
independent and identically distributed. Applied to a 1\,Hz electrical
stream, an i.i.d.\ clustering assigns each second in isolation and
consequently produces rapid, physically impossible state alternation
around cluster boundaries---a failure mode we quantify as the
\emph{flicker rate} in Section~\ref{sec:results}. Hidden Markov models
\cite{rabiner1989tutorial} address this by making the state sequence
itself the object of inference: the transition matrix encodes state
persistence, and Viterbi decoding returns the globally most probable
sequence rather than a concatenation of independent decisions. Fox
\textit{et al.} \cite{fox2011sticky} show that an explicit bias toward
self-transition (a ``sticky'' prior) reduces over-segmentation when the
emission distributions overlap, which is the regime encountered here.

It is worth stating precisely what this machinery can and cannot do,
because the point is routinely overstated. In a Viterbi decode the cost
of changing state is $\log(p_{\text{self}}/p_{\text{switch}})$, which is
a few nats even for extremely persistent transition matrices and is
bounded above however strong the prior is made. The emission term
carries no such bound: full-covariance Gaussians fitted to low-noise,
multi-channel electrical measurements can differ by hundreds of nats
between adjacent states. Where that holds, the transition model is
arithmetically unable to influence the decoded path and the HMM
degenerates toward per-sample assignment. Guaranteeing a minimum state
duration therefore requires modelling duration explicitly---a hidden
semi-Markov formulation \cite{yu2010hsmm}, or an equivalent
duration constraint imposed at decode time---rather than a heavier
transition prior. Section~\ref{sec:results} quantifies this for the
present data.

\subsubsection{Threshold-free decision rules}
Where a scalar decision boundary is unavoidable---for instance when
converting a continuous power-factor distribution into a binary
load/no-load verdict---Otsu's criterion \cite{otsu1979threshold}
selects the threshold that maximises between-class variance, removing
the arbitrariness of a hand-tuned cut-off. Agreement among several
independently derived partitions can then be aggregated using cluster
ensemble methods \cite{strehl2002cluster}.


% ---------------------------------------------------------------------
\subsection{Non-Intrusive Load Monitoring}
\label{sec:rw:nilm}

NILM originates with Hart \cite{hart1992nonintrusive}, who showed that
individual appliance activity can be recovered from an aggregate meter
by detecting steady-state step changes in real and reactive power.
Subsequent surveys \cite{zeifman2011nilm,schirmer2023nilmreview} trace
the field's progression from event-based signature matching to
probabilistic and, more recently, deep-learning formulations. Kolter
and Jaakkola \cite{kolter2012afamap} cast disaggregation as inference in
an additive factorial HMM, making explicit the connection between NILM
and the temporal state models of Section~\ref{sec:rw:states}. The deep
learning era is represented by Neural NILM \cite{kelly2015neuralnilm},
sequence-to-point regression \cite{zhang2018seq2point}, and denoising
autoencoders \cite{bonfigli2018denoising}. Standardised toolkits
\cite{batra2014nilmtk} and evaluation protocols
\cite{makonin2015performance} have made results across this line of
work broadly comparable.

Two properties of this literature are decisive for the present paper.
First, it is overwhelmingly \emph{residential}. The canonical benchmarks
---REDD \cite{kolter2011redd} and UK-DALE---consist of household
appliances whose signatures are sparse, short, and well separated.
Industrial machinery behaves very differently: loads are large,
continuously modulated by process throughput, dominated by induction
motors whose reactive power draw persists at zero mechanical output, and
governed by shift schedules rather than occupant behaviour. Holmegaard
and Kj\ae{}rgaard \cite{holmegaard2016industrial} evaluate established
NILM algorithms on industrial loads and report substantially degraded
performance, attributing it to precisely these differences. Martins
\textit{et al.} \cite{martins2018imdeld} apply a deep generative model
to industrial disaggregation and, in doing so, release the IMDELD
dataset \cite{imdeld2018dataset}---eleven meters recording eight heavy
machines in a Brazilian poultry-feed pellet plant at 1\,Hz across a
five-month window, of which the covered fraction differs by machine
(Section~\ref{sec:method:data}). IMDELD is the benchmark used throughout this
paper; its authors characterise the machines as three-state systems but
supply no per-timestamp state annotation, which is the reason the
present work must proceed without labels and validate its partition
against physics instead.

Second, and more fundamentally, the objective of NILM is
\emph{attribution}: given an aggregate signal, determine which appliance
consumed what. It is a retrospective accounting exercise. NILM does not
forecast, and it does not act. No NILM system, industrial or
residential, closes the loop from an inferred state to an actuation
decision. That step belongs to the two literatures reviewed next.


% ---------------------------------------------------------------------
\subsection{Industrial Energy Optimisation}
\label{sec:rw:industrial}

The manufacturing engineering community has independently established
that non-productive energy dominates the energy budget of discrete-part
production. Gutowski \textit{et al.} \cite{gutowski2006electrical} and
Dahmus and Gutowski \cite{dahmus2004environmental} show that the
load-independent baseline of a machine tool---spindle idling, hydraulics,
coolant pumps, control electronics---frequently exceeds the energy
actually delivered to the workpiece, and that this fraction rises as
automation increases. Kara and Li \cite{kara2011unit} formalise the
relationship as a unit-process model in which specific energy
consumption is inversely related to material removal rate, so that a
lightly loaded or idle machine is intrinsically inefficient. Devoldere
\textit{et al.} \cite{devoldere2007improvement} quantify the resulting
improvement potential and identify idle-mode elimination as the single
largest lever available without redesigning the machine. Duflou
\textit{et al.} \cite{duflou2012towards} synthesise this body of work
into a systems-level roadmap, placing operational measures---switching
equipment off when it is not needed---at the machine and cell levels of
their hierarchy.

On the management side, ISO~50001 \cite{iso50001} provides the
plan--do--check--act framework within which energy performance
indicators are defined and tracked, and ISO~14955-1 \cite{iso14955}
supplies the operating-mode taxonomy that makes standby energy a
measurable quantity in the first place.

The limitation of this literature is methodological rather than
conceptual. Its energy models are typically built from controlled
laboratory measurements of a single machine under prescribed process
parameters, and its recommendations are design guidelines and audit
procedures. It does not address how a facility, given only the
electrical measurements it already collects and no process-parameter
instrumentation, should determine \emph{which} of its machines are idle,
\emph{when}, and for \emph{how long}. That inference problem is left
open.


% ---------------------------------------------------------------------
\subsection{Predictive Shutdown and Decision Support}
\label{sec:rw:decision}

The operations research community has studied the shutdown decision
itself in considerable depth. Mouzon \textit{et al.}
\cite{mouzon2007operational} establish the central trade-off: turning a
machine off saves its idle power but incurs a restart energy penalty and
a restart delay, so shutdown is economical only when the idle period
exceeds a break-even duration determined by the ratio of restart energy
to idle power. They propose dispatching rules and a mathematical
programming formulation that exploit this, and extend the treatment to
joint energy--tardiness objectives \cite{mouzon2008framework}. Chen
\textit{et al.} \cite{chen2013schedule} develop schedule-based operation
policies for production lines, Sun and Li \cite{sun2013opportunity}
estimate real-time control opportunities from buffer state, Li and Sun
\cite{li2016energy} formulate dynamic energy control as a Markov
decision process, Frigerio and Matta \cite{frigerio2015energy} derive
control strategies for machine tools under stochastic job arrivals, and
He \textit{et al.} \cite{he2015energy} embed energy responsiveness in
job-shop machine selection and sequencing.

This literature is mathematically mature, and the break-even
formulation of Mouzon \textit{et al.} is the natural objective for the
decision layer of any standby-elimination system. Its universal
assumption, however, is that the idle interval is \emph{known}---either
because it is dictated by a deterministic production schedule, or
because it is drawn from an analytically tractable arrival process
calibrated by the modeller. The policies are evaluated in simulation or
on analytical queueing models. What is absent is any mechanism for
obtaining the idle-duration input from a machine's own measured
electrical behaviour, which is the only signal available in a retrofit
installation. In other words: this community can decide optimally given
a duration, and the NILM community can infer a state from
measurements, but no work connects the two.


% ---------------------------------------------------------------------
\subsection{Synthesis and Research Gap}
\label{sec:rw:gap}

Table~\ref{tab:related} compares representative works from the four
themes against the capabilities required for an autonomous
standby-elimination system: whether the method infers machine state
without labels, whether it models temporal structure explicitly, whether
it forecasts the \emph{duration} of a future non-productive interval,
whether it converts that forecast into an economically justified control
decision, and whether it is validated on real industrial measurements.

\begin{table*}[!t]
\centering
\caption{Comparison of representative prior work against the
capabilities required for autonomous standby elimination.
\CIRCLE~= provided, \LEFTcircle~= partially provided, \Circle~= not
provided, --~= not applicable.}
\label{tab:related}
\begin{threeparttable}
\small
\setlength{\tabcolsep}{4pt}
\begin{tabular}{@{}lllccccc@{}}
\toprule
\multirow{2}{*}{\textbf{Work}} &
\multirow{2}{*}{\textbf{Method family}} &
\multirow{2}{*}{\textbf{Data domain}} &
\textbf{Labels} &
\textbf{Temporal} &
\textbf{Duration} &
\textbf{Economic} &
\textbf{Real ind.} \\
 & & & \textbf{required} & \textbf{model} & \textbf{forecast} &
 \textbf{decision} & \textbf{data} \\
\midrule
\multicolumn{8}{@{}l}{\textit{Theme A --- machine state identification}}\\
Lloyd \cite{lloyd1982least}          & $k$-means                    & generic                & no  & \Circle    & \Circle & \Circle & \Circle \\
McLachlan \& Peel \cite{mclachlan2000finite} & Gaussian mixture     & generic                & no  & \Circle    & \Circle & \Circle & \Circle \\
Ester \textit{et al.} \cite{ester1996dbscan} & DBSCAN               & generic                & no  & \Circle    & \Circle & \Circle & \Circle \\
Fox \textit{et al.} \cite{fox2011sticky}     & sticky HDP-HMM       & audio                  & no  & \CIRCLE    & \Circle & \Circle & \Circle \\
\addlinespace
\multicolumn{8}{@{}l}{\textit{Theme B --- non-intrusive load monitoring}}\\
Hart \cite{hart1992nonintrusive}     & event/steady-state signature & residential            & yes\tnote{a} & \LEFTcircle & \Circle & \Circle & \Circle \\
Kolter \& Jaakkola \cite{kolter2012afamap} & additive factorial HMM & REDD (residential)     & no  & \CIRCLE    & \Circle & \Circle & \Circle \\
Kelly \& Knottenbelt \cite{kelly2015neuralnilm} & deep NN (dAE, LSTM) & UK-DALE (residential) & yes & \LEFTcircle & \Circle & \Circle & \Circle \\
Zhang \textit{et al.} \cite{zhang2018seq2point} & seq2point CNN     & UK-DALE, REDD          & yes & \LEFTcircle & \Circle & \Circle & \Circle \\
Holmegaard \& Kj\ae{}rgaard \cite{holmegaard2016industrial} & NILM algorithm evaluation & industrial & yes & \LEFTcircle & \Circle & \Circle & \CIRCLE \\
Martins \textit{et al.} \cite{martins2018imdeld} & deep generative disaggregation & IMDELD (industrial) & yes & \LEFTcircle & \Circle & \Circle & \CIRCLE \\
\addlinespace
\multicolumn{8}{@{}l}{\textit{Theme C --- industrial energy optimisation}}\\
Gutowski \textit{et al.} \cite{gutowski2006electrical} & thermodynamic/empirical process model & lab machine tools & -- & \Circle & \Circle & \Circle & \LEFTcircle \\
Kara \& Li \cite{kara2011unit}       & unit-process energy model    & lab machine tools      & --  & \Circle    & \Circle & \Circle & \LEFTcircle \\
Duflou \textit{et al.} \cite{duflou2012towards} & systems-level roadmap & survey             & --  & \Circle    & \Circle & \LEFTcircle & \LEFTcircle \\
\addlinespace
\multicolumn{8}{@{}l}{\textit{Theme D --- predictive shutdown and decision support}}\\
Mouzon \textit{et al.} \cite{mouzon2007operational} & dispatching rules, MILP & simulated single machine & -- & \Circle & \Circle & \CIRCLE & \Circle \\
Chen \textit{et al.} \cite{chen2013schedule} & schedule-based control  & simulated line     & --  & \LEFTcircle & \Circle & \CIRCLE & \Circle \\
Frigerio \& Matta \cite{frigerio2015energy} & stochastic control policy & analytical model  & --  & \CIRCLE    & \LEFTcircle\tnote{b} & \CIRCLE & \Circle \\
He \textit{et al.} \cite{he2015energy} & energy-responsive scheduling & job-shop model      & --  & \Circle    & \Circle & \CIRCLE & \Circle \\
\midrule
\textbf{This work}                   & \textbf{physics-gated GMM--HMM $\rightarrow$ duration forecast $\rightarrow$ optimisation} & \textbf{IMDELD (industrial)} & \textbf{no} & \CIRCLE & \CIRCLE & \CIRCLE & \CIRCLE \\
\bottomrule
\end{tabular}
\begin{tablenotes}[flushleft]\footnotesize
\item[a] Hart's formulation is unsupervised in its clustering step but
requires a manually curated appliance signature library to name the
recovered clusters.
\item[b] Idle duration is modelled as a known stochastic arrival
process calibrated by the analyst, not forecast from measured machine
behaviour.
\end{tablenotes}
\end{threeparttable}
\end{table*}

Reading Table~\ref{tab:related} column-wise makes the gap explicit.
Every capability required for autonomous standby elimination has been
solved somewhere, but the last column of any single row is always
empty:

\begin{itemize}
  \item \textbf{Label-free state inference with temporal structure}
  exists (Theme~A, and \cite{kolter2012afamap} in Theme~B), but is
  demonstrated on generic or residential data and terminates at the
  labelled sequence. Nothing is predicted and nothing is decided.

  \item \textbf{Industrial validation} exists
  (\cite{holmegaard2016industrial,martins2018imdeld}, Theme~B), but the
  task is disaggregation, the methods are supervised, and the output is
  a retrospective attribution of energy rather than an operating state
  usable for control.

  \item \textbf{Quantification of idle-energy waste} exists (Theme~C),
  but rests on controlled laboratory measurement of individual
  processes and yields design guidance, not an online inference
  procedure applicable to a running facility.

  \item \textbf{Optimal shutdown decisions} exist (Theme~D), and the
  break-even formulation is well understood, but every such policy
  takes the idle duration as a \emph{given} input---supplied by a
  deterministic schedule or an analyst-calibrated arrival
  process---and is evaluated in simulation.
\end{itemize}

\medskip
\noindent\fbox{\parbox{0.97\columnwidth}{%
\textbf{Research gap.} No existing work combines
(i)~label-free, temporally coherent inference of industrial machine
states from electrical measurements alone, (ii)~forecasting of the
\emph{duration} of the resulting non-productive intervals, and
(iii)~an economically optimised shutdown decision derived from that
forecast, validated end-to-end on real industrial measurement data.
The state-inference and decision-making capabilities have matured
independently and remain unconnected: the former produces labels no
controller consumes, the latter consumes durations no inference
procedure supplies.}}
\medskip

\noindent
Closing this gap requires more than concatenating existing components.
Three problems arise specifically at the interfaces. First, because no
ground-truth state annotation exists for industrial data, the inferred
partition cannot be validated by conventional supervised metrics, and an
alternative validation basis---independent physical and operational
evidence---must be constructed. Second, an i.i.d.\ clustering produces a
label sequence whose flicker makes duration forecasting meaningless,
so temporal coherence is a prerequisite for the forecasting stage rather
than a refinement of the inference stage. Third, a forecast duration
carries uncertainty that a deterministic break-even threshold ignores;
the decision layer must therefore be posed as a constrained
cost-minimisation problem in which the break-even point emerges as a
derived quantity rather than being supplied as a tuning parameter.
Section~\ref{sec:method} addresses each of these in turn.

%  Requires in the preamble:
%      \usepackage{booktabs}
%      \usepackage{threeparttable}   % for the table notes
%      \usepackage{multirow}
%      \usepackage{wasysym}          % \CIRCLE \LEFTcircle \Circle
%  Note: wasysym redefines \Box/\diamond; if that clashes with your
%  template, load it as \usepackage[nointegrals]{wasysym} or swap the
%  harvey-ball glyphs for {\ding{108}}/{\ding{109}} from pifont.
%  Bibliography: paper/references.bib
% =====================================================================

\section{Related Work}
\label{sec:related}

The problem addressed in this paper sits at the intersection of four
research communities that have so far developed largely in isolation:
(i)~unsupervised identification of machine operating states from
electrical measurements, (ii)~non-intrusive load monitoring (NILM),
(iii)~industrial energy optimisation at the process and system level,
and (iv)~predictive shutdown and decision support for manufacturing
equipment. We review each in turn, then summarise the resulting gap.


% ---------------------------------------------------------------------
\subsection{Machine State Identification}
\label{sec:rw:states}

The recognition that an electrically driven machine occupies a small
number of discrete operating regimes---typically \textsc{off},
\textsc{standby} (energised but not producing), and \textsc{working}---is
the foundation of essentially all energy-state analysis. ISO~14955-1
\cite{iso14955} formalises this decomposition as a design-time
requirement for energy-efficient machine tools, defining standby as the
operating mode in which a machine is ready for production but performs
no value-adding function. This definition is the one adopted throughout
the present work.

\subsubsection{Supervised approaches}
Where per-timestamp annotations are available, state identification
reduces to a standard classification problem, and support-vector
machines, random forests, and convolutional classifiers all perform
well. The practical obstacle in industrial deployments is that such
annotations essentially never exist: obtaining them requires either
instrumenting the machine's contactor and PLC directly or having an
operator observe each machine over an extended period. Both are
prohibitively expensive at facility scale, and neither transfers to a
new machine without repeating the effort. Supervised methods are
therefore poorly matched to the retrofit scenario that motivates this
paper.

\subsubsection{Unsupervised approaches}
Unsupervised state inference avoids the annotation cost entirely.
$k$-means \cite{lloyd1982least} is the simplest option but imposes
isotropic, equal-variance clusters, which is a poor model for
electrical states whose dispersion grows with load. Gaussian mixture
models (GMM) \cite{mclachlan2000finite}, fitted by
expectation--maximisation \cite{dempster1977em}, relax both assumptions:
each state receives its own mean and full covariance, and assignments
are soft, which matters precisely at the \textsc{standby}/\textsc{working}
boundary where the two regimes overlap. Density-based methods such as
DBSCAN \cite{ester1996dbscan} and HDBSCAN \cite{campello2015hdbscan}
make no parametric shape assumption and identify noise points
explicitly, while spectral clustering \cite{ng2001spectral} operates on
the affinity graph and can recover non-convex structure. Model order is
conventionally selected by the Bayesian information criterion
\cite{schwarz1978bic} or by silhouette analysis
\cite{rousseeuw1987silhouettes}.

Neither criterion selects a state count on records of the size
considered here, and this is worth stating because it determines the
design of the method proposed below. The BIC penalty grows as
$\tfrac{1}{2}p\log n$ while the likelihood gained by splitting a
non-Gaussian component grows linearly in $n$, so at $n>10^{6}$ samples
per machine the penalty is negligible and a mixture model keeps buying
components in order to approximate a distribution that is not a mixture
of $k$ Gaussians for any $k$. Scanning $k\in[2,8]$ over the eight
machines of Section~\ref{sec:multimachine}, BIC is still falling at the
boundary on seven of them, and the improvement from $k=4$ to $k=8$ is
$3.1$--$33.1\%$ of the criterion's own value. Silhouette behaves
similarly. Statistical order selection must therefore be supplemented
by an external, physical admissibility criterion rather than merely
confirmed by one --- which is the role of the validation battery of
Section~\ref{sec:method:accept_reject}, and the reason it is part of
the method here rather than an appendix to it.

\subsubsection{Temporal structure}
A limitation common to all of the above is that they treat samples as
independent and identically distributed. Applied to a 1\,Hz electrical
stream, an i.i.d.\ clustering assigns each second in isolation and
consequently produces rapid, physically impossible state alternation
around cluster boundaries---a failure mode we quantify as the
\emph{flicker rate} in Section~\ref{sec:results}. Hidden Markov models
\cite{rabiner1989tutorial} address this by making the state sequence
itself the object of inference: the transition matrix encodes state
persistence, and Viterbi decoding returns the globally most probable
sequence rather than a concatenation of independent decisions. Fox
\textit{et al.} \cite{fox2011sticky} show that an explicit bias toward
self-transition (a ``sticky'' prior) reduces over-segmentation when the
emission distributions overlap, which is the regime encountered here.

It is worth stating precisely what this machinery can and cannot do,
because the point is routinely overstated. In a Viterbi decode the cost
of changing state is $\log(p_{\text{self}}/p_{\text{switch}})$, which is
a few nats even for extremely persistent transition matrices and is
bounded above however strong the prior is made. The emission term
carries no such bound: full-covariance Gaussians fitted to low-noise,
multi-channel electrical measurements can differ by hundreds of nats
between adjacent states. Where that holds, the transition model is
arithmetically unable to influence the decoded path and the HMM
degenerates toward per-sample assignment. Guaranteeing a minimum state
duration therefore requires modelling duration explicitly---a hidden
semi-Markov formulation \cite{yu2010hsmm}, or an equivalent
duration constraint imposed at decode time---rather than a heavier
transition prior. Section~\ref{sec:results} quantifies this for the
present data.

\subsubsection{Threshold-free decision rules}
Where a scalar decision boundary is unavoidable---for instance when
converting a continuous power-factor distribution into a binary
load/no-load verdict---Otsu's criterion \cite{otsu1979threshold}
selects the threshold that maximises between-class variance, removing
the arbitrariness of a hand-tuned cut-off. Agreement among several
independently derived partitions can then be aggregated using cluster
ensemble methods \cite{strehl2002cluster}.


% ---------------------------------------------------------------------
\subsection{Non-Intrusive Load Monitoring}
\label{sec:rw:nilm}

NILM originates with Hart \cite{hart1992nonintrusive}, who showed that
individual appliance activity can be recovered from an aggregate meter
by detecting steady-state step changes in real and reactive power.
Subsequent surveys \cite{zeifman2011nilm,schirmer2023nilmreview} trace
the field's progression from event-based signature matching to
probabilistic and, more recently, deep-learning formulations. Kolter
and Jaakkola \cite{kolter2012afamap} cast disaggregation as inference in
an additive factorial HMM, making explicit the connection between NILM
and the temporal state models of Section~\ref{sec:rw:states}. The deep
learning era is represented by Neural NILM \cite{kelly2015neuralnilm},
sequence-to-point regression \cite{zhang2018seq2point}, and denoising
autoencoders \cite{bonfigli2018denoising}. Standardised toolkits
\cite{batra2014nilmtk} and evaluation protocols
\cite{makonin2015performance} have made results across this line of
work broadly comparable.

Two properties of this literature are decisive for the present paper.
First, it is overwhelmingly \emph{residential}. The canonical benchmarks
---REDD \cite{kolter2011redd} and UK-DALE---consist of household
appliances whose signatures are sparse, short, and well separated.
Industrial machinery behaves very differently: loads are large,
continuously modulated by process throughput, dominated by induction
motors whose reactive power draw persists at zero mechanical output, and
governed by shift schedules rather than occupant behaviour. Holmegaard
and Kj\ae{}rgaard \cite{holmegaard2016industrial} evaluate established
NILM algorithms on industrial loads and report substantially degraded
performance, attributing it to precisely these differences. Martins
\textit{et al.} \cite{martins2018imdeld} apply a deep generative model
to industrial disaggregation and, in doing so, release the IMDELD
dataset \cite{imdeld2018dataset}---eleven meters recording eight heavy
machines in a Brazilian poultry-feed pellet plant at 1\,Hz across a
five-month window, of which the covered fraction differs by machine
(Section~\ref{sec:method:data}). IMDELD is the benchmark used throughout this
paper; its authors characterise the machines as three-state systems but
supply no per-timestamp state annotation, which is the reason the
present work must proceed without labels and validate its partition
against physics instead.

Second, and more fundamentally, the objective of NILM is
\emph{attribution}: given an aggregate signal, determine which appliance
consumed what. It is a retrospective accounting exercise. NILM does not
forecast, and it does not act. No NILM system, industrial or
residential, closes the loop from an inferred state to an actuation
decision. That step belongs to the two literatures reviewed next.


% ---------------------------------------------------------------------
\subsection{Industrial Energy Optimisation}
\label{sec:rw:industrial}

The manufacturing engineering community has independently established
that non-productive energy dominates the energy budget of discrete-part
production. Gutowski \textit{et al.} \cite{gutowski2006electrical} and
Dahmus and Gutowski \cite{dahmus2004environmental} show that the
load-independent baseline of a machine tool---spindle idling, hydraulics,
coolant pumps, control electronics---frequently exceeds the energy
actually delivered to the workpiece, and that this fraction rises as
automation increases. Kara and Li \cite{kara2011unit} formalise the
relationship as a unit-process model in which specific energy
consumption is inversely related to material removal rate, so that a
lightly loaded or idle machine is intrinsically inefficient. Devoldere
\textit{et al.} \cite{devoldere2007improvement} quantify the resulting
improvement potential and identify idle-mode elimination as the single
largest lever available without redesigning the machine. Duflou
\textit{et al.} \cite{duflou2012towards} synthesise this body of work
into a systems-level roadmap, placing operational measures---switching
equipment off when it is not needed---at the machine and cell levels of
their hierarchy.

On the management side, ISO~50001 \cite{iso50001} provides the
plan--do--check--act framework within which energy performance
indicators are defined and tracked, and ISO~14955-1 \cite{iso14955}
supplies the operating-mode taxonomy that makes standby energy a
measurable quantity in the first place.

The limitation of this literature is methodological rather than
conceptual. Its energy models are typically built from controlled
laboratory measurements of a single machine under prescribed process
parameters, and its recommendations are design guidelines and audit
procedures. It does not address how a facility, given only the
electrical measurements it already collects and no process-parameter
instrumentation, should determine \emph{which} of its machines are idle,
\emph{when}, and for \emph{how long}. That inference problem is left
open.


% ---------------------------------------------------------------------
\subsection{Predictive Shutdown and Decision Support}
\label{sec:rw:decision}

The operations research community has studied the shutdown decision
itself in considerable depth. Mouzon \textit{et al.}
\cite{mouzon2007operational} establish the central trade-off: turning a
machine off saves its idle power but incurs a restart energy penalty and
a restart delay, so shutdown is economical only when the idle period
exceeds a break-even duration determined by the ratio of restart energy
to idle power. They propose dispatching rules and a mathematical
programming formulation that exploit this, and extend the treatment to
joint energy--tardiness objectives \cite{mouzon2008framework}. Chen
\textit{et al.} \cite{chen2013schedule} develop schedule-based operation
policies for production lines, Sun and Li \cite{sun2013opportunity}
estimate real-time control opportunities from buffer state, Li and Sun
\cite{li2016energy} formulate dynamic energy control as a Markov
decision process, Frigerio and Matta \cite{frigerio2015energy} derive
control strategies for machine tools under stochastic job arrivals, and
He \textit{et al.} \cite{he2015energy} embed energy responsiveness in
job-shop machine selection and sequencing.

This literature is mathematically mature, and the break-even
formulation of Mouzon \textit{et al.} is the natural objective for the
decision layer of any standby-elimination system. Its universal
assumption, however, is that the idle interval is \emph{known}---either
because it is dictated by a deterministic production schedule, or
because it is drawn from an analytically tractable arrival process
calibrated by the modeller. The policies are evaluated in simulation or
on analytical queueing models. What is absent is any mechanism for
obtaining the idle-duration input from a machine's own measured
electrical behaviour, which is the only signal available in a retrofit
installation. In other words: this community can decide optimally given
a duration, and the NILM community can infer a state from
measurements, but no work connects the two.


% ---------------------------------------------------------------------
\subsection{Synthesis and Research Gap}
\label{sec:rw:gap}

Table~\ref{tab:related} compares representative works from the four
themes against the capabilities required for an autonomous
standby-elimination system: whether the method infers machine state
without labels, whether it models temporal structure explicitly, whether
it forecasts the \emph{duration} of a future non-productive interval,
whether it converts that forecast into an economically justified control
decision, and whether it is validated on real industrial measurements.

\begin{table*}[!t]
\centering
\caption{Comparison of representative prior work against the
capabilities required for autonomous standby elimination.
\CIRCLE~= provided, \LEFTcircle~= partially provided, \Circle~= not
provided, --~= not applicable.}
\label{tab:related}
\begin{threeparttable}
\small
\setlength{\tabcolsep}{4pt}
\begin{tabular}{@{}lllccccc@{}}
\toprule
\multirow{2}{*}{\textbf{Work}} &
\multirow{2}{*}{\textbf{Method family}} &
\multirow{2}{*}{\textbf{Data domain}} &
\textbf{Labels} &
\textbf{Temporal} &
\textbf{Duration} &
\textbf{Economic} &
\textbf{Real ind.} \\
 & & & \textbf{required} & \textbf{model} & \textbf{forecast} &
 \textbf{decision} & \textbf{data} \\
\midrule
\multicolumn{8}{@{}l}{\textit{Theme A --- machine state identification}}\\
Lloyd \cite{lloyd1982least}          & $k$-means                    & generic                & no  & \Circle    & \Circle & \Circle & \Circle \\
McLachlan \& Peel \cite{mclachlan2000finite} & Gaussian mixture     & generic                & no  & \Circle    & \Circle & \Circle & \Circle \\
Ester \textit{et al.} \cite{ester1996dbscan} & DBSCAN               & generic                & no  & \Circle    & \Circle & \Circle & \Circle \\
Fox \textit{et al.} \cite{fox2011sticky}     & sticky HDP-HMM       & audio                  & no  & \CIRCLE    & \Circle & \Circle & \Circle \\
\addlinespace
\multicolumn{8}{@{}l}{\textit{Theme B --- non-intrusive load monitoring}}\\
Hart \cite{hart1992nonintrusive}     & event/steady-state signature & residential            & yes\tnote{a} & \LEFTcircle & \Circle & \Circle & \Circle \\
Kolter \& Jaakkola \cite{kolter2012afamap} & additive factorial HMM & REDD (residential)     & no  & \CIRCLE    & \Circle & \Circle & \Circle \\
Kelly \& Knottenbelt \cite{kelly2015neuralnilm} & deep NN (dAE, LSTM) & UK-DALE (residential) & yes & \LEFTcircle & \Circle & \Circle & \Circle \\
Zhang \textit{et al.} \cite{zhang2018seq2point} & seq2point CNN     & UK-DALE, REDD          & yes & \LEFTcircle & \Circle & \Circle & \Circle \\
Holmegaard \& Kj\ae{}rgaard \cite{holmegaard2016industrial} & NILM algorithm evaluation & industrial & yes & \LEFTcircle & \Circle & \Circle & \CIRCLE \\
Martins \textit{et al.} \cite{martins2018imdeld} & deep generative disaggregation & IMDELD (industrial) & yes & \LEFTcircle & \Circle & \Circle & \CIRCLE \\
\addlinespace
\multicolumn{8}{@{}l}{\textit{Theme C --- industrial energy optimisation}}\\
Gutowski \textit{et al.} \cite{gutowski2006electrical} & thermodynamic/empirical process model & lab machine tools & -- & \Circle & \Circle & \Circle & \LEFTcircle \\
Kara \& Li \cite{kara2011unit}       & unit-process energy model    & lab machine tools      & --  & \Circle    & \Circle & \Circle & \LEFTcircle \\
Duflou \textit{et al.} \cite{duflou2012towards} & systems-level roadmap & survey             & --  & \Circle    & \Circle & \LEFTcircle & \LEFTcircle \\
\addlinespace
\multicolumn{8}{@{}l}{\textit{Theme D --- predictive shutdown and decision support}}\\
Mouzon \textit{et al.} \cite{mouzon2007operational} & dispatching rules, MILP & simulated single machine & -- & \Circle & \Circle & \CIRCLE & \Circle \\
Chen \textit{et al.} \cite{chen2013schedule} & schedule-based control  & simulated line     & --  & \LEFTcircle & \Circle & \CIRCLE & \Circle \\
Frigerio \& Matta \cite{frigerio2015energy} & stochastic control policy & analytical model  & --  & \CIRCLE    & \LEFTcircle\tnote{b} & \CIRCLE & \Circle \\
He \textit{et al.} \cite{he2015energy} & energy-responsive scheduling & job-shop model      & --  & \Circle    & \Circle & \CIRCLE & \Circle \\
\midrule
\textbf{This work}                   & \textbf{physics-gated GMM--HMM $\rightarrow$ duration forecast $\rightarrow$ optimisation} & \textbf{IMDELD (industrial)} & \textbf{no} & \CIRCLE & \CIRCLE & \CIRCLE & \CIRCLE \\
\bottomrule
\end{tabular}
\begin{tablenotes}[flushleft]\footnotesize
\item[a] Hart's formulation is unsupervised in its clustering step but
requires a manually curated appliance signature library to name the
recovered clusters.
\item[b] Idle duration is modelled as a known stochastic arrival
process calibrated by the analyst, not forecast from measured machine
behaviour.
\end{tablenotes}
\end{threeparttable}
\end{table*}

Reading Table~\ref{tab:related} column-wise makes the gap explicit.
Every capability required for autonomous standby elimination has been
solved somewhere, but the last column of any single row is always
empty:

\begin{itemize}
  \item \textbf{Label-free state inference with temporal structure}
  exists (Theme~A, and \cite{kolter2012afamap} in Theme~B), but is
  demonstrated on generic or residential data and terminates at the
  labelled sequence. Nothing is predicted and nothing is decided.

  \item \textbf{Industrial validation} exists
  (\cite{holmegaard2016industrial,martins2018imdeld}, Theme~B), but the
  task is disaggregation, the methods are supervised, and the output is
  a retrospective attribution of energy rather than an operating state
  usable for control.

  \item \textbf{Quantification of idle-energy waste} exists (Theme~C),
  but rests on controlled laboratory measurement of individual
  processes and yields design guidance, not an online inference
  procedure applicable to a running facility.

  \item \textbf{Optimal shutdown decisions} exist (Theme~D), and the
  break-even formulation is well understood, but every such policy
  takes the idle duration as a \emph{given} input---supplied by a
  deterministic schedule or an analyst-calibrated arrival
  process---and is evaluated in simulation.
\end{itemize}

\medskip
\noindent\fbox{\parbox{0.97\columnwidth}{%
\textbf{Research gap.} No existing work combines
(i)~label-free, temporally coherent inference of industrial machine
states from electrical measurements alone, (ii)~forecasting of the
\emph{duration} of the resulting non-productive intervals, and
(iii)~an economically optimised shutdown decision derived from that
forecast, validated end-to-end on real industrial measurement data.
The state-inference and decision-making capabilities have matured
independently and remain unconnected: the former produces labels no
controller consumes, the latter consumes durations no inference
procedure supplies.}}
\medskip

\noindent
Closing this gap requires more than concatenating existing components.
Three problems arise specifically at the interfaces. First, because no
ground-truth state annotation exists for industrial data, the inferred
partition cannot be validated by conventional supervised metrics, and an
alternative validation basis---independent physical and operational
evidence---must be constructed. Second, an i.i.d.\ clustering produces a
label sequence whose flicker makes duration forecasting meaningless,
so temporal coherence is a prerequisite for the forecasting stage rather
than a refinement of the inference stage. Third, a forecast duration
carries uncertainty that a deterministic break-even threshold ignores;
the decision layer must therefore be posed as a constrained
cost-minimisation problem in which the break-even point emerges as a
derived quantity rather than being supplied as a tuning parameter.
Section~\ref{sec:method} addresses each of these in turn.

%  Requires in the preamble:
%      \usepackage{booktabs}
%      \usepackage{threeparttable}   % for the table notes
%      \usepackage{multirow}
%      \usepackage{wasysym}          % \CIRCLE \LEFTcircle \Circle
%  Note: wasysym redefines \Box/\diamond; if that clashes with your
%  template, load it as \usepackage[nointegrals]{wasysym} or swap the
%  harvey-ball glyphs for {\ding{108}}/{\ding{109}} from pifont.
%  Bibliography: paper/references.bib
% =====================================================================

\section{Related Work}
\label{sec:related}

The problem addressed in this paper sits at the intersection of four
research communities that have so far developed largely in isolation:
(i)~unsupervised identification of machine operating states from
electrical measurements, (ii)~non-intrusive load monitoring (NILM),
(iii)~industrial energy optimisation at the process and system level,
and (iv)~predictive shutdown and decision support for manufacturing
equipment. We review each in turn, then summarise the resulting gap.


% ---------------------------------------------------------------------
\subsection{Machine State Identification}
\label{sec:rw:states}

The recognition that an electrically driven machine occupies a small
number of discrete operating regimes---typically \textsc{off},
\textsc{standby} (energised but not producing), and \textsc{working}---is
the foundation of essentially all energy-state analysis. ISO~14955-1
\cite{iso14955} formalises this decomposition as a design-time
requirement for energy-efficient machine tools, defining standby as the
operating mode in which a machine is ready for production but performs
no value-adding function. This definition is the one adopted throughout
the present work.

\subsubsection{Supervised approaches}
Where per-timestamp annotations are available, state identification
reduces to a standard classification problem, and support-vector
machines, random forests, and convolutional classifiers all perform
well. The practical obstacle in industrial deployments is that such
annotations essentially never exist: obtaining them requires either
instrumenting the machine's contactor and PLC directly or having an
operator observe each machine over an extended period. Both are
prohibitively expensive at facility scale, and neither transfers to a
new machine without repeating the effort. Supervised methods are
therefore poorly matched to the retrofit scenario that motivates this
paper.

\subsubsection{Unsupervised approaches}
Unsupervised state inference avoids the annotation cost entirely.
$k$-means \cite{lloyd1982least} is the simplest option but imposes
isotropic, equal-variance clusters, which is a poor model for
electrical states whose dispersion grows with load. Gaussian mixture
models (GMM) \cite{mclachlan2000finite}, fitted by
expectation--maximisation \cite{dempster1977em}, relax both assumptions:
each state receives its own mean and full covariance, and assignments
are soft, which matters precisely at the \textsc{standby}/\textsc{working}
boundary where the two regimes overlap. Density-based methods such as
DBSCAN \cite{ester1996dbscan} and HDBSCAN \cite{campello2015hdbscan}
make no parametric shape assumption and identify noise points
explicitly, while spectral clustering \cite{ng2001spectral} operates on
the affinity graph and can recover non-convex structure. Model order is
conventionally selected by the Bayesian information criterion
\cite{schwarz1978bic} or by silhouette analysis
\cite{rousseeuw1987silhouettes}.

Neither criterion selects a state count on records of the size
considered here, and this is worth stating because it determines the
design of the method proposed below. The BIC penalty grows as
$\tfrac{1}{2}p\log n$ while the likelihood gained by splitting a
non-Gaussian component grows linearly in $n$, so at $n>10^{6}$ samples
per machine the penalty is negligible and a mixture model keeps buying
components in order to approximate a distribution that is not a mixture
of $k$ Gaussians for any $k$. Scanning $k\in[2,8]$ over the eight
machines of Section~\ref{sec:multimachine}, BIC is still falling at the
boundary on seven of them, and the improvement from $k=4$ to $k=8$ is
$3.1$--$33.1\%$ of the criterion's own value. Silhouette behaves
similarly. Statistical order selection must therefore be supplemented
by an external, physical admissibility criterion rather than merely
confirmed by one --- which is the role of the validation battery of
Section~\ref{sec:method:accept_reject}, and the reason it is part of
the method here rather than an appendix to it.

\subsubsection{Temporal structure}
A limitation common to all of the above is that they treat samples as
independent and identically distributed. Applied to a 1\,Hz electrical
stream, an i.i.d.\ clustering assigns each second in isolation and
consequently produces rapid, physically impossible state alternation
around cluster boundaries---a failure mode we quantify as the
\emph{flicker rate} in Section~\ref{sec:results}. Hidden Markov models
\cite{rabiner1989tutorial} address this by making the state sequence
itself the object of inference: the transition matrix encodes state
persistence, and Viterbi decoding returns the globally most probable
sequence rather than a concatenation of independent decisions. Fox
\textit{et al.} \cite{fox2011sticky} show that an explicit bias toward
self-transition (a ``sticky'' prior) reduces over-segmentation when the
emission distributions overlap, which is the regime encountered here.

It is worth stating precisely what this machinery can and cannot do,
because the point is routinely overstated. In a Viterbi decode the cost
of changing state is $\log(p_{\text{self}}/p_{\text{switch}})$, which is
a few nats even for extremely persistent transition matrices and is
bounded above however strong the prior is made. The emission term
carries no such bound: full-covariance Gaussians fitted to low-noise,
multi-channel electrical measurements can differ by hundreds of nats
between adjacent states. Where that holds, the transition model is
arithmetically unable to influence the decoded path and the HMM
degenerates toward per-sample assignment. Guaranteeing a minimum state
duration therefore requires modelling duration explicitly---a hidden
semi-Markov formulation \cite{yu2010hsmm}, or an equivalent
duration constraint imposed at decode time---rather than a heavier
transition prior. Section~\ref{sec:results} quantifies this for the
present data.

\subsubsection{Threshold-free decision rules}
Where a scalar decision boundary is unavoidable---for instance when
converting a continuous power-factor distribution into a binary
load/no-load verdict---Otsu's criterion \cite{otsu1979threshold}
selects the threshold that maximises between-class variance, removing
the arbitrariness of a hand-tuned cut-off. Agreement among several
independently derived partitions can then be aggregated using cluster
ensemble methods \cite{strehl2002cluster}.


% ---------------------------------------------------------------------
\subsection{Non-Intrusive Load Monitoring}
\label{sec:rw:nilm}

NILM originates with Hart \cite{hart1992nonintrusive}, who showed that
individual appliance activity can be recovered from an aggregate meter
by detecting steady-state step changes in real and reactive power.
Subsequent surveys \cite{zeifman2011nilm,schirmer2023nilmreview} trace
the field's progression from event-based signature matching to
probabilistic and, more recently, deep-learning formulations. Kolter
and Jaakkola \cite{kolter2012afamap} cast disaggregation as inference in
an additive factorial HMM, making explicit the connection between NILM
and the temporal state models of Section~\ref{sec:rw:states}. The deep
learning era is represented by Neural NILM \cite{kelly2015neuralnilm},
sequence-to-point regression \cite{zhang2018seq2point}, and denoising
autoencoders \cite{bonfigli2018denoising}. Standardised toolkits
\cite{batra2014nilmtk} and evaluation protocols
\cite{makonin2015performance} have made results across this line of
work broadly comparable.

Two properties of this literature are decisive for the present paper.
First, it is overwhelmingly \emph{residential}. The canonical benchmarks
---REDD \cite{kolter2011redd} and UK-DALE---consist of household
appliances whose signatures are sparse, short, and well separated.
Industrial machinery behaves very differently: loads are large,
continuously modulated by process throughput, dominated by induction
motors whose reactive power draw persists at zero mechanical output, and
governed by shift schedules rather than occupant behaviour. Holmegaard
and Kj\ae{}rgaard \cite{holmegaard2016industrial} evaluate established
NILM algorithms on industrial loads and report substantially degraded
performance, attributing it to precisely these differences. Martins
\textit{et al.} \cite{martins2018imdeld} apply a deep generative model
to industrial disaggregation and, in doing so, release the IMDELD
dataset \cite{imdeld2018dataset}---eleven meters recording eight heavy
machines in a Brazilian poultry-feed pellet plant at 1\,Hz across a
five-month window, of which the covered fraction differs by machine
(Section~\ref{sec:method:data}). IMDELD is the benchmark used throughout this
paper; its authors characterise the machines as three-state systems but
supply no per-timestamp state annotation, which is the reason the
present work must proceed without labels and validate its partition
against physics instead.

Second, and more fundamentally, the objective of NILM is
\emph{attribution}: given an aggregate signal, determine which appliance
consumed what. It is a retrospective accounting exercise. NILM does not
forecast, and it does not act. No NILM system, industrial or
residential, closes the loop from an inferred state to an actuation
decision. That step belongs to the two literatures reviewed next.


% ---------------------------------------------------------------------
\subsection{Industrial Energy Optimisation}
\label{sec:rw:industrial}

The manufacturing engineering community has independently established
that non-productive energy dominates the energy budget of discrete-part
production. Gutowski \textit{et al.} \cite{gutowski2006electrical} and
Dahmus and Gutowski \cite{dahmus2004environmental} show that the
load-independent baseline of a machine tool---spindle idling, hydraulics,
coolant pumps, control electronics---frequently exceeds the energy
actually delivered to the workpiece, and that this fraction rises as
automation increases. Kara and Li \cite{kara2011unit} formalise the
relationship as a unit-process model in which specific energy
consumption is inversely related to material removal rate, so that a
lightly loaded or idle machine is intrinsically inefficient. Devoldere
\textit{et al.} \cite{devoldere2007improvement} quantify the resulting
improvement potential and identify idle-mode elimination as the single
largest lever available without redesigning the machine. Duflou
\textit{et al.} \cite{duflou2012towards} synthesise this body of work
into a systems-level roadmap, placing operational measures---switching
equipment off when it is not needed---at the machine and cell levels of
their hierarchy.

On the management side, ISO~50001 \cite{iso50001} provides the
plan--do--check--act framework within which energy performance
indicators are defined and tracked, and ISO~14955-1 \cite{iso14955}
supplies the operating-mode taxonomy that makes standby energy a
measurable quantity in the first place.

The limitation of this literature is methodological rather than
conceptual. Its energy models are typically built from controlled
laboratory measurements of a single machine under prescribed process
parameters, and its recommendations are design guidelines and audit
procedures. It does not address how a facility, given only the
electrical measurements it already collects and no process-parameter
instrumentation, should determine \emph{which} of its machines are idle,
\emph{when}, and for \emph{how long}. That inference problem is left
open.


% ---------------------------------------------------------------------
\subsection{Predictive Shutdown and Decision Support}
\label{sec:rw:decision}

The operations research community has studied the shutdown decision
itself in considerable depth. Mouzon \textit{et al.}
\cite{mouzon2007operational} establish the central trade-off: turning a
machine off saves its idle power but incurs a restart energy penalty and
a restart delay, so shutdown is economical only when the idle period
exceeds a break-even duration determined by the ratio of restart energy
to idle power. They propose dispatching rules and a mathematical
programming formulation that exploit this, and extend the treatment to
joint energy--tardiness objectives \cite{mouzon2008framework}. Chen
\textit{et al.} \cite{chen2013schedule} develop schedule-based operation
policies for production lines, Sun and Li \cite{sun2013opportunity}
estimate real-time control opportunities from buffer state, Li and Sun
\cite{li2016energy} formulate dynamic energy control as a Markov
decision process, Frigerio and Matta \cite{frigerio2015energy} derive
control strategies for machine tools under stochastic job arrivals, and
He \textit{et al.} \cite{he2015energy} embed energy responsiveness in
job-shop machine selection and sequencing.

This literature is mathematically mature, and the break-even
formulation of Mouzon \textit{et al.} is the natural objective for the
decision layer of any standby-elimination system. Its universal
assumption, however, is that the idle interval is \emph{known}---either
because it is dictated by a deterministic production schedule, or
because it is drawn from an analytically tractable arrival process
calibrated by the modeller. The policies are evaluated in simulation or
on analytical queueing models. What is absent is any mechanism for
obtaining the idle-duration input from a machine's own measured
electrical behaviour, which is the only signal available in a retrofit
installation. In other words: this community can decide optimally given
a duration, and the NILM community can infer a state from
measurements, but no work connects the two.


% ---------------------------------------------------------------------
\subsection{Synthesis and Research Gap}
\label{sec:rw:gap}

Table~\ref{tab:related} compares representative works from the four
themes against the capabilities required for an autonomous
standby-elimination system: whether the method infers machine state
without labels, whether it models temporal structure explicitly, whether
it forecasts the \emph{duration} of a future non-productive interval,
whether it converts that forecast into an economically justified control
decision, and whether it is validated on real industrial measurements.

\begin{table*}[!t]
\centering
\caption{Comparison of representative prior work against the
capabilities required for autonomous standby elimination.
\CIRCLE~= provided, \LEFTcircle~= partially provided, \Circle~= not
provided, --~= not applicable.}
\label{tab:related}
\begin{threeparttable}
\small
\setlength{\tabcolsep}{4pt}
\begin{tabular}{@{}lllccccc@{}}
\toprule
\multirow{2}{*}{\textbf{Work}} &
\multirow{2}{*}{\textbf{Method family}} &
\multirow{2}{*}{\textbf{Data domain}} &
\textbf{Labels} &
\textbf{Temporal} &
\textbf{Duration} &
\textbf{Economic} &
\textbf{Real ind.} \\
 & & & \textbf{required} & \textbf{model} & \textbf{forecast} &
 \textbf{decision} & \textbf{data} \\
\midrule
\multicolumn{8}{@{}l}{\textit{Theme A --- machine state identification}}\\
Lloyd \cite{lloyd1982least}          & $k$-means                    & generic                & no  & \Circle    & \Circle & \Circle & \Circle \\
McLachlan \& Peel \cite{mclachlan2000finite} & Gaussian mixture     & generic                & no  & \Circle    & \Circle & \Circle & \Circle \\
Ester \textit{et al.} \cite{ester1996dbscan} & DBSCAN               & generic                & no  & \Circle    & \Circle & \Circle & \Circle \\
Fox \textit{et al.} \cite{fox2011sticky}     & sticky HDP-HMM       & audio                  & no  & \CIRCLE    & \Circle & \Circle & \Circle \\
\addlinespace
\multicolumn{8}{@{}l}{\textit{Theme B --- non-intrusive load monitoring}}\\
Hart \cite{hart1992nonintrusive}     & event/steady-state signature & residential            & yes\tnote{a} & \LEFTcircle & \Circle & \Circle & \Circle \\
Kolter \& Jaakkola \cite{kolter2012afamap} & additive factorial HMM & REDD (residential)     & no  & \CIRCLE    & \Circle & \Circle & \Circle \\
Kelly \& Knottenbelt \cite{kelly2015neuralnilm} & deep NN (dAE, LSTM) & UK-DALE (residential) & yes & \LEFTcircle & \Circle & \Circle & \Circle \\
Zhang \textit{et al.} \cite{zhang2018seq2point} & seq2point CNN     & UK-DALE, REDD          & yes & \LEFTcircle & \Circle & \Circle & \Circle \\
Holmegaard \& Kj\ae{}rgaard \cite{holmegaard2016industrial} & NILM algorithm evaluation & industrial & yes & \LEFTcircle & \Circle & \Circle & \CIRCLE \\
Martins \textit{et al.} \cite{martins2018imdeld} & deep generative disaggregation & IMDELD (industrial) & yes & \LEFTcircle & \Circle & \Circle & \CIRCLE \\
\addlinespace
\multicolumn{8}{@{}l}{\textit{Theme C --- industrial energy optimisation}}\\
Gutowski \textit{et al.} \cite{gutowski2006electrical} & thermodynamic/empirical process model & lab machine tools & -- & \Circle & \Circle & \Circle & \LEFTcircle \\
Kara \& Li \cite{kara2011unit}       & unit-process energy model    & lab machine tools      & --  & \Circle    & \Circle & \Circle & \LEFTcircle \\
Duflou \textit{et al.} \cite{duflou2012towards} & systems-level roadmap & survey             & --  & \Circle    & \Circle & \LEFTcircle & \LEFTcircle \\
\addlinespace
\multicolumn{8}{@{}l}{\textit{Theme D --- predictive shutdown and decision support}}\\
Mouzon \textit{et al.} \cite{mouzon2007operational} & dispatching rules, MILP & simulated single machine & -- & \Circle & \Circle & \CIRCLE & \Circle \\
Chen \textit{et al.} \cite{chen2013schedule} & schedule-based control  & simulated line     & --  & \LEFTcircle & \Circle & \CIRCLE & \Circle \\
Frigerio \& Matta \cite{frigerio2015energy} & stochastic control policy & analytical model  & --  & \CIRCLE    & \LEFTcircle\tnote{b} & \CIRCLE & \Circle \\
He \textit{et al.} \cite{he2015energy} & energy-responsive scheduling & job-shop model      & --  & \Circle    & \Circle & \CIRCLE & \Circle \\
\midrule
\textbf{This work}                   & \textbf{physics-gated GMM--HMM $\rightarrow$ duration forecast $\rightarrow$ optimisation} & \textbf{IMDELD (industrial)} & \textbf{no} & \CIRCLE & \CIRCLE & \CIRCLE & \CIRCLE \\
\bottomrule
\end{tabular}
\begin{tablenotes}[flushleft]\footnotesize
\item[a] Hart's formulation is unsupervised in its clustering step but
requires a manually curated appliance signature library to name the
recovered clusters.
\item[b] Idle duration is modelled as a known stochastic arrival
process calibrated by the analyst, not forecast from measured machine
behaviour.
\end{tablenotes}
\end{threeparttable}
\end{table*}

Reading Table~\ref{tab:related} column-wise makes the gap explicit.
Every capability required for autonomous standby elimination has been
solved somewhere, but the last column of any single row is always
empty:

\begin{itemize}
  \item \textbf{Label-free state inference with temporal structure}
  exists (Theme~A, and \cite{kolter2012afamap} in Theme~B), but is
  demonstrated on generic or residential data and terminates at the
  labelled sequence. Nothing is predicted and nothing is decided.

  \item \textbf{Industrial validation} exists
  (\cite{holmegaard2016industrial,martins2018imdeld}, Theme~B), but the
  task is disaggregation, the methods are supervised, and the output is
  a retrospective attribution of energy rather than an operating state
  usable for control.

  \item \textbf{Quantification of idle-energy waste} exists (Theme~C),
  but rests on controlled laboratory measurement of individual
  processes and yields design guidance, not an online inference
  procedure applicable to a running facility.

  \item \textbf{Optimal shutdown decisions} exist (Theme~D), and the
  break-even formulation is well understood, but every such policy
  takes the idle duration as a \emph{given} input---supplied by a
  deterministic schedule or an analyst-calibrated arrival
  process---and is evaluated in simulation.
\end{itemize}

\medskip
\noindent\fbox{\parbox{0.97\columnwidth}{%
\textbf{Research gap.} No existing work combines
(i)~label-free, temporally coherent inference of industrial machine
states from electrical measurements alone, (ii)~forecasting of the
\emph{duration} of the resulting non-productive intervals, and
(iii)~an economically optimised shutdown decision derived from that
forecast, validated end-to-end on real industrial measurement data.
The state-inference and decision-making capabilities have matured
independently and remain unconnected: the former produces labels no
controller consumes, the latter consumes durations no inference
procedure supplies.}}
\medskip

\noindent
Closing this gap requires more than concatenating existing components.
Three problems arise specifically at the interfaces. First, because no
ground-truth state annotation exists for industrial data, the inferred
partition cannot be validated by conventional supervised metrics, and an
alternative validation basis---independent physical and operational
evidence---must be constructed. Second, an i.i.d.\ clustering produces a
label sequence whose flicker makes duration forecasting meaningless,
so temporal coherence is a prerequisite for the forecasting stage rather
than a refinement of the inference stage. Third, a forecast duration
carries uncertainty that a deterministic break-even threshold ignores;
the decision layer must therefore be posed as a constrained
cost-minimisation problem in which the break-even point emerges as a
derived quantity rather than being supplied as a tuning parameter.
Section~\ref{sec:method} addresses each of these in turn.


% --- Framework, data, state inference, the two forecasters
% =====================================================================
%  method.tex  --  PHASE VII, Task 14C
%
%  Section 3.  The framework overview, the data, and the two forecasting
%  formulations.  The state-inference subsection lives in its own file
%  and is pulled in here.
%
%  Intended placement: Section 3 (\label{sec:method}), after related
%  work.
%
%  Drop-in usage: % =====================================================================
%  method.tex  --  PHASE VII, Task 14C
%
%  Section 3.  The framework overview, the data, and the two forecasting
%  formulations.  The state-inference subsection lives in its own file
%  and is pulled in here.
%
%  Intended placement: Section 3 (\label{sec:method}), after related
%  work.
%
%  Drop-in usage: % =====================================================================
%  method.tex  --  PHASE VII, Task 14C
%
%  Section 3.  The framework overview, the data, and the two forecasting
%  formulations.  The state-inference subsection lives in its own file
%  and is pulled in here.
%
%  Intended placement: Section 3 (\label{sec:method}), after related
%  work.
%
%  Drop-in usage: \input{method}
%  This file \inputs method_gmm_labeller.tex.
%  Bibliography: paper/references.bib
% =====================================================================

\section{Proposed Framework}
\label{sec:method}

\subsection{Overview}
\label{sec:method:overview}

The framework is a three-stage pipeline over a single input: the
per-phase voltage and current a standard retrofit meter already
records. Each stage is specified by what the next one needs.

\begin{enumerate}
\item \emph{State inference} (Section~\ref{sec:method:labelling}) maps
  the electrical record to a label per sample from
  $\{$OFF, STANDBY, WORKING, PEAK\_LOAD$\}$ without annotations, and is
  required to produce a sequence whose state durations are physically
  admissible --- not merely a per-sample partition --- because duration
  forecasting over a flickering label sequence is ill-posed.
\item \emph{Duration forecasting} (Sections~\ref{sec:forecasting}
  and~\ref{sec:decision:policies}) predicts how much longer the machine
  will remain non-productive, and is required to target the conditional
  \emph{mean} of that duration, because the decision rule below is
  affine in it.
\item \emph{Decision} (Section~\ref{sec:decision}) converts the forecast
  into a de-energise/keep-running action by minimising total cost
  subject to thermal and wear constraints, and returns the break-even
  duration as a derived quantity rather than consuming it as a
  parameter.
\end{enumerate}

The stages are separable in implementation and are evaluated separately
throughout, which is what makes the component ablation of
Section~\ref{sec:ablation} meaningful.

\subsection{Data}
\label{sec:method:data}

IMDELD~\cite{imdeld2018dataset,martins2018imdeld} publishes $1$\,Hz
three-phase measurements from eight machines of a Brazilian
poultry-feed pellet plant: two pelletizers, two milling machines, two
exhaust fans and two double-pole contactors. Six channels are used ---
active, reactive and apparent power, current, voltage and power factor
--- and no other information about the machines enters the models. The
plant's published calendar (weekday operation, a halt between 17:00 and
22:00 local time to avoid the peak tariff, and no weekend production)
is used \emph{only} as external validation evidence, never as a model
input to the state layer.

Coverage differs sharply by machine and is reported as measured rather
than as published. The pelletizer~I record spans $155$ days from
2017-10-30 to 2018-04-03 and contains $63.4$ days of data in $74$
contiguous segments; $31$ acquisition gaps account for the remainder,
the largest of them $30$ days. The two milling machines cover $12.2$ of
$43$ days in $1\,082$ and $580$ fragments. \textbf{Every annualised
quantity in this paper is scaled by covered time, not by span}, which
matters by a factor of $2.4$ on the five-month records and is the
reason the milling machines' figures are not annualised at all
(Section~\ref{sec:limitations:milling}).

\input{method_gmm_labeller}

\subsection{Two forecasting problems, not one}
\label{sec:method:forecasting}

The framework contains two forecasters. They are frequently conflated,
they consume different features, and no result in this paper reports
them in the same column.

\emph{The window state forecaster} takes a $600$\,s lookback of the
electrical record and predicts the state at each step of a $300$\,s
horizon. It is a sequence problem, it is what the baseline comparison of
Section~\ref{sec:forecasting} evaluates, and it is scored on windows.
Its design representation summarises each channel over the lookback ---
final observed value, sub-block means and standard deviations --- giving
$217$ columns for the tree, or the raw sequence for the neural
architectures.

\emph{The episode-duration forecaster} takes the state of an idle
episode at a decision epoch and predicts the \emph{remaining} idle time
$R$ in seconds. It is a regression problem, it is what the decision
layer of Section~\ref{sec:decision} consumes, and it is scored per
decision epoch. Its ten features are elapsed idle time, the current
position in the factory calendar, and the load history of the
production run that just ended.

Keeping them apart is not pedantry. Section~\ref{sec:explain} shows they
draw on almost disjoint information --- $11.6\%$ of the window model's
attribution is calendar against $82.4\%$ of the duration model's --- and
Section~\ref{sec:crossmachine} shows that one transfers between machines
and the other does not, for exactly that reason. Attributing the
economic result to the window model's $F_1$ would credit a model that
plays no part in producing it.

%  This file \inputs method_gmm_labeller.tex.
%  Bibliography: paper/references.bib
% =====================================================================

\section{Proposed Framework}
\label{sec:method}

\subsection{Overview}
\label{sec:method:overview}

The framework is a three-stage pipeline over a single input: the
per-phase voltage and current a standard retrofit meter already
records. Each stage is specified by what the next one needs.

\begin{enumerate}
\item \emph{State inference} (Section~\ref{sec:method:labelling}) maps
  the electrical record to a label per sample from
  $\{$OFF, STANDBY, WORKING, PEAK\_LOAD$\}$ without annotations, and is
  required to produce a sequence whose state durations are physically
  admissible --- not merely a per-sample partition --- because duration
  forecasting over a flickering label sequence is ill-posed.
\item \emph{Duration forecasting} (Sections~\ref{sec:forecasting}
  and~\ref{sec:decision:policies}) predicts how much longer the machine
  will remain non-productive, and is required to target the conditional
  \emph{mean} of that duration, because the decision rule below is
  affine in it.
\item \emph{Decision} (Section~\ref{sec:decision}) converts the forecast
  into a de-energise/keep-running action by minimising total cost
  subject to thermal and wear constraints, and returns the break-even
  duration as a derived quantity rather than consuming it as a
  parameter.
\end{enumerate}

The stages are separable in implementation and are evaluated separately
throughout, which is what makes the component ablation of
Section~\ref{sec:ablation} meaningful.

\subsection{Data}
\label{sec:method:data}

IMDELD~\cite{imdeld2018dataset,martins2018imdeld} publishes $1$\,Hz
three-phase measurements from eight machines of a Brazilian
poultry-feed pellet plant: two pelletizers, two milling machines, two
exhaust fans and two double-pole contactors. Six channels are used ---
active, reactive and apparent power, current, voltage and power factor
--- and no other information about the machines enters the models. The
plant's published calendar (weekday operation, a halt between 17:00 and
22:00 local time to avoid the peak tariff, and no weekend production)
is used \emph{only} as external validation evidence, never as a model
input to the state layer.

Coverage differs sharply by machine and is reported as measured rather
than as published. The pelletizer~I record spans $155$ days from
2017-10-30 to 2018-04-03 and contains $63.4$ days of data in $74$
contiguous segments; $31$ acquisition gaps account for the remainder,
the largest of them $30$ days. The two milling machines cover $12.2$ of
$43$ days in $1\,082$ and $580$ fragments. \textbf{Every annualised
quantity in this paper is scaled by covered time, not by span}, which
matters by a factor of $2.4$ on the five-month records and is the
reason the milling machines' figures are not annualised at all
(Section~\ref{sec:limitations:milling}).

% =====================================================================
%  method_gmm_labeller.tex  --  PHASE V CORRECTION 4 & 5
%
%  Drop-in description of the GMM-HMM labelling procedure, including
%  the accept/reject step that Phase V Task 12B showed is load-bearing
%  for reproducibility.  Insert into Section 3 (Method).
%
%  Numbers: all STANDBY-time totals are rounded to 2 significant
%  figures and stated as conditional on the physics battery passing,
%  per Phase V Task 11C (block-bootstrap CI ±15–25%) and
%  Task 12B (labelling CV 33.7% unconditional, 3.0% conditional).
% =====================================================================

\subsection{Machine-state inference}
\label{sec:method:labelling}

\subsubsection{Preprocessing}
\label{sec:method:preprocess}

Each machine's $1$\,Hz electrical record is preprocessed in three
deterministic steps, applied identically to every machine so that
differences in downstream results are attributable to the machines,
not to the preprocessing.

\begin{enumerate}
\item \emph{Spike masking.}  For each channel, rolling-MAD outliers
  (window $11$, threshold $5\times$MAD) are flagged and excluded
  from model fitting but retained in the record; removing them would
  destroy the row-count-to-seconds correspondence that episode
  durations depend on.

\item \emph{OFF-state denoising.}  Rows with active power below
  $5$\,W are smoothed with a rolling median (window $5$) to suppress
  sensor noise in the de-energised state without touching on-state
  dynamics.

\item \emph{Feature standardisation.}  Six electrical channels are
  log$_1$p-transformed (right-skewed channels: $P, Q, S, I$) and
  z-scored jointly before any model fitting.
\end{enumerate}

\subsubsection{Gaussian hidden Markov model}
\label{sec:method:ghmm}

Emission distributions are estimated by a Gaussian mixture model
(GMM, full covariance, $k = 4$ components) fitted on a uniform
random sample of $200{,}000$ rows.  A Gaussian HMM is then
warm-started from these emissions with $\texttt{params}=\texttt{``t''}$,
so only the transition matrix is learned from $64$ contiguous
two-hour blocks spread across the record; this decouples emission
shape (a property of the marginal distribution) from transition
dynamics (a property of the sequence) and avoids the local optima
that Baum--Welch finds when initialised randomly on this
dataset~\cite{fox2011sticky}.  Viterbi decoding is applied per
contiguous segment to prevent the decoder from asserting a transition
across an acquisition gap.  A minimum-dwell constraint of $10$\,s
is applied post-decode to remove label flicker without reshaping
the partition (Section~\ref{sec:results:task4}).

\subsubsection{Physics-validation accept/reject loop}
\label{sec:method:accept_reject}

\emph{This step is load-bearing for reproducibility and was not
explicit in the Phase I--IV description; it is stated here as a
result of Phase~V Task~12B.}

After decoding, a battery of ten physics tests (T1--T10:
power-factor separation, no-load current ratio, power conservation,
dwell plausibility, cross-machine replication, and factory-schedule
agreement) is run on a strided sample of the labelled record and
aggregated into a physics score $s \in [0, 1]$.  If $s < 0.90$,
the model is re-fitted with a different seed (incremented by a
fixed prime) and the battery is re-evaluated, up to three attempts.

The motivation is empirical.  Phase~V Task~12B fitted the labeller at
ten different seeds on identical preprocessed data from pelletizer~I
and found that three seeds (30\%) converge to a wrong cluster:

\begin{itemize}
\item seed~6: the ``STANDBY'' cluster absorbed a loaded running state
  (mean power $12{,}961$\,W instead of the correct $\sim3{,}750$\,W);
\item seeds~8 and~9: the ``STANDBY'' cluster absorbed the OFF state
  (mean power $0.8$--$1.1$\,W).
\end{itemize}

The physics battery scored those three seeds at $0.60$--$0.80$ and the
seven correct seeds at $1.00$, making the accept/reject decision exact
on this record.  Conditioning on a passing seed collapses the
STANDBY-hour coefficient of variation from $33.7\%$ (all ten seeds)
to $3.0\%$ (passing seeds), and the total reported throughout this
paper---$92 \pm 3$\,h over the 155-day record---is the mean of the
seven passing seeds.

\emph{All STANDBY-time totals in this paper are therefore stated as
conditional on the physics battery passing} (i.e., $s \ge 0.90$)
\emph{and are rounded to two significant figures}, because the
block-bootstrap sampling uncertainty (Phase~V Task~11C, one-week
blocks) is $\pm 15$--$25\%$ and the conditional labelling
uncertainty is $\pm 3\%$.  Three significant figures are not
supported by the data.

The accept/reject loop is implemented in \texttt{src/labelling.py}
(\texttt{label\_record}, \texttt{physics\_score\_threshold}=0.90,
\texttt{max\_refit\_attempts}=3) and its output is stored in the
model-info dictionary (\texttt{physics\_validation\_attempts},
\texttt{physics\_validation\_seed}, \texttt{physics\_score}) for
audit and reproduction.


\subsection{Two forecasting problems, not one}
\label{sec:method:forecasting}

The framework contains two forecasters. They are frequently conflated,
they consume different features, and no result in this paper reports
them in the same column.

\emph{The window state forecaster} takes a $600$\,s lookback of the
electrical record and predicts the state at each step of a $300$\,s
horizon. It is a sequence problem, it is what the baseline comparison of
Section~\ref{sec:forecasting} evaluates, and it is scored on windows.
Its design representation summarises each channel over the lookback ---
final observed value, sub-block means and standard deviations --- giving
$217$ columns for the tree, or the raw sequence for the neural
architectures.

\emph{The episode-duration forecaster} takes the state of an idle
episode at a decision epoch and predicts the \emph{remaining} idle time
$R$ in seconds. It is a regression problem, it is what the decision
layer of Section~\ref{sec:decision} consumes, and it is scored per
decision epoch. Its ten features are elapsed idle time, the current
position in the factory calendar, and the load history of the
production run that just ended.

Keeping them apart is not pedantry. Section~\ref{sec:explain} shows they
draw on almost disjoint information --- $11.6\%$ of the window model's
attribution is calendar against $82.4\%$ of the duration model's --- and
Section~\ref{sec:crossmachine} shows that one transfers between machines
and the other does not, for exactly that reason. Attributing the
economic result to the window model's $F_1$ would credit a model that
plays no part in producing it.

%  This file \inputs method_gmm_labeller.tex.
%  Bibliography: paper/references.bib
% =====================================================================

\section{Proposed Framework}
\label{sec:method}

\subsection{Overview}
\label{sec:method:overview}

The framework is a three-stage pipeline over a single input: the
per-phase voltage and current a standard retrofit meter already
records. Each stage is specified by what the next one needs.

\begin{enumerate}
\item \emph{State inference} (Section~\ref{sec:method:labelling}) maps
  the electrical record to a label per sample from
  $\{$OFF, STANDBY, WORKING, PEAK\_LOAD$\}$ without annotations, and is
  required to produce a sequence whose state durations are physically
  admissible --- not merely a per-sample partition --- because duration
  forecasting over a flickering label sequence is ill-posed.
\item \emph{Duration forecasting} (Sections~\ref{sec:forecasting}
  and~\ref{sec:decision:policies}) predicts how much longer the machine
  will remain non-productive, and is required to target the conditional
  \emph{mean} of that duration, because the decision rule below is
  affine in it.
\item \emph{Decision} (Section~\ref{sec:decision}) converts the forecast
  into a de-energise/keep-running action by minimising total cost
  subject to thermal and wear constraints, and returns the break-even
  duration as a derived quantity rather than consuming it as a
  parameter.
\end{enumerate}

The stages are separable in implementation and are evaluated separately
throughout, which is what makes the component ablation of
Section~\ref{sec:ablation} meaningful.

\subsection{Data}
\label{sec:method:data}

IMDELD~\cite{imdeld2018dataset,martins2018imdeld} publishes $1$\,Hz
three-phase measurements from eight machines of a Brazilian
poultry-feed pellet plant: two pelletizers, two milling machines, two
exhaust fans and two double-pole contactors. Six channels are used ---
active, reactive and apparent power, current, voltage and power factor
--- and no other information about the machines enters the models. The
plant's published calendar (weekday operation, a halt between 17:00 and
22:00 local time to avoid the peak tariff, and no weekend production)
is used \emph{only} as external validation evidence, never as a model
input to the state layer.

Coverage differs sharply by machine and is reported as measured rather
than as published. The pelletizer~I record spans $155$ days from
2017-10-30 to 2018-04-03 and contains $63.4$ days of data in $74$
contiguous segments; $31$ acquisition gaps account for the remainder,
the largest of them $30$ days. The two milling machines cover $12.2$ of
$43$ days in $1\,082$ and $580$ fragments. \textbf{Every annualised
quantity in this paper is scaled by covered time, not by span}, which
matters by a factor of $2.4$ on the five-month records and is the
reason the milling machines' figures are not annualised at all
(Section~\ref{sec:limitations:milling}).

% =====================================================================
%  method_gmm_labeller.tex  --  PHASE V CORRECTION 4 & 5
%
%  Drop-in description of the GMM-HMM labelling procedure, including
%  the accept/reject step that Phase V Task 12B showed is load-bearing
%  for reproducibility.  Insert into Section 3 (Method).
%
%  Numbers: all STANDBY-time totals are rounded to 2 significant
%  figures and stated as conditional on the physics battery passing,
%  per Phase V Task 11C (block-bootstrap CI ±15–25%) and
%  Task 12B (labelling CV 33.7% unconditional, 3.0% conditional).
% =====================================================================

\subsection{Machine-state inference}
\label{sec:method:labelling}

\subsubsection{Preprocessing}
\label{sec:method:preprocess}

Each machine's $1$\,Hz electrical record is preprocessed in three
deterministic steps, applied identically to every machine so that
differences in downstream results are attributable to the machines,
not to the preprocessing.

\begin{enumerate}
\item \emph{Spike masking.}  For each channel, rolling-MAD outliers
  (window $11$, threshold $5\times$MAD) are flagged and excluded
  from model fitting but retained in the record; removing them would
  destroy the row-count-to-seconds correspondence that episode
  durations depend on.

\item \emph{OFF-state denoising.}  Rows with active power below
  $5$\,W are smoothed with a rolling median (window $5$) to suppress
  sensor noise in the de-energised state without touching on-state
  dynamics.

\item \emph{Feature standardisation.}  Six electrical channels are
  log$_1$p-transformed (right-skewed channels: $P, Q, S, I$) and
  z-scored jointly before any model fitting.
\end{enumerate}

\subsubsection{Gaussian hidden Markov model}
\label{sec:method:ghmm}

Emission distributions are estimated by a Gaussian mixture model
(GMM, full covariance, $k = 4$ components) fitted on a uniform
random sample of $200{,}000$ rows.  A Gaussian HMM is then
warm-started from these emissions with $\texttt{params}=\texttt{``t''}$,
so only the transition matrix is learned from $64$ contiguous
two-hour blocks spread across the record; this decouples emission
shape (a property of the marginal distribution) from transition
dynamics (a property of the sequence) and avoids the local optima
that Baum--Welch finds when initialised randomly on this
dataset~\cite{fox2011sticky}.  Viterbi decoding is applied per
contiguous segment to prevent the decoder from asserting a transition
across an acquisition gap.  A minimum-dwell constraint of $10$\,s
is applied post-decode to remove label flicker without reshaping
the partition (Section~\ref{sec:results:task4}).

\subsubsection{Physics-validation accept/reject loop}
\label{sec:method:accept_reject}

\emph{This step is load-bearing for reproducibility and was not
explicit in the Phase I--IV description; it is stated here as a
result of Phase~V Task~12B.}

After decoding, a battery of ten physics tests (T1--T10:
power-factor separation, no-load current ratio, power conservation,
dwell plausibility, cross-machine replication, and factory-schedule
agreement) is run on a strided sample of the labelled record and
aggregated into a physics score $s \in [0, 1]$.  If $s < 0.90$,
the model is re-fitted with a different seed (incremented by a
fixed prime) and the battery is re-evaluated, up to three attempts.

The motivation is empirical.  Phase~V Task~12B fitted the labeller at
ten different seeds on identical preprocessed data from pelletizer~I
and found that three seeds (30\%) converge to a wrong cluster:

\begin{itemize}
\item seed~6: the ``STANDBY'' cluster absorbed a loaded running state
  (mean power $12{,}961$\,W instead of the correct $\sim3{,}750$\,W);
\item seeds~8 and~9: the ``STANDBY'' cluster absorbed the OFF state
  (mean power $0.8$--$1.1$\,W).
\end{itemize}

The physics battery scored those three seeds at $0.60$--$0.80$ and the
seven correct seeds at $1.00$, making the accept/reject decision exact
on this record.  Conditioning on a passing seed collapses the
STANDBY-hour coefficient of variation from $33.7\%$ (all ten seeds)
to $3.0\%$ (passing seeds), and the total reported throughout this
paper---$92 \pm 3$\,h over the 155-day record---is the mean of the
seven passing seeds.

\emph{All STANDBY-time totals in this paper are therefore stated as
conditional on the physics battery passing} (i.e., $s \ge 0.90$)
\emph{and are rounded to two significant figures}, because the
block-bootstrap sampling uncertainty (Phase~V Task~11C, one-week
blocks) is $\pm 15$--$25\%$ and the conditional labelling
uncertainty is $\pm 3\%$.  Three significant figures are not
supported by the data.

The accept/reject loop is implemented in \texttt{src/labelling.py}
(\texttt{label\_record}, \texttt{physics\_score\_threshold}=0.90,
\texttt{max\_refit\_attempts}=3) and its output is stored in the
model-info dictionary (\texttt{physics\_validation\_attempts},
\texttt{physics\_validation\_seed}, \texttt{physics\_score}) for
audit and reproduction.


\subsection{Two forecasting problems, not one}
\label{sec:method:forecasting}

The framework contains two forecasters. They are frequently conflated,
they consume different features, and no result in this paper reports
them in the same column.

\emph{The window state forecaster} takes a $600$\,s lookback of the
electrical record and predicts the state at each step of a $300$\,s
horizon. It is a sequence problem, it is what the baseline comparison of
Section~\ref{sec:forecasting} evaluates, and it is scored on windows.
Its design representation summarises each channel over the lookback ---
final observed value, sub-block means and standard deviations --- giving
$217$ columns for the tree, or the raw sequence for the neural
architectures.

\emph{The episode-duration forecaster} takes the state of an idle
episode at a decision epoch and predicts the \emph{remaining} idle time
$R$ in seconds. It is a regression problem, it is what the decision
layer of Section~\ref{sec:decision} consumes, and it is scored per
decision epoch. Its ten features are elapsed idle time, the current
position in the factory calendar, and the load history of the
production run that just ended.

Keeping them apart is not pedantry. Section~\ref{sec:explain} shows they
draw on almost disjoint information --- $11.6\%$ of the window model's
attribution is calendar against $82.4\%$ of the duration model's --- and
Section~\ref{sec:crossmachine} shows that one transfers between machines
and the other does not, for exactly that reason. Attributing the
economic result to the window model's $F_1$ would credit a model that
plays no part in producing it.


% --- State-layer results, flicker, validation
% =====================================================================
%  state_layer.tex  --  PHASE II, Tasks 3 and 4 (+ 4b, 4c)
%
%  Results for the state-inference stage: what separates STANDBY, how
%  the decoding rule compares against clustering baselines, why a
%  transition prior cannot remove flicker, and how a partition is
%  validated with no ground truth.
%
%  Intended placement: Section 4 (\label{sec:states}), immediately after
%  the method.  Also carries \label{sec:results} (it is the first
%  results section) and \label{sec:validation}.
%
%  Drop-in usage: % =====================================================================
%  state_layer.tex  --  PHASE II, Tasks 3 and 4 (+ 4b, 4c)
%
%  Results for the state-inference stage: what separates STANDBY, how
%  the decoding rule compares against clustering baselines, why a
%  transition prior cannot remove flicker, and how a partition is
%  validated with no ground truth.
%
%  Intended placement: Section 4 (\label{sec:states}), immediately after
%  the method.  Also carries \label{sec:results} (it is the first
%  results section) and \label{sec:validation}.
%
%  Drop-in usage: % =====================================================================
%  state_layer.tex  --  PHASE II, Tasks 3 and 4 (+ 4b, 4c)
%
%  Results for the state-inference stage: what separates STANDBY, how
%  the decoding rule compares against clustering baselines, why a
%  transition prior cannot remove flicker, and how a partition is
%  validated with no ground truth.
%
%  Intended placement: Section 4 (\label{sec:states}), immediately after
%  the method.  Also carries \label{sec:results} (it is the first
%  results section) and \label{sec:validation}.
%
%  Drop-in usage: \input{state_layer}
%  Figures are referenced by bare name; add to the preamble
%      \graphicspath{{../outputs/phase2/figures/pelletizer-I/}}
%  Bibliography: paper/references.bib
%  Numbers below are generated by
%      python run_phase2.py
%  and stored in outputs/phase2/task3/<machine>/task3_results.json,
%  outputs/phase2/task4/<machine>/{comparison_table.csv,
%  flicker_diagnosis.json,sticky_prior_sweep.csv,min_dwell_decode.csv}.
% =====================================================================

\section{The State Layer}
\label{sec:states}
\label{sec:results}

This section reports what the state-inference stage of
Section~\ref{sec:method:labelling} produces on pelletizer~I: which
measured quantities separate the states, how the choice of clustering
algorithm compares with the choice of decoding rule, and how a partition
is validated when there is nothing to validate it against.

\subsection{What separates STANDBY}
\label{sec:states:features}

Twenty-four features are derived from the six measured channels in four
groups: raw electrical (5), derived electrical (4 --- power factor,
$Q/P$ ratio, load factor, apparent-power residual), rolling statistics
(8, over a $300$\,s window), and temporal (7). Rolling windows are
\emph{trailing}, not centred, and are computed within contiguous
segments so that none spans an acquisition gap.

\textbf{The target is STANDBY against $\{$WORKING, PEAK\_LOAD$\}$ on
energised rows only.} Separating OFF from everything else is trivial and
would dominate any importance ranking computed over all four classes;
separating an energised idle machine from a lightly loaded running one
is the problem the framework actually has to solve. A four-class
ranking was also computed and is not reported here, because it measures
the OFF boundary and invites the two to be averaged.

\begin{table}[t]
\centering
\caption{Feature-group importance for the STANDBY-versus-productive
split, by three measures with different failure modes. Each column is
normalised to its own total.}
\label{tab:feature-groups}
\small
\begin{tabular}{@{}lrrr@{}}
\toprule
Group & mutual info. & split gain & SHAP \\
\midrule
raw electrical      & $0.350$ & $\mathbf{0.560}$ & $\mathbf{0.537}$ \\
derived electrical  & $0.250$ & $0.438$ & $0.310$ \\
rolling statistics  & $\mathbf{0.388}$ & $0.001$ & $0.116$ \\
temporal            & $0.012$ & $0.002$ & $0.036$ \\
\bottomrule
\end{tabular}
\end{table}

Three readings of Table~\ref{tab:feature-groups} matter for the rest of
the paper.

\textbf{The physics argument is empirically confirmed.}
\texttt{power\_factor} and \texttt{q\_p\_ratio} rank second and third by
consensus over the three measures, behind \texttt{active\_power} and
ahead of \texttt{current}, and the four derived electrical features
carry $31$--$44\%$ of total importance. Methods~1 and~2 of the
validation battery (Section~\ref{sec:validation}) were previously
justified by induction-motor theory alone; they now have a measured
basis on the machine's own data. Section~\ref{sec:explain:state}
strengthens this: under Shapley attribution power factor overtakes
active power outright.

\textbf{Temporal features are nearly irrelevant to this split}
($1$--$4\%$). That is the desirable outcome, not a null one: a state
model leaning on \texttt{hour\_of\_day} would be reproducing the shift
schedule rather than the machine's electrical condition, and would not
transfer to a facility with a different calendar. It is also the exact
opposite of what the \emph{decision} layer's forecaster does, which
draws $82\%$ of its attribution from the calendar
(Section~\ref{sec:crossmachine:decision}) --- the two stages of the
framework depend on almost disjoint information, and only the state
stage is physics-driven.

\textbf{High mutual information with near-zero split gain is the
redundancy signature.} Rolling statistics are individually informative
(MI $0.388$, the largest in the column) and add nothing once the tree
already holds the instantaneous value (gain $0.001$). Reporting either
measure alone would misdescribe them.

\paragraph{A caveat that must not be dropped.}
The surrogate used for the gain and SHAP columns reproduces its targets
to $99.9\%$ accuracy. \emph{That number is not a validation result.} The
targets are the state model's own labels and the state model was fitted
on a subset of these same features, so the surrogate is largely
re-deriving a deterministic function of its own inputs. The rankings are
meaningful; the accuracy is evidence of nothing. Validation is
Section~\ref{sec:validation}'s job.

\subsection{The emission model barely matters; the decoding rule does}
\label{sec:results:task4}

Six state-inference configurations are compared on one representation
and one protocol: $64$ contiguous two-hour blocks spread across the
record and never crossing a gap, with alternate blocks forming a fitting
set and an evaluation set of $230\,400$ rows each. DBSCAN and spectral
clustering are transductive, so their labels are extended to the
evaluation blocks by one-nearest-neighbour, applied identically to both.
Because no ground truth exists, the comparison is scored on flicker,
median dwell, the physics battery, the balanced schedule score, and the
STANDBY hour total --- never on accuracy.

\begin{table}[t]
\centering
\caption{State-inference comparison on the evaluation blocks. Flicker is
the fraction of state runs shorter than the $10$\,s physical
plausibility floor; the binary column recomputes it on the
STANDBY-versus-rest sequence the decision layer consumes. No column is
interpretable on its own --- see the spectral row.}
\label{tab:state-detection}
\small
\setlength{\tabcolsep}{4pt}
\begin{tabular}{@{}lrrrrrr@{}}
\toprule
Method & flick. & flick. & dwell & phys. & sched. & STANDBY \\
       & \%     & bin.\% & (s)   &       & \%      & (h) \\
\midrule
$k$-means \cite{lloyd1982least}      & $78.2$ & $49.1$ & $2$ & $5/5$ & $84.4$ & $4.38$ \\
DBSCAN \cite{ester1996dbscan}        & $62.5$ & $50.3$ & $4$ & $5/5$ & $84.4$ & $4.34$ \\
Spectral \cite{ng2001spectral}       & $98.2$ & $98.8$ & $2$ & $\mathbf{4/5}$ & $\mathbf{89.0}$ & $\mathbf{15.01}$ \\
GMM, i.i.d.\ \cite{mclachlan2000finite} & $66.9$ & $53.6$ & $3$ & $5/5$ & $84.5$ & $4.17$ \\
GMM--HMM \cite{rabiner1989tutorial}  & $44.5$ & $44.7$ & $13$ & $5/5$ & $84.5$ & $4.16$ \\
\textbf{GMM--HMM + min-dwell}        & $\mathbf{0.00}$ & $\mathbf{0.00}$ & $\mathbf{92.5}$ & $5/5$ & $84.5$ & $4.20$ \\
\bottomrule
\end{tabular}
\end{table}

\begin{figure*}[t]
\centering
\includegraphics[width=\textwidth]{fig_p7_state_detection}
\caption{The same six configurations as
Table~\ref{tab:state-detection}, plotted. \textbf{(A)} Flicker against
median dwell. The three filled points share one set of GMM emissions
and differ only in how those emissions are decoded; the arrows follow
i.i.d.\ assignment $\rightarrow$ Viterbi $\rightarrow$ Viterbi under a
minimum-dwell floor. The three squares are alternative emission models
decoded i.i.d.\ Changing the emission model moves the result within one
cluster; changing the decoding rule moves it two orders of magnitude in
dwell and to zero in flicker. \textbf{(B)} Inferred STANDBY hours.
Spectral clustering assigns $3.6\times$ every other method's total while
scoring the highest factory-schedule agreement of any row in the table,
which is why the battery is argued for in preference to any single
ground-truth-free metric. Hatching marks a failed physics check in both
panels.}
\label{fig:state-detection}
\end{figure*}

\textbf{Choosing a better clustering algorithm buys almost nothing;
choosing a better sequence model buys everything} ---
Figure~\ref{fig:state-detection}(A) is the table read as a picture, and
the two families separate along different axes. $k$-means, DBSCAN and
the GMM all pass $5/5$ physics checks, agree with the factory schedule
to within $0.2$ points, and find $4.17$--$4.38$ STANDBY hours on the
evaluation blocks. Switching from i.i.d.\ assignment to Viterbi decoding
over the \emph{same} emissions cuts flicker from $66.9\%$ to $44.5\%$
and quadruples the median dwell; adding the minimum-dwell constraint
takes flicker to zero and the median dwell to $92.5$\,s. The
emission-model rows span $62$--$78\%$ of flicker; the decoding rows span
$0$--$67\%$.

The comparison is clean in the one place it needs to be: the
\textsc{gmm} and \textsc{gmm--hmm} rows differ \emph{only} in the
decoding rule, because the HMM's emissions are warm-started from the
fitted GMM and held fixed (Section~\ref{sec:method:ghmm}). Fitting HMM
emissions from scratch by Baum--Welch converges to a materially worse
optimum on this record --- its ``STANDBY'' absorbs part of the
productive range at $10.3$\,kW instead of finding the no-load band at
$3.67$\,kW, and it fails the variance-ordering check --- which would
have confounded the value of temporal decoding with EM initialisation
luck.

\textbf{Spectral clustering is the one genuine failure, and it is
instructive} (Figure~\ref{fig:state-detection}(B)). It assigns $15.01$
STANDBY hours, $3.6\times$ every other
method, with $98.8\%$ binary flicker and a no-load current ratio of
$0.001$ against the acceptance band $[0.25,0.50]$. It also scores the
\emph{highest} schedule agreement in the table ($89.0\%$): a partition
that over-assigns non-productive states automatically looks good on the
factory-closed window. That is a weakness of the schedule metric in
isolation and the reason the STANDBY hour total is printed beside it.
\textbf{No single ground-truth-free metric can be quoted alone}, which
is the argument for the battery rather than for a favourite score.

\paragraph{A protocol correction worth recording.}
DBSCAN chooses its own cluster count and returned $28$. Naming the four
lowest-power clusters with the four canonical state names put all four
``states'' inside the machine's low-power range and produced an apparent
total failure --- $0.06$ STANDBY hours, $98.9\%$ flicker --- that was
entirely an artefact of the mapping. Cluster means are now grouped into
four size-weighted power bands, so sub-clusters of one physical state
are re-merged, and DBSCAN passes $5/5$ with $4.34$ STANDBY hours. Any
comparison against a method that over-segments needs this step, and
without it this paper would have reported a baseline failure that does
not exist.

\subsection{Flicker: why a transition prior cannot remove it}
\label{sec:flicker}

The expectation going in was that Viterbi decoding under a sticky prior
would take flicker to a few per cent. It does not, and the reason is
arithmetic rather than empirical.

\textbf{The sticky prior as conventionally set is numerically inert.}
Setting the Dirichlet pseudo-count to the target dwell in samples gives
$\kappa=30$ against roughly $576\,000$ observed transitions. A flat
prior and $\kappa=30$ give flicker $64.10\%$ and $64.09\%$ and minimum
self-transition $0.9677$ and $0.9678$: the prior is present in the
model, cited in its documentation, and has no measurable effect.

\textbf{Strengthening it does not help either.} Scaling the prior to ten
times the total evidence ($\kappa=575\,920$, minimum self-transition
$0.9994$, implied mean dwell $1\,605$\,s) still leaves flicker at
$50.2\%$.

\textbf{The bound is structural.} In a Viterbi decode the cost of
changing state is $\log(p_{\text{self}}/p_{\text{switch}})$, measured
here at $9.23$ nats at its largest and bounded above however strong the
prior is made. The emission log-likelihood gap between the best and
second-best state has median $27.8$ nats and $90$th percentile
$41\,701$ nats. \textbf{At $55.7\%$ of timesteps the emission term
exceeds the largest penalty the transition matrix can apply}, so the
transition model is arithmetically incapable of changing the decision
there and the HMM degenerates toward per-sample assignment. This is a
general property of applying an HMM to low-noise, multi-channel
measurements, not a property of this record.

\textbf{What fixes it is an explicit duration constraint} --- the
discrete analogue of a hidden semi-Markov duration
model~\cite{yu2010hsmm}, imposed at decode time. The floor is set equal
to the flicker threshold itself, so it removes exactly what the metric
defines as implausible and nothing more.
Table~\ref{tab:min-dwell} sweeps the floor over an order of magnitude.

\begin{table}[t]
\centering
\caption{Minimum-dwell sweep. Flicker goes to zero at no cost to the
physics checks, the schedule agreement, or the STANDBY total. This
sweep is run on the flicker-diagnosis decode rather than on the
evaluation blocks of Table~\ref{tab:state-detection}, which is why its
absolute physics score and STANDBY total differ from that table's; the
comparison of interest is down the column, not across the two tables.}
\label{tab:min-dwell}
\small
\begin{tabular}{@{}rrrrr@{}}
\toprule
min.\ dwell & flicker \% & median dwell (s) & physics & STANDBY (h) \\
\midrule
$0$\,s  & $64.09$ & $4$   & $4/5$ & $5.34$ \\
$10$\,s & $\mathbf{0.00}$ & $153$ & $4/5$ & $5.34$ \\
$30$\,s & $0.00$ & $327.5$ & $4/5$ & $5.37$ \\
$60$\,s & $0.00$ & $436.5$ & $4/5$ & $5.40$ \\
\bottomrule
\end{tabular}
\end{table}

\begin{figure}[t]
\centering
\includegraphics[width=\columnwidth]{fig_p7_min_dwell}
\caption{The minimum-dwell sweep of Table~\ref{tab:min-dwell}.
\textbf{(A)} Flicker under both definitions collapses to $0.00\%$ at the
first non-zero floor and stays there. \textbf{(B)} The median dwell
rises by two orders of magnitude while the inferred STANDBY total is
flat. The right-hand axis is anchored at zero deliberately: on an
auto-scaled axis the $1\%$ drift across the whole sweep would render as
a pronounced trend and say the opposite of what the data says.}
\label{fig:min-dwell}
\end{figure}

The constraint is removing noise, not reshaping the partition
(Figure~\ref{fig:min-dwell}): the
STANDBY total moves by $0.00$\,h from $0$\,s to $10$\,s and by $1\%$ out
to $60$\,s, and no physics check changes verdict. Its cost is measured
directly --- $0.19\%$ of rows relabelled on this machine, $0.15$--$6.1\%$
across the plant (Section~\ref{sec:multimachine:general}) --- and it is
the component the ablation of Section~\ref{sec:ablation:state} credits
with removing flicker, in preference to the HMM.

\subsection{Validation without ground truth}
\label{sec:validation}

IMDELD supplies no per-timestamp state annotation, and obtaining one
would require instrumenting each machine's contactor. The partition
therefore cannot be scored by any supervised metric, and the alternative
adopted here is to require agreement with evidence the model does not
have access to.

Three independent bodies of evidence are used, and they fail in
different ways, which is the property that makes their agreement
informative.

\textbf{(i) Circuit physics.} The five checks C1--C5 of
Section~\ref{sec:method:accept_reject}, of which C3 (power-factor
separation) and C4 (no-load current ratio) are the two that do not
depend on power magnitude and therefore the two that can distinguish an
energised idle machine from a lightly loaded one. Pelletizer~I passes
$5/5$, with a power-factor separation of $0.72$ and a no-load current
ratio of $0.36$ --- inside the $[0.25,0.50]$ band that induction-motor
theory predicts and that no part of the fitting procedure was told
about.

\textbf{(ii) Model-level diagnostics.} A ten-test harness (T1--T10)
covering data sanity, EM convergence, model order, component separation
at $2\sigma$, assignment confidence, ghost components, physical
thresholds, variance ordering, temporal alignment with nights and
weekends, and agreement with the published plant calendar. T3 is
reported as a \emph{diagnostic} rather than a selector: as
Section~\ref{sec:multimachine:general} shows, BIC does not identify $k$
on records of this size, and the physics constraints do the selecting.

\textbf{(iii) Independently derived classifiers.} Two threshold rules
constructed from the same raw channels but with no reference to the
mixture model: a power-factor classifier and a magnetising-current-ratio
classifier, each with its threshold chosen by Otsu's
criterion~\cite{otsu1979threshold} rather than by hand, so that neither
is tuned to agree. On pelletizer~I the Otsu power-factor threshold falls
at $0.474$ and the resulting classifier agrees with the state model on
$\mathbf{99.4\%}$ of energised rows; the current-ratio classifier agrees
on $97.3\%$. Agreement across the eight machines is $85$--$99.7\%$, with
one instructive exception at $55.1\%$ on the machine whose physics
checks fail hardest (Section~\ref{sec:multimachine:states}). Verdicts
are aggregated by consensus~\cite{strehl2002cluster}.

Finally, the plant's own calendar provides an external operational
check: the facility halts every weekday between 17:00 and 22:00 local
time to avoid the peak tariff, and does not run at weekends. On the full
record $76.8\%$ of readings inside that closed window are labelled
non-productive. The converse is reported alongside it, because a
degenerate partition calling everything OFF would otherwise score
$100\%$.

\textbf{The battery is not decoration, and Phase~V measured what it is
worth.} Re-fitting the labeller at ten seeds moves the headline STANDBY
total by a coefficient of variation of $33.7\%$, and three seeds in ten
converge on a cluster that is physically a different state altogether.
The battery scores exactly those three below threshold and the other
seven at $1.00$, and conditioning on a passing score collapses the
spread to $3.0\%$ (Section~\ref{sec:significance:seeds}). That is why
the accept/reject loop is part of the method rather than an appendix to
it, and why every state-time total in this paper is quoted conditional
on it.

%  Figures are referenced by bare name; add to the preamble
%      \graphicspath{{../outputs/phase2/figures/pelletizer-I/}}
%  Bibliography: paper/references.bib
%  Numbers below are generated by
%      python run_phase2.py
%  and stored in outputs/phase2/task3/<machine>/task3_results.json,
%  outputs/phase2/task4/<machine>/{comparison_table.csv,
%  flicker_diagnosis.json,sticky_prior_sweep.csv,min_dwell_decode.csv}.
% =====================================================================

\section{The State Layer}
\label{sec:states}
\label{sec:results}

This section reports what the state-inference stage of
Section~\ref{sec:method:labelling} produces on pelletizer~I: which
measured quantities separate the states, how the choice of clustering
algorithm compares with the choice of decoding rule, and how a partition
is validated when there is nothing to validate it against.

\subsection{What separates STANDBY}
\label{sec:states:features}

Twenty-four features are derived from the six measured channels in four
groups: raw electrical (5), derived electrical (4 --- power factor,
$Q/P$ ratio, load factor, apparent-power residual), rolling statistics
(8, over a $300$\,s window), and temporal (7). Rolling windows are
\emph{trailing}, not centred, and are computed within contiguous
segments so that none spans an acquisition gap.

\textbf{The target is STANDBY against $\{$WORKING, PEAK\_LOAD$\}$ on
energised rows only.} Separating OFF from everything else is trivial and
would dominate any importance ranking computed over all four classes;
separating an energised idle machine from a lightly loaded running one
is the problem the framework actually has to solve. A four-class
ranking was also computed and is not reported here, because it measures
the OFF boundary and invites the two to be averaged.

\begin{table}[t]
\centering
\caption{Feature-group importance for the STANDBY-versus-productive
split, by three measures with different failure modes. Each column is
normalised to its own total.}
\label{tab:feature-groups}
\small
\begin{tabular}{@{}lrrr@{}}
\toprule
Group & mutual info. & split gain & SHAP \\
\midrule
raw electrical      & $0.350$ & $\mathbf{0.560}$ & $\mathbf{0.537}$ \\
derived electrical  & $0.250$ & $0.438$ & $0.310$ \\
rolling statistics  & $\mathbf{0.388}$ & $0.001$ & $0.116$ \\
temporal            & $0.012$ & $0.002$ & $0.036$ \\
\bottomrule
\end{tabular}
\end{table}

Three readings of Table~\ref{tab:feature-groups} matter for the rest of
the paper.

\textbf{The physics argument is empirically confirmed.}
\texttt{power\_factor} and \texttt{q\_p\_ratio} rank second and third by
consensus over the three measures, behind \texttt{active\_power} and
ahead of \texttt{current}, and the four derived electrical features
carry $31$--$44\%$ of total importance. Methods~1 and~2 of the
validation battery (Section~\ref{sec:validation}) were previously
justified by induction-motor theory alone; they now have a measured
basis on the machine's own data. Section~\ref{sec:explain:state}
strengthens this: under Shapley attribution power factor overtakes
active power outright.

\textbf{Temporal features are nearly irrelevant to this split}
($1$--$4\%$). That is the desirable outcome, not a null one: a state
model leaning on \texttt{hour\_of\_day} would be reproducing the shift
schedule rather than the machine's electrical condition, and would not
transfer to a facility with a different calendar. It is also the exact
opposite of what the \emph{decision} layer's forecaster does, which
draws $82\%$ of its attribution from the calendar
(Section~\ref{sec:crossmachine:decision}) --- the two stages of the
framework depend on almost disjoint information, and only the state
stage is physics-driven.

\textbf{High mutual information with near-zero split gain is the
redundancy signature.} Rolling statistics are individually informative
(MI $0.388$, the largest in the column) and add nothing once the tree
already holds the instantaneous value (gain $0.001$). Reporting either
measure alone would misdescribe them.

\paragraph{A caveat that must not be dropped.}
The surrogate used for the gain and SHAP columns reproduces its targets
to $99.9\%$ accuracy. \emph{That number is not a validation result.} The
targets are the state model's own labels and the state model was fitted
on a subset of these same features, so the surrogate is largely
re-deriving a deterministic function of its own inputs. The rankings are
meaningful; the accuracy is evidence of nothing. Validation is
Section~\ref{sec:validation}'s job.

\subsection{The emission model barely matters; the decoding rule does}
\label{sec:results:task4}

Six state-inference configurations are compared on one representation
and one protocol: $64$ contiguous two-hour blocks spread across the
record and never crossing a gap, with alternate blocks forming a fitting
set and an evaluation set of $230\,400$ rows each. DBSCAN and spectral
clustering are transductive, so their labels are extended to the
evaluation blocks by one-nearest-neighbour, applied identically to both.
Because no ground truth exists, the comparison is scored on flicker,
median dwell, the physics battery, the balanced schedule score, and the
STANDBY hour total --- never on accuracy.

\begin{table}[t]
\centering
\caption{State-inference comparison on the evaluation blocks. Flicker is
the fraction of state runs shorter than the $10$\,s physical
plausibility floor; the binary column recomputes it on the
STANDBY-versus-rest sequence the decision layer consumes. No column is
interpretable on its own --- see the spectral row.}
\label{tab:state-detection}
\small
\setlength{\tabcolsep}{4pt}
\begin{tabular}{@{}lrrrrrr@{}}
\toprule
Method & flick. & flick. & dwell & phys. & sched. & STANDBY \\
       & \%     & bin.\% & (s)   &       & \%      & (h) \\
\midrule
$k$-means \cite{lloyd1982least}      & $78.2$ & $49.1$ & $2$ & $5/5$ & $84.4$ & $4.38$ \\
DBSCAN \cite{ester1996dbscan}        & $62.5$ & $50.3$ & $4$ & $5/5$ & $84.4$ & $4.34$ \\
Spectral \cite{ng2001spectral}       & $98.2$ & $98.8$ & $2$ & $\mathbf{4/5}$ & $\mathbf{89.0}$ & $\mathbf{15.01}$ \\
GMM, i.i.d.\ \cite{mclachlan2000finite} & $66.9$ & $53.6$ & $3$ & $5/5$ & $84.5$ & $4.17$ \\
GMM--HMM \cite{rabiner1989tutorial}  & $44.5$ & $44.7$ & $13$ & $5/5$ & $84.5$ & $4.16$ \\
\textbf{GMM--HMM + min-dwell}        & $\mathbf{0.00}$ & $\mathbf{0.00}$ & $\mathbf{92.5}$ & $5/5$ & $84.5$ & $4.20$ \\
\bottomrule
\end{tabular}
\end{table}

\begin{figure*}[t]
\centering
\includegraphics[width=\textwidth]{fig_p7_state_detection}
\caption{The same six configurations as
Table~\ref{tab:state-detection}, plotted. \textbf{(A)} Flicker against
median dwell. The three filled points share one set of GMM emissions
and differ only in how those emissions are decoded; the arrows follow
i.i.d.\ assignment $\rightarrow$ Viterbi $\rightarrow$ Viterbi under a
minimum-dwell floor. The three squares are alternative emission models
decoded i.i.d.\ Changing the emission model moves the result within one
cluster; changing the decoding rule moves it two orders of magnitude in
dwell and to zero in flicker. \textbf{(B)} Inferred STANDBY hours.
Spectral clustering assigns $3.6\times$ every other method's total while
scoring the highest factory-schedule agreement of any row in the table,
which is why the battery is argued for in preference to any single
ground-truth-free metric. Hatching marks a failed physics check in both
panels.}
\label{fig:state-detection}
\end{figure*}

\textbf{Choosing a better clustering algorithm buys almost nothing;
choosing a better sequence model buys everything} ---
Figure~\ref{fig:state-detection}(A) is the table read as a picture, and
the two families separate along different axes. $k$-means, DBSCAN and
the GMM all pass $5/5$ physics checks, agree with the factory schedule
to within $0.2$ points, and find $4.17$--$4.38$ STANDBY hours on the
evaluation blocks. Switching from i.i.d.\ assignment to Viterbi decoding
over the \emph{same} emissions cuts flicker from $66.9\%$ to $44.5\%$
and quadruples the median dwell; adding the minimum-dwell constraint
takes flicker to zero and the median dwell to $92.5$\,s. The
emission-model rows span $62$--$78\%$ of flicker; the decoding rows span
$0$--$67\%$.

The comparison is clean in the one place it needs to be: the
\textsc{gmm} and \textsc{gmm--hmm} rows differ \emph{only} in the
decoding rule, because the HMM's emissions are warm-started from the
fitted GMM and held fixed (Section~\ref{sec:method:ghmm}). Fitting HMM
emissions from scratch by Baum--Welch converges to a materially worse
optimum on this record --- its ``STANDBY'' absorbs part of the
productive range at $10.3$\,kW instead of finding the no-load band at
$3.67$\,kW, and it fails the variance-ordering check --- which would
have confounded the value of temporal decoding with EM initialisation
luck.

\textbf{Spectral clustering is the one genuine failure, and it is
instructive} (Figure~\ref{fig:state-detection}(B)). It assigns $15.01$
STANDBY hours, $3.6\times$ every other
method, with $98.8\%$ binary flicker and a no-load current ratio of
$0.001$ against the acceptance band $[0.25,0.50]$. It also scores the
\emph{highest} schedule agreement in the table ($89.0\%$): a partition
that over-assigns non-productive states automatically looks good on the
factory-closed window. That is a weakness of the schedule metric in
isolation and the reason the STANDBY hour total is printed beside it.
\textbf{No single ground-truth-free metric can be quoted alone}, which
is the argument for the battery rather than for a favourite score.

\paragraph{A protocol correction worth recording.}
DBSCAN chooses its own cluster count and returned $28$. Naming the four
lowest-power clusters with the four canonical state names put all four
``states'' inside the machine's low-power range and produced an apparent
total failure --- $0.06$ STANDBY hours, $98.9\%$ flicker --- that was
entirely an artefact of the mapping. Cluster means are now grouped into
four size-weighted power bands, so sub-clusters of one physical state
are re-merged, and DBSCAN passes $5/5$ with $4.34$ STANDBY hours. Any
comparison against a method that over-segments needs this step, and
without it this paper would have reported a baseline failure that does
not exist.

\subsection{Flicker: why a transition prior cannot remove it}
\label{sec:flicker}

The expectation going in was that Viterbi decoding under a sticky prior
would take flicker to a few per cent. It does not, and the reason is
arithmetic rather than empirical.

\textbf{The sticky prior as conventionally set is numerically inert.}
Setting the Dirichlet pseudo-count to the target dwell in samples gives
$\kappa=30$ against roughly $576\,000$ observed transitions. A flat
prior and $\kappa=30$ give flicker $64.10\%$ and $64.09\%$ and minimum
self-transition $0.9677$ and $0.9678$: the prior is present in the
model, cited in its documentation, and has no measurable effect.

\textbf{Strengthening it does not help either.} Scaling the prior to ten
times the total evidence ($\kappa=575\,920$, minimum self-transition
$0.9994$, implied mean dwell $1\,605$\,s) still leaves flicker at
$50.2\%$.

\textbf{The bound is structural.} In a Viterbi decode the cost of
changing state is $\log(p_{\text{self}}/p_{\text{switch}})$, measured
here at $9.23$ nats at its largest and bounded above however strong the
prior is made. The emission log-likelihood gap between the best and
second-best state has median $27.8$ nats and $90$th percentile
$41\,701$ nats. \textbf{At $55.7\%$ of timesteps the emission term
exceeds the largest penalty the transition matrix can apply}, so the
transition model is arithmetically incapable of changing the decision
there and the HMM degenerates toward per-sample assignment. This is a
general property of applying an HMM to low-noise, multi-channel
measurements, not a property of this record.

\textbf{What fixes it is an explicit duration constraint} --- the
discrete analogue of a hidden semi-Markov duration
model~\cite{yu2010hsmm}, imposed at decode time. The floor is set equal
to the flicker threshold itself, so it removes exactly what the metric
defines as implausible and nothing more.
Table~\ref{tab:min-dwell} sweeps the floor over an order of magnitude.

\begin{table}[t]
\centering
\caption{Minimum-dwell sweep. Flicker goes to zero at no cost to the
physics checks, the schedule agreement, or the STANDBY total. This
sweep is run on the flicker-diagnosis decode rather than on the
evaluation blocks of Table~\ref{tab:state-detection}, which is why its
absolute physics score and STANDBY total differ from that table's; the
comparison of interest is down the column, not across the two tables.}
\label{tab:min-dwell}
\small
\begin{tabular}{@{}rrrrr@{}}
\toprule
min.\ dwell & flicker \% & median dwell (s) & physics & STANDBY (h) \\
\midrule
$0$\,s  & $64.09$ & $4$   & $4/5$ & $5.34$ \\
$10$\,s & $\mathbf{0.00}$ & $153$ & $4/5$ & $5.34$ \\
$30$\,s & $0.00$ & $327.5$ & $4/5$ & $5.37$ \\
$60$\,s & $0.00$ & $436.5$ & $4/5$ & $5.40$ \\
\bottomrule
\end{tabular}
\end{table}

\begin{figure}[t]
\centering
\includegraphics[width=\columnwidth]{fig_p7_min_dwell}
\caption{The minimum-dwell sweep of Table~\ref{tab:min-dwell}.
\textbf{(A)} Flicker under both definitions collapses to $0.00\%$ at the
first non-zero floor and stays there. \textbf{(B)} The median dwell
rises by two orders of magnitude while the inferred STANDBY total is
flat. The right-hand axis is anchored at zero deliberately: on an
auto-scaled axis the $1\%$ drift across the whole sweep would render as
a pronounced trend and say the opposite of what the data says.}
\label{fig:min-dwell}
\end{figure}

The constraint is removing noise, not reshaping the partition
(Figure~\ref{fig:min-dwell}): the
STANDBY total moves by $0.00$\,h from $0$\,s to $10$\,s and by $1\%$ out
to $60$\,s, and no physics check changes verdict. Its cost is measured
directly --- $0.19\%$ of rows relabelled on this machine, $0.15$--$6.1\%$
across the plant (Section~\ref{sec:multimachine:general}) --- and it is
the component the ablation of Section~\ref{sec:ablation:state} credits
with removing flicker, in preference to the HMM.

\subsection{Validation without ground truth}
\label{sec:validation}

IMDELD supplies no per-timestamp state annotation, and obtaining one
would require instrumenting each machine's contactor. The partition
therefore cannot be scored by any supervised metric, and the alternative
adopted here is to require agreement with evidence the model does not
have access to.

Three independent bodies of evidence are used, and they fail in
different ways, which is the property that makes their agreement
informative.

\textbf{(i) Circuit physics.} The five checks C1--C5 of
Section~\ref{sec:method:accept_reject}, of which C3 (power-factor
separation) and C4 (no-load current ratio) are the two that do not
depend on power magnitude and therefore the two that can distinguish an
energised idle machine from a lightly loaded one. Pelletizer~I passes
$5/5$, with a power-factor separation of $0.72$ and a no-load current
ratio of $0.36$ --- inside the $[0.25,0.50]$ band that induction-motor
theory predicts and that no part of the fitting procedure was told
about.

\textbf{(ii) Model-level diagnostics.} A ten-test harness (T1--T10)
covering data sanity, EM convergence, model order, component separation
at $2\sigma$, assignment confidence, ghost components, physical
thresholds, variance ordering, temporal alignment with nights and
weekends, and agreement with the published plant calendar. T3 is
reported as a \emph{diagnostic} rather than a selector: as
Section~\ref{sec:multimachine:general} shows, BIC does not identify $k$
on records of this size, and the physics constraints do the selecting.

\textbf{(iii) Independently derived classifiers.} Two threshold rules
constructed from the same raw channels but with no reference to the
mixture model: a power-factor classifier and a magnetising-current-ratio
classifier, each with its threshold chosen by Otsu's
criterion~\cite{otsu1979threshold} rather than by hand, so that neither
is tuned to agree. On pelletizer~I the Otsu power-factor threshold falls
at $0.474$ and the resulting classifier agrees with the state model on
$\mathbf{99.4\%}$ of energised rows; the current-ratio classifier agrees
on $97.3\%$. Agreement across the eight machines is $85$--$99.7\%$, with
one instructive exception at $55.1\%$ on the machine whose physics
checks fail hardest (Section~\ref{sec:multimachine:states}). Verdicts
are aggregated by consensus~\cite{strehl2002cluster}.

Finally, the plant's own calendar provides an external operational
check: the facility halts every weekday between 17:00 and 22:00 local
time to avoid the peak tariff, and does not run at weekends. On the full
record $76.8\%$ of readings inside that closed window are labelled
non-productive. The converse is reported alongside it, because a
degenerate partition calling everything OFF would otherwise score
$100\%$.

\textbf{The battery is not decoration, and Phase~V measured what it is
worth.} Re-fitting the labeller at ten seeds moves the headline STANDBY
total by a coefficient of variation of $33.7\%$, and three seeds in ten
converge on a cluster that is physically a different state altogether.
The battery scores exactly those three below threshold and the other
seven at $1.00$, and conditioning on a passing score collapses the
spread to $3.0\%$ (Section~\ref{sec:significance:seeds}). That is why
the accept/reject loop is part of the method rather than an appendix to
it, and why every state-time total in this paper is quoted conditional
on it.

%  Figures are referenced by bare name; add to the preamble
%      \graphicspath{{../outputs/phase2/figures/pelletizer-I/}}
%  Bibliography: paper/references.bib
%  Numbers below are generated by
%      python run_phase2.py
%  and stored in outputs/phase2/task3/<machine>/task3_results.json,
%  outputs/phase2/task4/<machine>/{comparison_table.csv,
%  flicker_diagnosis.json,sticky_prior_sweep.csv,min_dwell_decode.csv}.
% =====================================================================

\section{The State Layer}
\label{sec:states}
\label{sec:results}

This section reports what the state-inference stage of
Section~\ref{sec:method:labelling} produces on pelletizer~I: which
measured quantities separate the states, how the choice of clustering
algorithm compares with the choice of decoding rule, and how a partition
is validated when there is nothing to validate it against.

\subsection{What separates STANDBY}
\label{sec:states:features}

Twenty-four features are derived from the six measured channels in four
groups: raw electrical (5), derived electrical (4 --- power factor,
$Q/P$ ratio, load factor, apparent-power residual), rolling statistics
(8, over a $300$\,s window), and temporal (7). Rolling windows are
\emph{trailing}, not centred, and are computed within contiguous
segments so that none spans an acquisition gap.

\textbf{The target is STANDBY against $\{$WORKING, PEAK\_LOAD$\}$ on
energised rows only.} Separating OFF from everything else is trivial and
would dominate any importance ranking computed over all four classes;
separating an energised idle machine from a lightly loaded running one
is the problem the framework actually has to solve. A four-class
ranking was also computed and is not reported here, because it measures
the OFF boundary and invites the two to be averaged.

\begin{table}[t]
\centering
\caption{Feature-group importance for the STANDBY-versus-productive
split, by three measures with different failure modes. Each column is
normalised to its own total.}
\label{tab:feature-groups}
\small
\begin{tabular}{@{}lrrr@{}}
\toprule
Group & mutual info. & split gain & SHAP \\
\midrule
raw electrical      & $0.350$ & $\mathbf{0.560}$ & $\mathbf{0.537}$ \\
derived electrical  & $0.250$ & $0.438$ & $0.310$ \\
rolling statistics  & $\mathbf{0.388}$ & $0.001$ & $0.116$ \\
temporal            & $0.012$ & $0.002$ & $0.036$ \\
\bottomrule
\end{tabular}
\end{table}

Three readings of Table~\ref{tab:feature-groups} matter for the rest of
the paper.

\textbf{The physics argument is empirically confirmed.}
\texttt{power\_factor} and \texttt{q\_p\_ratio} rank second and third by
consensus over the three measures, behind \texttt{active\_power} and
ahead of \texttt{current}, and the four derived electrical features
carry $31$--$44\%$ of total importance. Methods~1 and~2 of the
validation battery (Section~\ref{sec:validation}) were previously
justified by induction-motor theory alone; they now have a measured
basis on the machine's own data. Section~\ref{sec:explain:state}
strengthens this: under Shapley attribution power factor overtakes
active power outright.

\textbf{Temporal features are nearly irrelevant to this split}
($1$--$4\%$). That is the desirable outcome, not a null one: a state
model leaning on \texttt{hour\_of\_day} would be reproducing the shift
schedule rather than the machine's electrical condition, and would not
transfer to a facility with a different calendar. It is also the exact
opposite of what the \emph{decision} layer's forecaster does, which
draws $82\%$ of its attribution from the calendar
(Section~\ref{sec:crossmachine:decision}) --- the two stages of the
framework depend on almost disjoint information, and only the state
stage is physics-driven.

\textbf{High mutual information with near-zero split gain is the
redundancy signature.} Rolling statistics are individually informative
(MI $0.388$, the largest in the column) and add nothing once the tree
already holds the instantaneous value (gain $0.001$). Reporting either
measure alone would misdescribe them.

\paragraph{A caveat that must not be dropped.}
The surrogate used for the gain and SHAP columns reproduces its targets
to $99.9\%$ accuracy. \emph{That number is not a validation result.} The
targets are the state model's own labels and the state model was fitted
on a subset of these same features, so the surrogate is largely
re-deriving a deterministic function of its own inputs. The rankings are
meaningful; the accuracy is evidence of nothing. Validation is
Section~\ref{sec:validation}'s job.

\subsection{The emission model barely matters; the decoding rule does}
\label{sec:results:task4}

Six state-inference configurations are compared on one representation
and one protocol: $64$ contiguous two-hour blocks spread across the
record and never crossing a gap, with alternate blocks forming a fitting
set and an evaluation set of $230\,400$ rows each. DBSCAN and spectral
clustering are transductive, so their labels are extended to the
evaluation blocks by one-nearest-neighbour, applied identically to both.
Because no ground truth exists, the comparison is scored on flicker,
median dwell, the physics battery, the balanced schedule score, and the
STANDBY hour total --- never on accuracy.

\begin{table}[t]
\centering
\caption{State-inference comparison on the evaluation blocks. Flicker is
the fraction of state runs shorter than the $10$\,s physical
plausibility floor; the binary column recomputes it on the
STANDBY-versus-rest sequence the decision layer consumes. No column is
interpretable on its own --- see the spectral row.}
\label{tab:state-detection}
\small
\setlength{\tabcolsep}{4pt}
\begin{tabular}{@{}lrrrrrr@{}}
\toprule
Method & flick. & flick. & dwell & phys. & sched. & STANDBY \\
       & \%     & bin.\% & (s)   &       & \%      & (h) \\
\midrule
$k$-means \cite{lloyd1982least}      & $78.2$ & $49.1$ & $2$ & $5/5$ & $84.4$ & $4.38$ \\
DBSCAN \cite{ester1996dbscan}        & $62.5$ & $50.3$ & $4$ & $5/5$ & $84.4$ & $4.34$ \\
Spectral \cite{ng2001spectral}       & $98.2$ & $98.8$ & $2$ & $\mathbf{4/5}$ & $\mathbf{89.0}$ & $\mathbf{15.01}$ \\
GMM, i.i.d.\ \cite{mclachlan2000finite} & $66.9$ & $53.6$ & $3$ & $5/5$ & $84.5$ & $4.17$ \\
GMM--HMM \cite{rabiner1989tutorial}  & $44.5$ & $44.7$ & $13$ & $5/5$ & $84.5$ & $4.16$ \\
\textbf{GMM--HMM + min-dwell}        & $\mathbf{0.00}$ & $\mathbf{0.00}$ & $\mathbf{92.5}$ & $5/5$ & $84.5$ & $4.20$ \\
\bottomrule
\end{tabular}
\end{table}

\begin{figure*}[t]
\centering
\includegraphics[width=\textwidth]{fig_p7_state_detection}
\caption{The same six configurations as
Table~\ref{tab:state-detection}, plotted. \textbf{(A)} Flicker against
median dwell. The three filled points share one set of GMM emissions
and differ only in how those emissions are decoded; the arrows follow
i.i.d.\ assignment $\rightarrow$ Viterbi $\rightarrow$ Viterbi under a
minimum-dwell floor. The three squares are alternative emission models
decoded i.i.d.\ Changing the emission model moves the result within one
cluster; changing the decoding rule moves it two orders of magnitude in
dwell and to zero in flicker. \textbf{(B)} Inferred STANDBY hours.
Spectral clustering assigns $3.6\times$ every other method's total while
scoring the highest factory-schedule agreement of any row in the table,
which is why the battery is argued for in preference to any single
ground-truth-free metric. Hatching marks a failed physics check in both
panels.}
\label{fig:state-detection}
\end{figure*}

\textbf{Choosing a better clustering algorithm buys almost nothing;
choosing a better sequence model buys everything} ---
Figure~\ref{fig:state-detection}(A) is the table read as a picture, and
the two families separate along different axes. $k$-means, DBSCAN and
the GMM all pass $5/5$ physics checks, agree with the factory schedule
to within $0.2$ points, and find $4.17$--$4.38$ STANDBY hours on the
evaluation blocks. Switching from i.i.d.\ assignment to Viterbi decoding
over the \emph{same} emissions cuts flicker from $66.9\%$ to $44.5\%$
and quadruples the median dwell; adding the minimum-dwell constraint
takes flicker to zero and the median dwell to $92.5$\,s. The
emission-model rows span $62$--$78\%$ of flicker; the decoding rows span
$0$--$67\%$.

The comparison is clean in the one place it needs to be: the
\textsc{gmm} and \textsc{gmm--hmm} rows differ \emph{only} in the
decoding rule, because the HMM's emissions are warm-started from the
fitted GMM and held fixed (Section~\ref{sec:method:ghmm}). Fitting HMM
emissions from scratch by Baum--Welch converges to a materially worse
optimum on this record --- its ``STANDBY'' absorbs part of the
productive range at $10.3$\,kW instead of finding the no-load band at
$3.67$\,kW, and it fails the variance-ordering check --- which would
have confounded the value of temporal decoding with EM initialisation
luck.

\textbf{Spectral clustering is the one genuine failure, and it is
instructive} (Figure~\ref{fig:state-detection}(B)). It assigns $15.01$
STANDBY hours, $3.6\times$ every other
method, with $98.8\%$ binary flicker and a no-load current ratio of
$0.001$ against the acceptance band $[0.25,0.50]$. It also scores the
\emph{highest} schedule agreement in the table ($89.0\%$): a partition
that over-assigns non-productive states automatically looks good on the
factory-closed window. That is a weakness of the schedule metric in
isolation and the reason the STANDBY hour total is printed beside it.
\textbf{No single ground-truth-free metric can be quoted alone}, which
is the argument for the battery rather than for a favourite score.

\paragraph{A protocol correction worth recording.}
DBSCAN chooses its own cluster count and returned $28$. Naming the four
lowest-power clusters with the four canonical state names put all four
``states'' inside the machine's low-power range and produced an apparent
total failure --- $0.06$ STANDBY hours, $98.9\%$ flicker --- that was
entirely an artefact of the mapping. Cluster means are now grouped into
four size-weighted power bands, so sub-clusters of one physical state
are re-merged, and DBSCAN passes $5/5$ with $4.34$ STANDBY hours. Any
comparison against a method that over-segments needs this step, and
without it this paper would have reported a baseline failure that does
not exist.

\subsection{Flicker: why a transition prior cannot remove it}
\label{sec:flicker}

The expectation going in was that Viterbi decoding under a sticky prior
would take flicker to a few per cent. It does not, and the reason is
arithmetic rather than empirical.

\textbf{The sticky prior as conventionally set is numerically inert.}
Setting the Dirichlet pseudo-count to the target dwell in samples gives
$\kappa=30$ against roughly $576\,000$ observed transitions. A flat
prior and $\kappa=30$ give flicker $64.10\%$ and $64.09\%$ and minimum
self-transition $0.9677$ and $0.9678$: the prior is present in the
model, cited in its documentation, and has no measurable effect.

\textbf{Strengthening it does not help either.} Scaling the prior to ten
times the total evidence ($\kappa=575\,920$, minimum self-transition
$0.9994$, implied mean dwell $1\,605$\,s) still leaves flicker at
$50.2\%$.

\textbf{The bound is structural.} In a Viterbi decode the cost of
changing state is $\log(p_{\text{self}}/p_{\text{switch}})$, measured
here at $9.23$ nats at its largest and bounded above however strong the
prior is made. The emission log-likelihood gap between the best and
second-best state has median $27.8$ nats and $90$th percentile
$41\,701$ nats. \textbf{At $55.7\%$ of timesteps the emission term
exceeds the largest penalty the transition matrix can apply}, so the
transition model is arithmetically incapable of changing the decision
there and the HMM degenerates toward per-sample assignment. This is a
general property of applying an HMM to low-noise, multi-channel
measurements, not a property of this record.

\textbf{What fixes it is an explicit duration constraint} --- the
discrete analogue of a hidden semi-Markov duration
model~\cite{yu2010hsmm}, imposed at decode time. The floor is set equal
to the flicker threshold itself, so it removes exactly what the metric
defines as implausible and nothing more.
Table~\ref{tab:min-dwell} sweeps the floor over an order of magnitude.

\begin{table}[t]
\centering
\caption{Minimum-dwell sweep. Flicker goes to zero at no cost to the
physics checks, the schedule agreement, or the STANDBY total. This
sweep is run on the flicker-diagnosis decode rather than on the
evaluation blocks of Table~\ref{tab:state-detection}, which is why its
absolute physics score and STANDBY total differ from that table's; the
comparison of interest is down the column, not across the two tables.}
\label{tab:min-dwell}
\small
\begin{tabular}{@{}rrrrr@{}}
\toprule
min.\ dwell & flicker \% & median dwell (s) & physics & STANDBY (h) \\
\midrule
$0$\,s  & $64.09$ & $4$   & $4/5$ & $5.34$ \\
$10$\,s & $\mathbf{0.00}$ & $153$ & $4/5$ & $5.34$ \\
$30$\,s & $0.00$ & $327.5$ & $4/5$ & $5.37$ \\
$60$\,s & $0.00$ & $436.5$ & $4/5$ & $5.40$ \\
\bottomrule
\end{tabular}
\end{table}

\begin{figure}[t]
\centering
\includegraphics[width=\columnwidth]{fig_p7_min_dwell}
\caption{The minimum-dwell sweep of Table~\ref{tab:min-dwell}.
\textbf{(A)} Flicker under both definitions collapses to $0.00\%$ at the
first non-zero floor and stays there. \textbf{(B)} The median dwell
rises by two orders of magnitude while the inferred STANDBY total is
flat. The right-hand axis is anchored at zero deliberately: on an
auto-scaled axis the $1\%$ drift across the whole sweep would render as
a pronounced trend and say the opposite of what the data says.}
\label{fig:min-dwell}
\end{figure}

The constraint is removing noise, not reshaping the partition
(Figure~\ref{fig:min-dwell}): the
STANDBY total moves by $0.00$\,h from $0$\,s to $10$\,s and by $1\%$ out
to $60$\,s, and no physics check changes verdict. Its cost is measured
directly --- $0.19\%$ of rows relabelled on this machine, $0.15$--$6.1\%$
across the plant (Section~\ref{sec:multimachine:general}) --- and it is
the component the ablation of Section~\ref{sec:ablation:state} credits
with removing flicker, in preference to the HMM.

\subsection{Validation without ground truth}
\label{sec:validation}

IMDELD supplies no per-timestamp state annotation, and obtaining one
would require instrumenting each machine's contactor. The partition
therefore cannot be scored by any supervised metric, and the alternative
adopted here is to require agreement with evidence the model does not
have access to.

Three independent bodies of evidence are used, and they fail in
different ways, which is the property that makes their agreement
informative.

\textbf{(i) Circuit physics.} The five checks C1--C5 of
Section~\ref{sec:method:accept_reject}, of which C3 (power-factor
separation) and C4 (no-load current ratio) are the two that do not
depend on power magnitude and therefore the two that can distinguish an
energised idle machine from a lightly loaded one. Pelletizer~I passes
$5/5$, with a power-factor separation of $0.72$ and a no-load current
ratio of $0.36$ --- inside the $[0.25,0.50]$ band that induction-motor
theory predicts and that no part of the fitting procedure was told
about.

\textbf{(ii) Model-level diagnostics.} A ten-test harness (T1--T10)
covering data sanity, EM convergence, model order, component separation
at $2\sigma$, assignment confidence, ghost components, physical
thresholds, variance ordering, temporal alignment with nights and
weekends, and agreement with the published plant calendar. T3 is
reported as a \emph{diagnostic} rather than a selector: as
Section~\ref{sec:multimachine:general} shows, BIC does not identify $k$
on records of this size, and the physics constraints do the selecting.

\textbf{(iii) Independently derived classifiers.} Two threshold rules
constructed from the same raw channels but with no reference to the
mixture model: a power-factor classifier and a magnetising-current-ratio
classifier, each with its threshold chosen by Otsu's
criterion~\cite{otsu1979threshold} rather than by hand, so that neither
is tuned to agree. On pelletizer~I the Otsu power-factor threshold falls
at $0.474$ and the resulting classifier agrees with the state model on
$\mathbf{99.4\%}$ of energised rows; the current-ratio classifier agrees
on $97.3\%$. Agreement across the eight machines is $85$--$99.7\%$, with
one instructive exception at $55.1\%$ on the machine whose physics
checks fail hardest (Section~\ref{sec:multimachine:states}). Verdicts
are aggregated by consensus~\cite{strehl2002cluster}.

Finally, the plant's own calendar provides an external operational
check: the facility halts every weekday between 17:00 and 22:00 local
time to avoid the peak tariff, and does not run at weekends. On the full
record $76.8\%$ of readings inside that closed window are labelled
non-productive. The converse is reported alongside it, because a
degenerate partition calling everything OFF would otherwise score
$100\%$.

\textbf{The battery is not decoration, and Phase~V measured what it is
worth.} Re-fitting the labeller at ten seeds moves the headline STANDBY
total by a coefficient of variation of $33.7\%$, and three seeds in ten
converge on a cluster that is physically a different state altogether.
The battery scores exactly those three below threshold and the other
seven at $1.00$, and conditioning on a passing score collapses the
spread to $3.0\%$ (Section~\ref{sec:significance:seeds}). That is why
the accept/reject loop is part of the method rather than an appendix to
it, and why every state-time total in this paper is quoted conditional
on it.


% --- Forecasting comparison against an untrained reference
% =====================================================================
%  forecasting.tex  --  PHASE II Task 5 + PHASE V Task 12A
%
%  The forecasting comparison: seven models on one protocol, reported
%  as mean +/- sd over seeds, anchored on an untrained reference.
%
%  Intended placement: Section 5 (\label{sec:forecasting}), after the
%  state layer and before the decision layer.
%
%  Drop-in usage: % =====================================================================
%  forecasting.tex  --  PHASE II Task 5 + PHASE V Task 12A
%
%  The forecasting comparison: seven models on one protocol, reported
%  as mean +/- sd over seeds, anchored on an untrained reference.
%
%  Intended placement: Section 5 (\label{sec:forecasting}), after the
%  state layer and before the decision layer.
%
%  Drop-in usage: % =====================================================================
%  forecasting.tex  --  PHASE II Task 5 + PHASE V Task 12A
%
%  The forecasting comparison: seven models on one protocol, reported
%  as mean +/- sd over seeds, anchored on an untrained reference.
%
%  Intended placement: Section 5 (\label{sec:forecasting}), after the
%  state layer and before the decision layer.
%
%  Drop-in usage: \input{forecasting}
%  Figures are referenced by bare name; add to the preamble
%      \graphicspath{{../outputs/phase2/figures/pelletizer-I/}
%                    {../outputs/phase6/figures/}}
%  Bibliography: paper/references.bib
%  Numbers below are generated by
%      python -m experiments.task5_forecasting_baselines   (single seed)
%      python run_phase5.py                                (seeds)
%  and stored in
%      outputs/phase2/task5/<machine>/comparison_table.csv
%      outputs/phase5/task12/<machine>/forecasting/seed_table.csv
%  Significance testing of every contrast quoted here is
%  Section~\ref{sec:significance}; attribution is
%  Section~\ref{sec:explain:window}.
% =====================================================================

\section{Forecasting the Non-Productive Horizon}
\label{sec:forecasting}

The state layer of Section~\ref{sec:states} produces a label per second.
A shutdown decision consumes something else: how much longer the machine
will remain non-productive. This section evaluates that forecasting step
as a supervised sequence task over the inferred labels, under one
protocol, against a reference that does no training at all.

\subsection{Protocol}
\label{sec:forecasting:protocol}

The labelled record is decimated $10\times$ to $0.1$\,Hz and cut into
windows of $600$\,s lookback and $300$\,s horizon at a stride of
$120$\,s. The split is chronological ($70/15/15$) and windows are built
\emph{inside} each split, so no window straddles a split boundary and no
future sample can enter a training window. Class weights are balanced
identically for every model. The test set is $5\,449$ windows, on which
the mean true STANDBY content is $20.1$\,s of a $300$\,s horizon and
only $13.5\%$ of windows contain any STANDBY at all.

Two metrics are reported. STANDBY $F_1$ is the per-step classification
quality on the minority class the decision layer cares about. The
per-window \emph{STANDBY-seconds MAE} --- the absolute error in how many
of the next $300$\,s are predicted to be standby --- is the quantity
closer to what a duration-consuming policy would use, and it is the
metric on which the significance tests of
Section~\ref{sec:significance} are pre-registered.

\paragraph{Seeds are reported, not a single run.}
Every stochastic model is trained at multiple seeds and the table
reports mean\,$\pm$\,sd: $15$ seeds for the gradient-boosted tree, $5$
for each neural architecture, $1$ for the deterministic reference. This
is not a formality here. The spread between a good and a bad seed of one
architecture reaches $0.166$ of $F_1$, which is larger than the entire
spread between architectures, so a single-seed table of this comparison
is not a ranking (Figure~\ref{fig:seed-stability}).

\paragraph{The reference is not a trained model.}
\textsc{persistence} continues the last observed state across the whole
horizon. It has no parameters and no fitting step. A comparison among
trained models can establish which architecture is better; only a
comparison against not training can establish that the forecasting stage
belongs in the framework at all. This row was absent from the
comparison as first run, and adding it changed the conclusion.

\subsection{Results}
\label{sec:forecasting:results}

\begin{table}[t]
\centering
\caption{Forecasting comparison on $5\,449$ held-out windows of
pelletizer~I, mean\,$\pm$\,sd over seeds. Lower MAE is better. The
untrained reference is deterministic. Ordered by MAE.}
\label{tab:forecasting}
\small
\setlength{\tabcolsep}{4pt}
\begin{tabular}{@{}lrrrr@{}}
\toprule
Model & seeds & STANDBY $F_1$ & MAE (s) & params \\
\midrule
XGBoost \cite{chen2016xgboost}          & 15 & $\mathbf{0.689 \pm 0.003}$ & $\mathbf{11.50 \pm 0.11}$ & 236 trees \\
\textsc{persistence} (untrained)        &  1 & $0.681$                    & $12.95$                   & $\mathbf{0}$ \\
Transformer \cite{nie2023patchtst}      &  5 & $0.595 \pm 0.032$          & $17.45 \pm 3.13$          & $310\,488$ \\
GRU \cite{cho2014gru}                   &  5 & $0.599 \pm 0.073$          & $17.68 \pm 5.59$          & $121\,604$ \\
TCN \cite{bai2018tcn}                   &  5 & $0.588 \pm 0.050$          & $18.66 \pm 3.29$          & $131\,332$ \\
Vanilla LSTM \cite{hochreiter1997lstm}  &  5 & $0.578 \pm 0.036$          & $19.12 \pm 3.01$          & $93\,816$ \\
Seq2Seq LSTM \cite{sutskever2014seq2seq}&  5 & $0.561 \pm 0.064$          & $20.90 \pm 4.74$          & $161\,028$ \\
\bottomrule
\end{tabular}
\end{table}

Table~\ref{tab:forecasting} orders the seven models by the error metric
the decision layer would consume. Three results follow, and the first
two are not the ones this comparison was designed to produce.

\textbf{Every neural architecture loses to the untrained reference.}
Persistence reaches $12.95$\,s of per-window MAE against $17.45$--$20.90$\,s
for the five sequence models, and $F_1 = 0.681$ against
$0.561$--$0.599$. The margins are $4.5$--$8.0$\,s of MAE, all of them
above the $1$\,s threshold fixed in advance as the smallest difference
that could change a decision, and all significant at
$p_{\text{Holm}}\le6\times10^{-31}$ (Section~\ref{sec:significance}). The
sequence-to-sequence LSTM --- the architecture this framework originally
proposed for this stage --- is the worst of the seven, losing to
persistence by $7.95$\,s of MAE while costing $161$k parameters and
$865\pm75$\,s of training.

\textbf{One trained model clears the reference, and it is the tree.}
XGBoost reaches $11.50$\,s, $1.45$\,s better than persistence at
$p_{\text{Holm}}=7\times10^{-4}$: statistically detectable, above the
practical threshold, and the only comparison in the family where
training beats not training. It is also the cheapest of the trained
models at inference ($0.18$\,ms per window against $0.51$--$1.16$\,ms
for the neural models). \textbf{The gradient-boosted tree is therefore the
forecasting component of the proposed framework}, and the sequence
models are reported as what they are: compared baselines that did not
earn the stage.

\textbf{The tree is also an order of magnitude more stable.} Its $F_1$
standard deviation over $15$ seeds is $0.003$; the neural models' over
$5$ seeds is $0.032$--$0.073$, a factor of $13$--$29$
(Figure~\ref{fig:seed-stability}). The practical consequence is visible
in the same figure: seed~$0$ is the best of five draws for four of the
five neural architectures, so a single-seed protocol --- which is what
the first run of this comparison used --- reports each of them at close
to its most favourable outcome. That protocol was sound and its runs
reproduce bit-for-bit; it simply cannot support a ranking.

\begin{figure}[t]
  \centering
  \includegraphics[width=\linewidth]{fig_p6_seed_stability}
  \caption{STANDBY $F_1$ by training seed for all seven forecasters.
  The tree's $15$ seeds span $0.009$ of $F_1$; the neural
  architectures' five seeds span $0.081$--$0.166$, wider than the gap
  between architectures. The dashed line is the untrained reference.}
  \label{fig:seed-stability}
\end{figure}

\subsection{Why the untrained reference is so hard to beat}
\label{sec:forecasting:why}

The result is not an artefact of an unfair protocol, and the explanation
is available from two independent directions.

The first is the operating regime. The horizon is $300$\,s while the
median state dwell after the minimum-dwell constraint is $92.5$\,s with
a heavy tail --- STANDBY runs last minutes to hours
(Section~\ref{sec:results:task4}). ``Nothing changes in the next five
minutes'' is therefore right most of the time, and a forecaster deployed
in this regime has to beat that, not a strawman.

The second is measured on the winning model itself. Attributing
XGBoost's predictions (Section~\ref{sec:explain:window}) puts $43.5\%$
of the attribution on the final observed value of the lookback window,
against $26.4\%$ on block means and $26.3\%$ on block standard
deviations. A model whose largest single input is the state the machine
is in \emph{now} is a refinement of persistence, and cannot be expected
to beat it by much. It does not: $1.45$\,s. The margin is real, and it
is small for a reason the attribution makes explicit rather than leaving
to speculation.

\paragraph{A caveat for deployment elsewhere.}
The same attribution shows that supply voltage is the largest single
channel contributor to this model, at $20.0\%$ across its nine design
columns and $11.8\%$ on the last observed value alone. Supply voltage is
not a state variable of the machine: it sags under load and recovers
when the drive unloads, so it is a real but indirect indicator, and one
tied to this site's point of common coupling and transformer impedance.
\textbf{A deployment on another site should not expect that contribution
to transfer}, and the cross-dataset failure reported in
Section~\ref{sec:crossmachine:dataset} is consistent with it. Split gain
ranks voltage tenth of the twenty-five channels and gives no such
warning, which is one more reason the attribution is reported.

\paragraph{Scope.}
This comparison is one machine, one horizon, and one decimation.
Section~\ref{sec:crossmachine} repeats the training across all eight
machines and asks whether any of it transfers; the answer bears directly
on how much of Table~\ref{tab:forecasting} is a property of the method
and how much of pelletizer~I.

%  Figures are referenced by bare name; add to the preamble
%      \graphicspath{{../outputs/phase2/figures/pelletizer-I/}
%                    {../outputs/phase6/figures/}}
%  Bibliography: paper/references.bib
%  Numbers below are generated by
%      python -m experiments.task5_forecasting_baselines   (single seed)
%      python run_phase5.py                                (seeds)
%  and stored in
%      outputs/phase2/task5/<machine>/comparison_table.csv
%      outputs/phase5/task12/<machine>/forecasting/seed_table.csv
%  Significance testing of every contrast quoted here is
%  Section~\ref{sec:significance}; attribution is
%  Section~\ref{sec:explain:window}.
% =====================================================================

\section{Forecasting the Non-Productive Horizon}
\label{sec:forecasting}

The state layer of Section~\ref{sec:states} produces a label per second.
A shutdown decision consumes something else: how much longer the machine
will remain non-productive. This section evaluates that forecasting step
as a supervised sequence task over the inferred labels, under one
protocol, against a reference that does no training at all.

\subsection{Protocol}
\label{sec:forecasting:protocol}

The labelled record is decimated $10\times$ to $0.1$\,Hz and cut into
windows of $600$\,s lookback and $300$\,s horizon at a stride of
$120$\,s. The split is chronological ($70/15/15$) and windows are built
\emph{inside} each split, so no window straddles a split boundary and no
future sample can enter a training window. Class weights are balanced
identically for every model. The test set is $5\,449$ windows, on which
the mean true STANDBY content is $20.1$\,s of a $300$\,s horizon and
only $13.5\%$ of windows contain any STANDBY at all.

Two metrics are reported. STANDBY $F_1$ is the per-step classification
quality on the minority class the decision layer cares about. The
per-window \emph{STANDBY-seconds MAE} --- the absolute error in how many
of the next $300$\,s are predicted to be standby --- is the quantity
closer to what a duration-consuming policy would use, and it is the
metric on which the significance tests of
Section~\ref{sec:significance} are pre-registered.

\paragraph{Seeds are reported, not a single run.}
Every stochastic model is trained at multiple seeds and the table
reports mean\,$\pm$\,sd: $15$ seeds for the gradient-boosted tree, $5$
for each neural architecture, $1$ for the deterministic reference. This
is not a formality here. The spread between a good and a bad seed of one
architecture reaches $0.166$ of $F_1$, which is larger than the entire
spread between architectures, so a single-seed table of this comparison
is not a ranking (Figure~\ref{fig:seed-stability}).

\paragraph{The reference is not a trained model.}
\textsc{persistence} continues the last observed state across the whole
horizon. It has no parameters and no fitting step. A comparison among
trained models can establish which architecture is better; only a
comparison against not training can establish that the forecasting stage
belongs in the framework at all. This row was absent from the
comparison as first run, and adding it changed the conclusion.

\subsection{Results}
\label{sec:forecasting:results}

\begin{table}[t]
\centering
\caption{Forecasting comparison on $5\,449$ held-out windows of
pelletizer~I, mean\,$\pm$\,sd over seeds. Lower MAE is better. The
untrained reference is deterministic. Ordered by MAE.}
\label{tab:forecasting}
\small
\setlength{\tabcolsep}{4pt}
\begin{tabular}{@{}lrrrr@{}}
\toprule
Model & seeds & STANDBY $F_1$ & MAE (s) & params \\
\midrule
XGBoost \cite{chen2016xgboost}          & 15 & $\mathbf{0.689 \pm 0.003}$ & $\mathbf{11.50 \pm 0.11}$ & 236 trees \\
\textsc{persistence} (untrained)        &  1 & $0.681$                    & $12.95$                   & $\mathbf{0}$ \\
Transformer \cite{nie2023patchtst}      &  5 & $0.595 \pm 0.032$          & $17.45 \pm 3.13$          & $310\,488$ \\
GRU \cite{cho2014gru}                   &  5 & $0.599 \pm 0.073$          & $17.68 \pm 5.59$          & $121\,604$ \\
TCN \cite{bai2018tcn}                   &  5 & $0.588 \pm 0.050$          & $18.66 \pm 3.29$          & $131\,332$ \\
Vanilla LSTM \cite{hochreiter1997lstm}  &  5 & $0.578 \pm 0.036$          & $19.12 \pm 3.01$          & $93\,816$ \\
Seq2Seq LSTM \cite{sutskever2014seq2seq}&  5 & $0.561 \pm 0.064$          & $20.90 \pm 4.74$          & $161\,028$ \\
\bottomrule
\end{tabular}
\end{table}

Table~\ref{tab:forecasting} orders the seven models by the error metric
the decision layer would consume. Three results follow, and the first
two are not the ones this comparison was designed to produce.

\textbf{Every neural architecture loses to the untrained reference.}
Persistence reaches $12.95$\,s of per-window MAE against $17.45$--$20.90$\,s
for the five sequence models, and $F_1 = 0.681$ against
$0.561$--$0.599$. The margins are $4.5$--$8.0$\,s of MAE, all of them
above the $1$\,s threshold fixed in advance as the smallest difference
that could change a decision, and all significant at
$p_{\text{Holm}}\le6\times10^{-31}$ (Section~\ref{sec:significance}). The
sequence-to-sequence LSTM --- the architecture this framework originally
proposed for this stage --- is the worst of the seven, losing to
persistence by $7.95$\,s of MAE while costing $161$k parameters and
$865\pm75$\,s of training.

\textbf{One trained model clears the reference, and it is the tree.}
XGBoost reaches $11.50$\,s, $1.45$\,s better than persistence at
$p_{\text{Holm}}=7\times10^{-4}$: statistically detectable, above the
practical threshold, and the only comparison in the family where
training beats not training. It is also the cheapest of the trained
models at inference ($0.18$\,ms per window against $0.51$--$1.16$\,ms
for the neural models). \textbf{The gradient-boosted tree is therefore the
forecasting component of the proposed framework}, and the sequence
models are reported as what they are: compared baselines that did not
earn the stage.

\textbf{The tree is also an order of magnitude more stable.} Its $F_1$
standard deviation over $15$ seeds is $0.003$; the neural models' over
$5$ seeds is $0.032$--$0.073$, a factor of $13$--$29$
(Figure~\ref{fig:seed-stability}). The practical consequence is visible
in the same figure: seed~$0$ is the best of five draws for four of the
five neural architectures, so a single-seed protocol --- which is what
the first run of this comparison used --- reports each of them at close
to its most favourable outcome. That protocol was sound and its runs
reproduce bit-for-bit; it simply cannot support a ranking.

\begin{figure}[t]
  \centering
  \includegraphics[width=\linewidth]{fig_p6_seed_stability}
  \caption{STANDBY $F_1$ by training seed for all seven forecasters.
  The tree's $15$ seeds span $0.009$ of $F_1$; the neural
  architectures' five seeds span $0.081$--$0.166$, wider than the gap
  between architectures. The dashed line is the untrained reference.}
  \label{fig:seed-stability}
\end{figure}

\subsection{Why the untrained reference is so hard to beat}
\label{sec:forecasting:why}

The result is not an artefact of an unfair protocol, and the explanation
is available from two independent directions.

The first is the operating regime. The horizon is $300$\,s while the
median state dwell after the minimum-dwell constraint is $92.5$\,s with
a heavy tail --- STANDBY runs last minutes to hours
(Section~\ref{sec:results:task4}). ``Nothing changes in the next five
minutes'' is therefore right most of the time, and a forecaster deployed
in this regime has to beat that, not a strawman.

The second is measured on the winning model itself. Attributing
XGBoost's predictions (Section~\ref{sec:explain:window}) puts $43.5\%$
of the attribution on the final observed value of the lookback window,
against $26.4\%$ on block means and $26.3\%$ on block standard
deviations. A model whose largest single input is the state the machine
is in \emph{now} is a refinement of persistence, and cannot be expected
to beat it by much. It does not: $1.45$\,s. The margin is real, and it
is small for a reason the attribution makes explicit rather than leaving
to speculation.

\paragraph{A caveat for deployment elsewhere.}
The same attribution shows that supply voltage is the largest single
channel contributor to this model, at $20.0\%$ across its nine design
columns and $11.8\%$ on the last observed value alone. Supply voltage is
not a state variable of the machine: it sags under load and recovers
when the drive unloads, so it is a real but indirect indicator, and one
tied to this site's point of common coupling and transformer impedance.
\textbf{A deployment on another site should not expect that contribution
to transfer}, and the cross-dataset failure reported in
Section~\ref{sec:crossmachine:dataset} is consistent with it. Split gain
ranks voltage tenth of the twenty-five channels and gives no such
warning, which is one more reason the attribution is reported.

\paragraph{Scope.}
This comparison is one machine, one horizon, and one decimation.
Section~\ref{sec:crossmachine} repeats the training across all eight
machines and asks whether any of it transfers; the answer bears directly
on how much of Table~\ref{tab:forecasting} is a property of the method
and how much of pelletizer~I.

%  Figures are referenced by bare name; add to the preamble
%      \graphicspath{{../outputs/phase2/figures/pelletizer-I/}
%                    {../outputs/phase6/figures/}}
%  Bibliography: paper/references.bib
%  Numbers below are generated by
%      python -m experiments.task5_forecasting_baselines   (single seed)
%      python run_phase5.py                                (seeds)
%  and stored in
%      outputs/phase2/task5/<machine>/comparison_table.csv
%      outputs/phase5/task12/<machine>/forecasting/seed_table.csv
%  Significance testing of every contrast quoted here is
%  Section~\ref{sec:significance}; attribution is
%  Section~\ref{sec:explain:window}.
% =====================================================================

\section{Forecasting the Non-Productive Horizon}
\label{sec:forecasting}

The state layer of Section~\ref{sec:states} produces a label per second.
A shutdown decision consumes something else: how much longer the machine
will remain non-productive. This section evaluates that forecasting step
as a supervised sequence task over the inferred labels, under one
protocol, against a reference that does no training at all.

\subsection{Protocol}
\label{sec:forecasting:protocol}

The labelled record is decimated $10\times$ to $0.1$\,Hz and cut into
windows of $600$\,s lookback and $300$\,s horizon at a stride of
$120$\,s. The split is chronological ($70/15/15$) and windows are built
\emph{inside} each split, so no window straddles a split boundary and no
future sample can enter a training window. Class weights are balanced
identically for every model. The test set is $5\,449$ windows, on which
the mean true STANDBY content is $20.1$\,s of a $300$\,s horizon and
only $13.5\%$ of windows contain any STANDBY at all.

Two metrics are reported. STANDBY $F_1$ is the per-step classification
quality on the minority class the decision layer cares about. The
per-window \emph{STANDBY-seconds MAE} --- the absolute error in how many
of the next $300$\,s are predicted to be standby --- is the quantity
closer to what a duration-consuming policy would use, and it is the
metric on which the significance tests of
Section~\ref{sec:significance} are pre-registered.

\paragraph{Seeds are reported, not a single run.}
Every stochastic model is trained at multiple seeds and the table
reports mean\,$\pm$\,sd: $15$ seeds for the gradient-boosted tree, $5$
for each neural architecture, $1$ for the deterministic reference. This
is not a formality here. The spread between a good and a bad seed of one
architecture reaches $0.166$ of $F_1$, which is larger than the entire
spread between architectures, so a single-seed table of this comparison
is not a ranking (Figure~\ref{fig:seed-stability}).

\paragraph{The reference is not a trained model.}
\textsc{persistence} continues the last observed state across the whole
horizon. It has no parameters and no fitting step. A comparison among
trained models can establish which architecture is better; only a
comparison against not training can establish that the forecasting stage
belongs in the framework at all. This row was absent from the
comparison as first run, and adding it changed the conclusion.

\subsection{Results}
\label{sec:forecasting:results}

\begin{table}[t]
\centering
\caption{Forecasting comparison on $5\,449$ held-out windows of
pelletizer~I, mean\,$\pm$\,sd over seeds. Lower MAE is better. The
untrained reference is deterministic. Ordered by MAE.}
\label{tab:forecasting}
\small
\setlength{\tabcolsep}{4pt}
\begin{tabular}{@{}lrrrr@{}}
\toprule
Model & seeds & STANDBY $F_1$ & MAE (s) & params \\
\midrule
XGBoost \cite{chen2016xgboost}          & 15 & $\mathbf{0.689 \pm 0.003}$ & $\mathbf{11.50 \pm 0.11}$ & 236 trees \\
\textsc{persistence} (untrained)        &  1 & $0.681$                    & $12.95$                   & $\mathbf{0}$ \\
Transformer \cite{nie2023patchtst}      &  5 & $0.595 \pm 0.032$          & $17.45 \pm 3.13$          & $310\,488$ \\
GRU \cite{cho2014gru}                   &  5 & $0.599 \pm 0.073$          & $17.68 \pm 5.59$          & $121\,604$ \\
TCN \cite{bai2018tcn}                   &  5 & $0.588 \pm 0.050$          & $18.66 \pm 3.29$          & $131\,332$ \\
Vanilla LSTM \cite{hochreiter1997lstm}  &  5 & $0.578 \pm 0.036$          & $19.12 \pm 3.01$          & $93\,816$ \\
Seq2Seq LSTM \cite{sutskever2014seq2seq}&  5 & $0.561 \pm 0.064$          & $20.90 \pm 4.74$          & $161\,028$ \\
\bottomrule
\end{tabular}
\end{table}

Table~\ref{tab:forecasting} orders the seven models by the error metric
the decision layer would consume. Three results follow, and the first
two are not the ones this comparison was designed to produce.

\textbf{Every neural architecture loses to the untrained reference.}
Persistence reaches $12.95$\,s of per-window MAE against $17.45$--$20.90$\,s
for the five sequence models, and $F_1 = 0.681$ against
$0.561$--$0.599$. The margins are $4.5$--$8.0$\,s of MAE, all of them
above the $1$\,s threshold fixed in advance as the smallest difference
that could change a decision, and all significant at
$p_{\text{Holm}}\le6\times10^{-31}$ (Section~\ref{sec:significance}). The
sequence-to-sequence LSTM --- the architecture this framework originally
proposed for this stage --- is the worst of the seven, losing to
persistence by $7.95$\,s of MAE while costing $161$k parameters and
$865\pm75$\,s of training.

\textbf{One trained model clears the reference, and it is the tree.}
XGBoost reaches $11.50$\,s, $1.45$\,s better than persistence at
$p_{\text{Holm}}=7\times10^{-4}$: statistically detectable, above the
practical threshold, and the only comparison in the family where
training beats not training. It is also the cheapest of the trained
models at inference ($0.18$\,ms per window against $0.51$--$1.16$\,ms
for the neural models). \textbf{The gradient-boosted tree is therefore the
forecasting component of the proposed framework}, and the sequence
models are reported as what they are: compared baselines that did not
earn the stage.

\textbf{The tree is also an order of magnitude more stable.} Its $F_1$
standard deviation over $15$ seeds is $0.003$; the neural models' over
$5$ seeds is $0.032$--$0.073$, a factor of $13$--$29$
(Figure~\ref{fig:seed-stability}). The practical consequence is visible
in the same figure: seed~$0$ is the best of five draws for four of the
five neural architectures, so a single-seed protocol --- which is what
the first run of this comparison used --- reports each of them at close
to its most favourable outcome. That protocol was sound and its runs
reproduce bit-for-bit; it simply cannot support a ranking.

\begin{figure}[t]
  \centering
  \includegraphics[width=\linewidth]{fig_p6_seed_stability}
  \caption{STANDBY $F_1$ by training seed for all seven forecasters.
  The tree's $15$ seeds span $0.009$ of $F_1$; the neural
  architectures' five seeds span $0.081$--$0.166$, wider than the gap
  between architectures. The dashed line is the untrained reference.}
  \label{fig:seed-stability}
\end{figure}

\subsection{Why the untrained reference is so hard to beat}
\label{sec:forecasting:why}

The result is not an artefact of an unfair protocol, and the explanation
is available from two independent directions.

The first is the operating regime. The horizon is $300$\,s while the
median state dwell after the minimum-dwell constraint is $92.5$\,s with
a heavy tail --- STANDBY runs last minutes to hours
(Section~\ref{sec:results:task4}). ``Nothing changes in the next five
minutes'' is therefore right most of the time, and a forecaster deployed
in this regime has to beat that, not a strawman.

The second is measured on the winning model itself. Attributing
XGBoost's predictions (Section~\ref{sec:explain:window}) puts $43.5\%$
of the attribution on the final observed value of the lookback window,
against $26.4\%$ on block means and $26.3\%$ on block standard
deviations. A model whose largest single input is the state the machine
is in \emph{now} is a refinement of persistence, and cannot be expected
to beat it by much. It does not: $1.45$\,s. The margin is real, and it
is small for a reason the attribution makes explicit rather than leaving
to speculation.

\paragraph{A caveat for deployment elsewhere.}
The same attribution shows that supply voltage is the largest single
channel contributor to this model, at $20.0\%$ across its nine design
columns and $11.8\%$ on the last observed value alone. Supply voltage is
not a state variable of the machine: it sags under load and recovers
when the drive unloads, so it is a real but indirect indicator, and one
tied to this site's point of common coupling and transformer impedance.
\textbf{A deployment on another site should not expect that contribution
to transfer}, and the cross-dataset failure reported in
Section~\ref{sec:crossmachine:dataset} is consistent with it. Split gain
ranks voltage tenth of the twenty-five channels and gives no such
warning, which is one more reason the attribution is reported.

\paragraph{Scope.}
This comparison is one machine, one horizon, and one decimation.
Section~\ref{sec:crossmachine} repeats the training across all eight
machines and asks whether any of it transfers; the answer bears directly
on how much of Table~\ref{tab:forecasting} is a property of the method
and how much of pelletizer~I.


% --- Optimisation-based decision support
% =====================================================================
%  decision_layer.tex  --  PHASE III, Tasks 6 and 7
%
%  The optimisation-based shutdown decision layer: problem statement,
%  the reduction that makes it exactly solvable, the policy family,
%  and the empirical results including the economic sensitivity.
%
%  Intended placement: Section 5 (\label{sec:decision}), after the
%  forecasting section and before the sustainability analysis.
%
%  Drop-in usage: % =====================================================================
%  decision_layer.tex  --  PHASE III, Tasks 6 and 7
%
%  The optimisation-based shutdown decision layer: problem statement,
%  the reduction that makes it exactly solvable, the policy family,
%  and the empirical results including the economic sensitivity.
%
%  Intended placement: Section 5 (\label{sec:decision}), after the
%  forecasting section and before the sustainability analysis.
%
%  Drop-in usage: % =====================================================================
%  decision_layer.tex  --  PHASE III, Tasks 6 and 7
%
%  The optimisation-based shutdown decision layer: problem statement,
%  the reduction that makes it exactly solvable, the policy family,
%  and the empirical results including the economic sensitivity.
%
%  Intended placement: Section 5 (\label{sec:decision}), after the
%  forecasting section and before the sustainability analysis.
%
%  Drop-in usage: \input{decision_layer}
%  Bibliography: paper/references.bib
%  Numbers below are generated by  python run_phase3.py  and stored in
%  outputs/phase3/task6/<machine>/task6_results.json
%             and outputs/phase3/task7/<machine>/task7_results.json
% =====================================================================

\section{Optimisation-Based Shutdown Decision Support}
\label{sec:decision}

\subsection{Problem statement}
\label{sec:decision:problem}

Let $x(t)\in\{0,1\}$ indicate that the machine is energised at time $t$.
Over a horizon the plant incurs standby energy cost while idle and
energised, a restart cost each time the machine is re-energised, and a
production-delay cost because a cold machine reaches production more
slowly than a warm one. The shutdown policy is the solution of

\begin{equation}
\label{eq:objective}
\min_{x(\cdot)}\;
C_{\text{total}}
= c_{\text{sb}}\,h_{\text{standby}}
+ c_{\text{re}}\,N_{\text{restarts}}
+ c_{\text{dl}}\,h_{\text{delay}}
\end{equation}
subject to
\begin{align}
h_{\text{off}}^{(i)} &\ge h_{\min}
&&\text{(motor thermal cool-down)} \label{eq:c1}\\
N_{\text{restarts}}^{(d)} &\le N_{\max}
&&\text{(mechanical wear, per day $d$)} \label{eq:c2}\\
x(t) &\in \{0,1\}. \label{eq:c3}
\end{align}

The coefficients are
$c_{\text{sb}} = \pi\,P_{\text{sb}}$, with $\pi$ the electricity tariff
and $P_{\text{sb}}$ the mean standby power taken from the inferred
STANDBY component of Section~\ref{sec:method}; and
$c_{\text{re}} = \pi E_{r} + \ell + c_{\text{dl}} T_{d}$, with $E_{r}$
the restart energy, $\ell$ any non-energy per-cycle cost, and $T_{d}$
the restart delay. Section~\ref{sec:decision:coeff} measures $P_{\text{sb}}$,
$E_r$ and $T_d$ from the record itself.

\paragraph{The break-even duration is derived, not supplied.}
Setting the saving from de-energising an idle interval of length $D$ to
zero gives
\begin{equation}
\label{eq:breakeven}
D^{\star} \;=\; \frac{c_{\text{re}}}{c_{\text{sb}}}
\;=\; \frac{\pi E_{r} + \ell + c_{\text{dl}} T_{d}}{\pi\,P_{\text{sb}}} .
\end{equation}
This replaces the fixed threshold used in the conventional formulation,
in which $D^{\star}$ is an operator-supplied constant. Equation
\eqref{eq:breakeven} also carries a consequence that is easy to miss:
\emph{when the restart cost is purely electrical} ($\ell = 0$,
$c_{\text{dl}} = 0$) \emph{the tariff cancels}, leaving
$D^{\star} = E_{r}/P_{\text{sb}}$, and the break-even duration is
independent of electricity price. Price enters only through the
non-energy terms. We verify this numerically in
Section~\ref{sec:decision:sensitivity} and it governs how the
sensitivity analysis must be presented.

\subsection{Reduction to one decision per idle episode}
\label{sec:decision:reduction}

Although $x(t)$ is defined at every sample, $C_{\text{total}}$ changes
value only at the boundaries of an \emph{idle episode} --- a maximal
interval of non-productive operation bounded by production. Within one
episode the standby cost is linear and non-decreasing in the time spent
energised, so the saving from de-energising after $\tau$ seconds of an
episode of realised length $D$ is
\begin{equation}
\label{eq:epsave}
S(D,\tau) =
\begin{cases}
c_{\text{sb}}\,(D-\tau) - c_{\text{re}}, & \tau < D,\\
0, & \text{otherwise,}
\end{cases}
\end{equation}
which is decreasing in $\tau$. Any episode that is de-energised at all
is therefore optimally de-energised at its onset, and the problem
reduces to choosing one scalar $\tau_i$ per episode.

\begin{proposition}[Exact offline optimum]
\label{prop:oracle}
With realised durations known, the optimum of
\eqref{eq:objective}--\eqref{eq:c3} is obtained by selecting, within
each day, the $N_{\max}$ episodes of largest positive $S(D_i,0)$ among
those with $D_i \ge h_{\min}$.
\end{proposition}

\noindent
The objective is modular over episodes and the only coupling is the
per-day cardinality constraint \eqref{eq:c2}, so greedy selection is
optimal. Consequently the oracle bound reported below is the exact
solution of the stated problem rather than a heuristic for it, and no
dynamic programme over the $5.5\times10^{6}$ samples is required.

\subsection{Measuring the restart coefficients}
\label{sec:decision:coeff}

Rather than assume $E_r$ and $T_d$, we estimate them from the record.
Every resumption of production is classified by the state the machine
resumed \emph{from} --- cold (OFF) or warm (STANDBY) --- and the cost
attributable to having de-energised is the difference between the two.
Charging the entire cold-start ramp to the shutdown decision would bill
it for a ramp the machine would have paid in any case.

On pelletizer-I the machine reaches 90\% of production power in a median
$574$\,s from cold ($n=126$) against $308$\,s from warm ($n=682$), giving
\begin{equation}
T_{d} = 266\ \text{s}\qquad (\text{95\% CI } [140,\,351]),
\end{equation}
by percentile bootstrap. Energy integrated over a common $1800$\,s
window from the moment production resumes gives
$-6.21$\,kWh (95\% CI $[-7.99,\,-4.64]$; $n=109$ cold, $2526$ warm), and
a milestone-based estimator gives $-1.78$\,kWh.

Both energy estimates are negative and significantly so. This is not an
artefact: a cold start reaches production more slowly, so within any
fixed window it has produced less and therefore drawn less. The energy
difference and the delay are two views of one physical fact, and summing
a negative energy term with a positive delay term would double-count it.
We therefore report that \emph{this machine has no separately
identifiable electrical restart penalty --- its restart cost is a
production-delay cost} --- set $E_r = 0$ in the baseline, and sweep $E_r$
over the range reported in the literature in
Section~\ref{sec:decision:sensitivity}, which also covers machines for
which the penalty is real.

\subsection{Policies}
\label{sec:decision:policies}

Online, $D$ is unknown when the decision must be made. Paying a small
cost repeatedly or a large one-off cost that ends the payments is
exactly the ski-rental (spin-block) structure, and the classical
guarantee applies: de-energising once accumulated standby cost equals
the restart cost is $2$-competitive against the offline optimum for
\emph{any} duration distribution~\cite{karlin1994competitive,karlin1988competitive}.
The same structure underlies online device power-down~\cite{irani2003online}.
This --- not ``never shut down'' --- is the baseline a forecast-driven
policy must beat.

We compare: \textsc{observed}, the plant's own behaviour read off the
inferred labels; \textsc{never} and \textsc{immediate}, the two trivial
policies; \textsc{static}, the conventional rule that forecasts the
duration once at onset and compares it with $D^{\star}$;
\textsc{ski-rental}, the forecast-free elapsed-time rule; the proposed
\textsc{forecast-opt}; and \textsc{oracle} from
Proposition~\ref{prop:oracle}.

\paragraph{Proposed policy.}
Every $60$\,s while the machine remains idle, the rule re-decides using
a forecast of the \emph{remaining} idle time $R$:
\begin{equation}
\label{eq:rule}
\text{de-energise} \iff
\underbrace{c_{\text{sb}}\,\mathbb{E}[R \mid \mathcal{F}_\tau] > c_{\text{re}}}_{\text{economics}}
\;\wedge\;
\underbrace{\Pr(R \ge h_{\min} \mid \mathcal{F}_\tau) \ge 1-\alpha}_{\text{feasibility}} ,
\end{equation}
where $\mathcal{F}_\tau$ is the information available at elapsed time
$\tau$: elapsed idle time, the \emph{current} position in the factory
calendar, and the load history of the production run that just ended.

Two aspects of \eqref{eq:rule} are deliberate. First, it is
\emph{sequential} rather than a one-shot classification at onset --- the
moment when the least information is available --- so each additional
idle minute is used as evidence. Second, the two conditions are not
redundant. Because $S$ is affine in $R$ by \eqref{eq:epsave}, the
unconstrained risk-neutral rule depends on the predictive distribution
only through its \emph{mean}, and a mean regression suffices. Constraint
\eqref{eq:c1} breaks that affinity: a shutdown whose idle period proves
shorter than $h_{\min}$ is infeasible rather than merely unprofitable,
and no expectation over $R$ expresses this. Feasibility is therefore
imposed as a chance constraint~\cite{charnes1959chance} at a stated
confidence $1-\alpha = 0.90$, using a separate classifier for
$\Pr(R\ge h_{\min})$ (AUC $0.939$, Brier $0.102$). Setting $\alpha=1$
recovers the pure risk-neutral rule and isolates the constraint's
contribution.

\paragraph{The estimand matters more than the model.}
Idle durations span four orders of magnitude, so fitting
$\log(1+R)$ is the natural choice --- and it is the wrong one here.
Back-transforming a log-scale fit estimates the conditional
\emph{median}, and Duan's smearing correction~\cite{duan1983smearing}
repairs it with a single global constant that cannot absorb a
covariate-dependent bias. At episode onset the conditional median is
$7$\,s while the conditional mean exceeds $2\,000$\,s. Since
\eqref{eq:rule} consumes the mean, the estimand must be elicited by the
loss used to fit it~\cite{gneiting2011making}; we therefore fit the raw
scale under squared error and retain the log-scale model as an ablation.

\subsection{Experimental protocol}
\label{sec:decision:protocol}

Spike rows are retained: removing $19.5\%$ of a $1$\,Hz series destroys
the correspondence between row counts and seconds on which every episode
duration depends. Episodes touching one of the record's $31$
acquisition gaps are excluded ($51$ episodes, $187.9$\,h), because their
true extent is unknown and counting them would manufacture savings out
of missing data. This leaves $3\,224$ usable episodes spanning
$474.9$\,h of idle time. The split is \emph{chronological} ($70/30$):
unlike the unsupervised comparison of Section~\ref{sec:results}, a
decision policy is deployed forward in time, and a shuffled split would
let it learn the future duty cycle. The per-day restart cap is enforced
at evaluation time, identically for every policy. All annualisation uses
the $63.4$ days the record actually covers, not its $154$-day span.

\paragraph{The status quo is the baseline.}
Of the $474.9$ usable idle hours the plant already de-energises for
$388.1$, leaving $86.8$\,h energised. Savings quoted against a
``never shut down'' hypothetical would therefore be savings the plant has
already banked; every figure below is quoted against \textsc{observed}.

\paragraph{Scenarios.}
Every coefficient is measurable except the valuation of lost production,
and $D^{\star}$ is linear in it, so asserting a value would silently
determine the conclusion. Scenarios are instead indexed by where
$D^{\star}$ falls within the observed idle-duration distribution, and the
delay valuation each implies is reported as an \emph{output}
(Table~\ref{tab:decision:scenarios}).

\begin{table}[t]
\centering
\caption{Economic scenarios, indexed by the derived break-even point.
The implied production-delay valuation is an output of
Eq.~\eqref{eq:breakeven}, not an assumption.}
\label{tab:decision:scenarios}
\begin{tabular}{lrrrr}
\hline
Scenario & $D^{\star}$ & Implied $c_{\text{dl}}$ & $c_{\text{re}}$ & Episodes above \\
         & (min) & (USD/h) & (USD) & \\
\hline
S1 low       &  10 &  1.01 & 0.075 & 228 \\
S2 moderate  &  60 &  6.09 & 0.450 &  53 \\
S3 high      & 240 & 24.34 & 1.799 &  43 \\
\hline
\end{tabular}
\end{table}

\subsection{Results}
\label{sec:decision:results}

Table~\ref{tab:decision:policies} reports annual saving against the
status quo and the share of the oracle head-room captured, on $968$
held-out episodes spanning $78.3$ days.

\begin{table}[t]
\centering
\caption{Policy comparison on held-out episodes: annual saving versus
the plant's own behaviour (USD/yr) and share of the attainable head-room
(\%). The three middle rows are ablations of the proposed policy.}
\label{tab:decision:policies}
\begin{tabular}{lrrrrrr}
\hline
& \multicolumn{2}{c}{S1 (10\,min)} & \multicolumn{2}{c}{S2 (60\,min)} & \multicolumn{2}{c}{S3 (240\,min)}\\
Policy & USD/yr & \% & USD/yr & \% & USD/yr & \% \\
\hline
Never shut down            & $-215$ & $-530$ & $-133$ & $-190$ & $+163$ & $65$ \\
Observed (status quo)      & $0$ & $0$ & $0$ & $0$ & $0$ & $0$ \\
Immediate shutdown         & $-237$ & $-584$ & $-316$ & $-451$ & $-599$ & $-241$ \\
Static break-even          & $-98$ & $-242$ & $-46$ & $-65$ & $+175$ & $70$ \\
Ski-rental (no forecast)   & $-23$ & $-56$ & $+24$ & $34$ & $+106$ & $43$ \\
\quad ablation: median estimand      & $-23$ & $-56$ & $-14$ & $-20$ & $+180$ & $72$ \\
\quad ablation: no chance constraint & $-73$ & $-180$ & $+6$ & $8$ & $+173$ & $69$ \\
\textbf{Forecast optimisation}       & $\mathbf{-3}$ & $\mathbf{-8}$ & $\mathbf{+39}$ & $\mathbf{55}$ & $\mathbf{+190}$ & $\mathbf{76}$ \\
Oracle (offline optimum)   & $+41$ & $100$ & $+70$ & $100$ & $+249$ & $100$ \\
\hline
\end{tabular}
\end{table}

\paragraph{These figures are conditional on the labelling.}
Table~\ref{tab:decision:policies} is computed on the state labels of
Section~\ref{sec:states}, in which the Viterbi path is left
unconstrained. Re-deriving the same table from labels carrying the
minimum-dwell constraint of Section~\ref{sec:flicker}---the same
machine, the same policy, the same coefficients, and a test window cut
at the same instant---changes the S2 column materially: the attainable
head-room falls from $71$\,USD/yr to $38$\,USD/yr and the proposed
policy from $+41$\,USD/yr ($57\%$ of the oracle) to
$+2.9 \pm 0.6$\,USD/yr ($7.8 \pm 1.6\%$; mean and standard deviation
over ten random seeds from Phase~V Task~12C). The physical quantities
are unmoved (idle $474.9$ vs $474.1$\,h; energised idle $86.8$ vs
$87.8$\,h; standby power $3\,747$ vs $4\,208$\,W); what moves is the
\emph{partition of idle time into episodes}, which is the unit the
decision layer bids on: enforcing a minimum dwell merges $3\,224$
episodes into $1\,312$ and raises the median episode from $12$ to
$90$\,s. Sub-ten-second ``opportunities'' created by label flicker are
not shutdown opportunities, so the smaller figure is the defensible
one, and both are reported here rather than the more favourable one
alone. The seed-to-seed standard deviation ($0.6$\,USD/yr) is small
relative to the bootstrap sampling interval ($\pm 50$\,USD/yr over
18 factory days; Phase~V Task~11B), confirming that model-fitting
variance does not explain the uncertainty.

The proposed policy captures $55$--$76\%$ of the attainable head-room in
S2 and S3 and beats the forecast-free ski-rental rule in all three
scenarios. Both ablations are material: replacing the mean estimand with
the log-scale median costs $53$\,USD/yr in S2, and removing the chance
constraint costs $33$\,USD/yr in S2 while raising minimum-off violations
from $2$ to $19$ there and from $2$ to $54$ in S1. The policy does
\emph{not} beat the status quo in S1: when the
restart cost is low the winning move is to de-energise aggressively for
long blocks, which the plant already does, and forecast noise costs more
than selectivity gains.

\subsection{Economic sensitivity}
\label{sec:decision:sensitivity}

\paragraph{Break-even is tariff-invariant absent non-energy costs.}
Sweeping the tariff over a $12\times$ range ($0.04$--$0.50$\,USD/kWh)
at $E_r = 10$\,kWh with $\ell = c_{\text{dl}} = 0$ leaves $D^{\star}$
constant at $2.669$\,h, matching $E_r/P_{\text{sb}}$ to seven
significant figures; introducing a $1.00$\,USD non-energy restart cost
makes it span $3.20$--$9.34$\,h. A tariff-versus-break-even map is
therefore uninformative for a purely electrical restart cost, and we
present the plane over restart energy and delay valuation instead, with
tariff as a third swept axis.

\begin{figure*}[t]
\centering
\includegraphics[width=\textwidth]{fig_p7_economic_plane}
\caption{The economic plane at a tariff of $0.12$\,USD/kWh, over the two
coefficients that carry information once break-even is shown to be
tariff-invariant. \textbf{(A)} The derived break-even duration
$D^{\star}$, an output of the optimisation rather than an input
threshold. \textbf{(B)} Annual saving under the proposed policy and
\textbf{(C)} energy saved over the test window, both on a diverging
scale with the zero contour drawn. \textbf{Panels B and C disagree in
sign over part of the plane}: where restarts are expensive the
profitable advice is to stop de-energising, so currency and energy move
in opposite directions and neither can be reported without the other.}
\label{fig:economic-plane}
\end{figure*}

\paragraph{Regimes.}
Figure~\ref{fig:economic-plane} maps the plane.
Over $256$ combinations of tariff, restart energy and delay valuation the
proposed policy is profitable in $249$, but only $171$ also reduce
energy: where restarts are expensive the cheapest advice is to
\emph{stop} de-energising, so currency saved and energy saved carry
opposite signs, and the two must be reported together. The regime
boundaries are
\begin{itemize}
\item the policy overtakes the status quo from $D^{\star} \approx 0.27$\,h;
\item ``never shut down'' becomes optimal from $D^{\star}\approx 27.5$\,h,
      just beyond the longest observed idle episode ($24.1$\,h);
\item the advantage over ski-rental \emph{grows} with $D^{\star}$
      ($+3.1$\,USD in S2, $+17.9$\,USD in S3 over the test window).
\end{itemize}
The last point is the practical message: when waiting is cheap a clock is
almost as good as a forecast, and forecasting earns its place precisely
where restarts are expensive.

\paragraph{Payback.}
At the measured coefficients and $D^{\star}=1$\,h the policy saves
$39$\,USD/yr per machine on the unconstrained labels of
Table~\ref{tab:decision:policies} and $2.9 \pm 0.6$\,USD/yr
(mean $\pm$ sd over ten seeds; Phase~V Task~12C) under the
minimum-dwell labelling, so a $5\,000$\,USD retrofit would pay back in
$130$ years on the first figure and not at all on the second, and a
three-year payback demands a per-machine cost below $116$\,USD in the
most favourable of the two. This follows directly from the plant already de-energising
for the long gaps: only $87$\,h of energised idle remain over $63$ days
of record, some $325$\,kWh. On pelletizer-I the framework's value is
therefore in decision quality --- it captures most of the attainable
head-room while avoiding the $20$ constraint-violating shutdowns the
status quo commits, and, in the headline scenario, its $29$ loss-making
ones --- rather than in a
large absolute saving. Section~\ref{sec:multimachine} examines whether
any other machine in the facility has a larger energised-idle fraction
and so changes this conclusion; none does.

%  Bibliography: paper/references.bib
%  Numbers below are generated by  python run_phase3.py  and stored in
%  outputs/phase3/task6/<machine>/task6_results.json
%             and outputs/phase3/task7/<machine>/task7_results.json
% =====================================================================

\section{Optimisation-Based Shutdown Decision Support}
\label{sec:decision}

\subsection{Problem statement}
\label{sec:decision:problem}

Let $x(t)\in\{0,1\}$ indicate that the machine is energised at time $t$.
Over a horizon the plant incurs standby energy cost while idle and
energised, a restart cost each time the machine is re-energised, and a
production-delay cost because a cold machine reaches production more
slowly than a warm one. The shutdown policy is the solution of

\begin{equation}
\label{eq:objective}
\min_{x(\cdot)}\;
C_{\text{total}}
= c_{\text{sb}}\,h_{\text{standby}}
+ c_{\text{re}}\,N_{\text{restarts}}
+ c_{\text{dl}}\,h_{\text{delay}}
\end{equation}
subject to
\begin{align}
h_{\text{off}}^{(i)} &\ge h_{\min}
&&\text{(motor thermal cool-down)} \label{eq:c1}\\
N_{\text{restarts}}^{(d)} &\le N_{\max}
&&\text{(mechanical wear, per day $d$)} \label{eq:c2}\\
x(t) &\in \{0,1\}. \label{eq:c3}
\end{align}

The coefficients are
$c_{\text{sb}} = \pi\,P_{\text{sb}}$, with $\pi$ the electricity tariff
and $P_{\text{sb}}$ the mean standby power taken from the inferred
STANDBY component of Section~\ref{sec:method}; and
$c_{\text{re}} = \pi E_{r} + \ell + c_{\text{dl}} T_{d}$, with $E_{r}$
the restart energy, $\ell$ any non-energy per-cycle cost, and $T_{d}$
the restart delay. Section~\ref{sec:decision:coeff} measures $P_{\text{sb}}$,
$E_r$ and $T_d$ from the record itself.

\paragraph{The break-even duration is derived, not supplied.}
Setting the saving from de-energising an idle interval of length $D$ to
zero gives
\begin{equation}
\label{eq:breakeven}
D^{\star} \;=\; \frac{c_{\text{re}}}{c_{\text{sb}}}
\;=\; \frac{\pi E_{r} + \ell + c_{\text{dl}} T_{d}}{\pi\,P_{\text{sb}}} .
\end{equation}
This replaces the fixed threshold used in the conventional formulation,
in which $D^{\star}$ is an operator-supplied constant. Equation
\eqref{eq:breakeven} also carries a consequence that is easy to miss:
\emph{when the restart cost is purely electrical} ($\ell = 0$,
$c_{\text{dl}} = 0$) \emph{the tariff cancels}, leaving
$D^{\star} = E_{r}/P_{\text{sb}}$, and the break-even duration is
independent of electricity price. Price enters only through the
non-energy terms. We verify this numerically in
Section~\ref{sec:decision:sensitivity} and it governs how the
sensitivity analysis must be presented.

\subsection{Reduction to one decision per idle episode}
\label{sec:decision:reduction}

Although $x(t)$ is defined at every sample, $C_{\text{total}}$ changes
value only at the boundaries of an \emph{idle episode} --- a maximal
interval of non-productive operation bounded by production. Within one
episode the standby cost is linear and non-decreasing in the time spent
energised, so the saving from de-energising after $\tau$ seconds of an
episode of realised length $D$ is
\begin{equation}
\label{eq:epsave}
S(D,\tau) =
\begin{cases}
c_{\text{sb}}\,(D-\tau) - c_{\text{re}}, & \tau < D,\\
0, & \text{otherwise,}
\end{cases}
\end{equation}
which is decreasing in $\tau$. Any episode that is de-energised at all
is therefore optimally de-energised at its onset, and the problem
reduces to choosing one scalar $\tau_i$ per episode.

\begin{proposition}[Exact offline optimum]
\label{prop:oracle}
With realised durations known, the optimum of
\eqref{eq:objective}--\eqref{eq:c3} is obtained by selecting, within
each day, the $N_{\max}$ episodes of largest positive $S(D_i,0)$ among
those with $D_i \ge h_{\min}$.
\end{proposition}

\noindent
The objective is modular over episodes and the only coupling is the
per-day cardinality constraint \eqref{eq:c2}, so greedy selection is
optimal. Consequently the oracle bound reported below is the exact
solution of the stated problem rather than a heuristic for it, and no
dynamic programme over the $5.5\times10^{6}$ samples is required.

\subsection{Measuring the restart coefficients}
\label{sec:decision:coeff}

Rather than assume $E_r$ and $T_d$, we estimate them from the record.
Every resumption of production is classified by the state the machine
resumed \emph{from} --- cold (OFF) or warm (STANDBY) --- and the cost
attributable to having de-energised is the difference between the two.
Charging the entire cold-start ramp to the shutdown decision would bill
it for a ramp the machine would have paid in any case.

On pelletizer-I the machine reaches 90\% of production power in a median
$574$\,s from cold ($n=126$) against $308$\,s from warm ($n=682$), giving
\begin{equation}
T_{d} = 266\ \text{s}\qquad (\text{95\% CI } [140,\,351]),
\end{equation}
by percentile bootstrap. Energy integrated over a common $1800$\,s
window from the moment production resumes gives
$-6.21$\,kWh (95\% CI $[-7.99,\,-4.64]$; $n=109$ cold, $2526$ warm), and
a milestone-based estimator gives $-1.78$\,kWh.

Both energy estimates are negative and significantly so. This is not an
artefact: a cold start reaches production more slowly, so within any
fixed window it has produced less and therefore drawn less. The energy
difference and the delay are two views of one physical fact, and summing
a negative energy term with a positive delay term would double-count it.
We therefore report that \emph{this machine has no separately
identifiable electrical restart penalty --- its restart cost is a
production-delay cost} --- set $E_r = 0$ in the baseline, and sweep $E_r$
over the range reported in the literature in
Section~\ref{sec:decision:sensitivity}, which also covers machines for
which the penalty is real.

\subsection{Policies}
\label{sec:decision:policies}

Online, $D$ is unknown when the decision must be made. Paying a small
cost repeatedly or a large one-off cost that ends the payments is
exactly the ski-rental (spin-block) structure, and the classical
guarantee applies: de-energising once accumulated standby cost equals
the restart cost is $2$-competitive against the offline optimum for
\emph{any} duration distribution~\cite{karlin1994competitive,karlin1988competitive}.
The same structure underlies online device power-down~\cite{irani2003online}.
This --- not ``never shut down'' --- is the baseline a forecast-driven
policy must beat.

We compare: \textsc{observed}, the plant's own behaviour read off the
inferred labels; \textsc{never} and \textsc{immediate}, the two trivial
policies; \textsc{static}, the conventional rule that forecasts the
duration once at onset and compares it with $D^{\star}$;
\textsc{ski-rental}, the forecast-free elapsed-time rule; the proposed
\textsc{forecast-opt}; and \textsc{oracle} from
Proposition~\ref{prop:oracle}.

\paragraph{Proposed policy.}
Every $60$\,s while the machine remains idle, the rule re-decides using
a forecast of the \emph{remaining} idle time $R$:
\begin{equation}
\label{eq:rule}
\text{de-energise} \iff
\underbrace{c_{\text{sb}}\,\mathbb{E}[R \mid \mathcal{F}_\tau] > c_{\text{re}}}_{\text{economics}}
\;\wedge\;
\underbrace{\Pr(R \ge h_{\min} \mid \mathcal{F}_\tau) \ge 1-\alpha}_{\text{feasibility}} ,
\end{equation}
where $\mathcal{F}_\tau$ is the information available at elapsed time
$\tau$: elapsed idle time, the \emph{current} position in the factory
calendar, and the load history of the production run that just ended.

Two aspects of \eqref{eq:rule} are deliberate. First, it is
\emph{sequential} rather than a one-shot classification at onset --- the
moment when the least information is available --- so each additional
idle minute is used as evidence. Second, the two conditions are not
redundant. Because $S$ is affine in $R$ by \eqref{eq:epsave}, the
unconstrained risk-neutral rule depends on the predictive distribution
only through its \emph{mean}, and a mean regression suffices. Constraint
\eqref{eq:c1} breaks that affinity: a shutdown whose idle period proves
shorter than $h_{\min}$ is infeasible rather than merely unprofitable,
and no expectation over $R$ expresses this. Feasibility is therefore
imposed as a chance constraint~\cite{charnes1959chance} at a stated
confidence $1-\alpha = 0.90$, using a separate classifier for
$\Pr(R\ge h_{\min})$ (AUC $0.939$, Brier $0.102$). Setting $\alpha=1$
recovers the pure risk-neutral rule and isolates the constraint's
contribution.

\paragraph{The estimand matters more than the model.}
Idle durations span four orders of magnitude, so fitting
$\log(1+R)$ is the natural choice --- and it is the wrong one here.
Back-transforming a log-scale fit estimates the conditional
\emph{median}, and Duan's smearing correction~\cite{duan1983smearing}
repairs it with a single global constant that cannot absorb a
covariate-dependent bias. At episode onset the conditional median is
$7$\,s while the conditional mean exceeds $2\,000$\,s. Since
\eqref{eq:rule} consumes the mean, the estimand must be elicited by the
loss used to fit it~\cite{gneiting2011making}; we therefore fit the raw
scale under squared error and retain the log-scale model as an ablation.

\subsection{Experimental protocol}
\label{sec:decision:protocol}

Spike rows are retained: removing $19.5\%$ of a $1$\,Hz series destroys
the correspondence between row counts and seconds on which every episode
duration depends. Episodes touching one of the record's $31$
acquisition gaps are excluded ($51$ episodes, $187.9$\,h), because their
true extent is unknown and counting them would manufacture savings out
of missing data. This leaves $3\,224$ usable episodes spanning
$474.9$\,h of idle time. The split is \emph{chronological} ($70/30$):
unlike the unsupervised comparison of Section~\ref{sec:results}, a
decision policy is deployed forward in time, and a shuffled split would
let it learn the future duty cycle. The per-day restart cap is enforced
at evaluation time, identically for every policy. All annualisation uses
the $63.4$ days the record actually covers, not its $154$-day span.

\paragraph{The status quo is the baseline.}
Of the $474.9$ usable idle hours the plant already de-energises for
$388.1$, leaving $86.8$\,h energised. Savings quoted against a
``never shut down'' hypothetical would therefore be savings the plant has
already banked; every figure below is quoted against \textsc{observed}.

\paragraph{Scenarios.}
Every coefficient is measurable except the valuation of lost production,
and $D^{\star}$ is linear in it, so asserting a value would silently
determine the conclusion. Scenarios are instead indexed by where
$D^{\star}$ falls within the observed idle-duration distribution, and the
delay valuation each implies is reported as an \emph{output}
(Table~\ref{tab:decision:scenarios}).

\begin{table}[t]
\centering
\caption{Economic scenarios, indexed by the derived break-even point.
The implied production-delay valuation is an output of
Eq.~\eqref{eq:breakeven}, not an assumption.}
\label{tab:decision:scenarios}
\begin{tabular}{lrrrr}
\hline
Scenario & $D^{\star}$ & Implied $c_{\text{dl}}$ & $c_{\text{re}}$ & Episodes above \\
         & (min) & (USD/h) & (USD) & \\
\hline
S1 low       &  10 &  1.01 & 0.075 & 228 \\
S2 moderate  &  60 &  6.09 & 0.450 &  53 \\
S3 high      & 240 & 24.34 & 1.799 &  43 \\
\hline
\end{tabular}
\end{table}

\subsection{Results}
\label{sec:decision:results}

Table~\ref{tab:decision:policies} reports annual saving against the
status quo and the share of the oracle head-room captured, on $968$
held-out episodes spanning $78.3$ days.

\begin{table}[t]
\centering
\caption{Policy comparison on held-out episodes: annual saving versus
the plant's own behaviour (USD/yr) and share of the attainable head-room
(\%). The three middle rows are ablations of the proposed policy.}
\label{tab:decision:policies}
\begin{tabular}{lrrrrrr}
\hline
& \multicolumn{2}{c}{S1 (10\,min)} & \multicolumn{2}{c}{S2 (60\,min)} & \multicolumn{2}{c}{S3 (240\,min)}\\
Policy & USD/yr & \% & USD/yr & \% & USD/yr & \% \\
\hline
Never shut down            & $-215$ & $-530$ & $-133$ & $-190$ & $+163$ & $65$ \\
Observed (status quo)      & $0$ & $0$ & $0$ & $0$ & $0$ & $0$ \\
Immediate shutdown         & $-237$ & $-584$ & $-316$ & $-451$ & $-599$ & $-241$ \\
Static break-even          & $-98$ & $-242$ & $-46$ & $-65$ & $+175$ & $70$ \\
Ski-rental (no forecast)   & $-23$ & $-56$ & $+24$ & $34$ & $+106$ & $43$ \\
\quad ablation: median estimand      & $-23$ & $-56$ & $-14$ & $-20$ & $+180$ & $72$ \\
\quad ablation: no chance constraint & $-73$ & $-180$ & $+6$ & $8$ & $+173$ & $69$ \\
\textbf{Forecast optimisation}       & $\mathbf{-3}$ & $\mathbf{-8}$ & $\mathbf{+39}$ & $\mathbf{55}$ & $\mathbf{+190}$ & $\mathbf{76}$ \\
Oracle (offline optimum)   & $+41$ & $100$ & $+70$ & $100$ & $+249$ & $100$ \\
\hline
\end{tabular}
\end{table}

\paragraph{These figures are conditional on the labelling.}
Table~\ref{tab:decision:policies} is computed on the state labels of
Section~\ref{sec:states}, in which the Viterbi path is left
unconstrained. Re-deriving the same table from labels carrying the
minimum-dwell constraint of Section~\ref{sec:flicker}---the same
machine, the same policy, the same coefficients, and a test window cut
at the same instant---changes the S2 column materially: the attainable
head-room falls from $71$\,USD/yr to $38$\,USD/yr and the proposed
policy from $+41$\,USD/yr ($57\%$ of the oracle) to
$+2.9 \pm 0.6$\,USD/yr ($7.8 \pm 1.6\%$; mean and standard deviation
over ten random seeds from Phase~V Task~12C). The physical quantities
are unmoved (idle $474.9$ vs $474.1$\,h; energised idle $86.8$ vs
$87.8$\,h; standby power $3\,747$ vs $4\,208$\,W); what moves is the
\emph{partition of idle time into episodes}, which is the unit the
decision layer bids on: enforcing a minimum dwell merges $3\,224$
episodes into $1\,312$ and raises the median episode from $12$ to
$90$\,s. Sub-ten-second ``opportunities'' created by label flicker are
not shutdown opportunities, so the smaller figure is the defensible
one, and both are reported here rather than the more favourable one
alone. The seed-to-seed standard deviation ($0.6$\,USD/yr) is small
relative to the bootstrap sampling interval ($\pm 50$\,USD/yr over
18 factory days; Phase~V Task~11B), confirming that model-fitting
variance does not explain the uncertainty.

The proposed policy captures $55$--$76\%$ of the attainable head-room in
S2 and S3 and beats the forecast-free ski-rental rule in all three
scenarios. Both ablations are material: replacing the mean estimand with
the log-scale median costs $53$\,USD/yr in S2, and removing the chance
constraint costs $33$\,USD/yr in S2 while raising minimum-off violations
from $2$ to $19$ there and from $2$ to $54$ in S1. The policy does
\emph{not} beat the status quo in S1: when the
restart cost is low the winning move is to de-energise aggressively for
long blocks, which the plant already does, and forecast noise costs more
than selectivity gains.

\subsection{Economic sensitivity}
\label{sec:decision:sensitivity}

\paragraph{Break-even is tariff-invariant absent non-energy costs.}
Sweeping the tariff over a $12\times$ range ($0.04$--$0.50$\,USD/kWh)
at $E_r = 10$\,kWh with $\ell = c_{\text{dl}} = 0$ leaves $D^{\star}$
constant at $2.669$\,h, matching $E_r/P_{\text{sb}}$ to seven
significant figures; introducing a $1.00$\,USD non-energy restart cost
makes it span $3.20$--$9.34$\,h. A tariff-versus-break-even map is
therefore uninformative for a purely electrical restart cost, and we
present the plane over restart energy and delay valuation instead, with
tariff as a third swept axis.

\begin{figure*}[t]
\centering
\includegraphics[width=\textwidth]{fig_p7_economic_plane}
\caption{The economic plane at a tariff of $0.12$\,USD/kWh, over the two
coefficients that carry information once break-even is shown to be
tariff-invariant. \textbf{(A)} The derived break-even duration
$D^{\star}$, an output of the optimisation rather than an input
threshold. \textbf{(B)} Annual saving under the proposed policy and
\textbf{(C)} energy saved over the test window, both on a diverging
scale with the zero contour drawn. \textbf{Panels B and C disagree in
sign over part of the plane}: where restarts are expensive the
profitable advice is to stop de-energising, so currency and energy move
in opposite directions and neither can be reported without the other.}
\label{fig:economic-plane}
\end{figure*}

\paragraph{Regimes.}
Figure~\ref{fig:economic-plane} maps the plane.
Over $256$ combinations of tariff, restart energy and delay valuation the
proposed policy is profitable in $249$, but only $171$ also reduce
energy: where restarts are expensive the cheapest advice is to
\emph{stop} de-energising, so currency saved and energy saved carry
opposite signs, and the two must be reported together. The regime
boundaries are
\begin{itemize}
\item the policy overtakes the status quo from $D^{\star} \approx 0.27$\,h;
\item ``never shut down'' becomes optimal from $D^{\star}\approx 27.5$\,h,
      just beyond the longest observed idle episode ($24.1$\,h);
\item the advantage over ski-rental \emph{grows} with $D^{\star}$
      ($+3.1$\,USD in S2, $+17.9$\,USD in S3 over the test window).
\end{itemize}
The last point is the practical message: when waiting is cheap a clock is
almost as good as a forecast, and forecasting earns its place precisely
where restarts are expensive.

\paragraph{Payback.}
At the measured coefficients and $D^{\star}=1$\,h the policy saves
$39$\,USD/yr per machine on the unconstrained labels of
Table~\ref{tab:decision:policies} and $2.9 \pm 0.6$\,USD/yr
(mean $\pm$ sd over ten seeds; Phase~V Task~12C) under the
minimum-dwell labelling, so a $5\,000$\,USD retrofit would pay back in
$130$ years on the first figure and not at all on the second, and a
three-year payback demands a per-machine cost below $116$\,USD in the
most favourable of the two. This follows directly from the plant already de-energising
for the long gaps: only $87$\,h of energised idle remain over $63$ days
of record, some $325$\,kWh. On pelletizer-I the framework's value is
therefore in decision quality --- it captures most of the attainable
head-room while avoiding the $20$ constraint-violating shutdowns the
status quo commits, and, in the headline scenario, its $29$ loss-making
ones --- rather than in a
large absolute saving. Section~\ref{sec:multimachine} examines whether
any other machine in the facility has a larger energised-idle fraction
and so changes this conclusion; none does.

%  Bibliography: paper/references.bib
%  Numbers below are generated by  python run_phase3.py  and stored in
%  outputs/phase3/task6/<machine>/task6_results.json
%             and outputs/phase3/task7/<machine>/task7_results.json
% =====================================================================

\section{Optimisation-Based Shutdown Decision Support}
\label{sec:decision}

\subsection{Problem statement}
\label{sec:decision:problem}

Let $x(t)\in\{0,1\}$ indicate that the machine is energised at time $t$.
Over a horizon the plant incurs standby energy cost while idle and
energised, a restart cost each time the machine is re-energised, and a
production-delay cost because a cold machine reaches production more
slowly than a warm one. The shutdown policy is the solution of

\begin{equation}
\label{eq:objective}
\min_{x(\cdot)}\;
C_{\text{total}}
= c_{\text{sb}}\,h_{\text{standby}}
+ c_{\text{re}}\,N_{\text{restarts}}
+ c_{\text{dl}}\,h_{\text{delay}}
\end{equation}
subject to
\begin{align}
h_{\text{off}}^{(i)} &\ge h_{\min}
&&\text{(motor thermal cool-down)} \label{eq:c1}\\
N_{\text{restarts}}^{(d)} &\le N_{\max}
&&\text{(mechanical wear, per day $d$)} \label{eq:c2}\\
x(t) &\in \{0,1\}. \label{eq:c3}
\end{align}

The coefficients are
$c_{\text{sb}} = \pi\,P_{\text{sb}}$, with $\pi$ the electricity tariff
and $P_{\text{sb}}$ the mean standby power taken from the inferred
STANDBY component of Section~\ref{sec:method}; and
$c_{\text{re}} = \pi E_{r} + \ell + c_{\text{dl}} T_{d}$, with $E_{r}$
the restart energy, $\ell$ any non-energy per-cycle cost, and $T_{d}$
the restart delay. Section~\ref{sec:decision:coeff} measures $P_{\text{sb}}$,
$E_r$ and $T_d$ from the record itself.

\paragraph{The break-even duration is derived, not supplied.}
Setting the saving from de-energising an idle interval of length $D$ to
zero gives
\begin{equation}
\label{eq:breakeven}
D^{\star} \;=\; \frac{c_{\text{re}}}{c_{\text{sb}}}
\;=\; \frac{\pi E_{r} + \ell + c_{\text{dl}} T_{d}}{\pi\,P_{\text{sb}}} .
\end{equation}
This replaces the fixed threshold used in the conventional formulation,
in which $D^{\star}$ is an operator-supplied constant. Equation
\eqref{eq:breakeven} also carries a consequence that is easy to miss:
\emph{when the restart cost is purely electrical} ($\ell = 0$,
$c_{\text{dl}} = 0$) \emph{the tariff cancels}, leaving
$D^{\star} = E_{r}/P_{\text{sb}}$, and the break-even duration is
independent of electricity price. Price enters only through the
non-energy terms. We verify this numerically in
Section~\ref{sec:decision:sensitivity} and it governs how the
sensitivity analysis must be presented.

\subsection{Reduction to one decision per idle episode}
\label{sec:decision:reduction}

Although $x(t)$ is defined at every sample, $C_{\text{total}}$ changes
value only at the boundaries of an \emph{idle episode} --- a maximal
interval of non-productive operation bounded by production. Within one
episode the standby cost is linear and non-decreasing in the time spent
energised, so the saving from de-energising after $\tau$ seconds of an
episode of realised length $D$ is
\begin{equation}
\label{eq:epsave}
S(D,\tau) =
\begin{cases}
c_{\text{sb}}\,(D-\tau) - c_{\text{re}}, & \tau < D,\\
0, & \text{otherwise,}
\end{cases}
\end{equation}
which is decreasing in $\tau$. Any episode that is de-energised at all
is therefore optimally de-energised at its onset, and the problem
reduces to choosing one scalar $\tau_i$ per episode.

\begin{proposition}[Exact offline optimum]
\label{prop:oracle}
With realised durations known, the optimum of
\eqref{eq:objective}--\eqref{eq:c3} is obtained by selecting, within
each day, the $N_{\max}$ episodes of largest positive $S(D_i,0)$ among
those with $D_i \ge h_{\min}$.
\end{proposition}

\noindent
The objective is modular over episodes and the only coupling is the
per-day cardinality constraint \eqref{eq:c2}, so greedy selection is
optimal. Consequently the oracle bound reported below is the exact
solution of the stated problem rather than a heuristic for it, and no
dynamic programme over the $5.5\times10^{6}$ samples is required.

\subsection{Measuring the restart coefficients}
\label{sec:decision:coeff}

Rather than assume $E_r$ and $T_d$, we estimate them from the record.
Every resumption of production is classified by the state the machine
resumed \emph{from} --- cold (OFF) or warm (STANDBY) --- and the cost
attributable to having de-energised is the difference between the two.
Charging the entire cold-start ramp to the shutdown decision would bill
it for a ramp the machine would have paid in any case.

On pelletizer-I the machine reaches 90\% of production power in a median
$574$\,s from cold ($n=126$) against $308$\,s from warm ($n=682$), giving
\begin{equation}
T_{d} = 266\ \text{s}\qquad (\text{95\% CI } [140,\,351]),
\end{equation}
by percentile bootstrap. Energy integrated over a common $1800$\,s
window from the moment production resumes gives
$-6.21$\,kWh (95\% CI $[-7.99,\,-4.64]$; $n=109$ cold, $2526$ warm), and
a milestone-based estimator gives $-1.78$\,kWh.

Both energy estimates are negative and significantly so. This is not an
artefact: a cold start reaches production more slowly, so within any
fixed window it has produced less and therefore drawn less. The energy
difference and the delay are two views of one physical fact, and summing
a negative energy term with a positive delay term would double-count it.
We therefore report that \emph{this machine has no separately
identifiable electrical restart penalty --- its restart cost is a
production-delay cost} --- set $E_r = 0$ in the baseline, and sweep $E_r$
over the range reported in the literature in
Section~\ref{sec:decision:sensitivity}, which also covers machines for
which the penalty is real.

\subsection{Policies}
\label{sec:decision:policies}

Online, $D$ is unknown when the decision must be made. Paying a small
cost repeatedly or a large one-off cost that ends the payments is
exactly the ski-rental (spin-block) structure, and the classical
guarantee applies: de-energising once accumulated standby cost equals
the restart cost is $2$-competitive against the offline optimum for
\emph{any} duration distribution~\cite{karlin1994competitive,karlin1988competitive}.
The same structure underlies online device power-down~\cite{irani2003online}.
This --- not ``never shut down'' --- is the baseline a forecast-driven
policy must beat.

We compare: \textsc{observed}, the plant's own behaviour read off the
inferred labels; \textsc{never} and \textsc{immediate}, the two trivial
policies; \textsc{static}, the conventional rule that forecasts the
duration once at onset and compares it with $D^{\star}$;
\textsc{ski-rental}, the forecast-free elapsed-time rule; the proposed
\textsc{forecast-opt}; and \textsc{oracle} from
Proposition~\ref{prop:oracle}.

\paragraph{Proposed policy.}
Every $60$\,s while the machine remains idle, the rule re-decides using
a forecast of the \emph{remaining} idle time $R$:
\begin{equation}
\label{eq:rule}
\text{de-energise} \iff
\underbrace{c_{\text{sb}}\,\mathbb{E}[R \mid \mathcal{F}_\tau] > c_{\text{re}}}_{\text{economics}}
\;\wedge\;
\underbrace{\Pr(R \ge h_{\min} \mid \mathcal{F}_\tau) \ge 1-\alpha}_{\text{feasibility}} ,
\end{equation}
where $\mathcal{F}_\tau$ is the information available at elapsed time
$\tau$: elapsed idle time, the \emph{current} position in the factory
calendar, and the load history of the production run that just ended.

Two aspects of \eqref{eq:rule} are deliberate. First, it is
\emph{sequential} rather than a one-shot classification at onset --- the
moment when the least information is available --- so each additional
idle minute is used as evidence. Second, the two conditions are not
redundant. Because $S$ is affine in $R$ by \eqref{eq:epsave}, the
unconstrained risk-neutral rule depends on the predictive distribution
only through its \emph{mean}, and a mean regression suffices. Constraint
\eqref{eq:c1} breaks that affinity: a shutdown whose idle period proves
shorter than $h_{\min}$ is infeasible rather than merely unprofitable,
and no expectation over $R$ expresses this. Feasibility is therefore
imposed as a chance constraint~\cite{charnes1959chance} at a stated
confidence $1-\alpha = 0.90$, using a separate classifier for
$\Pr(R\ge h_{\min})$ (AUC $0.939$, Brier $0.102$). Setting $\alpha=1$
recovers the pure risk-neutral rule and isolates the constraint's
contribution.

\paragraph{The estimand matters more than the model.}
Idle durations span four orders of magnitude, so fitting
$\log(1+R)$ is the natural choice --- and it is the wrong one here.
Back-transforming a log-scale fit estimates the conditional
\emph{median}, and Duan's smearing correction~\cite{duan1983smearing}
repairs it with a single global constant that cannot absorb a
covariate-dependent bias. At episode onset the conditional median is
$7$\,s while the conditional mean exceeds $2\,000$\,s. Since
\eqref{eq:rule} consumes the mean, the estimand must be elicited by the
loss used to fit it~\cite{gneiting2011making}; we therefore fit the raw
scale under squared error and retain the log-scale model as an ablation.

\subsection{Experimental protocol}
\label{sec:decision:protocol}

Spike rows are retained: removing $19.5\%$ of a $1$\,Hz series destroys
the correspondence between row counts and seconds on which every episode
duration depends. Episodes touching one of the record's $31$
acquisition gaps are excluded ($51$ episodes, $187.9$\,h), because their
true extent is unknown and counting them would manufacture savings out
of missing data. This leaves $3\,224$ usable episodes spanning
$474.9$\,h of idle time. The split is \emph{chronological} ($70/30$):
unlike the unsupervised comparison of Section~\ref{sec:results}, a
decision policy is deployed forward in time, and a shuffled split would
let it learn the future duty cycle. The per-day restart cap is enforced
at evaluation time, identically for every policy. All annualisation uses
the $63.4$ days the record actually covers, not its $154$-day span.

\paragraph{The status quo is the baseline.}
Of the $474.9$ usable idle hours the plant already de-energises for
$388.1$, leaving $86.8$\,h energised. Savings quoted against a
``never shut down'' hypothetical would therefore be savings the plant has
already banked; every figure below is quoted against \textsc{observed}.

\paragraph{Scenarios.}
Every coefficient is measurable except the valuation of lost production,
and $D^{\star}$ is linear in it, so asserting a value would silently
determine the conclusion. Scenarios are instead indexed by where
$D^{\star}$ falls within the observed idle-duration distribution, and the
delay valuation each implies is reported as an \emph{output}
(Table~\ref{tab:decision:scenarios}).

\begin{table}[t]
\centering
\caption{Economic scenarios, indexed by the derived break-even point.
The implied production-delay valuation is an output of
Eq.~\eqref{eq:breakeven}, not an assumption.}
\label{tab:decision:scenarios}
\begin{tabular}{lrrrr}
\hline
Scenario & $D^{\star}$ & Implied $c_{\text{dl}}$ & $c_{\text{re}}$ & Episodes above \\
         & (min) & (USD/h) & (USD) & \\
\hline
S1 low       &  10 &  1.01 & 0.075 & 228 \\
S2 moderate  &  60 &  6.09 & 0.450 &  53 \\
S3 high      & 240 & 24.34 & 1.799 &  43 \\
\hline
\end{tabular}
\end{table}

\subsection{Results}
\label{sec:decision:results}

Table~\ref{tab:decision:policies} reports annual saving against the
status quo and the share of the oracle head-room captured, on $968$
held-out episodes spanning $78.3$ days.

\begin{table}[t]
\centering
\caption{Policy comparison on held-out episodes: annual saving versus
the plant's own behaviour (USD/yr) and share of the attainable head-room
(\%). The three middle rows are ablations of the proposed policy.}
\label{tab:decision:policies}
\begin{tabular}{lrrrrrr}
\hline
& \multicolumn{2}{c}{S1 (10\,min)} & \multicolumn{2}{c}{S2 (60\,min)} & \multicolumn{2}{c}{S3 (240\,min)}\\
Policy & USD/yr & \% & USD/yr & \% & USD/yr & \% \\
\hline
Never shut down            & $-215$ & $-530$ & $-133$ & $-190$ & $+163$ & $65$ \\
Observed (status quo)      & $0$ & $0$ & $0$ & $0$ & $0$ & $0$ \\
Immediate shutdown         & $-237$ & $-584$ & $-316$ & $-451$ & $-599$ & $-241$ \\
Static break-even          & $-98$ & $-242$ & $-46$ & $-65$ & $+175$ & $70$ \\
Ski-rental (no forecast)   & $-23$ & $-56$ & $+24$ & $34$ & $+106$ & $43$ \\
\quad ablation: median estimand      & $-23$ & $-56$ & $-14$ & $-20$ & $+180$ & $72$ \\
\quad ablation: no chance constraint & $-73$ & $-180$ & $+6$ & $8$ & $+173$ & $69$ \\
\textbf{Forecast optimisation}       & $\mathbf{-3}$ & $\mathbf{-8}$ & $\mathbf{+39}$ & $\mathbf{55}$ & $\mathbf{+190}$ & $\mathbf{76}$ \\
Oracle (offline optimum)   & $+41$ & $100$ & $+70$ & $100$ & $+249$ & $100$ \\
\hline
\end{tabular}
\end{table}

\paragraph{These figures are conditional on the labelling.}
Table~\ref{tab:decision:policies} is computed on the state labels of
Section~\ref{sec:states}, in which the Viterbi path is left
unconstrained. Re-deriving the same table from labels carrying the
minimum-dwell constraint of Section~\ref{sec:flicker}---the same
machine, the same policy, the same coefficients, and a test window cut
at the same instant---changes the S2 column materially: the attainable
head-room falls from $71$\,USD/yr to $38$\,USD/yr and the proposed
policy from $+41$\,USD/yr ($57\%$ of the oracle) to
$+2.9 \pm 0.6$\,USD/yr ($7.8 \pm 1.6\%$; mean and standard deviation
over ten random seeds from Phase~V Task~12C). The physical quantities
are unmoved (idle $474.9$ vs $474.1$\,h; energised idle $86.8$ vs
$87.8$\,h; standby power $3\,747$ vs $4\,208$\,W); what moves is the
\emph{partition of idle time into episodes}, which is the unit the
decision layer bids on: enforcing a minimum dwell merges $3\,224$
episodes into $1\,312$ and raises the median episode from $12$ to
$90$\,s. Sub-ten-second ``opportunities'' created by label flicker are
not shutdown opportunities, so the smaller figure is the defensible
one, and both are reported here rather than the more favourable one
alone. The seed-to-seed standard deviation ($0.6$\,USD/yr) is small
relative to the bootstrap sampling interval ($\pm 50$\,USD/yr over
18 factory days; Phase~V Task~11B), confirming that model-fitting
variance does not explain the uncertainty.

The proposed policy captures $55$--$76\%$ of the attainable head-room in
S2 and S3 and beats the forecast-free ski-rental rule in all three
scenarios. Both ablations are material: replacing the mean estimand with
the log-scale median costs $53$\,USD/yr in S2, and removing the chance
constraint costs $33$\,USD/yr in S2 while raising minimum-off violations
from $2$ to $19$ there and from $2$ to $54$ in S1. The policy does
\emph{not} beat the status quo in S1: when the
restart cost is low the winning move is to de-energise aggressively for
long blocks, which the plant already does, and forecast noise costs more
than selectivity gains.

\subsection{Economic sensitivity}
\label{sec:decision:sensitivity}

\paragraph{Break-even is tariff-invariant absent non-energy costs.}
Sweeping the tariff over a $12\times$ range ($0.04$--$0.50$\,USD/kWh)
at $E_r = 10$\,kWh with $\ell = c_{\text{dl}} = 0$ leaves $D^{\star}$
constant at $2.669$\,h, matching $E_r/P_{\text{sb}}$ to seven
significant figures; introducing a $1.00$\,USD non-energy restart cost
makes it span $3.20$--$9.34$\,h. A tariff-versus-break-even map is
therefore uninformative for a purely electrical restart cost, and we
present the plane over restart energy and delay valuation instead, with
tariff as a third swept axis.

\begin{figure*}[t]
\centering
\includegraphics[width=\textwidth]{fig_p7_economic_plane}
\caption{The economic plane at a tariff of $0.12$\,USD/kWh, over the two
coefficients that carry information once break-even is shown to be
tariff-invariant. \textbf{(A)} The derived break-even duration
$D^{\star}$, an output of the optimisation rather than an input
threshold. \textbf{(B)} Annual saving under the proposed policy and
\textbf{(C)} energy saved over the test window, both on a diverging
scale with the zero contour drawn. \textbf{Panels B and C disagree in
sign over part of the plane}: where restarts are expensive the
profitable advice is to stop de-energising, so currency and energy move
in opposite directions and neither can be reported without the other.}
\label{fig:economic-plane}
\end{figure*}

\paragraph{Regimes.}
Figure~\ref{fig:economic-plane} maps the plane.
Over $256$ combinations of tariff, restart energy and delay valuation the
proposed policy is profitable in $249$, but only $171$ also reduce
energy: where restarts are expensive the cheapest advice is to
\emph{stop} de-energising, so currency saved and energy saved carry
opposite signs, and the two must be reported together. The regime
boundaries are
\begin{itemize}
\item the policy overtakes the status quo from $D^{\star} \approx 0.27$\,h;
\item ``never shut down'' becomes optimal from $D^{\star}\approx 27.5$\,h,
      just beyond the longest observed idle episode ($24.1$\,h);
\item the advantage over ski-rental \emph{grows} with $D^{\star}$
      ($+3.1$\,USD in S2, $+17.9$\,USD in S3 over the test window).
\end{itemize}
The last point is the practical message: when waiting is cheap a clock is
almost as good as a forecast, and forecasting earns its place precisely
where restarts are expensive.

\paragraph{Payback.}
At the measured coefficients and $D^{\star}=1$\,h the policy saves
$39$\,USD/yr per machine on the unconstrained labels of
Table~\ref{tab:decision:policies} and $2.9 \pm 0.6$\,USD/yr
(mean $\pm$ sd over ten seeds; Phase~V Task~12C) under the
minimum-dwell labelling, so a $5\,000$\,USD retrofit would pay back in
$130$ years on the first figure and not at all on the second, and a
three-year payback demands a per-machine cost below $116$\,USD in the
most favourable of the two. This follows directly from the plant already de-energising
for the long gaps: only $87$\,h of energised idle remain over $63$ days
of record, some $325$\,kWh. On pelletizer-I the framework's value is
therefore in decision quality --- it captures most of the attainable
head-room while avoiding the $20$ constraint-violating shutdowns the
status quo commits, and, in the headline scenario, its $29$ loss-making
ones --- rather than in a
large absolute saving. Section~\ref{sec:multimachine} examines whether
any other machine in the facility has a larger energised-idle fraction
and so changes this conclusion; none does.


% --- Significance testing and repeated runs
% =====================================================================
%  significance.tex  --  PHASE V, Tasks 11 and 12
%
%  Significance testing and repeated runs: what the resampling unit is,
%  which differences survive, which designs cannot detect anything, and
%  what ten labelling seeds do to the headline state-time total.
%
%  Intended placement: Section 7 (\label{sec:significance}), AFTER both
%  result sections it tests -- the forecasting comparison and the
%  decision layer -- so that a reader has seen both tables before being
%  shown what survives resampling.  The decision layer's own intervals
%  are reported here rather than in decision_layer.tex, because the
%  resampling design they depend on is stated here.
%
%  Drop-in usage: % =====================================================================
%  significance.tex  --  PHASE V, Tasks 11 and 12
%
%  Significance testing and repeated runs: what the resampling unit is,
%  which differences survive, which designs cannot detect anything, and
%  what ten labelling seeds do to the headline state-time total.
%
%  Intended placement: Section 7 (\label{sec:significance}), AFTER both
%  result sections it tests -- the forecasting comparison and the
%  decision layer -- so that a reader has seen both tables before being
%  shown what survives resampling.  The decision layer's own intervals
%  are reported here rather than in decision_layer.tex, because the
%  resampling design they depend on is stated here.
%
%  Drop-in usage: % =====================================================================
%  significance.tex  --  PHASE V, Tasks 11 and 12
%
%  Significance testing and repeated runs: what the resampling unit is,
%  which differences survive, which designs cannot detect anything, and
%  what ten labelling seeds do to the headline state-time total.
%
%  Intended placement: Section 7 (\label{sec:significance}), AFTER both
%  result sections it tests -- the forecasting comparison and the
%  decision layer -- so that a reader has seen both tables before being
%  shown what survives resampling.  The decision layer's own intervals
%  are reported here rather than in decision_layer.tex, because the
%  resampling design they depend on is stated here.
%
%  Drop-in usage: \input{significance}
%  Figures are referenced by bare name; add to the preamble
%      \graphicspath{{../outputs/phase6/figures/}}
%  Bibliography: paper/references.bib
%  Numbers below are generated by  python run_phase5.py  and stored in
%    outputs/phase5/task11/pelletizer-I/forecasting_tests.json
%    outputs/phase5/task11/pelletizer-I/decision_bootstrap_S*.json
%    outputs/phase5/task11/cross_machine_tests.json
%    outputs/phase5/task11/state_block_bootstrap.csv
%    outputs/phase5/task12/pelletizer-I/{forecasting,labelling,decision}/
% =====================================================================

\section{Statistical Analysis}
\label{sec:significance}

Every headline quantity in the preceding sections is a point estimate.
This section reports what happens to those estimates under resampling
and under repeated fitting, and --- for two of the designs used here ---
what they are incapable of detecting in the first place.

\subsection{Four decisions that determine every interval below}
\label{sec:significance:method}

\textbf{The resampling unit is never a row.} A $1$\,Hz industrial record
is massively autocorrelated, and an i.i.d.\ bootstrap over rows would
treat $5.5$ million readings as $5.5$ million independent observations.
Measured on a synthetic AR(1) series with $\rho=0.999$, the i.i.d.\
bootstrap returns a standard error $25\times$ too small. Three
dependence-aware units are used instead, each chosen from the structure
of the quantity being estimated:

\begin{center}
\small
\begin{tabular}{@{}lll@{}}
\toprule
quantity & unit & why \\
\midrule
state-time totals & one-week moving blocks~\cite{kunsch1989jackknife}
  & duty cycle closes over a week \\
decision savings & factory days
  & the restart cap is per calendar day \\
forecast errors & test windows, \emph{paired}
  & one problem solved twice \\
\bottomrule
\end{tabular}
\end{center}

\textbf{A day drawn twice counts as two days.} In the decision bootstrap
the resampled copies are re-keyed. Leaving the label unchanged would let
two copies of one day share a single day's restart budget, silently
tightening the very constraint under evaluation.

\textbf{Every policy is scored on the same resample.} Drawing
independent resamples per policy would break the pairing and inflate the
interval of every contrast --- and the contrasts, not the individual
intervals, are what the decision claim rests on
(Section~\ref{sec:significance:decision}).

\textbf{Statistical and economic significance are separate columns.} A
paired test over $5\,449$ windows detects differences far below anything
that changes a decision, so every comparison carries a $p$-value, an
effect size, and a threshold fixed before the tests were run: $1$\,s of
per-window MAE, $0.02$ of $F_1$, $5$\,USD/yr. A result is allowed to be
\emph{detectable but immaterial}, and two of the twenty-one forecasting
contrasts are. Multiplicity is controlled by Holm's step-down
procedure~\cite{holm1979simple}, which achieves the same family-wise
error rate as Bonferroni and is uniformly more powerful;
Benjamini--Hochberg~\cite{benjamini1995controlling} is reported beside
it for the exploratory families, and both are shown next to the raw
$p$-value so the cost of the correction stays visible.

\subsection{Forecasting}
\label{sec:significance:forecast}

Seven models, $5\,449$ identical held-out windows. Every stored
prediction array was re-scored and matched against its recorded metrics
before any test ran, and all seven sit on a byte-identical vector of
targets, which is what makes the pairing legitimate.

\paragraph{The pre-registered test.}
The comparison named in advance was the proposed Seq2Seq LSTM against
the untrained reference. It resolves decisively, and in the direction
opposite to the one it was written to check
(Table~\ref{tab:sig:headline}). The design was set up to ask whether a
small \emph{positive} gap was real; there is no positive gap. The $F_1$
interval is computed on seed~$0$ of each model, which is the
\emph{best} of the sequence model's five seeds --- and even granting it
its most favourable draw the interval excludes zero in persistence's
favour. On the seed mean the gap is five times wider.

\begin{table}[t]
\centering
\caption{The pre-registered comparison: proposed sequence model versus
the untrained reference, on $5\,449$ paired windows.}
\label{tab:sig:headline}
\small
\begin{tabular}{@{}lrrr@{}}
\toprule
quantity & Seq2Seq & \textsc{persistence} & $p$ \\
\midrule
per-window MAE (s)        & $20.90$ & $12.95$ & $5.7\times10^{-44}$ \\
horizon-step accuracy     & \multicolumn{2}{c}{$-0.0574$ (worse)} & $1.5\times10^{-113}$ \\
STANDBY $F_1$, seed mean  & $0.561$ & $0.681$ & --- \\
STANDBY $F_1$, seed 0     & $0.660$ & $0.681$ & CI $[-0.042,-0.003]$ \\
\bottomrule
\end{tabular}
\end{table}

\paragraph{The family.}
Of the $21$ pairwise contrasts, $19$ are significant after Holm
correction and $17$ of those also clear the $1$\,s practical threshold.
Two conclusions survive both filters. \emph{Persistence beats all five
neural models}, at $p_{\text{Holm}}\le6\times10^{-31}$ and by $4.5$--$8.0$\,s of
MAE. \emph{XGBoost beats persistence}, by $1.45$\,s at
$p_{\text{Holm}}=7.0\times10^{-4}$ with rank-biserial correlation
$0.24$ --- the single contrast in the entire family in which training
beats not training.

\paragraph{Friedman over seeds, and how little it can see.}
Blocking on the five seeds and treating the six trained models as
treatments gives $\chi^2=12.77$, $p=0.026$, Kendall's $W=0.511$
\cite{friedman1937use,demsar2006statistical}. XGBoost takes mean rank
$1.00$ --- it is the best model on every seed block --- and Seq2Seq
takes $5.00$, the worst. But the Nemenyi critical difference is $3.37$
on a $1$--$6$ rank scale, so \textbf{the only pair this diagram
separates is XGBoost against Seq2Seq} ($4.00$ apart) and everything else
is joined (Figure~\ref{fig:cd}). Measured power at this design is $0.14$
for a $1$-sd effect and $0.53$ at $2$\,sd. The seed-blocked comparison
is nearly blind, and the confirmatory conclusions above do not rest on
it: those are paired over $5\,449$ windows, not over $5$ seeds.

\begin{figure}[t]
  \centering
  \includegraphics[width=\linewidth]{fig_p6_cd_diagram}
  \caption{Critical-difference diagram over five seed blocks. Models
  joined by a bar are not separable at $\alpha=0.05$. The bar spans
  almost the whole scale because $n=5$ blocks buys a critical difference
  of $3.37$ ranks.}
  \label{fig:cd}
\end{figure}

\subsection{The decision layer}
\label{sec:significance:decision}

Day-block bootstrap, $2\,000$ draws, on the minimum-dwell labelling:
$394$ held-out episodes over \textbf{18 factory days} spanning $73.6$
calendar days.

\begin{table}[t]
\centering
\caption{Annual saving against the plant's own behaviour (USD/yr per
machine), with $95\%$ day-block bootstrap intervals. Every policy is
scored on the same resample.}
\label{tab:sig:decision}
\small
\setlength{\tabcolsep}{3pt}
\begin{tabular}{@{}lrrr@{}}
\toprule
Policy & S1 ($10$\,min) & S2 ($60$\,min) & S3 ($240$\,min) \\
\midrule
never shut down   & $-219.8$\,{\footnotesize$[-334,-132]$} & $-157.1$\,{\footnotesize$[-276,-65]$} & $+68.6$\,{\footnotesize$[-84,+213]$} \\
immediate         & $-219.7$\,{\footnotesize$[-338,-126]$} & $-307.5$\,{\footnotesize$[-430,-212]$} & $-623.4$\,{\footnotesize$[-774,-478]$} \\
static break-even & $-166.7$\,{\footnotesize$[-290,-71]$}  & $-66.0$\,{\footnotesize$[-196,+24]$}  & $+6.3$\,{\footnotesize$[-170,+168]$} \\
ski-rental        & $-46.7$\,{\footnotesize$[-93,-7]$}     & $+0.0$\,{\footnotesize$[-25,+29]$}    & $+30.4$\,{\footnotesize$[-79,+155]$} \\
\textbf{proposed} & $\mathbf{-28.5}$\,{\footnotesize$[-79,+1]$} & $\mathbf{+2.9}$\,{\footnotesize$[-56,+46]$} & $\mathbf{+94.5}$\,{\footnotesize$[-58,+237]$} \\
oracle            & $+25.5$\,{\footnotesize$[+15,+38]$}    & $+37.6$\,{\footnotesize$[+17,+61]$}   & $+160.8$\,{\footnotesize$[+65,+266]$} \\
\bottomrule
\end{tabular}
\end{table}

Table~\ref{tab:sig:decision} has to be read twice, and the two readings
give different answers.

\textbf{Read row by row, no policy but the oracle demonstrates a
saving.} Every other interval against the status quo spans zero,
including the proposed policy's headline $+2.9$ USD/yr inside
$[-55.9,+46.2]$. The point estimate the earlier sections quote is
indistinguishable from zero, and always was; the bootstrap supplies the
interval that makes it visible.

\textbf{Read as paired contrasts, one result survives}
(Table~\ref{tab:sig:contrasts}). Because every policy is scored on the
same resample, a difference between two rows is far better determined
than either row: the pairing removes the day-to-day variance that makes
the individual intervals wide. The proposed policy beats the static
break-even rule it was built to replace by $+68.9$ USD/yr
($p=0.001$, interval excluding zero, above the $5$ USD/yr economic
threshold) in S2 and by $+138$ and $+88$ USD/yr in S1 and S3. It does
\emph{not} beat forecast-free ski-rental, and it cannot be shown to beat
the plant's own behaviour.

\begin{table}[t]
\centering
\caption{Paired contrasts, same resample. The claim the decision layer
supports is the third row, and only the third row.}
\label{tab:sig:contrasts}
\small
\begin{tabular}{@{}lrrr@{}}
\toprule
Contrast & S1 & S2 & S3 \\
\midrule
proposed $-$ status quo   & $-28.5$ & $+2.9$ & $+94.5$ \\
                          & \footnotesize$p=0.072$ & \footnotesize$p=0.828$ & \footnotesize$p=0.196$ \\
proposed $-$ ski-rental   & $+18.2$ & $+2.8$ & $+64.1$ \\
                          & \footnotesize$p=0.553$ & \footnotesize$p=0.729$ & \footnotesize$p=0.212$ \\
\textbf{proposed $-$ static} & $\mathbf{+138.2}$ & $\mathbf{+68.9}$ & $\mathbf{+88.3}$ \\
                          & \footnotesize$\mathbf{p<0.001}$ & \footnotesize$\mathbf{p=0.001}$ & \footnotesize$\mathbf{p<0.001}$ \\
proposed $-$ oracle       & $-54.0$ & $-34.8$ & $-66.3$ \\
                          & \footnotesize$p<0.001$ & \footnotesize$p<0.001$ & \footnotesize$p<0.001$ \\
\bottomrule
\end{tabular}
\end{table}

\textbf{The oracle gap is also real}, in all three scenarios: the
head-room is genuine and mostly uncaptured --- the proposed policy takes
$7.6\%$ of it in S2.

\textbf{The width is the sample, not the estimator.} Eighteen factory
days is the binding constraint: the day-block bootstrap has eighteen
units to resample, and no refinement of the estimator narrows these
intervals. Any claim of economic benefit from this layer needs a longer
record, and this paper does not make one.
Figure~\ref{fig:decision-forest} shows all three scenarios at once,
with the intervals that cross zero drawn so that they cannot be
mistaken for the ones that do not.

\begin{figure}[t]
  \centering
  \includegraphics[width=\linewidth]{fig_p6_decision_forest}
  \caption{Decision-layer intervals in the three scenarios. Intervals
  crossing zero are drawn grey and hollow; only the oracle's excludes
  zero in every scenario.}
  \label{fig:decision-forest}
\end{figure}

\subsection{Cross-machine tests, and a design that cannot detect
anything}
\label{sec:significance:cross}

Three paired comparisons across the eight machines are reported in
Table~\ref{tab:sig:cross}, and the third column of that table is as
important as the second.

\begin{table}[t]
\centering
\caption{Paired signed-rank tests across machines. The $p$-value floor
is the smallest two-sided value the test can return at that $n$.}
\label{tab:sig:cross}
\small
\setlength{\tabcolsep}{4pt}
\begin{tabular}{@{}lrrrrr@{}}
\toprule
Comparison & $n$ & $\Delta F_1$ & $p$ & floor & power \\
\midrule
own vs \textsc{persistence}, all      & 8 & $+0.194$ & $0.016$ & $0.008$ & $0.41$ \\
own vs \textsc{persistence}, motors   & 4 & $+0.048$ & $0.250$ & $\mathbf{0.125}$ & $\mathbf{0.00}$ \\
local vs transferred                  & 8 & $+0.384$ & $\mathbf{0.008}$ & $0.008$ & $d=1.75$ \\
\bottomrule
\end{tabular}
\end{table}

\textbf{Training beats not training across the eight machines, but the
effect lives in the auxiliary plant.} The $+0.194$ mean difference is
dominated by exhaust-fan~I ($0.982$ against $0.288$) and dpc~I ($0.402$
against $0.012$), machines whose ``STANDBY'' cluster is really their off
state and is therefore trivially predictable once a model has seen it
(Section~\ref{sec:multimachine}). On the four motor-driven machines, the
ones on which the state model is physically valid, the difference is
$+0.048$ and cannot be tested.

\textbf{``Cannot be tested'' is meant literally.} With four paired
observations the smallest two-sided $p$-value the signed-rank test can
return is $2/2^{4}=0.125$, which is above $\alpha=0.05$. No data, and no
effect size however large, can produce a significant result at $n=4$.
\textbf{A non-significant verdict here is a statement about the design,
not about the models}, and it is reported in exactly those terms rather
than as ``no significant difference''.
Figure~\ref{fig:power} puts numbers on both designs: at $n=8$ the paired
test reaches $80\%$ power only above $\approx1.4$\,sd; at $n=4$ it never
does, at any effect size.

\textbf{The transfer degradation is real.} Locally fitted sequence
models beat transferred ones by $0.384$ of $F_1$ with $p$ at the floor
and $d=+1.75$. This is the one cross-machine claim of
Section~\ref{sec:crossmachine} that survives testing outright.

\begin{figure}[t]
  \centering
  \includegraphics[width=\linewidth]{fig_p6_power_curve}
  \caption{Power of the paired signed-rank test against effect size for
  $n=4$ and $n=8$. The marker is the effect actually observed on the
  four motor-driven machines ($d=0.56$), whose measured power is $0.00$
  --- readable off the axis rather than taken on trust.}
  \label{fig:power}
\end{figure}

\subsection{Repeated fitting, and the result that matters most}
\label{sec:significance:seeds}

\subsubsection{The labelling is not reproducible from the seed alone}

Re-fitting the state model at ten seeds on one deterministic
preprocessing of pelletizer~I gives a STANDBY total of
$84.6\pm28.5$\,h --- a coefficient of variation of $\mathbf{33.7\%}$ on
the headline state-time quantity of the paper. Three seeds in ten
converge on a cluster that is not standby at all:

\begin{center}
\small
\begin{tabular}{@{}rrrrl@{}}
\toprule
seed & physics & STANDBY & mean power & what it found \\
\midrule
6 & $0.80$ & $121.6$\,h & $12\,961$\,W & a loaded running state \\
8 & $0.60$ & $63.2$\,h  & $0.8$\,W     & the OFF state \\
9 & $0.60$ & $13.8$\,h  & $1.1$\,W     & the OFF state \\
\bottomrule
\end{tabular}
\end{center}

A ``standby'' cluster at $0.8$\,W drawing $0.06\%$ of rated current is
the machine switched off; one at $13$\,kW is the machine working. EM
initialisation decides which, and in $30\%$ of initialisations it
decides wrongly.

\textbf{What rescues the result is the physics battery, and this is the
constructive finding.} The tests of
Section~\ref{sec:method:accept_reject} score exactly those three seeds
at $0.60$--$0.80$ and the other seven at $1.00$, so the accept/reject
decision is exact on this record. Conditioning on a passing score
collapses the spread:

\begin{center}
\small
\begin{tabular}{@{}lrr@{}}
\toprule
 & all $10$ seeds & the $7$ that pass \\
\midrule
STANDBY hours   & $84.6\pm28.5$\,h & $\mathbf{92.4\pm2.8}$\,h \\
coefficient of variation & $\mathbf{33.7\%}$ & $\mathbf{3.0\%}$ \\
range           & $13.8$--$121.6$\,h & $87.3$--$94.2$\,h \\
STANDBY power   & $4\,570\pm3\,543$\,W & $4\,676\pm343$\,W \\
\bottomrule
\end{tabular}
\end{center}

\textbf{The pipeline is reproducible because of the physics validation,
not despite the clustering.} That is a stronger statement of this
paper's contribution than an unconditional total would be, and it is why
the accept/reject loop is part of the method in
Section~\ref{sec:method:accept_reject} rather than an appendix to it.
Flicker is $0.00\%$ at every seed and the adjusted Rand index against
seed~$0$ stays $\ge0.889$: the \emph{partition} is stable because OFF
and PEAK\_LOAD dominate the record, and it is the identity of the
STANDBY cluster specifically that moves.

\subsubsection{Model-fitting variance is not what makes the economics
uncertain}

Ten seeds of the identical decision configuration (S2, minimum-dwell
labels) give a proposed-policy saving of
$\mathbf{+2.92\pm0.59}$\,USD/yr, $7.8\pm1.6\%$ of oracle, on head-room
of $37.63$\,USD/yr with \emph{no} seed variance at all --- head-room
being a property of the episodes and the coefficients rather than of the
fit. The forecaster's own error is $8\,954\pm156$\,s, about $2.5$\,h on
episode durations, which is why the policy captures under $8\%$ of what
is available.

Set the two sources of uncertainty side by side: model fitting
contributes $\pm0.6$\,USD/yr, day-to-day sampling contributes
$\pm50$\,USD/yr. \textbf{The uncertainty in the economic claim is the
18-day sample, not the model}, and no amount of retraining addresses it.

\subsubsection{State-time totals}

The one-week moving-block bootstrap gives intervals of $\pm15$--$25\%$
of the point estimate on the five-month records
(Table~\ref{tab:multimachine}, per machine). Combined with the $\pm3\%$
conditional labelling spread above, this fixes the precision at which
every state-time number in this paper is quoted: \textbf{two significant
figures, and conditional on the physics battery passing}. Three
significant figures assert a precision neither interval supports.

%  Figures are referenced by bare name; add to the preamble
%      \graphicspath{{../outputs/phase6/figures/}}
%  Bibliography: paper/references.bib
%  Numbers below are generated by  python run_phase5.py  and stored in
%    outputs/phase5/task11/pelletizer-I/forecasting_tests.json
%    outputs/phase5/task11/pelletizer-I/decision_bootstrap_S*.json
%    outputs/phase5/task11/cross_machine_tests.json
%    outputs/phase5/task11/state_block_bootstrap.csv
%    outputs/phase5/task12/pelletizer-I/{forecasting,labelling,decision}/
% =====================================================================

\section{Statistical Analysis}
\label{sec:significance}

Every headline quantity in the preceding sections is a point estimate.
This section reports what happens to those estimates under resampling
and under repeated fitting, and --- for two of the designs used here ---
what they are incapable of detecting in the first place.

\subsection{Four decisions that determine every interval below}
\label{sec:significance:method}

\textbf{The resampling unit is never a row.} A $1$\,Hz industrial record
is massively autocorrelated, and an i.i.d.\ bootstrap over rows would
treat $5.5$ million readings as $5.5$ million independent observations.
Measured on a synthetic AR(1) series with $\rho=0.999$, the i.i.d.\
bootstrap returns a standard error $25\times$ too small. Three
dependence-aware units are used instead, each chosen from the structure
of the quantity being estimated:

\begin{center}
\small
\begin{tabular}{@{}lll@{}}
\toprule
quantity & unit & why \\
\midrule
state-time totals & one-week moving blocks~\cite{kunsch1989jackknife}
  & duty cycle closes over a week \\
decision savings & factory days
  & the restart cap is per calendar day \\
forecast errors & test windows, \emph{paired}
  & one problem solved twice \\
\bottomrule
\end{tabular}
\end{center}

\textbf{A day drawn twice counts as two days.} In the decision bootstrap
the resampled copies are re-keyed. Leaving the label unchanged would let
two copies of one day share a single day's restart budget, silently
tightening the very constraint under evaluation.

\textbf{Every policy is scored on the same resample.} Drawing
independent resamples per policy would break the pairing and inflate the
interval of every contrast --- and the contrasts, not the individual
intervals, are what the decision claim rests on
(Section~\ref{sec:significance:decision}).

\textbf{Statistical and economic significance are separate columns.} A
paired test over $5\,449$ windows detects differences far below anything
that changes a decision, so every comparison carries a $p$-value, an
effect size, and a threshold fixed before the tests were run: $1$\,s of
per-window MAE, $0.02$ of $F_1$, $5$\,USD/yr. A result is allowed to be
\emph{detectable but immaterial}, and two of the twenty-one forecasting
contrasts are. Multiplicity is controlled by Holm's step-down
procedure~\cite{holm1979simple}, which achieves the same family-wise
error rate as Bonferroni and is uniformly more powerful;
Benjamini--Hochberg~\cite{benjamini1995controlling} is reported beside
it for the exploratory families, and both are shown next to the raw
$p$-value so the cost of the correction stays visible.

\subsection{Forecasting}
\label{sec:significance:forecast}

Seven models, $5\,449$ identical held-out windows. Every stored
prediction array was re-scored and matched against its recorded metrics
before any test ran, and all seven sit on a byte-identical vector of
targets, which is what makes the pairing legitimate.

\paragraph{The pre-registered test.}
The comparison named in advance was the proposed Seq2Seq LSTM against
the untrained reference. It resolves decisively, and in the direction
opposite to the one it was written to check
(Table~\ref{tab:sig:headline}). The design was set up to ask whether a
small \emph{positive} gap was real; there is no positive gap. The $F_1$
interval is computed on seed~$0$ of each model, which is the
\emph{best} of the sequence model's five seeds --- and even granting it
its most favourable draw the interval excludes zero in persistence's
favour. On the seed mean the gap is five times wider.

\begin{table}[t]
\centering
\caption{The pre-registered comparison: proposed sequence model versus
the untrained reference, on $5\,449$ paired windows.}
\label{tab:sig:headline}
\small
\begin{tabular}{@{}lrrr@{}}
\toprule
quantity & Seq2Seq & \textsc{persistence} & $p$ \\
\midrule
per-window MAE (s)        & $20.90$ & $12.95$ & $5.7\times10^{-44}$ \\
horizon-step accuracy     & \multicolumn{2}{c}{$-0.0574$ (worse)} & $1.5\times10^{-113}$ \\
STANDBY $F_1$, seed mean  & $0.561$ & $0.681$ & --- \\
STANDBY $F_1$, seed 0     & $0.660$ & $0.681$ & CI $[-0.042,-0.003]$ \\
\bottomrule
\end{tabular}
\end{table}

\paragraph{The family.}
Of the $21$ pairwise contrasts, $19$ are significant after Holm
correction and $17$ of those also clear the $1$\,s practical threshold.
Two conclusions survive both filters. \emph{Persistence beats all five
neural models}, at $p_{\text{Holm}}\le6\times10^{-31}$ and by $4.5$--$8.0$\,s of
MAE. \emph{XGBoost beats persistence}, by $1.45$\,s at
$p_{\text{Holm}}=7.0\times10^{-4}$ with rank-biserial correlation
$0.24$ --- the single contrast in the entire family in which training
beats not training.

\paragraph{Friedman over seeds, and how little it can see.}
Blocking on the five seeds and treating the six trained models as
treatments gives $\chi^2=12.77$, $p=0.026$, Kendall's $W=0.511$
\cite{friedman1937use,demsar2006statistical}. XGBoost takes mean rank
$1.00$ --- it is the best model on every seed block --- and Seq2Seq
takes $5.00$, the worst. But the Nemenyi critical difference is $3.37$
on a $1$--$6$ rank scale, so \textbf{the only pair this diagram
separates is XGBoost against Seq2Seq} ($4.00$ apart) and everything else
is joined (Figure~\ref{fig:cd}). Measured power at this design is $0.14$
for a $1$-sd effect and $0.53$ at $2$\,sd. The seed-blocked comparison
is nearly blind, and the confirmatory conclusions above do not rest on
it: those are paired over $5\,449$ windows, not over $5$ seeds.

\begin{figure}[t]
  \centering
  \includegraphics[width=\linewidth]{fig_p6_cd_diagram}
  \caption{Critical-difference diagram over five seed blocks. Models
  joined by a bar are not separable at $\alpha=0.05$. The bar spans
  almost the whole scale because $n=5$ blocks buys a critical difference
  of $3.37$ ranks.}
  \label{fig:cd}
\end{figure}

\subsection{The decision layer}
\label{sec:significance:decision}

Day-block bootstrap, $2\,000$ draws, on the minimum-dwell labelling:
$394$ held-out episodes over \textbf{18 factory days} spanning $73.6$
calendar days.

\begin{table}[t]
\centering
\caption{Annual saving against the plant's own behaviour (USD/yr per
machine), with $95\%$ day-block bootstrap intervals. Every policy is
scored on the same resample.}
\label{tab:sig:decision}
\small
\setlength{\tabcolsep}{3pt}
\begin{tabular}{@{}lrrr@{}}
\toprule
Policy & S1 ($10$\,min) & S2 ($60$\,min) & S3 ($240$\,min) \\
\midrule
never shut down   & $-219.8$\,{\footnotesize$[-334,-132]$} & $-157.1$\,{\footnotesize$[-276,-65]$} & $+68.6$\,{\footnotesize$[-84,+213]$} \\
immediate         & $-219.7$\,{\footnotesize$[-338,-126]$} & $-307.5$\,{\footnotesize$[-430,-212]$} & $-623.4$\,{\footnotesize$[-774,-478]$} \\
static break-even & $-166.7$\,{\footnotesize$[-290,-71]$}  & $-66.0$\,{\footnotesize$[-196,+24]$}  & $+6.3$\,{\footnotesize$[-170,+168]$} \\
ski-rental        & $-46.7$\,{\footnotesize$[-93,-7]$}     & $+0.0$\,{\footnotesize$[-25,+29]$}    & $+30.4$\,{\footnotesize$[-79,+155]$} \\
\textbf{proposed} & $\mathbf{-28.5}$\,{\footnotesize$[-79,+1]$} & $\mathbf{+2.9}$\,{\footnotesize$[-56,+46]$} & $\mathbf{+94.5}$\,{\footnotesize$[-58,+237]$} \\
oracle            & $+25.5$\,{\footnotesize$[+15,+38]$}    & $+37.6$\,{\footnotesize$[+17,+61]$}   & $+160.8$\,{\footnotesize$[+65,+266]$} \\
\bottomrule
\end{tabular}
\end{table}

Table~\ref{tab:sig:decision} has to be read twice, and the two readings
give different answers.

\textbf{Read row by row, no policy but the oracle demonstrates a
saving.} Every other interval against the status quo spans zero,
including the proposed policy's headline $+2.9$ USD/yr inside
$[-55.9,+46.2]$. The point estimate the earlier sections quote is
indistinguishable from zero, and always was; the bootstrap supplies the
interval that makes it visible.

\textbf{Read as paired contrasts, one result survives}
(Table~\ref{tab:sig:contrasts}). Because every policy is scored on the
same resample, a difference between two rows is far better determined
than either row: the pairing removes the day-to-day variance that makes
the individual intervals wide. The proposed policy beats the static
break-even rule it was built to replace by $+68.9$ USD/yr
($p=0.001$, interval excluding zero, above the $5$ USD/yr economic
threshold) in S2 and by $+138$ and $+88$ USD/yr in S1 and S3. It does
\emph{not} beat forecast-free ski-rental, and it cannot be shown to beat
the plant's own behaviour.

\begin{table}[t]
\centering
\caption{Paired contrasts, same resample. The claim the decision layer
supports is the third row, and only the third row.}
\label{tab:sig:contrasts}
\small
\begin{tabular}{@{}lrrr@{}}
\toprule
Contrast & S1 & S2 & S3 \\
\midrule
proposed $-$ status quo   & $-28.5$ & $+2.9$ & $+94.5$ \\
                          & \footnotesize$p=0.072$ & \footnotesize$p=0.828$ & \footnotesize$p=0.196$ \\
proposed $-$ ski-rental   & $+18.2$ & $+2.8$ & $+64.1$ \\
                          & \footnotesize$p=0.553$ & \footnotesize$p=0.729$ & \footnotesize$p=0.212$ \\
\textbf{proposed $-$ static} & $\mathbf{+138.2}$ & $\mathbf{+68.9}$ & $\mathbf{+88.3}$ \\
                          & \footnotesize$\mathbf{p<0.001}$ & \footnotesize$\mathbf{p=0.001}$ & \footnotesize$\mathbf{p<0.001}$ \\
proposed $-$ oracle       & $-54.0$ & $-34.8$ & $-66.3$ \\
                          & \footnotesize$p<0.001$ & \footnotesize$p<0.001$ & \footnotesize$p<0.001$ \\
\bottomrule
\end{tabular}
\end{table}

\textbf{The oracle gap is also real}, in all three scenarios: the
head-room is genuine and mostly uncaptured --- the proposed policy takes
$7.6\%$ of it in S2.

\textbf{The width is the sample, not the estimator.} Eighteen factory
days is the binding constraint: the day-block bootstrap has eighteen
units to resample, and no refinement of the estimator narrows these
intervals. Any claim of economic benefit from this layer needs a longer
record, and this paper does not make one.
Figure~\ref{fig:decision-forest} shows all three scenarios at once,
with the intervals that cross zero drawn so that they cannot be
mistaken for the ones that do not.

\begin{figure}[t]
  \centering
  \includegraphics[width=\linewidth]{fig_p6_decision_forest}
  \caption{Decision-layer intervals in the three scenarios. Intervals
  crossing zero are drawn grey and hollow; only the oracle's excludes
  zero in every scenario.}
  \label{fig:decision-forest}
\end{figure}

\subsection{Cross-machine tests, and a design that cannot detect
anything}
\label{sec:significance:cross}

Three paired comparisons across the eight machines are reported in
Table~\ref{tab:sig:cross}, and the third column of that table is as
important as the second.

\begin{table}[t]
\centering
\caption{Paired signed-rank tests across machines. The $p$-value floor
is the smallest two-sided value the test can return at that $n$.}
\label{tab:sig:cross}
\small
\setlength{\tabcolsep}{4pt}
\begin{tabular}{@{}lrrrrr@{}}
\toprule
Comparison & $n$ & $\Delta F_1$ & $p$ & floor & power \\
\midrule
own vs \textsc{persistence}, all      & 8 & $+0.194$ & $0.016$ & $0.008$ & $0.41$ \\
own vs \textsc{persistence}, motors   & 4 & $+0.048$ & $0.250$ & $\mathbf{0.125}$ & $\mathbf{0.00}$ \\
local vs transferred                  & 8 & $+0.384$ & $\mathbf{0.008}$ & $0.008$ & $d=1.75$ \\
\bottomrule
\end{tabular}
\end{table}

\textbf{Training beats not training across the eight machines, but the
effect lives in the auxiliary plant.} The $+0.194$ mean difference is
dominated by exhaust-fan~I ($0.982$ against $0.288$) and dpc~I ($0.402$
against $0.012$), machines whose ``STANDBY'' cluster is really their off
state and is therefore trivially predictable once a model has seen it
(Section~\ref{sec:multimachine}). On the four motor-driven machines, the
ones on which the state model is physically valid, the difference is
$+0.048$ and cannot be tested.

\textbf{``Cannot be tested'' is meant literally.} With four paired
observations the smallest two-sided $p$-value the signed-rank test can
return is $2/2^{4}=0.125$, which is above $\alpha=0.05$. No data, and no
effect size however large, can produce a significant result at $n=4$.
\textbf{A non-significant verdict here is a statement about the design,
not about the models}, and it is reported in exactly those terms rather
than as ``no significant difference''.
Figure~\ref{fig:power} puts numbers on both designs: at $n=8$ the paired
test reaches $80\%$ power only above $\approx1.4$\,sd; at $n=4$ it never
does, at any effect size.

\textbf{The transfer degradation is real.} Locally fitted sequence
models beat transferred ones by $0.384$ of $F_1$ with $p$ at the floor
and $d=+1.75$. This is the one cross-machine claim of
Section~\ref{sec:crossmachine} that survives testing outright.

\begin{figure}[t]
  \centering
  \includegraphics[width=\linewidth]{fig_p6_power_curve}
  \caption{Power of the paired signed-rank test against effect size for
  $n=4$ and $n=8$. The marker is the effect actually observed on the
  four motor-driven machines ($d=0.56$), whose measured power is $0.00$
  --- readable off the axis rather than taken on trust.}
  \label{fig:power}
\end{figure}

\subsection{Repeated fitting, and the result that matters most}
\label{sec:significance:seeds}

\subsubsection{The labelling is not reproducible from the seed alone}

Re-fitting the state model at ten seeds on one deterministic
preprocessing of pelletizer~I gives a STANDBY total of
$84.6\pm28.5$\,h --- a coefficient of variation of $\mathbf{33.7\%}$ on
the headline state-time quantity of the paper. Three seeds in ten
converge on a cluster that is not standby at all:

\begin{center}
\small
\begin{tabular}{@{}rrrrl@{}}
\toprule
seed & physics & STANDBY & mean power & what it found \\
\midrule
6 & $0.80$ & $121.6$\,h & $12\,961$\,W & a loaded running state \\
8 & $0.60$ & $63.2$\,h  & $0.8$\,W     & the OFF state \\
9 & $0.60$ & $13.8$\,h  & $1.1$\,W     & the OFF state \\
\bottomrule
\end{tabular}
\end{center}

A ``standby'' cluster at $0.8$\,W drawing $0.06\%$ of rated current is
the machine switched off; one at $13$\,kW is the machine working. EM
initialisation decides which, and in $30\%$ of initialisations it
decides wrongly.

\textbf{What rescues the result is the physics battery, and this is the
constructive finding.} The tests of
Section~\ref{sec:method:accept_reject} score exactly those three seeds
at $0.60$--$0.80$ and the other seven at $1.00$, so the accept/reject
decision is exact on this record. Conditioning on a passing score
collapses the spread:

\begin{center}
\small
\begin{tabular}{@{}lrr@{}}
\toprule
 & all $10$ seeds & the $7$ that pass \\
\midrule
STANDBY hours   & $84.6\pm28.5$\,h & $\mathbf{92.4\pm2.8}$\,h \\
coefficient of variation & $\mathbf{33.7\%}$ & $\mathbf{3.0\%}$ \\
range           & $13.8$--$121.6$\,h & $87.3$--$94.2$\,h \\
STANDBY power   & $4\,570\pm3\,543$\,W & $4\,676\pm343$\,W \\
\bottomrule
\end{tabular}
\end{center}

\textbf{The pipeline is reproducible because of the physics validation,
not despite the clustering.} That is a stronger statement of this
paper's contribution than an unconditional total would be, and it is why
the accept/reject loop is part of the method in
Section~\ref{sec:method:accept_reject} rather than an appendix to it.
Flicker is $0.00\%$ at every seed and the adjusted Rand index against
seed~$0$ stays $\ge0.889$: the \emph{partition} is stable because OFF
and PEAK\_LOAD dominate the record, and it is the identity of the
STANDBY cluster specifically that moves.

\subsubsection{Model-fitting variance is not what makes the economics
uncertain}

Ten seeds of the identical decision configuration (S2, minimum-dwell
labels) give a proposed-policy saving of
$\mathbf{+2.92\pm0.59}$\,USD/yr, $7.8\pm1.6\%$ of oracle, on head-room
of $37.63$\,USD/yr with \emph{no} seed variance at all --- head-room
being a property of the episodes and the coefficients rather than of the
fit. The forecaster's own error is $8\,954\pm156$\,s, about $2.5$\,h on
episode durations, which is why the policy captures under $8\%$ of what
is available.

Set the two sources of uncertainty side by side: model fitting
contributes $\pm0.6$\,USD/yr, day-to-day sampling contributes
$\pm50$\,USD/yr. \textbf{The uncertainty in the economic claim is the
18-day sample, not the model}, and no amount of retraining addresses it.

\subsubsection{State-time totals}

The one-week moving-block bootstrap gives intervals of $\pm15$--$25\%$
of the point estimate on the five-month records
(Table~\ref{tab:multimachine}, per machine). Combined with the $\pm3\%$
conditional labelling spread above, this fixes the precision at which
every state-time number in this paper is quoted: \textbf{two significant
figures, and conditional on the physics battery passing}. Three
significant figures assert a precision neither interval supports.

%  Figures are referenced by bare name; add to the preamble
%      \graphicspath{{../outputs/phase6/figures/}}
%  Bibliography: paper/references.bib
%  Numbers below are generated by  python run_phase5.py  and stored in
%    outputs/phase5/task11/pelletizer-I/forecasting_tests.json
%    outputs/phase5/task11/pelletizer-I/decision_bootstrap_S*.json
%    outputs/phase5/task11/cross_machine_tests.json
%    outputs/phase5/task11/state_block_bootstrap.csv
%    outputs/phase5/task12/pelletizer-I/{forecasting,labelling,decision}/
% =====================================================================

\section{Statistical Analysis}
\label{sec:significance}

Every headline quantity in the preceding sections is a point estimate.
This section reports what happens to those estimates under resampling
and under repeated fitting, and --- for two of the designs used here ---
what they are incapable of detecting in the first place.

\subsection{Four decisions that determine every interval below}
\label{sec:significance:method}

\textbf{The resampling unit is never a row.} A $1$\,Hz industrial record
is massively autocorrelated, and an i.i.d.\ bootstrap over rows would
treat $5.5$ million readings as $5.5$ million independent observations.
Measured on a synthetic AR(1) series with $\rho=0.999$, the i.i.d.\
bootstrap returns a standard error $25\times$ too small. Three
dependence-aware units are used instead, each chosen from the structure
of the quantity being estimated:

\begin{center}
\small
\begin{tabular}{@{}lll@{}}
\toprule
quantity & unit & why \\
\midrule
state-time totals & one-week moving blocks~\cite{kunsch1989jackknife}
  & duty cycle closes over a week \\
decision savings & factory days
  & the restart cap is per calendar day \\
forecast errors & test windows, \emph{paired}
  & one problem solved twice \\
\bottomrule
\end{tabular}
\end{center}

\textbf{A day drawn twice counts as two days.} In the decision bootstrap
the resampled copies are re-keyed. Leaving the label unchanged would let
two copies of one day share a single day's restart budget, silently
tightening the very constraint under evaluation.

\textbf{Every policy is scored on the same resample.} Drawing
independent resamples per policy would break the pairing and inflate the
interval of every contrast --- and the contrasts, not the individual
intervals, are what the decision claim rests on
(Section~\ref{sec:significance:decision}).

\textbf{Statistical and economic significance are separate columns.} A
paired test over $5\,449$ windows detects differences far below anything
that changes a decision, so every comparison carries a $p$-value, an
effect size, and a threshold fixed before the tests were run: $1$\,s of
per-window MAE, $0.02$ of $F_1$, $5$\,USD/yr. A result is allowed to be
\emph{detectable but immaterial}, and two of the twenty-one forecasting
contrasts are. Multiplicity is controlled by Holm's step-down
procedure~\cite{holm1979simple}, which achieves the same family-wise
error rate as Bonferroni and is uniformly more powerful;
Benjamini--Hochberg~\cite{benjamini1995controlling} is reported beside
it for the exploratory families, and both are shown next to the raw
$p$-value so the cost of the correction stays visible.

\subsection{Forecasting}
\label{sec:significance:forecast}

Seven models, $5\,449$ identical held-out windows. Every stored
prediction array was re-scored and matched against its recorded metrics
before any test ran, and all seven sit on a byte-identical vector of
targets, which is what makes the pairing legitimate.

\paragraph{The pre-registered test.}
The comparison named in advance was the proposed Seq2Seq LSTM against
the untrained reference. It resolves decisively, and in the direction
opposite to the one it was written to check
(Table~\ref{tab:sig:headline}). The design was set up to ask whether a
small \emph{positive} gap was real; there is no positive gap. The $F_1$
interval is computed on seed~$0$ of each model, which is the
\emph{best} of the sequence model's five seeds --- and even granting it
its most favourable draw the interval excludes zero in persistence's
favour. On the seed mean the gap is five times wider.

\begin{table}[t]
\centering
\caption{The pre-registered comparison: proposed sequence model versus
the untrained reference, on $5\,449$ paired windows.}
\label{tab:sig:headline}
\small
\begin{tabular}{@{}lrrr@{}}
\toprule
quantity & Seq2Seq & \textsc{persistence} & $p$ \\
\midrule
per-window MAE (s)        & $20.90$ & $12.95$ & $5.7\times10^{-44}$ \\
horizon-step accuracy     & \multicolumn{2}{c}{$-0.0574$ (worse)} & $1.5\times10^{-113}$ \\
STANDBY $F_1$, seed mean  & $0.561$ & $0.681$ & --- \\
STANDBY $F_1$, seed 0     & $0.660$ & $0.681$ & CI $[-0.042,-0.003]$ \\
\bottomrule
\end{tabular}
\end{table}

\paragraph{The family.}
Of the $21$ pairwise contrasts, $19$ are significant after Holm
correction and $17$ of those also clear the $1$\,s practical threshold.
Two conclusions survive both filters. \emph{Persistence beats all five
neural models}, at $p_{\text{Holm}}\le6\times10^{-31}$ and by $4.5$--$8.0$\,s of
MAE. \emph{XGBoost beats persistence}, by $1.45$\,s at
$p_{\text{Holm}}=7.0\times10^{-4}$ with rank-biserial correlation
$0.24$ --- the single contrast in the entire family in which training
beats not training.

\paragraph{Friedman over seeds, and how little it can see.}
Blocking on the five seeds and treating the six trained models as
treatments gives $\chi^2=12.77$, $p=0.026$, Kendall's $W=0.511$
\cite{friedman1937use,demsar2006statistical}. XGBoost takes mean rank
$1.00$ --- it is the best model on every seed block --- and Seq2Seq
takes $5.00$, the worst. But the Nemenyi critical difference is $3.37$
on a $1$--$6$ rank scale, so \textbf{the only pair this diagram
separates is XGBoost against Seq2Seq} ($4.00$ apart) and everything else
is joined (Figure~\ref{fig:cd}). Measured power at this design is $0.14$
for a $1$-sd effect and $0.53$ at $2$\,sd. The seed-blocked comparison
is nearly blind, and the confirmatory conclusions above do not rest on
it: those are paired over $5\,449$ windows, not over $5$ seeds.

\begin{figure}[t]
  \centering
  \includegraphics[width=\linewidth]{fig_p6_cd_diagram}
  \caption{Critical-difference diagram over five seed blocks. Models
  joined by a bar are not separable at $\alpha=0.05$. The bar spans
  almost the whole scale because $n=5$ blocks buys a critical difference
  of $3.37$ ranks.}
  \label{fig:cd}
\end{figure}

\subsection{The decision layer}
\label{sec:significance:decision}

Day-block bootstrap, $2\,000$ draws, on the minimum-dwell labelling:
$394$ held-out episodes over \textbf{18 factory days} spanning $73.6$
calendar days.

\begin{table}[t]
\centering
\caption{Annual saving against the plant's own behaviour (USD/yr per
machine), with $95\%$ day-block bootstrap intervals. Every policy is
scored on the same resample.}
\label{tab:sig:decision}
\small
\setlength{\tabcolsep}{3pt}
\begin{tabular}{@{}lrrr@{}}
\toprule
Policy & S1 ($10$\,min) & S2 ($60$\,min) & S3 ($240$\,min) \\
\midrule
never shut down   & $-219.8$\,{\footnotesize$[-334,-132]$} & $-157.1$\,{\footnotesize$[-276,-65]$} & $+68.6$\,{\footnotesize$[-84,+213]$} \\
immediate         & $-219.7$\,{\footnotesize$[-338,-126]$} & $-307.5$\,{\footnotesize$[-430,-212]$} & $-623.4$\,{\footnotesize$[-774,-478]$} \\
static break-even & $-166.7$\,{\footnotesize$[-290,-71]$}  & $-66.0$\,{\footnotesize$[-196,+24]$}  & $+6.3$\,{\footnotesize$[-170,+168]$} \\
ski-rental        & $-46.7$\,{\footnotesize$[-93,-7]$}     & $+0.0$\,{\footnotesize$[-25,+29]$}    & $+30.4$\,{\footnotesize$[-79,+155]$} \\
\textbf{proposed} & $\mathbf{-28.5}$\,{\footnotesize$[-79,+1]$} & $\mathbf{+2.9}$\,{\footnotesize$[-56,+46]$} & $\mathbf{+94.5}$\,{\footnotesize$[-58,+237]$} \\
oracle            & $+25.5$\,{\footnotesize$[+15,+38]$}    & $+37.6$\,{\footnotesize$[+17,+61]$}   & $+160.8$\,{\footnotesize$[+65,+266]$} \\
\bottomrule
\end{tabular}
\end{table}

Table~\ref{tab:sig:decision} has to be read twice, and the two readings
give different answers.

\textbf{Read row by row, no policy but the oracle demonstrates a
saving.} Every other interval against the status quo spans zero,
including the proposed policy's headline $+2.9$ USD/yr inside
$[-55.9,+46.2]$. The point estimate the earlier sections quote is
indistinguishable from zero, and always was; the bootstrap supplies the
interval that makes it visible.

\textbf{Read as paired contrasts, one result survives}
(Table~\ref{tab:sig:contrasts}). Because every policy is scored on the
same resample, a difference between two rows is far better determined
than either row: the pairing removes the day-to-day variance that makes
the individual intervals wide. The proposed policy beats the static
break-even rule it was built to replace by $+68.9$ USD/yr
($p=0.001$, interval excluding zero, above the $5$ USD/yr economic
threshold) in S2 and by $+138$ and $+88$ USD/yr in S1 and S3. It does
\emph{not} beat forecast-free ski-rental, and it cannot be shown to beat
the plant's own behaviour.

\begin{table}[t]
\centering
\caption{Paired contrasts, same resample. The claim the decision layer
supports is the third row, and only the third row.}
\label{tab:sig:contrasts}
\small
\begin{tabular}{@{}lrrr@{}}
\toprule
Contrast & S1 & S2 & S3 \\
\midrule
proposed $-$ status quo   & $-28.5$ & $+2.9$ & $+94.5$ \\
                          & \footnotesize$p=0.072$ & \footnotesize$p=0.828$ & \footnotesize$p=0.196$ \\
proposed $-$ ski-rental   & $+18.2$ & $+2.8$ & $+64.1$ \\
                          & \footnotesize$p=0.553$ & \footnotesize$p=0.729$ & \footnotesize$p=0.212$ \\
\textbf{proposed $-$ static} & $\mathbf{+138.2}$ & $\mathbf{+68.9}$ & $\mathbf{+88.3}$ \\
                          & \footnotesize$\mathbf{p<0.001}$ & \footnotesize$\mathbf{p=0.001}$ & \footnotesize$\mathbf{p<0.001}$ \\
proposed $-$ oracle       & $-54.0$ & $-34.8$ & $-66.3$ \\
                          & \footnotesize$p<0.001$ & \footnotesize$p<0.001$ & \footnotesize$p<0.001$ \\
\bottomrule
\end{tabular}
\end{table}

\textbf{The oracle gap is also real}, in all three scenarios: the
head-room is genuine and mostly uncaptured --- the proposed policy takes
$7.6\%$ of it in S2.

\textbf{The width is the sample, not the estimator.} Eighteen factory
days is the binding constraint: the day-block bootstrap has eighteen
units to resample, and no refinement of the estimator narrows these
intervals. Any claim of economic benefit from this layer needs a longer
record, and this paper does not make one.
Figure~\ref{fig:decision-forest} shows all three scenarios at once,
with the intervals that cross zero drawn so that they cannot be
mistaken for the ones that do not.

\begin{figure}[t]
  \centering
  \includegraphics[width=\linewidth]{fig_p6_decision_forest}
  \caption{Decision-layer intervals in the three scenarios. Intervals
  crossing zero are drawn grey and hollow; only the oracle's excludes
  zero in every scenario.}
  \label{fig:decision-forest}
\end{figure}

\subsection{Cross-machine tests, and a design that cannot detect
anything}
\label{sec:significance:cross}

Three paired comparisons across the eight machines are reported in
Table~\ref{tab:sig:cross}, and the third column of that table is as
important as the second.

\begin{table}[t]
\centering
\caption{Paired signed-rank tests across machines. The $p$-value floor
is the smallest two-sided value the test can return at that $n$.}
\label{tab:sig:cross}
\small
\setlength{\tabcolsep}{4pt}
\begin{tabular}{@{}lrrrrr@{}}
\toprule
Comparison & $n$ & $\Delta F_1$ & $p$ & floor & power \\
\midrule
own vs \textsc{persistence}, all      & 8 & $+0.194$ & $0.016$ & $0.008$ & $0.41$ \\
own vs \textsc{persistence}, motors   & 4 & $+0.048$ & $0.250$ & $\mathbf{0.125}$ & $\mathbf{0.00}$ \\
local vs transferred                  & 8 & $+0.384$ & $\mathbf{0.008}$ & $0.008$ & $d=1.75$ \\
\bottomrule
\end{tabular}
\end{table}

\textbf{Training beats not training across the eight machines, but the
effect lives in the auxiliary plant.} The $+0.194$ mean difference is
dominated by exhaust-fan~I ($0.982$ against $0.288$) and dpc~I ($0.402$
against $0.012$), machines whose ``STANDBY'' cluster is really their off
state and is therefore trivially predictable once a model has seen it
(Section~\ref{sec:multimachine}). On the four motor-driven machines, the
ones on which the state model is physically valid, the difference is
$+0.048$ and cannot be tested.

\textbf{``Cannot be tested'' is meant literally.} With four paired
observations the smallest two-sided $p$-value the signed-rank test can
return is $2/2^{4}=0.125$, which is above $\alpha=0.05$. No data, and no
effect size however large, can produce a significant result at $n=4$.
\textbf{A non-significant verdict here is a statement about the design,
not about the models}, and it is reported in exactly those terms rather
than as ``no significant difference''.
Figure~\ref{fig:power} puts numbers on both designs: at $n=8$ the paired
test reaches $80\%$ power only above $\approx1.4$\,sd; at $n=4$ it never
does, at any effect size.

\textbf{The transfer degradation is real.} Locally fitted sequence
models beat transferred ones by $0.384$ of $F_1$ with $p$ at the floor
and $d=+1.75$. This is the one cross-machine claim of
Section~\ref{sec:crossmachine} that survives testing outright.

\begin{figure}[t]
  \centering
  \includegraphics[width=\linewidth]{fig_p6_power_curve}
  \caption{Power of the paired signed-rank test against effect size for
  $n=4$ and $n=8$. The marker is the effect actually observed on the
  four motor-driven machines ($d=0.56$), whose measured power is $0.00$
  --- readable off the axis rather than taken on trust.}
  \label{fig:power}
\end{figure}

\subsection{Repeated fitting, and the result that matters most}
\label{sec:significance:seeds}

\subsubsection{The labelling is not reproducible from the seed alone}

Re-fitting the state model at ten seeds on one deterministic
preprocessing of pelletizer~I gives a STANDBY total of
$84.6\pm28.5$\,h --- a coefficient of variation of $\mathbf{33.7\%}$ on
the headline state-time quantity of the paper. Three seeds in ten
converge on a cluster that is not standby at all:

\begin{center}
\small
\begin{tabular}{@{}rrrrl@{}}
\toprule
seed & physics & STANDBY & mean power & what it found \\
\midrule
6 & $0.80$ & $121.6$\,h & $12\,961$\,W & a loaded running state \\
8 & $0.60$ & $63.2$\,h  & $0.8$\,W     & the OFF state \\
9 & $0.60$ & $13.8$\,h  & $1.1$\,W     & the OFF state \\
\bottomrule
\end{tabular}
\end{center}

A ``standby'' cluster at $0.8$\,W drawing $0.06\%$ of rated current is
the machine switched off; one at $13$\,kW is the machine working. EM
initialisation decides which, and in $30\%$ of initialisations it
decides wrongly.

\textbf{What rescues the result is the physics battery, and this is the
constructive finding.} The tests of
Section~\ref{sec:method:accept_reject} score exactly those three seeds
at $0.60$--$0.80$ and the other seven at $1.00$, so the accept/reject
decision is exact on this record. Conditioning on a passing score
collapses the spread:

\begin{center}
\small
\begin{tabular}{@{}lrr@{}}
\toprule
 & all $10$ seeds & the $7$ that pass \\
\midrule
STANDBY hours   & $84.6\pm28.5$\,h & $\mathbf{92.4\pm2.8}$\,h \\
coefficient of variation & $\mathbf{33.7\%}$ & $\mathbf{3.0\%}$ \\
range           & $13.8$--$121.6$\,h & $87.3$--$94.2$\,h \\
STANDBY power   & $4\,570\pm3\,543$\,W & $4\,676\pm343$\,W \\
\bottomrule
\end{tabular}
\end{center}

\textbf{The pipeline is reproducible because of the physics validation,
not despite the clustering.} That is a stronger statement of this
paper's contribution than an unconditional total would be, and it is why
the accept/reject loop is part of the method in
Section~\ref{sec:method:accept_reject} rather than an appendix to it.
Flicker is $0.00\%$ at every seed and the adjusted Rand index against
seed~$0$ stays $\ge0.889$: the \emph{partition} is stable because OFF
and PEAK\_LOAD dominate the record, and it is the identity of the
STANDBY cluster specifically that moves.

\subsubsection{Model-fitting variance is not what makes the economics
uncertain}

Ten seeds of the identical decision configuration (S2, minimum-dwell
labels) give a proposed-policy saving of
$\mathbf{+2.92\pm0.59}$\,USD/yr, $7.8\pm1.6\%$ of oracle, on head-room
of $37.63$\,USD/yr with \emph{no} seed variance at all --- head-room
being a property of the episodes and the coefficients rather than of the
fit. The forecaster's own error is $8\,954\pm156$\,s, about $2.5$\,h on
episode durations, which is why the policy captures under $8\%$ of what
is available.

Set the two sources of uncertainty side by side: model fitting
contributes $\pm0.6$\,USD/yr, day-to-day sampling contributes
$\pm50$\,USD/yr. \textbf{The uncertainty in the economic claim is the
18-day sample, not the model}, and no amount of retraining addresses it.

\subsubsection{State-time totals}

The one-week moving-block bootstrap gives intervals of $\pm15$--$25\%$
of the point estimate on the five-month records
(Table~\ref{tab:multimachine}, per machine). Combined with the $\pm3\%$
conditional labelling spread above, this fixes the precision at which
every state-time number in this paper is quoted: \textbf{two significant
figures, and conditional on the physics battery passing}. Three
significant figures assert a precision neither interval supports.


% --- The eight machines
% =====================================================================
%  multimachine.tex  --  PHASE IV, Task 8
%
%  The eight-machine characterisation: where the state definition holds,
%  where it fails, how it says so, and which machines have a decision
%  problem at all.
%
%  Intended placement: Section 8 (\label{sec:multimachine}), after the
%  decision layer.
%
%  Drop-in usage: % =====================================================================
%  multimachine.tex  --  PHASE IV, Task 8
%
%  The eight-machine characterisation: where the state definition holds,
%  where it fails, how it says so, and which machines have a decision
%  problem at all.
%
%  Intended placement: Section 8 (\label{sec:multimachine}), after the
%  decision layer.
%
%  Drop-in usage: % =====================================================================
%  multimachine.tex  --  PHASE IV, Task 8
%
%  The eight-machine characterisation: where the state definition holds,
%  where it fails, how it says so, and which machines have a decision
%  problem at all.
%
%  Intended placement: Section 8 (\label{sec:multimachine}), after the
%  decision layer.
%
%  Drop-in usage: \input{multimachine}
%  Figures are referenced by bare name; add to the preamble
%      \graphicspath{{../outputs/phase4/figures/}{../outputs/phase6/figures/}}
%  Bibliography: paper/references.bib
%  Numbers below are generated by  python run_phase4.py --only task8
%  and stored in outputs/phase4/task8/table_*.csv, with the intervals
%  from outputs/phase5/task11/state_block_bootstrap.csv.
% =====================================================================

\section{The Eight Machines}
\label{sec:multimachine}

Everything to this point is pelletizer~I. This section applies the
identical procedure --- one preprocessing, one labeller, $k=4$ fixed,
one physics battery --- to all eight IMDELD machines, and asks two
questions of the result: on which machines is the state definition
physically valid, and on which is there a shutdown decision to make at
all. The two answers are not the same, and neither follows the
prediction that motivated the experiment.

\subsection{Uniform labelling}
\label{sec:multimachine:labelling}

The Phase~I--III labels exist only for pelletizer~I and were produced by
a script whose configuration has since moved, so comparing that export
against seven fresh labellings would confound machine differences with
procedure differences. All eight records are therefore re-labelled
identically. Three departures from the original implementation were
required, each verified rather than assumed.

\emph{Blank currents.} The two milling-machine meters leave the current
channel empty on $36.1\%$ and $36.3\%$ of rows --- matching their OFF
shares almost exactly, i.e.\ the meter blanks the channel whenever the
machine is de-energised. On every such row the apparent power is exactly
$0$\,VA with voltage present, and $S=VI$ then admits only $I=0$\,A, so
the cells are filled with zero and a unit test asserts that every filled
cell met that condition. Leaving them missing makes a GMM
unfittable; dropping them would delete a third of each record precisely
where the machine is off, biasing every state-time total.

\emph{Timestamps and storage.} IMDELD spans a daylight-saving change, so
the timestamp column carries mixed UTC offsets; splitting the naive part
from the offset is exact and ten times faster than the fallback parse.
Records are stored as \texttt{float32}, whose one measurable consequence
is that about two rows per $10^{5}$ fall on the other side of the
spike threshold.

On pelletizer~I --- the one machine both labellings cover --- the
uniform labeller reproduces the pipeline's own export at $98.88\%$ row
agreement, $99.84\%$ on STANDBY-versus-rest, ARI $0.978$, and $89$\,h of
STANDBY against the export's $87$\,h. A number in this section that
differs from an earlier one is therefore a consequence of the analysis,
not of the labels.

Both of those figures are single-seed. The total this paper quotes for
pelletizer~I is the physics-conditional mean over passing seeds,
$92\pm3$\,h (Section~\ref{sec:method:accept_reject}); the single-seed
values sit inside that interval, and the per-machine rows of
Table~\ref{tab:multimachine} are reported per covered day with a
sampling interval for the same reason.

\subsection{Where the state definition holds}
\label{sec:multimachine:states}

\begin{table*}[!t]
\centering
\caption{Per-machine characterisation. STANDBY hours per covered day
carry a $95\%$ one-week moving-block bootstrap interval
(Section~\ref{sec:significance}); $n_{b}$ is the number of blocks the
record supports. $P_{\text{sb}}$ is the mean power of the inferred
STANDBY component. The no-load current ratio has an acceptance band of
$[0.25,0.50]$ from induction-motor theory. M1 is the fraction of
energised rows on which the independent Otsu power-factor classifier
agrees with the state model.}
\label{tab:multimachine}
\small
\begin{tabular}{@{}lrrrrrrrr@{}}
\toprule
Machine & cov.\ (d) & STANDBY (h/d) & $n_{b}$ & $P_{\text{sb}}$ & physics & PF sep. & $I$ ratio & M1 \\
\midrule
\multicolumn{9}{@{}l}{\textit{Motor-driven --- the state definition applies}}\\
Pelletizer I    & $63.4$ & $1.40$\,{\footnotesize$[1.15,1.63]$} & $10$ & $4.2$\,kW & $\mathbf{5/5}$ & $0.72$ & $0.36$ & $99.4\%$ \\
Pelletizer II   & $62.4$ & $1.85$\,{\footnotesize$[1.47,2.28]$} & $9$  & $4.5$\,kW & $\mathbf{5/5}$ & $0.71$ & $0.36$ & $99.7\%$ \\
Milling I       & $12.2$ & $5.52$\,{\footnotesize$[4.58,6.23]$} & $\mathbf{2}$ & $4.6$\,kW & $\mathbf{5/5}$ & $0.61$ & $0.43$ & $98.1\%$ \\
Milling II      & $12.2$ & $5.27$\,{\footnotesize$[4.44,5.89]$} & $\mathbf{2}$ & $5.4$\,kW & $\mathbf{5/5}$ & $0.52$ & $0.33$ & $97.7\%$ \\
\addlinespace
\multicolumn{9}{@{}l}{\textit{Auxiliary --- no energised-idle state exists to find}}\\
Exhaust fan I   & $63.7$ & $8.73$\,{\footnotesize$[7.31,10.03]$} & $10$ & $0.3$\,W & $4/5$ & $0.91$ & $\mathbf{0.008}$ & $96.6\%$ \\
Exhaust fan II  & $62.3$ & $1.45$\,{\footnotesize$[1.09,2.07]$}  & $9$  & $0.2$\,W & $3/5$ & $0.62$ & $\mathbf{0.002}$ & $\mathbf{55.1\%}$ \\
Contactor I     & $63.7$ & $0.90$\,{\footnotesize$[0.71,1.12]$}  & $10$ & $0.3$\,W & $3/5$ & $0.31$ & $\mathbf{0.009}$ & $90.3\%$ \\
Contactor II    & $62.5$ & $1.15$\,{\footnotesize$[0.85,1.42]$}  & $9$  & $9.3$\,W & $3/5$ & $\mathbf{-0.05}$ & $\mathbf{0.007}$ & $85.4\%$ \\
\bottomrule
\end{tabular}
\end{table*}

\textbf{The two pelletizers and the two milling machines pass all five
physics checks.} Their STANDBY components sit at $4.2$--$5.4$\,kW with a
power factor of $0.11$--$0.17$ against $0.64$--$0.65$ when loaded, and a
no-load current ratio inside the motor band $[0.25,0.50]$. That is the
textbook signature of an energised induction motor drawing magnetising
current, recovered independently on four machines from unlabelled data,
and it is the strongest evidence in this paper that the state definition
is physical rather than fitted.

\textbf{On the fans and contactors there is no energised-idle state to
find, and the battery says so.} Their ``STANDBY'' components sit at
$0.2$--$9.3$\,W drawing $0.09$--$0.14$\,A: the model is splitting the
OFF band in two because it was told to find four states. The current
ratio is $0.002$--$0.009$ against an acceptance band starting at
$0.25$ on all four, and contactor~II additionally shows a
\emph{negative} power-factor separation. \textbf{This is the physics
battery doing the job it was built for.} The framework's state
definition is specific to motor-driven loads, and on the machines where
it does not apply it reports that rather than returning four confident
states regardless.

Two mechanisms behind those failures are worth naming, because both are
limitations of the method rather than of the data.

\emph{Ranking states by active power can misname them.} Exhaust fan~I's
``OFF'' component ($0.8\%$ of the record) carries $6.6$\,A at $0$\,W and
power factor $0$ --- standing reactive current with no real power ---
while its ``STANDBY'' component ($36.4\%$) sits at $0.3$\,W and
$0.13$\,A and is the machine's true de-energised state. The canonical
names are assigned by mean active power, so a near-zero-power component
carrying reactive current ranks below the genuine off state and takes
its name. No hour total changes, since both are non-productive, but it
is why the current-ratio check fails in the direction it does. The same
machine returns the largest forbidden-transition probability in the
plant ($3.4\times10^{-3}$, against $1.9\times10^{-4}$ on the four motor
machines), which is an independent measurement arriving at the same
conclusion (Section~\ref{sec:explain:transitions}).

\emph{The independent classifier fails where the model has nothing to
lock onto.} Otsu-thresholded power factor agrees with the state model on
$85$--$99.7\%$ of energised rows on seven machines and on $\mathbf{55.1\%}$
on exhaust fan~II --- the machine whose physics checks fail hardest.
Method~1's usefulness as an external check is thereby confirmed
\emph{and} its failure mode identified.

\subsection{Two results that generalise}
\label{sec:multimachine:general}

\textbf{The minimum-dwell result holds on all eight records.} Flicker
before the constraint ranges from $26.7\%$ (milling~II) to $77.7\%$
(exhaust fan~II), median $50.6\%$; after it, $0.00\%$ on seven machines
and $0.17\%$ on milling~I, at a cost of relabelling $0.15$--$6.1\%$ of
rows. What Section~\ref{sec:results:task4} established on one machine
--- that a transition prior cannot control flicker and an explicit
dwell constraint can --- now holds across duty cycles differing by an
order of magnitude.

\textbf{BIC does not select a state count here.} Widening the scan to
$k\in[2,8]$ returns $k=8$ on seven machines and $k=7$ on pelletizer~II,
with the criterion still falling at the boundary in seven of eight
cases; the improvement from $k=4$ to $k=8$ is $3.1$--$33.1\%$ of the
criterion's own value. This is a result about the criterion, not the
data (Section~\ref{sec:rw:states}). Fixing $k=4$ everywhere for
comparability is a choice, and this is the honest statement of what it
costs: \textbf{the physics vetoes are not a belt-and-braces addition on
top of BIC --- on records of this size they are what is actually doing
the selecting.}

\subsection{Which machines are hard, and why}
\label{sec:multimachine:difficulty}

Difficulty is measured on three axes: separability of the no-load/load
boundary (Cohen's $d$ and Bhattacharyya distance on log power), record
quality (coverage, spike rate, segment count), and agreement with the
independent classifier.

The ordering is not the one predicted. The expectation was that the
low-power auxiliary machines would be hardest. What the data show is
that the fans and contactors have the \emph{largest} separability
figures in the plant --- Cohen's $d$ of $9.3$ to $75.0$ against the
pelletizers' $4.3$ to $7.2$ --- while failing the physics checks,
because the separation being measured there is OFF against running, not
no-load against loaded. \textbf{Separability and validity are close to
orthogonal on this data}, which is the argument for keeping an external
physics battery rather than an internal cluster-quality score:
silhouette, BIC and Cohen's $d$ are all maximised by the partitions that
are physically wrong.

The two milling machines are hard for an unrelated reason. Their records
cover $12.2$ of $43$ days in $1\,082$ and $580$ fragments, against
$66$--$74$ fragments for the five-month records. They support two
one-week bootstrap blocks, which is not enough to form a usable
interval, and the intervals in Table~\ref{tab:multimachine} are printed
for those rows only to make that visible.
\textbf{Their $5.5$\,h/day --- the largest per-day standby figure among
the motor-driven machines --- must not be annualised}, and
Figure~\ref{fig:standby-hours} prints the block count inside each bar so
the figure cannot be read as the largest opportunity in the plant.

\begin{figure}[t]
  \centering
  \includegraphics[width=\linewidth]{fig_p6_standby_hours}
  \caption{STANDBY hours per covered day with $95\%$ weekly-block
  intervals. Bars are marked according to whether the machine passes the
  physics battery, and the number of weekly blocks is printed inside
  each bar; hatching marks an interval that rests on too few blocks to
  be usable.}
  \label{fig:standby-hours}
\end{figure}

\subsection{The decision layer exists on three machines of eight}
\label{sec:multimachine:decision}

A decision problem requires a measurable standby power and an observable
cold restart. Five machines fail an explicit degeneracy test rather than
an inspection: four because standby power is below $1$\,W, so the
standby cost per second is numerically zero, the break-even point is
infinite, and every currency figure would be an artefact of dividing by
zero; and pelletizer~II because in five months it is \emph{never once
observed resuming production from OFF}, so the cold-versus-warm contrast
cannot be measured at all. Two machines return a \emph{negative} restart
delay --- cold starts reaching production faster than warm ones --- which
is not a plausible physical reading and is a further symptom of the same
problem. The three machines that survive the test are the subject of
Table~\ref{tab:multimachine:decision}.

\begin{table}[t]
\centering
\caption{Machines on which the shutdown decision is well posed.
Head-room is the oracle saving; the proposed policy's figure is quoted
against the plant's own behaviour.}
\label{tab:multimachine:decision}
\small
\setlength{\tabcolsep}{4pt}
\begin{tabular}{@{}lrrrr@{}}
\toprule
Machine & energised idle & head-room & forecast MAE & proposed \\
        & (h, \% of idle) & (USD/yr) & (s) & (USD/yr) \\
\midrule
Pelletizer I & $88$ ($18.5\%$) & $37.6$ & $9\,034$ & $+2.9$ \\
Milling I    & $67$ ($48.8\%$) & $24.9$ & $18\,205$ & $\mathbf{-1\,041}$ \\
Milling II   & $64$ ($47.6\%$) & $25.2$ & $18\,440$ & $\mathbf{-1\,256}$ \\
\bottomrule
\end{tabular}
\end{table}

\textbf{The conjecture that motivated this experiment is refuted.} The
open question after the single-machine study was whether some other
IMDELD machine has enough recoverable standby to fund a retrofit ---
higher standby power, a larger energised-idle fraction, or a plant that
does not already shut down manually. None does. The three machines with
a well-posed problem have head-rooms of $37.6$, $24.9$ and
$25.2$\,USD/yr, so \textbf{the facility's total attainable saving from
standby recovery, summed over every machine where the question is even
meaningful, is under $90$\,USD/yr}. The framework's contribution on this
facility is decision quality, not recovered energy, and this paper says
so in its abstract rather than at the end of its discussion.

\textbf{The milling machines are worse than that, and the negative
result is reported as it stands.} The proposed policy \emph{loses}
$1\,041$ and $1\,256$\,USD/yr there against a head-room of
$25$\,USD/yr. The cause is visible in the same table: a duration
forecaster with MAE $18\,200$\,s and Spearman $0.34$--$0.42$, against
$9\,034$\,s and $0.58$ on pelletizer~I, fitted on a record broken into
$1\,082$ fragments. It over-predicts remaining idle time and fires
shutdowns the realised episode does not support. The decision layer
inherits its forecaster's quality, and on a record this fragmented the
forecaster is not good enough to act on --- which is a deployment
criterion, not merely an observation.

%  Figures are referenced by bare name; add to the preamble
%      \graphicspath{{../outputs/phase4/figures/}{../outputs/phase6/figures/}}
%  Bibliography: paper/references.bib
%  Numbers below are generated by  python run_phase4.py --only task8
%  and stored in outputs/phase4/task8/table_*.csv, with the intervals
%  from outputs/phase5/task11/state_block_bootstrap.csv.
% =====================================================================

\section{The Eight Machines}
\label{sec:multimachine}

Everything to this point is pelletizer~I. This section applies the
identical procedure --- one preprocessing, one labeller, $k=4$ fixed,
one physics battery --- to all eight IMDELD machines, and asks two
questions of the result: on which machines is the state definition
physically valid, and on which is there a shutdown decision to make at
all. The two answers are not the same, and neither follows the
prediction that motivated the experiment.

\subsection{Uniform labelling}
\label{sec:multimachine:labelling}

The Phase~I--III labels exist only for pelletizer~I and were produced by
a script whose configuration has since moved, so comparing that export
against seven fresh labellings would confound machine differences with
procedure differences. All eight records are therefore re-labelled
identically. Three departures from the original implementation were
required, each verified rather than assumed.

\emph{Blank currents.} The two milling-machine meters leave the current
channel empty on $36.1\%$ and $36.3\%$ of rows --- matching their OFF
shares almost exactly, i.e.\ the meter blanks the channel whenever the
machine is de-energised. On every such row the apparent power is exactly
$0$\,VA with voltage present, and $S=VI$ then admits only $I=0$\,A, so
the cells are filled with zero and a unit test asserts that every filled
cell met that condition. Leaving them missing makes a GMM
unfittable; dropping them would delete a third of each record precisely
where the machine is off, biasing every state-time total.

\emph{Timestamps and storage.} IMDELD spans a daylight-saving change, so
the timestamp column carries mixed UTC offsets; splitting the naive part
from the offset is exact and ten times faster than the fallback parse.
Records are stored as \texttt{float32}, whose one measurable consequence
is that about two rows per $10^{5}$ fall on the other side of the
spike threshold.

On pelletizer~I --- the one machine both labellings cover --- the
uniform labeller reproduces the pipeline's own export at $98.88\%$ row
agreement, $99.84\%$ on STANDBY-versus-rest, ARI $0.978$, and $89$\,h of
STANDBY against the export's $87$\,h. A number in this section that
differs from an earlier one is therefore a consequence of the analysis,
not of the labels.

Both of those figures are single-seed. The total this paper quotes for
pelletizer~I is the physics-conditional mean over passing seeds,
$92\pm3$\,h (Section~\ref{sec:method:accept_reject}); the single-seed
values sit inside that interval, and the per-machine rows of
Table~\ref{tab:multimachine} are reported per covered day with a
sampling interval for the same reason.

\subsection{Where the state definition holds}
\label{sec:multimachine:states}

\begin{table*}[!t]
\centering
\caption{Per-machine characterisation. STANDBY hours per covered day
carry a $95\%$ one-week moving-block bootstrap interval
(Section~\ref{sec:significance}); $n_{b}$ is the number of blocks the
record supports. $P_{\text{sb}}$ is the mean power of the inferred
STANDBY component. The no-load current ratio has an acceptance band of
$[0.25,0.50]$ from induction-motor theory. M1 is the fraction of
energised rows on which the independent Otsu power-factor classifier
agrees with the state model.}
\label{tab:multimachine}
\small
\begin{tabular}{@{}lrrrrrrrr@{}}
\toprule
Machine & cov.\ (d) & STANDBY (h/d) & $n_{b}$ & $P_{\text{sb}}$ & physics & PF sep. & $I$ ratio & M1 \\
\midrule
\multicolumn{9}{@{}l}{\textit{Motor-driven --- the state definition applies}}\\
Pelletizer I    & $63.4$ & $1.40$\,{\footnotesize$[1.15,1.63]$} & $10$ & $4.2$\,kW & $\mathbf{5/5}$ & $0.72$ & $0.36$ & $99.4\%$ \\
Pelletizer II   & $62.4$ & $1.85$\,{\footnotesize$[1.47,2.28]$} & $9$  & $4.5$\,kW & $\mathbf{5/5}$ & $0.71$ & $0.36$ & $99.7\%$ \\
Milling I       & $12.2$ & $5.52$\,{\footnotesize$[4.58,6.23]$} & $\mathbf{2}$ & $4.6$\,kW & $\mathbf{5/5}$ & $0.61$ & $0.43$ & $98.1\%$ \\
Milling II      & $12.2$ & $5.27$\,{\footnotesize$[4.44,5.89]$} & $\mathbf{2}$ & $5.4$\,kW & $\mathbf{5/5}$ & $0.52$ & $0.33$ & $97.7\%$ \\
\addlinespace
\multicolumn{9}{@{}l}{\textit{Auxiliary --- no energised-idle state exists to find}}\\
Exhaust fan I   & $63.7$ & $8.73$\,{\footnotesize$[7.31,10.03]$} & $10$ & $0.3$\,W & $4/5$ & $0.91$ & $\mathbf{0.008}$ & $96.6\%$ \\
Exhaust fan II  & $62.3$ & $1.45$\,{\footnotesize$[1.09,2.07]$}  & $9$  & $0.2$\,W & $3/5$ & $0.62$ & $\mathbf{0.002}$ & $\mathbf{55.1\%}$ \\
Contactor I     & $63.7$ & $0.90$\,{\footnotesize$[0.71,1.12]$}  & $10$ & $0.3$\,W & $3/5$ & $0.31$ & $\mathbf{0.009}$ & $90.3\%$ \\
Contactor II    & $62.5$ & $1.15$\,{\footnotesize$[0.85,1.42]$}  & $9$  & $9.3$\,W & $3/5$ & $\mathbf{-0.05}$ & $\mathbf{0.007}$ & $85.4\%$ \\
\bottomrule
\end{tabular}
\end{table*}

\textbf{The two pelletizers and the two milling machines pass all five
physics checks.} Their STANDBY components sit at $4.2$--$5.4$\,kW with a
power factor of $0.11$--$0.17$ against $0.64$--$0.65$ when loaded, and a
no-load current ratio inside the motor band $[0.25,0.50]$. That is the
textbook signature of an energised induction motor drawing magnetising
current, recovered independently on four machines from unlabelled data,
and it is the strongest evidence in this paper that the state definition
is physical rather than fitted.

\textbf{On the fans and contactors there is no energised-idle state to
find, and the battery says so.} Their ``STANDBY'' components sit at
$0.2$--$9.3$\,W drawing $0.09$--$0.14$\,A: the model is splitting the
OFF band in two because it was told to find four states. The current
ratio is $0.002$--$0.009$ against an acceptance band starting at
$0.25$ on all four, and contactor~II additionally shows a
\emph{negative} power-factor separation. \textbf{This is the physics
battery doing the job it was built for.} The framework's state
definition is specific to motor-driven loads, and on the machines where
it does not apply it reports that rather than returning four confident
states regardless.

Two mechanisms behind those failures are worth naming, because both are
limitations of the method rather than of the data.

\emph{Ranking states by active power can misname them.} Exhaust fan~I's
``OFF'' component ($0.8\%$ of the record) carries $6.6$\,A at $0$\,W and
power factor $0$ --- standing reactive current with no real power ---
while its ``STANDBY'' component ($36.4\%$) sits at $0.3$\,W and
$0.13$\,A and is the machine's true de-energised state. The canonical
names are assigned by mean active power, so a near-zero-power component
carrying reactive current ranks below the genuine off state and takes
its name. No hour total changes, since both are non-productive, but it
is why the current-ratio check fails in the direction it does. The same
machine returns the largest forbidden-transition probability in the
plant ($3.4\times10^{-3}$, against $1.9\times10^{-4}$ on the four motor
machines), which is an independent measurement arriving at the same
conclusion (Section~\ref{sec:explain:transitions}).

\emph{The independent classifier fails where the model has nothing to
lock onto.} Otsu-thresholded power factor agrees with the state model on
$85$--$99.7\%$ of energised rows on seven machines and on $\mathbf{55.1\%}$
on exhaust fan~II --- the machine whose physics checks fail hardest.
Method~1's usefulness as an external check is thereby confirmed
\emph{and} its failure mode identified.

\subsection{Two results that generalise}
\label{sec:multimachine:general}

\textbf{The minimum-dwell result holds on all eight records.} Flicker
before the constraint ranges from $26.7\%$ (milling~II) to $77.7\%$
(exhaust fan~II), median $50.6\%$; after it, $0.00\%$ on seven machines
and $0.17\%$ on milling~I, at a cost of relabelling $0.15$--$6.1\%$ of
rows. What Section~\ref{sec:results:task4} established on one machine
--- that a transition prior cannot control flicker and an explicit
dwell constraint can --- now holds across duty cycles differing by an
order of magnitude.

\textbf{BIC does not select a state count here.} Widening the scan to
$k\in[2,8]$ returns $k=8$ on seven machines and $k=7$ on pelletizer~II,
with the criterion still falling at the boundary in seven of eight
cases; the improvement from $k=4$ to $k=8$ is $3.1$--$33.1\%$ of the
criterion's own value. This is a result about the criterion, not the
data (Section~\ref{sec:rw:states}). Fixing $k=4$ everywhere for
comparability is a choice, and this is the honest statement of what it
costs: \textbf{the physics vetoes are not a belt-and-braces addition on
top of BIC --- on records of this size they are what is actually doing
the selecting.}

\subsection{Which machines are hard, and why}
\label{sec:multimachine:difficulty}

Difficulty is measured on three axes: separability of the no-load/load
boundary (Cohen's $d$ and Bhattacharyya distance on log power), record
quality (coverage, spike rate, segment count), and agreement with the
independent classifier.

The ordering is not the one predicted. The expectation was that the
low-power auxiliary machines would be hardest. What the data show is
that the fans and contactors have the \emph{largest} separability
figures in the plant --- Cohen's $d$ of $9.3$ to $75.0$ against the
pelletizers' $4.3$ to $7.2$ --- while failing the physics checks,
because the separation being measured there is OFF against running, not
no-load against loaded. \textbf{Separability and validity are close to
orthogonal on this data}, which is the argument for keeping an external
physics battery rather than an internal cluster-quality score:
silhouette, BIC and Cohen's $d$ are all maximised by the partitions that
are physically wrong.

The two milling machines are hard for an unrelated reason. Their records
cover $12.2$ of $43$ days in $1\,082$ and $580$ fragments, against
$66$--$74$ fragments for the five-month records. They support two
one-week bootstrap blocks, which is not enough to form a usable
interval, and the intervals in Table~\ref{tab:multimachine} are printed
for those rows only to make that visible.
\textbf{Their $5.5$\,h/day --- the largest per-day standby figure among
the motor-driven machines --- must not be annualised}, and
Figure~\ref{fig:standby-hours} prints the block count inside each bar so
the figure cannot be read as the largest opportunity in the plant.

\begin{figure}[t]
  \centering
  \includegraphics[width=\linewidth]{fig_p6_standby_hours}
  \caption{STANDBY hours per covered day with $95\%$ weekly-block
  intervals. Bars are marked according to whether the machine passes the
  physics battery, and the number of weekly blocks is printed inside
  each bar; hatching marks an interval that rests on too few blocks to
  be usable.}
  \label{fig:standby-hours}
\end{figure}

\subsection{The decision layer exists on three machines of eight}
\label{sec:multimachine:decision}

A decision problem requires a measurable standby power and an observable
cold restart. Five machines fail an explicit degeneracy test rather than
an inspection: four because standby power is below $1$\,W, so the
standby cost per second is numerically zero, the break-even point is
infinite, and every currency figure would be an artefact of dividing by
zero; and pelletizer~II because in five months it is \emph{never once
observed resuming production from OFF}, so the cold-versus-warm contrast
cannot be measured at all. Two machines return a \emph{negative} restart
delay --- cold starts reaching production faster than warm ones --- which
is not a plausible physical reading and is a further symptom of the same
problem. The three machines that survive the test are the subject of
Table~\ref{tab:multimachine:decision}.

\begin{table}[t]
\centering
\caption{Machines on which the shutdown decision is well posed.
Head-room is the oracle saving; the proposed policy's figure is quoted
against the plant's own behaviour.}
\label{tab:multimachine:decision}
\small
\setlength{\tabcolsep}{4pt}
\begin{tabular}{@{}lrrrr@{}}
\toprule
Machine & energised idle & head-room & forecast MAE & proposed \\
        & (h, \% of idle) & (USD/yr) & (s) & (USD/yr) \\
\midrule
Pelletizer I & $88$ ($18.5\%$) & $37.6$ & $9\,034$ & $+2.9$ \\
Milling I    & $67$ ($48.8\%$) & $24.9$ & $18\,205$ & $\mathbf{-1\,041}$ \\
Milling II   & $64$ ($47.6\%$) & $25.2$ & $18\,440$ & $\mathbf{-1\,256}$ \\
\bottomrule
\end{tabular}
\end{table}

\textbf{The conjecture that motivated this experiment is refuted.} The
open question after the single-machine study was whether some other
IMDELD machine has enough recoverable standby to fund a retrofit ---
higher standby power, a larger energised-idle fraction, or a plant that
does not already shut down manually. None does. The three machines with
a well-posed problem have head-rooms of $37.6$, $24.9$ and
$25.2$\,USD/yr, so \textbf{the facility's total attainable saving from
standby recovery, summed over every machine where the question is even
meaningful, is under $90$\,USD/yr}. The framework's contribution on this
facility is decision quality, not recovered energy, and this paper says
so in its abstract rather than at the end of its discussion.

\textbf{The milling machines are worse than that, and the negative
result is reported as it stands.} The proposed policy \emph{loses}
$1\,041$ and $1\,256$\,USD/yr there against a head-room of
$25$\,USD/yr. The cause is visible in the same table: a duration
forecaster with MAE $18\,200$\,s and Spearman $0.34$--$0.42$, against
$9\,034$\,s and $0.58$ on pelletizer~I, fitted on a record broken into
$1\,082$ fragments. It over-predicts remaining idle time and fires
shutdowns the realised episode does not support. The decision layer
inherits its forecaster's quality, and on a record this fragmented the
forecaster is not good enough to act on --- which is a deployment
criterion, not merely an observation.

%  Figures are referenced by bare name; add to the preamble
%      \graphicspath{{../outputs/phase4/figures/}{../outputs/phase6/figures/}}
%  Bibliography: paper/references.bib
%  Numbers below are generated by  python run_phase4.py --only task8
%  and stored in outputs/phase4/task8/table_*.csv, with the intervals
%  from outputs/phase5/task11/state_block_bootstrap.csv.
% =====================================================================

\section{The Eight Machines}
\label{sec:multimachine}

Everything to this point is pelletizer~I. This section applies the
identical procedure --- one preprocessing, one labeller, $k=4$ fixed,
one physics battery --- to all eight IMDELD machines, and asks two
questions of the result: on which machines is the state definition
physically valid, and on which is there a shutdown decision to make at
all. The two answers are not the same, and neither follows the
prediction that motivated the experiment.

\subsection{Uniform labelling}
\label{sec:multimachine:labelling}

The Phase~I--III labels exist only for pelletizer~I and were produced by
a script whose configuration has since moved, so comparing that export
against seven fresh labellings would confound machine differences with
procedure differences. All eight records are therefore re-labelled
identically. Three departures from the original implementation were
required, each verified rather than assumed.

\emph{Blank currents.} The two milling-machine meters leave the current
channel empty on $36.1\%$ and $36.3\%$ of rows --- matching their OFF
shares almost exactly, i.e.\ the meter blanks the channel whenever the
machine is de-energised. On every such row the apparent power is exactly
$0$\,VA with voltage present, and $S=VI$ then admits only $I=0$\,A, so
the cells are filled with zero and a unit test asserts that every filled
cell met that condition. Leaving them missing makes a GMM
unfittable; dropping them would delete a third of each record precisely
where the machine is off, biasing every state-time total.

\emph{Timestamps and storage.} IMDELD spans a daylight-saving change, so
the timestamp column carries mixed UTC offsets; splitting the naive part
from the offset is exact and ten times faster than the fallback parse.
Records are stored as \texttt{float32}, whose one measurable consequence
is that about two rows per $10^{5}$ fall on the other side of the
spike threshold.

On pelletizer~I --- the one machine both labellings cover --- the
uniform labeller reproduces the pipeline's own export at $98.88\%$ row
agreement, $99.84\%$ on STANDBY-versus-rest, ARI $0.978$, and $89$\,h of
STANDBY against the export's $87$\,h. A number in this section that
differs from an earlier one is therefore a consequence of the analysis,
not of the labels.

Both of those figures are single-seed. The total this paper quotes for
pelletizer~I is the physics-conditional mean over passing seeds,
$92\pm3$\,h (Section~\ref{sec:method:accept_reject}); the single-seed
values sit inside that interval, and the per-machine rows of
Table~\ref{tab:multimachine} are reported per covered day with a
sampling interval for the same reason.

\subsection{Where the state definition holds}
\label{sec:multimachine:states}

\begin{table*}[!t]
\centering
\caption{Per-machine characterisation. STANDBY hours per covered day
carry a $95\%$ one-week moving-block bootstrap interval
(Section~\ref{sec:significance}); $n_{b}$ is the number of blocks the
record supports. $P_{\text{sb}}$ is the mean power of the inferred
STANDBY component. The no-load current ratio has an acceptance band of
$[0.25,0.50]$ from induction-motor theory. M1 is the fraction of
energised rows on which the independent Otsu power-factor classifier
agrees with the state model.}
\label{tab:multimachine}
\small
\begin{tabular}{@{}lrrrrrrrr@{}}
\toprule
Machine & cov.\ (d) & STANDBY (h/d) & $n_{b}$ & $P_{\text{sb}}$ & physics & PF sep. & $I$ ratio & M1 \\
\midrule
\multicolumn{9}{@{}l}{\textit{Motor-driven --- the state definition applies}}\\
Pelletizer I    & $63.4$ & $1.40$\,{\footnotesize$[1.15,1.63]$} & $10$ & $4.2$\,kW & $\mathbf{5/5}$ & $0.72$ & $0.36$ & $99.4\%$ \\
Pelletizer II   & $62.4$ & $1.85$\,{\footnotesize$[1.47,2.28]$} & $9$  & $4.5$\,kW & $\mathbf{5/5}$ & $0.71$ & $0.36$ & $99.7\%$ \\
Milling I       & $12.2$ & $5.52$\,{\footnotesize$[4.58,6.23]$} & $\mathbf{2}$ & $4.6$\,kW & $\mathbf{5/5}$ & $0.61$ & $0.43$ & $98.1\%$ \\
Milling II      & $12.2$ & $5.27$\,{\footnotesize$[4.44,5.89]$} & $\mathbf{2}$ & $5.4$\,kW & $\mathbf{5/5}$ & $0.52$ & $0.33$ & $97.7\%$ \\
\addlinespace
\multicolumn{9}{@{}l}{\textit{Auxiliary --- no energised-idle state exists to find}}\\
Exhaust fan I   & $63.7$ & $8.73$\,{\footnotesize$[7.31,10.03]$} & $10$ & $0.3$\,W & $4/5$ & $0.91$ & $\mathbf{0.008}$ & $96.6\%$ \\
Exhaust fan II  & $62.3$ & $1.45$\,{\footnotesize$[1.09,2.07]$}  & $9$  & $0.2$\,W & $3/5$ & $0.62$ & $\mathbf{0.002}$ & $\mathbf{55.1\%}$ \\
Contactor I     & $63.7$ & $0.90$\,{\footnotesize$[0.71,1.12]$}  & $10$ & $0.3$\,W & $3/5$ & $0.31$ & $\mathbf{0.009}$ & $90.3\%$ \\
Contactor II    & $62.5$ & $1.15$\,{\footnotesize$[0.85,1.42]$}  & $9$  & $9.3$\,W & $3/5$ & $\mathbf{-0.05}$ & $\mathbf{0.007}$ & $85.4\%$ \\
\bottomrule
\end{tabular}
\end{table*}

\textbf{The two pelletizers and the two milling machines pass all five
physics checks.} Their STANDBY components sit at $4.2$--$5.4$\,kW with a
power factor of $0.11$--$0.17$ against $0.64$--$0.65$ when loaded, and a
no-load current ratio inside the motor band $[0.25,0.50]$. That is the
textbook signature of an energised induction motor drawing magnetising
current, recovered independently on four machines from unlabelled data,
and it is the strongest evidence in this paper that the state definition
is physical rather than fitted.

\textbf{On the fans and contactors there is no energised-idle state to
find, and the battery says so.} Their ``STANDBY'' components sit at
$0.2$--$9.3$\,W drawing $0.09$--$0.14$\,A: the model is splitting the
OFF band in two because it was told to find four states. The current
ratio is $0.002$--$0.009$ against an acceptance band starting at
$0.25$ on all four, and contactor~II additionally shows a
\emph{negative} power-factor separation. \textbf{This is the physics
battery doing the job it was built for.} The framework's state
definition is specific to motor-driven loads, and on the machines where
it does not apply it reports that rather than returning four confident
states regardless.

Two mechanisms behind those failures are worth naming, because both are
limitations of the method rather than of the data.

\emph{Ranking states by active power can misname them.} Exhaust fan~I's
``OFF'' component ($0.8\%$ of the record) carries $6.6$\,A at $0$\,W and
power factor $0$ --- standing reactive current with no real power ---
while its ``STANDBY'' component ($36.4\%$) sits at $0.3$\,W and
$0.13$\,A and is the machine's true de-energised state. The canonical
names are assigned by mean active power, so a near-zero-power component
carrying reactive current ranks below the genuine off state and takes
its name. No hour total changes, since both are non-productive, but it
is why the current-ratio check fails in the direction it does. The same
machine returns the largest forbidden-transition probability in the
plant ($3.4\times10^{-3}$, against $1.9\times10^{-4}$ on the four motor
machines), which is an independent measurement arriving at the same
conclusion (Section~\ref{sec:explain:transitions}).

\emph{The independent classifier fails where the model has nothing to
lock onto.} Otsu-thresholded power factor agrees with the state model on
$85$--$99.7\%$ of energised rows on seven machines and on $\mathbf{55.1\%}$
on exhaust fan~II --- the machine whose physics checks fail hardest.
Method~1's usefulness as an external check is thereby confirmed
\emph{and} its failure mode identified.

\subsection{Two results that generalise}
\label{sec:multimachine:general}

\textbf{The minimum-dwell result holds on all eight records.} Flicker
before the constraint ranges from $26.7\%$ (milling~II) to $77.7\%$
(exhaust fan~II), median $50.6\%$; after it, $0.00\%$ on seven machines
and $0.17\%$ on milling~I, at a cost of relabelling $0.15$--$6.1\%$ of
rows. What Section~\ref{sec:results:task4} established on one machine
--- that a transition prior cannot control flicker and an explicit
dwell constraint can --- now holds across duty cycles differing by an
order of magnitude.

\textbf{BIC does not select a state count here.} Widening the scan to
$k\in[2,8]$ returns $k=8$ on seven machines and $k=7$ on pelletizer~II,
with the criterion still falling at the boundary in seven of eight
cases; the improvement from $k=4$ to $k=8$ is $3.1$--$33.1\%$ of the
criterion's own value. This is a result about the criterion, not the
data (Section~\ref{sec:rw:states}). Fixing $k=4$ everywhere for
comparability is a choice, and this is the honest statement of what it
costs: \textbf{the physics vetoes are not a belt-and-braces addition on
top of BIC --- on records of this size they are what is actually doing
the selecting.}

\subsection{Which machines are hard, and why}
\label{sec:multimachine:difficulty}

Difficulty is measured on three axes: separability of the no-load/load
boundary (Cohen's $d$ and Bhattacharyya distance on log power), record
quality (coverage, spike rate, segment count), and agreement with the
independent classifier.

The ordering is not the one predicted. The expectation was that the
low-power auxiliary machines would be hardest. What the data show is
that the fans and contactors have the \emph{largest} separability
figures in the plant --- Cohen's $d$ of $9.3$ to $75.0$ against the
pelletizers' $4.3$ to $7.2$ --- while failing the physics checks,
because the separation being measured there is OFF against running, not
no-load against loaded. \textbf{Separability and validity are close to
orthogonal on this data}, which is the argument for keeping an external
physics battery rather than an internal cluster-quality score:
silhouette, BIC and Cohen's $d$ are all maximised by the partitions that
are physically wrong.

The two milling machines are hard for an unrelated reason. Their records
cover $12.2$ of $43$ days in $1\,082$ and $580$ fragments, against
$66$--$74$ fragments for the five-month records. They support two
one-week bootstrap blocks, which is not enough to form a usable
interval, and the intervals in Table~\ref{tab:multimachine} are printed
for those rows only to make that visible.
\textbf{Their $5.5$\,h/day --- the largest per-day standby figure among
the motor-driven machines --- must not be annualised}, and
Figure~\ref{fig:standby-hours} prints the block count inside each bar so
the figure cannot be read as the largest opportunity in the plant.

\begin{figure}[t]
  \centering
  \includegraphics[width=\linewidth]{fig_p6_standby_hours}
  \caption{STANDBY hours per covered day with $95\%$ weekly-block
  intervals. Bars are marked according to whether the machine passes the
  physics battery, and the number of weekly blocks is printed inside
  each bar; hatching marks an interval that rests on too few blocks to
  be usable.}
  \label{fig:standby-hours}
\end{figure}

\subsection{The decision layer exists on three machines of eight}
\label{sec:multimachine:decision}

A decision problem requires a measurable standby power and an observable
cold restart. Five machines fail an explicit degeneracy test rather than
an inspection: four because standby power is below $1$\,W, so the
standby cost per second is numerically zero, the break-even point is
infinite, and every currency figure would be an artefact of dividing by
zero; and pelletizer~II because in five months it is \emph{never once
observed resuming production from OFF}, so the cold-versus-warm contrast
cannot be measured at all. Two machines return a \emph{negative} restart
delay --- cold starts reaching production faster than warm ones --- which
is not a plausible physical reading and is a further symptom of the same
problem. The three machines that survive the test are the subject of
Table~\ref{tab:multimachine:decision}.

\begin{table}[t]
\centering
\caption{Machines on which the shutdown decision is well posed.
Head-room is the oracle saving; the proposed policy's figure is quoted
against the plant's own behaviour.}
\label{tab:multimachine:decision}
\small
\setlength{\tabcolsep}{4pt}
\begin{tabular}{@{}lrrrr@{}}
\toprule
Machine & energised idle & head-room & forecast MAE & proposed \\
        & (h, \% of idle) & (USD/yr) & (s) & (USD/yr) \\
\midrule
Pelletizer I & $88$ ($18.5\%$) & $37.6$ & $9\,034$ & $+2.9$ \\
Milling I    & $67$ ($48.8\%$) & $24.9$ & $18\,205$ & $\mathbf{-1\,041}$ \\
Milling II   & $64$ ($47.6\%$) & $25.2$ & $18\,440$ & $\mathbf{-1\,256}$ \\
\bottomrule
\end{tabular}
\end{table}

\textbf{The conjecture that motivated this experiment is refuted.} The
open question after the single-machine study was whether some other
IMDELD machine has enough recoverable standby to fund a retrofit ---
higher standby power, a larger energised-idle fraction, or a plant that
does not already shut down manually. None does. The three machines with
a well-posed problem have head-rooms of $37.6$, $24.9$ and
$25.2$\,USD/yr, so \textbf{the facility's total attainable saving from
standby recovery, summed over every machine where the question is even
meaningful, is under $90$\,USD/yr}. The framework's contribution on this
facility is decision quality, not recovered energy, and this paper says
so in its abstract rather than at the end of its discussion.

\textbf{The milling machines are worse than that, and the negative
result is reported as it stands.} The proposed policy \emph{loses}
$1\,041$ and $1\,256$\,USD/yr there against a head-room of
$25$\,USD/yr. The cause is visible in the same table: a duration
forecaster with MAE $18\,200$\,s and Spearman $0.34$--$0.42$, against
$9\,034$\,s and $0.58$ on pelletizer~I, fitted on a record broken into
$1\,082$ fragments. It over-predicts remaining idle time and fires
shutdowns the realised episode does not support. The decision layer
inherits its forecaster's quality, and on a record this fragmented the
forecaster is not good enough to act on --- which is a deployment
criterion, not merely an observation.


% --- Transfer between machines and between sites
% =====================================================================
%  crossmachine.tex  --  PHASE IV, Tasks 9 and 10
%
%  Transfer: every machine pair in both directions for both forecasting
%  layers, and a second dataset from a different site including a
%  negative control.
%
%  Intended placement: Section 9 (\label{sec:crossmachine}), after the
%  eight-machine characterisation.
%
%  Drop-in usage: % =====================================================================
%  crossmachine.tex  --  PHASE IV, Tasks 9 and 10
%
%  Transfer: every machine pair in both directions for both forecasting
%  layers, and a second dataset from a different site including a
%  negative control.
%
%  Intended placement: Section 9 (\label{sec:crossmachine}), after the
%  eight-machine characterisation.
%
%  Drop-in usage: % =====================================================================
%  crossmachine.tex  --  PHASE IV, Tasks 9 and 10
%
%  Transfer: every machine pair in both directions for both forecasting
%  layers, and a second dataset from a different site including a
%  negative control.
%
%  Intended placement: Section 9 (\label{sec:crossmachine}), after the
%  eight-machine characterisation.
%
%  Drop-in usage: \input{crossmachine}
%  Figures are referenced by bare name; add to the preamble
%      \graphicspath{{../outputs/phase4/figures/}}
%  Bibliography: paper/references.bib
%  Numbers below are generated by  python run_phase4.py  and stored in
%    outputs/phase4/task9/part_a/{transfer_matrix.csv,part_a_summary.json}
%    outputs/phase4/task9/part_b/{transfer_matrix.csv,part_b_summary.json}
%    outputs/phase4/task10/{table_cross_dataset.csv,task10_results.json}
%
%  NOTE.  The calendar-share figure in Sec. "What the forecaster is
%  actually using" is the Shapley attribution (82.4%), NOT the split
%  gain (96.7%).  An earlier draft quoted 94.5% from gain; gain is
%  demonstrably misranking this model.  See explainability.tex
%  Sec.~\ref{sec:explain:decision} and PHASE6_NOTES correction 1.
% =====================================================================

\section{Transfer Between Machines and Between Sites}
\label{sec:crossmachine}

A framework fitted per machine is only as useful as the effort of
fitting it. This section asks what can be moved: every ordered pair of
the eight machines, in both directions, for each of the two forecasting
layers separately, and then a record from a different site and
acquisition system entirely. The two layers give opposite answers, and
the contrast between them is the result.

\subsection{The duration forecaster transfers almost perfectly}
\label{sec:crossmachine:decision}

The decision layer's remaining-duration forecaster is fitted on each
machine and evaluated on every machine, using the target's own
economics: $64$ cells, in two input variants.

\textbf{Transfer costs almost nothing.} Median MAE is $11\,118$\,s
within-machine, $11\,471$\,s between machines of the same type, and
$11\,653$\,s across types --- a $4.8\%$ degradation from the hardest
transfer available. On two targets a \emph{foreign} model is the better
one.

That is not a success. A model that moves this freely between a $40$\,kW
pelletizer and a contactor is not carrying machine-specific structure,
and the obvious hypothesis --- that the residual gap is a units mismatch
--- is testable and false. Rescaling the one absolute-power feature by
the target's own median productive power, a constant available from
unlabelled history, changes median cross-type MAE from $11\,653$\,s to
$11\,597$\,s ($0.5\%$) and \emph{lowers} median cross-type Spearman
correlation from $0.357$ to $0.287$.

\subsubsection*{What the forecaster is actually using}

The attribution of Section~\ref{sec:explain:decision} answers this
directly rather than by inference from the transfer matrix. Seven
calendar features carry $\mathbf{82.4\%}$ of the mean $|\phi|$ of the
duration forecaster; the machine's own production history carries
$13.3\%$ and elapsed idle time $4.3\%$.

\textbf{This figure supersedes the split-gain importance an earlier
version of this analysis reported.} Measured by split gain the same
features take $96.7\%$, and gain is demonstrably misranking this model:
the binary \texttt{cur\_is\_weekend} takes $63.2\%$ of the gain and
$5.1\%$ of the attribution, because it is split on near the root of a
partition that later splits then refine. The two measures agree only at
Spearman $\rho=0.72$. The corrected statement is that the forecaster is
\emph{four fifths} plant calendar, not nineteen twentieths --- and the
conclusion is unchanged and better supported: a model that is $82\%$
plant timetable transfers between any two machines sharing a shift
pattern, and to nothing else. The machine's own history at $13.3\%$ is
not negligible, so ``the forecaster is a clock'' is too strong; ``the
forecaster is four fifths clock'' is right.

\subsubsection*{In currency, on the one target where the problem is well
posed}

MAE is not what an operator buys. Table~\ref{tab:transfer:currency}
therefore re-scores the same transfer on the only machine in the plant
where the decision problem is well posed and the head-room is not
numerical noise.

\begin{table}[t]
\centering
\caption{Every source machine's duration forecaster evaluated on
pelletizer~I under that machine's own economics (S2), without
adaptation. Head-room is $37.6$\,USD/yr.}
\label{tab:transfer:currency}
\small
\begin{tabular}{@{}llrrc@{}}
\toprule
Source & relation & USD/yr & \% oracle & $>$ ski-rental \\
\midrule
Pelletizer II  & same type & $\mathbf{+8.07}$ & $21.4$ & yes \\
Exhaust fan I  & cross     & $+6.74$ & $17.9$ & yes \\
Pelletizer I   & itself    & $+2.86$ & $7.6$  & yes \\
Contactor II   & cross     & $-4.71$ & $-12.5$ & no \\
Contactor I    & cross     & $-5.80$ & $-15.4$ & no \\
Exhaust fan II & cross     & $-25.4$ & $-67.4$ & no \\
Milling II     & cross     & $-38.1$ & $-101$  & no \\
Milling I      & cross     & $-45.6$ & $-121$  & no \\
\bottomrule
\end{tabular}
\end{table}

The peer machine's forecaster beats the machine's own ($+8.07$ against
$+2.86$). With $1\,312$ episodes and a $30\%$ test split there is no
basis for calling that a real ordering, and it is not claimed as one:
the honest reading is that same-type transfer costs nothing measurable
and that both sit close to the forecast-free floor. Cross-type transfer
is a different matter --- five of six cross-type sources lose money on
this target, and across the whole matrix only $5.6\%$ of cross-type
pairs beat ski-rental against $33\%$ of within-machine pairs.

\textbf{The deployable conclusion} is that a duration forecaster may be
moved between machines of the same type without refitting, must not be
moved between types, and on this plant is worth little either way
against a rule that consults only a clock.

\subsection{The sequence forecaster does not transfer at all}
\label{sec:crossmachine:sequence}

The window state model is trained once per machine under the protocol of
Section~\ref{sec:forecasting} and evaluated on every machine's held-out
windows, in two scaler variants, with a persistence reference computed
on each target. Table~\ref{tab:transfer:sequence} summarises the
resulting $8\times8$ matrix by relation.

\begin{table}[t]
\centering
\caption{Sequence-model transfer, median over the $8\times8$ matrix.
The persistence row needs no training on any machine.}
\label{tab:transfer:sequence}
\small
\begin{tabular}{@{}lrrr@{}}
\toprule
 & macro $F_1$ & $F_1$ (STANDBY) & MAE (s) \\
\midrule
trained on the target      & $\mathbf{0.695}$ & $\mathbf{0.539}$ & $\mathbf{16.4}$ \\
same-type transfer         & $0.511$ & $0.479$ & $36.7$ \\
cross-type (source scaler) & $0.286$ & $0.076$ & $70.8$ \\
cross-type (target scaler) & $0.374$ & $0.149$ & $44.2$ \\
\textsc{persistence}       & $0.629$ & $0.463$ & $28.8$ \\
\bottomrule
\end{tabular}
\end{table}

\begin{figure*}[t]
\centering
\includegraphics[width=\textwidth]{fig_p7_transfer}
\caption{The two transfer matrices on one axis convention: rows are the
machine a model was fitted on, columns the machine it was evaluated on,
and the outlined diagonal is the locally fitted model.
\textbf{(A)} The decision layer's duration forecaster, raw variant --- no
target information of any kind, the honest zero-shot condition. Values
vary down each column far less than across the row, i.e.\ the error is a
property of the target and barely of the source. \textbf{(B)} The
sequence state model under its source scaler: a bright diagonal on a
dark field. \textbf{What transfers here is not machine-specific, and
what is machine-specific does not transfer.}}
\label{fig:transfer}
\end{figure*}

\textbf{The contrast with the previous subsection is the point}, and
Figure~\ref{fig:transfer} is that contrast in one frame. The
duration forecaster moved between machines at a $4.8\%$ cost in MAE; the
sequence model loses $86\%$ of its STANDBY $F_1$ crossing machine types.
The two results are consistent rather than contradictory, and the
attributions of Section~\ref{sec:explain} say why: what the sequence
model learns \emph{is} machine-specific --- the electrical shape of one
motor's duty cycle, $11.6\%$ calendar --- while what the duration
forecaster learns is the plant calendar at $82.4\%$, which every machine
in the facility shares. \textbf{What transfers here is not
machine-specific, and what is machine-specific does not transfer.}
The paired test over the eight machines puts the transfer degradation at
$0.384$ of $F_1$, $p=0.008$ at the floor of the test, $d=+1.75$
(Section~\ref{sec:significance:cross}) --- the one cross-machine claim
here that survives testing outright.

Re-standardising to the target helps cross-type transfer ($F_1$
$0.076\rightarrow0.149$) and hurts same-type transfer
($0.479\rightarrow0.305$): a source model applied to a machine of the
same type is better off keeping the scale it was trained on, since the
two units genuinely have similar ratings, while across types it is
better off being told the new machine's scale. Neither variant comes
close to a locally fitted model.

\begin{table}[t]
\centering
\caption{Per-machine sequence forecasting, STANDBY $F_1$. The
\textsc{persistence} column requires no training.}
\label{tab:transfer:persistence}
\small
\begin{tabular}{@{}lrrr@{}}
\toprule
Target & own model & \textsc{persistence} & transferred (med.) \\
\midrule
Pelletizer I   & $0.660$ & $0.650$ & $0.20$ \\
Pelletizer II  & $0.607$ & $0.432$ & $0.43$ \\
Milling I      & $0.515$ & $0.504$ & $0.20$ \\
Milling II     & $0.489$ & $\mathbf{0.493}$ & $0.25$ \\
Exhaust fan I  & $\mathbf{0.982}$ & $0.288$ & $0.11$ \\
Exhaust fan II & $0.477$ & $0.211$ & $0.01$ \\
Contactor I    & $0.402$ & $0.012$ & $0.08$ \\
Contactor II   & $0.563$ & $0.552$ & $0.12$ \\
\bottomrule
\end{tabular}
\end{table}

Table~\ref{tab:transfer:persistence} carries the uncomfortable result of
this section. \textbf{On pelletizer~I --- the machine the rest of this
paper is built on --- the sequence model beats ``assume nothing changes
for the next five minutes'' by $0.010$ of $F_1$}, and on milling~II
persistence wins outright. The trained model earns its place only on the
machines whose STANDBY state is easy for the wrong reason: exhaust
fan~I, where STANDBY \emph{is} the off state
(Section~\ref{sec:multimachine:states}) and $0.98$ is achievable. Only
$18$--$21\%$ of transferred pairs beat the target's own persistence
baseline, against $87.5\%$ of locally trained models. Averaged over the
eight machines the training gain is $+0.194$ of $F_1$ and significant,
but it is carried by the auxiliary plant; on the four motor-driven
machines it is $+0.048$ and, at $n=4$, untestable by construction
(Section~\ref{sec:significance:cross}).

\subsection{A second site, and a control that the battery nearly passes}
\label{sec:crossmachine:dataset}

The eight machines share a plant, a calendar and an acquisition system.
To test the framework outside all three, it is applied unchanged to two
single-channel active-power records at $5$\,s resolution from a
different site: a CNC machining centre ($348$ covered days) and the AC
output of a photovoltaic inverter ($86$ covered days). \textbf{The
second is a negative control} --- a generator, not a driven load, for
which no shutdown decision exists --- included because a validation
battery is only worth quoting if it can reject something.
Table~\ref{tab:crossdataset} reports what the unmodified pipeline
returns on each.

\begin{table}[t]
\centering
\caption{Cross-dataset application. Two of the five physics checks
require current and power-factor channels and cannot be evaluated on a
single-channel record.}
\label{tab:crossdataset}
\small
\begin{tabular}{@{}lrr@{}}
\toprule
 & CNC (machine) & PV inverter (control) \\
\midrule
physics checks evaluable & $3/5$ & $3/5$ \\
checks passed            & $\mathbf{3/3}$ & $2/3$ \\
flicker                  & $0.00\%$ & $0.00\%$ \\
median dwell             & $120$\,s & $30$\,s \\
inferred STANDBY         & $186$\,h @ $179$\,W & $42$\,h @ $444$\,W \\
productive at night      & $7.9\%$ & $\mathbf{0.0\%}$ \\
productive at midday     & $34.2\%$ & $54.1\%$ \\
\bottomrule
\end{tabular}
\end{table}

\textbf{The pipeline runs unchanged on both and returns a
plausible-looking four-state partition for both.} On the CNC that is the
right answer: OFF at $0$\,W, STANDBY at $179$\,W, WORKING at $2.24$\,kW,
PEAK at $5.70$\,kW, with the diurnal profile of a day-shift machine.

\textbf{On a solar inverter it is the wrong answer, and the battery
nearly fails to say so.} Power-factor separation and no-load current
ratio are undefined without current and power-factor channels --- and
those are precisely the two checks that do not depend on power
magnitude, i.e.\ the two that test whether STANDBY has been separated
from a lightly loaded WORKING state rather than merely from a
lower-power one. Of the three that remain the control passes power
ordering and OFF-near-zero, and fails only variance ordering, by a $2\%$
margin ($\sigma_{\text{OFF}}=234$\,W against
$\sigma_{\text{STANDBY}}=229$\,W). The framework reports $42$\,h of
recoverable ``standby'' on a device that cannot be shut down.

\begin{figure*}[t]
\centering
\includegraphics[width=\textwidth]{fig_p7_cross_dataset}
\caption{Why the negative control nearly passes. \textbf{(A)} The five
physics checks on the four-channel reference machine and on the two
single-channel second-site records. C3 and C4 are the only two checks
independent of power magnitude --- hence the only two that can separate
energised idle from a lightly loaded productive state --- and they are
exactly the two that a single-channel record cannot evaluate. The
control fails on the one remaining discriminating check, and by a $2\%$
margin. \textbf{(B)} What the battery is left with: the photovoltaic
inverter is productive $0.0\%$ of the time at night, which no driven
load is. That the rejection rests on a diurnal signal rather than a
physical one is the finding.}
\label{fig:crossdataset-checks}
\end{figure*}

\textbf{This is the central cross-dataset finding, and it is a
limitation of the method rather than of the data}
(Figure~\ref{fig:crossdataset-checks}). On a single-channel
record the state model still produces states and the decision layer
still consumes them, but the physics validation has lost most of its
discriminating power --- and it is the physics validation that
Section~\ref{sec:significance:seeds} showed to be load-bearing for
reproducibility.

Nor does the diurnal profile rescue it. Both records are daytime-active,
and on the midday figure the inverter looks like the \emph{better}
machine ($54.1\%$ productive against the CNC's $34.2\%$); only the night
figure separates them, and then by $7.9$ points. What actually differs
is the \emph{shape} --- the CNC's profile is a plateau with sharp edges,
the inverter's a smooth bell centred on solar noon --- and no test in the
battery looks at shape. A single-channel deployment therefore needs an
external admissibility check, \emph{is this a driven load?}, which the
framework as it stands does not provide.

Two further results follow.

\textbf{The decision layer cannot be instantiated on either record.}
Both return a negative cold-versus-warm restart delay, so the same guard
that flagged five IMDELD machines
(Section~\ref{sec:multimachine:decision}) flags both of these. On the
CNC the reason is visible in the episode structure: $229$ usable idle
episodes with a median duration of $14.9$\,h, against pelletizer~I's
$90$\,s --- a machine that is off far more than it is on, and is rarely
observed restarting at all.

\textbf{Zero-shot transfer across sites fails outright.} The
pelletizer~I duration forecaster scores MAE $120\,955$\,s and Spearman
$\mathbf{-0.245}$ on the CNC record, against $60\,386$\,s and $+0.659$
for a locally fitted model; rescaling does not help ($-0.234$). A
\emph{negative} rank correlation means the transferred model orders the
CNC's idle episodes worse than chance. This is the predicted consequence
of Section~\ref{sec:crossmachine:decision}: what that forecaster carries
is the IMDELD plant calendar, and another site does not keep IMDELD's
shifts.

%  Figures are referenced by bare name; add to the preamble
%      \graphicspath{{../outputs/phase4/figures/}}
%  Bibliography: paper/references.bib
%  Numbers below are generated by  python run_phase4.py  and stored in
%    outputs/phase4/task9/part_a/{transfer_matrix.csv,part_a_summary.json}
%    outputs/phase4/task9/part_b/{transfer_matrix.csv,part_b_summary.json}
%    outputs/phase4/task10/{table_cross_dataset.csv,task10_results.json}
%
%  NOTE.  The calendar-share figure in Sec. "What the forecaster is
%  actually using" is the Shapley attribution (82.4%), NOT the split
%  gain (96.7%).  An earlier draft quoted 94.5% from gain; gain is
%  demonstrably misranking this model.  See explainability.tex
%  Sec.~\ref{sec:explain:decision} and PHASE6_NOTES correction 1.
% =====================================================================

\section{Transfer Between Machines and Between Sites}
\label{sec:crossmachine}

A framework fitted per machine is only as useful as the effort of
fitting it. This section asks what can be moved: every ordered pair of
the eight machines, in both directions, for each of the two forecasting
layers separately, and then a record from a different site and
acquisition system entirely. The two layers give opposite answers, and
the contrast between them is the result.

\subsection{The duration forecaster transfers almost perfectly}
\label{sec:crossmachine:decision}

The decision layer's remaining-duration forecaster is fitted on each
machine and evaluated on every machine, using the target's own
economics: $64$ cells, in two input variants.

\textbf{Transfer costs almost nothing.} Median MAE is $11\,118$\,s
within-machine, $11\,471$\,s between machines of the same type, and
$11\,653$\,s across types --- a $4.8\%$ degradation from the hardest
transfer available. On two targets a \emph{foreign} model is the better
one.

That is not a success. A model that moves this freely between a $40$\,kW
pelletizer and a contactor is not carrying machine-specific structure,
and the obvious hypothesis --- that the residual gap is a units mismatch
--- is testable and false. Rescaling the one absolute-power feature by
the target's own median productive power, a constant available from
unlabelled history, changes median cross-type MAE from $11\,653$\,s to
$11\,597$\,s ($0.5\%$) and \emph{lowers} median cross-type Spearman
correlation from $0.357$ to $0.287$.

\subsubsection*{What the forecaster is actually using}

The attribution of Section~\ref{sec:explain:decision} answers this
directly rather than by inference from the transfer matrix. Seven
calendar features carry $\mathbf{82.4\%}$ of the mean $|\phi|$ of the
duration forecaster; the machine's own production history carries
$13.3\%$ and elapsed idle time $4.3\%$.

\textbf{This figure supersedes the split-gain importance an earlier
version of this analysis reported.} Measured by split gain the same
features take $96.7\%$, and gain is demonstrably misranking this model:
the binary \texttt{cur\_is\_weekend} takes $63.2\%$ of the gain and
$5.1\%$ of the attribution, because it is split on near the root of a
partition that later splits then refine. The two measures agree only at
Spearman $\rho=0.72$. The corrected statement is that the forecaster is
\emph{four fifths} plant calendar, not nineteen twentieths --- and the
conclusion is unchanged and better supported: a model that is $82\%$
plant timetable transfers between any two machines sharing a shift
pattern, and to nothing else. The machine's own history at $13.3\%$ is
not negligible, so ``the forecaster is a clock'' is too strong; ``the
forecaster is four fifths clock'' is right.

\subsubsection*{In currency, on the one target where the problem is well
posed}

MAE is not what an operator buys. Table~\ref{tab:transfer:currency}
therefore re-scores the same transfer on the only machine in the plant
where the decision problem is well posed and the head-room is not
numerical noise.

\begin{table}[t]
\centering
\caption{Every source machine's duration forecaster evaluated on
pelletizer~I under that machine's own economics (S2), without
adaptation. Head-room is $37.6$\,USD/yr.}
\label{tab:transfer:currency}
\small
\begin{tabular}{@{}llrrc@{}}
\toprule
Source & relation & USD/yr & \% oracle & $>$ ski-rental \\
\midrule
Pelletizer II  & same type & $\mathbf{+8.07}$ & $21.4$ & yes \\
Exhaust fan I  & cross     & $+6.74$ & $17.9$ & yes \\
Pelletizer I   & itself    & $+2.86$ & $7.6$  & yes \\
Contactor II   & cross     & $-4.71$ & $-12.5$ & no \\
Contactor I    & cross     & $-5.80$ & $-15.4$ & no \\
Exhaust fan II & cross     & $-25.4$ & $-67.4$ & no \\
Milling II     & cross     & $-38.1$ & $-101$  & no \\
Milling I      & cross     & $-45.6$ & $-121$  & no \\
\bottomrule
\end{tabular}
\end{table}

The peer machine's forecaster beats the machine's own ($+8.07$ against
$+2.86$). With $1\,312$ episodes and a $30\%$ test split there is no
basis for calling that a real ordering, and it is not claimed as one:
the honest reading is that same-type transfer costs nothing measurable
and that both sit close to the forecast-free floor. Cross-type transfer
is a different matter --- five of six cross-type sources lose money on
this target, and across the whole matrix only $5.6\%$ of cross-type
pairs beat ski-rental against $33\%$ of within-machine pairs.

\textbf{The deployable conclusion} is that a duration forecaster may be
moved between machines of the same type without refitting, must not be
moved between types, and on this plant is worth little either way
against a rule that consults only a clock.

\subsection{The sequence forecaster does not transfer at all}
\label{sec:crossmachine:sequence}

The window state model is trained once per machine under the protocol of
Section~\ref{sec:forecasting} and evaluated on every machine's held-out
windows, in two scaler variants, with a persistence reference computed
on each target. Table~\ref{tab:transfer:sequence} summarises the
resulting $8\times8$ matrix by relation.

\begin{table}[t]
\centering
\caption{Sequence-model transfer, median over the $8\times8$ matrix.
The persistence row needs no training on any machine.}
\label{tab:transfer:sequence}
\small
\begin{tabular}{@{}lrrr@{}}
\toprule
 & macro $F_1$ & $F_1$ (STANDBY) & MAE (s) \\
\midrule
trained on the target      & $\mathbf{0.695}$ & $\mathbf{0.539}$ & $\mathbf{16.4}$ \\
same-type transfer         & $0.511$ & $0.479$ & $36.7$ \\
cross-type (source scaler) & $0.286$ & $0.076$ & $70.8$ \\
cross-type (target scaler) & $0.374$ & $0.149$ & $44.2$ \\
\textsc{persistence}       & $0.629$ & $0.463$ & $28.8$ \\
\bottomrule
\end{tabular}
\end{table}

\begin{figure*}[t]
\centering
\includegraphics[width=\textwidth]{fig_p7_transfer}
\caption{The two transfer matrices on one axis convention: rows are the
machine a model was fitted on, columns the machine it was evaluated on,
and the outlined diagonal is the locally fitted model.
\textbf{(A)} The decision layer's duration forecaster, raw variant --- no
target information of any kind, the honest zero-shot condition. Values
vary down each column far less than across the row, i.e.\ the error is a
property of the target and barely of the source. \textbf{(B)} The
sequence state model under its source scaler: a bright diagonal on a
dark field. \textbf{What transfers here is not machine-specific, and
what is machine-specific does not transfer.}}
\label{fig:transfer}
\end{figure*}

\textbf{The contrast with the previous subsection is the point}, and
Figure~\ref{fig:transfer} is that contrast in one frame. The
duration forecaster moved between machines at a $4.8\%$ cost in MAE; the
sequence model loses $86\%$ of its STANDBY $F_1$ crossing machine types.
The two results are consistent rather than contradictory, and the
attributions of Section~\ref{sec:explain} say why: what the sequence
model learns \emph{is} machine-specific --- the electrical shape of one
motor's duty cycle, $11.6\%$ calendar --- while what the duration
forecaster learns is the plant calendar at $82.4\%$, which every machine
in the facility shares. \textbf{What transfers here is not
machine-specific, and what is machine-specific does not transfer.}
The paired test over the eight machines puts the transfer degradation at
$0.384$ of $F_1$, $p=0.008$ at the floor of the test, $d=+1.75$
(Section~\ref{sec:significance:cross}) --- the one cross-machine claim
here that survives testing outright.

Re-standardising to the target helps cross-type transfer ($F_1$
$0.076\rightarrow0.149$) and hurts same-type transfer
($0.479\rightarrow0.305$): a source model applied to a machine of the
same type is better off keeping the scale it was trained on, since the
two units genuinely have similar ratings, while across types it is
better off being told the new machine's scale. Neither variant comes
close to a locally fitted model.

\begin{table}[t]
\centering
\caption{Per-machine sequence forecasting, STANDBY $F_1$. The
\textsc{persistence} column requires no training.}
\label{tab:transfer:persistence}
\small
\begin{tabular}{@{}lrrr@{}}
\toprule
Target & own model & \textsc{persistence} & transferred (med.) \\
\midrule
Pelletizer I   & $0.660$ & $0.650$ & $0.20$ \\
Pelletizer II  & $0.607$ & $0.432$ & $0.43$ \\
Milling I      & $0.515$ & $0.504$ & $0.20$ \\
Milling II     & $0.489$ & $\mathbf{0.493}$ & $0.25$ \\
Exhaust fan I  & $\mathbf{0.982}$ & $0.288$ & $0.11$ \\
Exhaust fan II & $0.477$ & $0.211$ & $0.01$ \\
Contactor I    & $0.402$ & $0.012$ & $0.08$ \\
Contactor II   & $0.563$ & $0.552$ & $0.12$ \\
\bottomrule
\end{tabular}
\end{table}

Table~\ref{tab:transfer:persistence} carries the uncomfortable result of
this section. \textbf{On pelletizer~I --- the machine the rest of this
paper is built on --- the sequence model beats ``assume nothing changes
for the next five minutes'' by $0.010$ of $F_1$}, and on milling~II
persistence wins outright. The trained model earns its place only on the
machines whose STANDBY state is easy for the wrong reason: exhaust
fan~I, where STANDBY \emph{is} the off state
(Section~\ref{sec:multimachine:states}) and $0.98$ is achievable. Only
$18$--$21\%$ of transferred pairs beat the target's own persistence
baseline, against $87.5\%$ of locally trained models. Averaged over the
eight machines the training gain is $+0.194$ of $F_1$ and significant,
but it is carried by the auxiliary plant; on the four motor-driven
machines it is $+0.048$ and, at $n=4$, untestable by construction
(Section~\ref{sec:significance:cross}).

\subsection{A second site, and a control that the battery nearly passes}
\label{sec:crossmachine:dataset}

The eight machines share a plant, a calendar and an acquisition system.
To test the framework outside all three, it is applied unchanged to two
single-channel active-power records at $5$\,s resolution from a
different site: a CNC machining centre ($348$ covered days) and the AC
output of a photovoltaic inverter ($86$ covered days). \textbf{The
second is a negative control} --- a generator, not a driven load, for
which no shutdown decision exists --- included because a validation
battery is only worth quoting if it can reject something.
Table~\ref{tab:crossdataset} reports what the unmodified pipeline
returns on each.

\begin{table}[t]
\centering
\caption{Cross-dataset application. Two of the five physics checks
require current and power-factor channels and cannot be evaluated on a
single-channel record.}
\label{tab:crossdataset}
\small
\begin{tabular}{@{}lrr@{}}
\toprule
 & CNC (machine) & PV inverter (control) \\
\midrule
physics checks evaluable & $3/5$ & $3/5$ \\
checks passed            & $\mathbf{3/3}$ & $2/3$ \\
flicker                  & $0.00\%$ & $0.00\%$ \\
median dwell             & $120$\,s & $30$\,s \\
inferred STANDBY         & $186$\,h @ $179$\,W & $42$\,h @ $444$\,W \\
productive at night      & $7.9\%$ & $\mathbf{0.0\%}$ \\
productive at midday     & $34.2\%$ & $54.1\%$ \\
\bottomrule
\end{tabular}
\end{table}

\textbf{The pipeline runs unchanged on both and returns a
plausible-looking four-state partition for both.} On the CNC that is the
right answer: OFF at $0$\,W, STANDBY at $179$\,W, WORKING at $2.24$\,kW,
PEAK at $5.70$\,kW, with the diurnal profile of a day-shift machine.

\textbf{On a solar inverter it is the wrong answer, and the battery
nearly fails to say so.} Power-factor separation and no-load current
ratio are undefined without current and power-factor channels --- and
those are precisely the two checks that do not depend on power
magnitude, i.e.\ the two that test whether STANDBY has been separated
from a lightly loaded WORKING state rather than merely from a
lower-power one. Of the three that remain the control passes power
ordering and OFF-near-zero, and fails only variance ordering, by a $2\%$
margin ($\sigma_{\text{OFF}}=234$\,W against
$\sigma_{\text{STANDBY}}=229$\,W). The framework reports $42$\,h of
recoverable ``standby'' on a device that cannot be shut down.

\begin{figure*}[t]
\centering
\includegraphics[width=\textwidth]{fig_p7_cross_dataset}
\caption{Why the negative control nearly passes. \textbf{(A)} The five
physics checks on the four-channel reference machine and on the two
single-channel second-site records. C3 and C4 are the only two checks
independent of power magnitude --- hence the only two that can separate
energised idle from a lightly loaded productive state --- and they are
exactly the two that a single-channel record cannot evaluate. The
control fails on the one remaining discriminating check, and by a $2\%$
margin. \textbf{(B)} What the battery is left with: the photovoltaic
inverter is productive $0.0\%$ of the time at night, which no driven
load is. That the rejection rests on a diurnal signal rather than a
physical one is the finding.}
\label{fig:crossdataset-checks}
\end{figure*}

\textbf{This is the central cross-dataset finding, and it is a
limitation of the method rather than of the data}
(Figure~\ref{fig:crossdataset-checks}). On a single-channel
record the state model still produces states and the decision layer
still consumes them, but the physics validation has lost most of its
discriminating power --- and it is the physics validation that
Section~\ref{sec:significance:seeds} showed to be load-bearing for
reproducibility.

Nor does the diurnal profile rescue it. Both records are daytime-active,
and on the midday figure the inverter looks like the \emph{better}
machine ($54.1\%$ productive against the CNC's $34.2\%$); only the night
figure separates them, and then by $7.9$ points. What actually differs
is the \emph{shape} --- the CNC's profile is a plateau with sharp edges,
the inverter's a smooth bell centred on solar noon --- and no test in the
battery looks at shape. A single-channel deployment therefore needs an
external admissibility check, \emph{is this a driven load?}, which the
framework as it stands does not provide.

Two further results follow.

\textbf{The decision layer cannot be instantiated on either record.}
Both return a negative cold-versus-warm restart delay, so the same guard
that flagged five IMDELD machines
(Section~\ref{sec:multimachine:decision}) flags both of these. On the
CNC the reason is visible in the episode structure: $229$ usable idle
episodes with a median duration of $14.9$\,h, against pelletizer~I's
$90$\,s --- a machine that is off far more than it is on, and is rarely
observed restarting at all.

\textbf{Zero-shot transfer across sites fails outright.} The
pelletizer~I duration forecaster scores MAE $120\,955$\,s and Spearman
$\mathbf{-0.245}$ on the CNC record, against $60\,386$\,s and $+0.659$
for a locally fitted model; rescaling does not help ($-0.234$). A
\emph{negative} rank correlation means the transferred model orders the
CNC's idle episodes worse than chance. This is the predicted consequence
of Section~\ref{sec:crossmachine:decision}: what that forecaster carries
is the IMDELD plant calendar, and another site does not keep IMDELD's
shifts.

%  Figures are referenced by bare name; add to the preamble
%      \graphicspath{{../outputs/phase4/figures/}}
%  Bibliography: paper/references.bib
%  Numbers below are generated by  python run_phase4.py  and stored in
%    outputs/phase4/task9/part_a/{transfer_matrix.csv,part_a_summary.json}
%    outputs/phase4/task9/part_b/{transfer_matrix.csv,part_b_summary.json}
%    outputs/phase4/task10/{table_cross_dataset.csv,task10_results.json}
%
%  NOTE.  The calendar-share figure in Sec. "What the forecaster is
%  actually using" is the Shapley attribution (82.4%), NOT the split
%  gain (96.7%).  An earlier draft quoted 94.5% from gain; gain is
%  demonstrably misranking this model.  See explainability.tex
%  Sec.~\ref{sec:explain:decision} and PHASE6_NOTES correction 1.
% =====================================================================

\section{Transfer Between Machines and Between Sites}
\label{sec:crossmachine}

A framework fitted per machine is only as useful as the effort of
fitting it. This section asks what can be moved: every ordered pair of
the eight machines, in both directions, for each of the two forecasting
layers separately, and then a record from a different site and
acquisition system entirely. The two layers give opposite answers, and
the contrast between them is the result.

\subsection{The duration forecaster transfers almost perfectly}
\label{sec:crossmachine:decision}

The decision layer's remaining-duration forecaster is fitted on each
machine and evaluated on every machine, using the target's own
economics: $64$ cells, in two input variants.

\textbf{Transfer costs almost nothing.} Median MAE is $11\,118$\,s
within-machine, $11\,471$\,s between machines of the same type, and
$11\,653$\,s across types --- a $4.8\%$ degradation from the hardest
transfer available. On two targets a \emph{foreign} model is the better
one.

That is not a success. A model that moves this freely between a $40$\,kW
pelletizer and a contactor is not carrying machine-specific structure,
and the obvious hypothesis --- that the residual gap is a units mismatch
--- is testable and false. Rescaling the one absolute-power feature by
the target's own median productive power, a constant available from
unlabelled history, changes median cross-type MAE from $11\,653$\,s to
$11\,597$\,s ($0.5\%$) and \emph{lowers} median cross-type Spearman
correlation from $0.357$ to $0.287$.

\subsubsection*{What the forecaster is actually using}

The attribution of Section~\ref{sec:explain:decision} answers this
directly rather than by inference from the transfer matrix. Seven
calendar features carry $\mathbf{82.4\%}$ of the mean $|\phi|$ of the
duration forecaster; the machine's own production history carries
$13.3\%$ and elapsed idle time $4.3\%$.

\textbf{This figure supersedes the split-gain importance an earlier
version of this analysis reported.} Measured by split gain the same
features take $96.7\%$, and gain is demonstrably misranking this model:
the binary \texttt{cur\_is\_weekend} takes $63.2\%$ of the gain and
$5.1\%$ of the attribution, because it is split on near the root of a
partition that later splits then refine. The two measures agree only at
Spearman $\rho=0.72$. The corrected statement is that the forecaster is
\emph{four fifths} plant calendar, not nineteen twentieths --- and the
conclusion is unchanged and better supported: a model that is $82\%$
plant timetable transfers between any two machines sharing a shift
pattern, and to nothing else. The machine's own history at $13.3\%$ is
not negligible, so ``the forecaster is a clock'' is too strong; ``the
forecaster is four fifths clock'' is right.

\subsubsection*{In currency, on the one target where the problem is well
posed}

MAE is not what an operator buys. Table~\ref{tab:transfer:currency}
therefore re-scores the same transfer on the only machine in the plant
where the decision problem is well posed and the head-room is not
numerical noise.

\begin{table}[t]
\centering
\caption{Every source machine's duration forecaster evaluated on
pelletizer~I under that machine's own economics (S2), without
adaptation. Head-room is $37.6$\,USD/yr.}
\label{tab:transfer:currency}
\small
\begin{tabular}{@{}llrrc@{}}
\toprule
Source & relation & USD/yr & \% oracle & $>$ ski-rental \\
\midrule
Pelletizer II  & same type & $\mathbf{+8.07}$ & $21.4$ & yes \\
Exhaust fan I  & cross     & $+6.74$ & $17.9$ & yes \\
Pelletizer I   & itself    & $+2.86$ & $7.6$  & yes \\
Contactor II   & cross     & $-4.71$ & $-12.5$ & no \\
Contactor I    & cross     & $-5.80$ & $-15.4$ & no \\
Exhaust fan II & cross     & $-25.4$ & $-67.4$ & no \\
Milling II     & cross     & $-38.1$ & $-101$  & no \\
Milling I      & cross     & $-45.6$ & $-121$  & no \\
\bottomrule
\end{tabular}
\end{table}

The peer machine's forecaster beats the machine's own ($+8.07$ against
$+2.86$). With $1\,312$ episodes and a $30\%$ test split there is no
basis for calling that a real ordering, and it is not claimed as one:
the honest reading is that same-type transfer costs nothing measurable
and that both sit close to the forecast-free floor. Cross-type transfer
is a different matter --- five of six cross-type sources lose money on
this target, and across the whole matrix only $5.6\%$ of cross-type
pairs beat ski-rental against $33\%$ of within-machine pairs.

\textbf{The deployable conclusion} is that a duration forecaster may be
moved between machines of the same type without refitting, must not be
moved between types, and on this plant is worth little either way
against a rule that consults only a clock.

\subsection{The sequence forecaster does not transfer at all}
\label{sec:crossmachine:sequence}

The window state model is trained once per machine under the protocol of
Section~\ref{sec:forecasting} and evaluated on every machine's held-out
windows, in two scaler variants, with a persistence reference computed
on each target. Table~\ref{tab:transfer:sequence} summarises the
resulting $8\times8$ matrix by relation.

\begin{table}[t]
\centering
\caption{Sequence-model transfer, median over the $8\times8$ matrix.
The persistence row needs no training on any machine.}
\label{tab:transfer:sequence}
\small
\begin{tabular}{@{}lrrr@{}}
\toprule
 & macro $F_1$ & $F_1$ (STANDBY) & MAE (s) \\
\midrule
trained on the target      & $\mathbf{0.695}$ & $\mathbf{0.539}$ & $\mathbf{16.4}$ \\
same-type transfer         & $0.511$ & $0.479$ & $36.7$ \\
cross-type (source scaler) & $0.286$ & $0.076$ & $70.8$ \\
cross-type (target scaler) & $0.374$ & $0.149$ & $44.2$ \\
\textsc{persistence}       & $0.629$ & $0.463$ & $28.8$ \\
\bottomrule
\end{tabular}
\end{table}

\begin{figure*}[t]
\centering
\includegraphics[width=\textwidth]{fig_p7_transfer}
\caption{The two transfer matrices on one axis convention: rows are the
machine a model was fitted on, columns the machine it was evaluated on,
and the outlined diagonal is the locally fitted model.
\textbf{(A)} The decision layer's duration forecaster, raw variant --- no
target information of any kind, the honest zero-shot condition. Values
vary down each column far less than across the row, i.e.\ the error is a
property of the target and barely of the source. \textbf{(B)} The
sequence state model under its source scaler: a bright diagonal on a
dark field. \textbf{What transfers here is not machine-specific, and
what is machine-specific does not transfer.}}
\label{fig:transfer}
\end{figure*}

\textbf{The contrast with the previous subsection is the point}, and
Figure~\ref{fig:transfer} is that contrast in one frame. The
duration forecaster moved between machines at a $4.8\%$ cost in MAE; the
sequence model loses $86\%$ of its STANDBY $F_1$ crossing machine types.
The two results are consistent rather than contradictory, and the
attributions of Section~\ref{sec:explain} say why: what the sequence
model learns \emph{is} machine-specific --- the electrical shape of one
motor's duty cycle, $11.6\%$ calendar --- while what the duration
forecaster learns is the plant calendar at $82.4\%$, which every machine
in the facility shares. \textbf{What transfers here is not
machine-specific, and what is machine-specific does not transfer.}
The paired test over the eight machines puts the transfer degradation at
$0.384$ of $F_1$, $p=0.008$ at the floor of the test, $d=+1.75$
(Section~\ref{sec:significance:cross}) --- the one cross-machine claim
here that survives testing outright.

Re-standardising to the target helps cross-type transfer ($F_1$
$0.076\rightarrow0.149$) and hurts same-type transfer
($0.479\rightarrow0.305$): a source model applied to a machine of the
same type is better off keeping the scale it was trained on, since the
two units genuinely have similar ratings, while across types it is
better off being told the new machine's scale. Neither variant comes
close to a locally fitted model.

\begin{table}[t]
\centering
\caption{Per-machine sequence forecasting, STANDBY $F_1$. The
\textsc{persistence} column requires no training.}
\label{tab:transfer:persistence}
\small
\begin{tabular}{@{}lrrr@{}}
\toprule
Target & own model & \textsc{persistence} & transferred (med.) \\
\midrule
Pelletizer I   & $0.660$ & $0.650$ & $0.20$ \\
Pelletizer II  & $0.607$ & $0.432$ & $0.43$ \\
Milling I      & $0.515$ & $0.504$ & $0.20$ \\
Milling II     & $0.489$ & $\mathbf{0.493}$ & $0.25$ \\
Exhaust fan I  & $\mathbf{0.982}$ & $0.288$ & $0.11$ \\
Exhaust fan II & $0.477$ & $0.211$ & $0.01$ \\
Contactor I    & $0.402$ & $0.012$ & $0.08$ \\
Contactor II   & $0.563$ & $0.552$ & $0.12$ \\
\bottomrule
\end{tabular}
\end{table}

Table~\ref{tab:transfer:persistence} carries the uncomfortable result of
this section. \textbf{On pelletizer~I --- the machine the rest of this
paper is built on --- the sequence model beats ``assume nothing changes
for the next five minutes'' by $0.010$ of $F_1$}, and on milling~II
persistence wins outright. The trained model earns its place only on the
machines whose STANDBY state is easy for the wrong reason: exhaust
fan~I, where STANDBY \emph{is} the off state
(Section~\ref{sec:multimachine:states}) and $0.98$ is achievable. Only
$18$--$21\%$ of transferred pairs beat the target's own persistence
baseline, against $87.5\%$ of locally trained models. Averaged over the
eight machines the training gain is $+0.194$ of $F_1$ and significant,
but it is carried by the auxiliary plant; on the four motor-driven
machines it is $+0.048$ and, at $n=4$, untestable by construction
(Section~\ref{sec:significance:cross}).

\subsection{A second site, and a control that the battery nearly passes}
\label{sec:crossmachine:dataset}

The eight machines share a plant, a calendar and an acquisition system.
To test the framework outside all three, it is applied unchanged to two
single-channel active-power records at $5$\,s resolution from a
different site: a CNC machining centre ($348$ covered days) and the AC
output of a photovoltaic inverter ($86$ covered days). \textbf{The
second is a negative control} --- a generator, not a driven load, for
which no shutdown decision exists --- included because a validation
battery is only worth quoting if it can reject something.
Table~\ref{tab:crossdataset} reports what the unmodified pipeline
returns on each.

\begin{table}[t]
\centering
\caption{Cross-dataset application. Two of the five physics checks
require current and power-factor channels and cannot be evaluated on a
single-channel record.}
\label{tab:crossdataset}
\small
\begin{tabular}{@{}lrr@{}}
\toprule
 & CNC (machine) & PV inverter (control) \\
\midrule
physics checks evaluable & $3/5$ & $3/5$ \\
checks passed            & $\mathbf{3/3}$ & $2/3$ \\
flicker                  & $0.00\%$ & $0.00\%$ \\
median dwell             & $120$\,s & $30$\,s \\
inferred STANDBY         & $186$\,h @ $179$\,W & $42$\,h @ $444$\,W \\
productive at night      & $7.9\%$ & $\mathbf{0.0\%}$ \\
productive at midday     & $34.2\%$ & $54.1\%$ \\
\bottomrule
\end{tabular}
\end{table}

\textbf{The pipeline runs unchanged on both and returns a
plausible-looking four-state partition for both.} On the CNC that is the
right answer: OFF at $0$\,W, STANDBY at $179$\,W, WORKING at $2.24$\,kW,
PEAK at $5.70$\,kW, with the diurnal profile of a day-shift machine.

\textbf{On a solar inverter it is the wrong answer, and the battery
nearly fails to say so.} Power-factor separation and no-load current
ratio are undefined without current and power-factor channels --- and
those are precisely the two checks that do not depend on power
magnitude, i.e.\ the two that test whether STANDBY has been separated
from a lightly loaded WORKING state rather than merely from a
lower-power one. Of the three that remain the control passes power
ordering and OFF-near-zero, and fails only variance ordering, by a $2\%$
margin ($\sigma_{\text{OFF}}=234$\,W against
$\sigma_{\text{STANDBY}}=229$\,W). The framework reports $42$\,h of
recoverable ``standby'' on a device that cannot be shut down.

\begin{figure*}[t]
\centering
\includegraphics[width=\textwidth]{fig_p7_cross_dataset}
\caption{Why the negative control nearly passes. \textbf{(A)} The five
physics checks on the four-channel reference machine and on the two
single-channel second-site records. C3 and C4 are the only two checks
independent of power magnitude --- hence the only two that can separate
energised idle from a lightly loaded productive state --- and they are
exactly the two that a single-channel record cannot evaluate. The
control fails on the one remaining discriminating check, and by a $2\%$
margin. \textbf{(B)} What the battery is left with: the photovoltaic
inverter is productive $0.0\%$ of the time at night, which no driven
load is. That the rejection rests on a diurnal signal rather than a
physical one is the finding.}
\label{fig:crossdataset-checks}
\end{figure*}

\textbf{This is the central cross-dataset finding, and it is a
limitation of the method rather than of the data}
(Figure~\ref{fig:crossdataset-checks}). On a single-channel
record the state model still produces states and the decision layer
still consumes them, but the physics validation has lost most of its
discriminating power --- and it is the physics validation that
Section~\ref{sec:significance:seeds} showed to be load-bearing for
reproducibility.

Nor does the diurnal profile rescue it. Both records are daytime-active,
and on the midday figure the inverter looks like the \emph{better}
machine ($54.1\%$ productive against the CNC's $34.2\%$); only the night
figure separates them, and then by $7.9$ points. What actually differs
is the \emph{shape} --- the CNC's profile is a plateau with sharp edges,
the inverter's a smooth bell centred on solar noon --- and no test in the
battery looks at shape. A single-channel deployment therefore needs an
external admissibility check, \emph{is this a driven load?}, which the
framework as it stands does not provide.

Two further results follow.

\textbf{The decision layer cannot be instantiated on either record.}
Both return a negative cold-versus-warm restart delay, so the same guard
that flagged five IMDELD machines
(Section~\ref{sec:multimachine:decision}) flags both of these. On the
CNC the reason is visible in the episode structure: $229$ usable idle
episodes with a median duration of $14.9$\,h, against pelletizer~I's
$90$\,s --- a machine that is off far more than it is on, and is rarely
observed restarting at all.

\textbf{Zero-shot transfer across sites fails outright.} The
pelletizer~I duration forecaster scores MAE $120\,955$\,s and Spearman
$\mathbf{-0.245}$ on the CNC record, against $60\,386$\,s and $+0.659$
for a locally fitted model; rescaling does not help ($-0.234$). A
\emph{negative} rank correlation means the transferred model orders the
CNC's idle episodes worse than chance. This is the predicted consequence
of Section~\ref{sec:crossmachine:decision}: what that forecaster carries
is the IMDELD plant calendar, and another site does not keep IMDELD's
shifts.


% --- Shapley attribution and transition structure
% =====================================================================
%  explainability.tex  --  PHASE VI, Tasks 13A-13C
%
%  Shapley attribution of the two forecasting components and of the
%  state layer, and the transition structure the temporal model learns.
%
%  Intended placement: Section 7 (\label{sec:explain}), after the
%  decision layer and before the ablation.
%
%  Drop-in usage: % =====================================================================
%  explainability.tex  --  PHASE VI, Tasks 13A-13C
%
%  Shapley attribution of the two forecasting components and of the
%  state layer, and the transition structure the temporal model learns.
%
%  Intended placement: Section 7 (\label{sec:explain}), after the
%  decision layer and before the ablation.
%
%  Drop-in usage: % =====================================================================
%  explainability.tex  --  PHASE VI, Tasks 13A-13C
%
%  Shapley attribution of the two forecasting components and of the
%  state layer, and the transition structure the temporal model learns.
%
%  Intended placement: Section 7 (\label{sec:explain}), after the
%  decision layer and before the ablation.
%
%  Drop-in usage: \input{explainability}
%  Figures are referenced by bare name; add to the preamble
%      \graphicspath{{../outputs/phase6/figures/}}
%  and use the PDF copies (both PDF and PNG are written).
%  Bibliography: paper/references.bib
%  Numbers below are generated by  python run_phase6.py  and stored in
%    outputs/phase6/task13a/<machine>/task13a_results.json
%    outputs/phase6/task13b/<machine>/task13b_results.json
%    outputs/phase6/task13c/task13c_results.json
%  Figures:
%    fig_task13a_decision_shap_phase4, fig_task13a_window_shap,
%    fig_task13b_state_shap_phase4, fig_task13c_transition_matrix,
%    fig_task13c_transition_graph
%
%  FORWARD REFERENCES.  This file cites four sections that do not yet
%  have a .tex of their own; they will resolve to ?? until those are
%  written, which is deliberate rather than an oversight.  The labels
%  expected are:
%    sec:forecasting    the Task 5 / Task 12A forecasting comparison
%    sec:significance   the Phase V significance testing
%    sec:multimachine   the Task 8 eight-machine analysis
%    sec:crossmachine   the Task 9 transfer study
%  Everything else resolves against method_gmm_labeller.tex,
%  related_work.tex and decision_layer.tex as they stand.
% =====================================================================

\section{What the Models Use}
\label{sec:explain}

The three stages of the framework are opaque in different ways. The state
layer assigns a label without saying which measured quantity carried the
decision; both forecasters report split-gain importances, which are known to
be biased toward high-cardinality splits and describe the fitted trees rather
than any individual prediction. This section replaces those importances with
Shapley attributions~\cite{lundberg2017shap}, which are additive, consistent,
and defined per prediction, and reports the transition structure the temporal
model learns.

One caveat governs everything below and is not repeated at each result. Every
target attributed here is derived from the pipeline's own GMM--HMM labels.
These attributions therefore explain what each model learned; they are not
evidence that the underlying state assignment is correct. Validation of the
state definition is the physics battery of
Section~\ref{sec:method:accept_reject} and the cross-machine evidence of
Section~\ref{sec:crossmachine}, not this section.

\subsection{Attribution method}
\label{sec:explain:method}

Every model attributed here is a gradient-boosted
tree~\cite{chen2016xgboost} --- the two forecasters, the chance-constraint
classifier, and the surrogate fitted to the state labels --- so exact tree SHAP
applies and no sampling approximation is needed. Attributions are
computed on held-out rows: $7\,440$ decision epochs for the duration
forecaster, $15\,000$ design rows for the window forecaster, and $20\,000$
energised rows for the state surrogate.

Shapley values satisfy $f(x)=\mathbb{E}[f]+\sum_j\phi_j(x)$, and we verify that
identity rather than assume it. The verification is performed on the model's
raw margin --- the quantity tree SHAP decomposes --- and on a \emph{relative}
scale: the duration forecaster's outputs are seconds and run to $10^{5}$,
while the tree traversal accumulates in single precision, so an absolute
tolerance written for probability outputs is not the right instrument. The
largest relative reconstruction error measured over the four forecaster
attributions is $2.0\times10^{-6}$, four orders of magnitude below the
smallest difference any claim below rests on; the state surrogate's
attributions satisfy the stricter absolute check unmodified.

\subsection{The decision layer's forecaster is a calendar model}
\label{sec:explain:decision}

Section~\ref{sec:crossmachine} reported that the remaining-duration forecaster
draws $94.5\%$ of its split gain from four calendar features, and inferred
from that a model carrying the plant timetable rather than machine physics.
The Shapley attribution confirms the conclusion and corrects the number.

Under the minimum-dwell labelling the seven calendar features carry $82.4\%$
of the mean $|\phi|$ against $96.7\%$ of the split gain; on the
Phase~I--III labelling the same quantities are $78.0\%$ and $96.2\%$. The
direction is the same on both labellings, and the gap between the two measures
is systematic rather than incidental: gain and Shapley agree only at
Spearman $\rho=0.72$ over the ten features ($\rho=0.43$ on the Phase~I--III
labelling), and the two that move furthest are the two the original claim
leaned on.
Figure~\ref{fig:shap-decision}B makes the disagreement visible. The binary
\texttt{cur\_is\_weekend} takes $63.2\%$ of the split gain and $5.1\%$ of the
attribution: it is split on early and often, near the root, on a partition that
subsequent splits then refine, so it accumulates gain out of proportion to the
change it makes to any prediction. Conversely \texttt{onset\_hour} carries
$17.5\%$ of the gain and $43.1\%$ of the attribution.

\textbf{The corrected statement of the finding is therefore that four fifths,
not nineteen twentieths, of the decision forecaster's behaviour is calendar
--- and the conclusion drawn from it in Section~\ref{sec:crossmachine} is
unaffected.} A model that is $82\%$ plant timetable still transfers between any
two machines sharing a shift pattern and still transfers to nothing else,
which is exactly what the transfer matrices measure. What the attribution adds
is that the machine's own history is not negligible: the previous production
run contributes $13.3\%$ of the attribution against $2.3\%$ of the gain, so the
forecaster is not purely a clock.

The chance-constraint classifier keys on something different, and the
difference is what justifies fitting it separately rather than deriving it
from the point forecast. Elapsed idle time is its single largest input
($33.1\%$ of attribution, rank 1) against sixth and $4.3\%$ for the duration
regressor. The feasibility question --- \emph{will this episode last at least
$h_{\min}$ longer?} --- depends on how long the episode has already run in a
way the conditional mean does not.

\begin{figure}[t]
  \centering
  \includegraphics[width=\linewidth]{fig_task13a_decision_shap_phase4}
  \caption{Shapley attribution of the remaining-idle-duration forecaster on
  the minimum-dwell labelling. (A)~Per-prediction attributions on $7\,440$
  held-out decision epochs; the horizontal axis is symmetric-log because
  attributions span $\pm10^{5}$~s while the bulk lies inside $\pm10^{3}$~s.
  Colour is the feature's own value. (B)~Mean $|\phi|$ against split gain; the
  dashed line is where the two measures would agree.}
  \label{fig:shap-decision}
\end{figure}

\subsection{The window forecaster is an electrical model, and mostly a
persistence one}
\label{sec:explain:window}

The window forecaster and the duration forecaster are separate models serving
separate consumers, and their attributions have almost nothing in common. The
window model draws $40.8\%$ of its attribution from the raw electrical
channels, $23.9\%$ from the rolling statistics and $20.1\%$ from the derived
electrical features, against $11.6\%$ for everything calendar
(Figure~\ref{fig:shap-window}A). Set beside the $82\%$ of the previous
subsection, this is the quantitative form of the observation in
Section~\ref{sec:crossmachine} that the two layers depend on nearly disjoint
information --- and it explains why one transfers between machines and the
other does not.

Two results here are worth stating plainly because neither is convenient.

\textbf{Voltage is the single largest contributor}, at $20.0\%$ of the
attribution across its nine design columns and $11.8\%$ on the final observed
value alone, against $3.0\%$ of the split gain. Supply voltage is not a state
variable of the machine; it sags under load and recovers when the drive
unloads, so it is a legitimate but indirect indicator, and one that would not
survive a change of supply transformer or a different point of common
coupling. A deployment on another site should expect this contribution not to
transfer.

\textbf{Almost half the attribution sits on the final observed value.}
Grouping the design columns by what they summarise
(Figure~\ref{fig:shap-window}B) gives $43.5\%$ to the last sample of the
lookback window, $26.4\%$ to block means and $26.3\%$ to block standard
deviations. A model whose largest single input is the state the machine is in
right now is a refinement of ``nothing changes'', and it cannot be expected to
beat that reference by much. It does not: the measured margin over untrained
persistence is $1.45$~s of per-window MAE
(Section~\ref{sec:significance}), which the attribution here now explains
rather than merely reports.

\begin{figure}[t]
  \centering
  \includegraphics[width=\linewidth]{fig_task13a_window_shap}
  \caption{Shapley attribution of the window state forecaster, refitted at
  seed~0 under the protocol of Section~\ref{sec:forecasting} and verified to
  reproduce its stored metrics exactly. (A)~The twelve largest measured
  channels, summed over their nine design columns each; hollow bars are the
  split gain. (B)~The same attribution grouped by what each design column
  summarises.}
  \label{fig:shap-window}
\end{figure}

\subsection{The state layer separates STANDBY on power factor}
\label{sec:explain:state}

Methods~1 and 2 of the five-method proof rest on a physical assertion: an
energised but unloaded induction motor is identified by its power factor and
its magnetising-current ratio, not by its power magnitude. That assertion is
testable against the labels the pipeline actually produced, by fitting a
gradient-boosted surrogate on them and attributing it.

On the hard split --- STANDBY against $\{$WORKING, PEAK\_LOAD$\}$, restricted
to energised rows so that the trivial OFF boundary cannot dominate --- the
attribution ranks \texttt{current} first ($17.8\%$), \texttt{power\_factor}
second ($16.9\%$) and \texttt{active\_power} third ($12.3\%$). The four
physics-named channels carry $30.9\%$ of the attribution between them.

The result is clearest where the two measures disagree. By split gain,
\texttt{active\_power} takes $60.3\%$ and \texttt{power\_factor} $5.1\%$; by
Shapley attribution the order reverses. \textbf{Power magnitude is what the
trees split on; power factor is what changes the predictions.} That is the
physics argument confirmed by the data rather than asserted from motor theory
--- and it is confirmed on the labelling the paper now quotes, not only on the
one the argument was originally made against. Agreement with the corresponding
Phase~II ranking on the earlier labelling is Spearman $\rho=0.833$
($p=4\times10^{-7}$) with four of five top features shared, so the
minimum-dwell re-labelling does not move the explanation.

Figure~\ref{fig:shap-state}A shows the mechanism directly: for
\texttt{power\_factor} and \texttt{current}, the attributions pushing towards
STANDBY are the low feature values, which is the magnetising-current signature
stated in Section~\ref{sec:method:accept_reject}, read off the model.

The surrogate reproduces its targets to $99.74\%$ accuracy. \emph{That number
is not a validation result and should not be read as one}: the targets are the
state model's own labels, so the accuracy measures only that a tree can
re-express a mixture model's decision boundary. Only the ranking is a result.

\begin{figure}[t]
  \centering
  \includegraphics[width=\linewidth]{fig_task13b_state_shap_phase4}
  \caption{Attribution of a gradient-boosted surrogate fitted on the
  GMM--HMM labels, for the STANDBY-versus-productive split on energised rows.
  (A)~Per-row attributions; blue is a low feature value.
  (B)~Three importance measures with independent failure modes, each
  normalised to its own total.}
  \label{fig:shap-state}
\end{figure}

\subsection{The transition structure}
\label{sec:explain:transitions}

The temporal model is claimed to learn a physically admissible state graph.
Figure~\ref{fig:transitions} reports three matrices that are commonly
conflated: the HMM's estimated transition matrix, the frequencies realised by
Viterbi decoding, and the frequencies after the minimum-dwell constraint.
Transitions are never counted across an acquisition gap.

\textbf{The prediction the state definition makes is that a machine cannot
move between de-energised and loaded in one second.} On pelletizer~I the
learned probability of OFF\,$\rightarrow$\,PEAK\_LOAD is exactly zero, and the
decoded sequence contains \emph{zero} such transitions in $5.47$~million
readings. OFF\,$\rightarrow$\,WORKING occurs $126$ times
($6.1\times10^{-5}$ of the outgoing mass from OFF), which at one-second
sampling is the expected consequence of a machine crossing the
STANDBY band faster than the sampling interval rather than a
violation of the ordering.

\textbf{The minimum-dwell constraint acts on the transitions, not only on the
labels.} WORKING\,$\rightarrow$\,OFF falls from $124$ events to $25$, and the
expected dwell in STANDBY rises from $96$~s to $246$~s. The constraint is
therefore doing what it was introduced to do --- removing runs too short to be
physical --- rather than merely relabelling rows.

It also introduces a small artefact that should be stated: three
PEAK\_LOAD\,$\rightarrow$\,OFF transitions appear after the constraint that
the decoder never produced. Absorbing a sub-threshold intermediate run into
its neighbours can join two states that were not adjacent in the decoded
sequence. Three events in $5.47$~million is $5\times10^{-7}$ of the record and
changes no reported quantity, but it is a property of the operator rather than
of the machine.

Across the plant the structure holds where the state model does and fails where
it does not. The largest forbidden-transition probability among the four
motor-driven machines is $1.9\times10^{-4}$ (pelletizer~I's
WORKING\,$\rightarrow$\,OFF); on exhaust-fan~I it is $3.4\times10^{-3}$,
eighteen times larger --- the same machine whose state names
Section~\ref{sec:multimachine} showed to be mis-assigned, arriving here through
an independent measurement.

\begin{figure}[t]
  \centering
  \includegraphics[width=\linewidth]{fig_task13c_transition_matrix}
  \caption{Transition structure on pelletizer~I. (A)~The HMM's estimated
  matrix. (B)~Frequencies counted from the Viterbi-decoded sequence.
  (C)~The same after the minimum-dwell constraint. Colour is logarithmic;
  parenthesised numbers are event counts for the transitions the state
  definition rules out.}
  \label{fig:transitions}
\end{figure}

%  Figures are referenced by bare name; add to the preamble
%      \graphicspath{{../outputs/phase6/figures/}}
%  and use the PDF copies (both PDF and PNG are written).
%  Bibliography: paper/references.bib
%  Numbers below are generated by  python run_phase6.py  and stored in
%    outputs/phase6/task13a/<machine>/task13a_results.json
%    outputs/phase6/task13b/<machine>/task13b_results.json
%    outputs/phase6/task13c/task13c_results.json
%  Figures:
%    fig_task13a_decision_shap_phase4, fig_task13a_window_shap,
%    fig_task13b_state_shap_phase4, fig_task13c_transition_matrix,
%    fig_task13c_transition_graph
%
%  FORWARD REFERENCES.  This file cites four sections that do not yet
%  have a .tex of their own; they will resolve to ?? until those are
%  written, which is deliberate rather than an oversight.  The labels
%  expected are:
%    sec:forecasting    the Task 5 / Task 12A forecasting comparison
%    sec:significance   the Phase V significance testing
%    sec:multimachine   the Task 8 eight-machine analysis
%    sec:crossmachine   the Task 9 transfer study
%  Everything else resolves against method_gmm_labeller.tex,
%  related_work.tex and decision_layer.tex as they stand.
% =====================================================================

\section{What the Models Use}
\label{sec:explain}

The three stages of the framework are opaque in different ways. The state
layer assigns a label without saying which measured quantity carried the
decision; both forecasters report split-gain importances, which are known to
be biased toward high-cardinality splits and describe the fitted trees rather
than any individual prediction. This section replaces those importances with
Shapley attributions~\cite{lundberg2017shap}, which are additive, consistent,
and defined per prediction, and reports the transition structure the temporal
model learns.

One caveat governs everything below and is not repeated at each result. Every
target attributed here is derived from the pipeline's own GMM--HMM labels.
These attributions therefore explain what each model learned; they are not
evidence that the underlying state assignment is correct. Validation of the
state definition is the physics battery of
Section~\ref{sec:method:accept_reject} and the cross-machine evidence of
Section~\ref{sec:crossmachine}, not this section.

\subsection{Attribution method}
\label{sec:explain:method}

Every model attributed here is a gradient-boosted
tree~\cite{chen2016xgboost} --- the two forecasters, the chance-constraint
classifier, and the surrogate fitted to the state labels --- so exact tree SHAP
applies and no sampling approximation is needed. Attributions are
computed on held-out rows: $7\,440$ decision epochs for the duration
forecaster, $15\,000$ design rows for the window forecaster, and $20\,000$
energised rows for the state surrogate.

Shapley values satisfy $f(x)=\mathbb{E}[f]+\sum_j\phi_j(x)$, and we verify that
identity rather than assume it. The verification is performed on the model's
raw margin --- the quantity tree SHAP decomposes --- and on a \emph{relative}
scale: the duration forecaster's outputs are seconds and run to $10^{5}$,
while the tree traversal accumulates in single precision, so an absolute
tolerance written for probability outputs is not the right instrument. The
largest relative reconstruction error measured over the four forecaster
attributions is $2.0\times10^{-6}$, four orders of magnitude below the
smallest difference any claim below rests on; the state surrogate's
attributions satisfy the stricter absolute check unmodified.

\subsection{The decision layer's forecaster is a calendar model}
\label{sec:explain:decision}

Section~\ref{sec:crossmachine} reported that the remaining-duration forecaster
draws $94.5\%$ of its split gain from four calendar features, and inferred
from that a model carrying the plant timetable rather than machine physics.
The Shapley attribution confirms the conclusion and corrects the number.

Under the minimum-dwell labelling the seven calendar features carry $82.4\%$
of the mean $|\phi|$ against $96.7\%$ of the split gain; on the
Phase~I--III labelling the same quantities are $78.0\%$ and $96.2\%$. The
direction is the same on both labellings, and the gap between the two measures
is systematic rather than incidental: gain and Shapley agree only at
Spearman $\rho=0.72$ over the ten features ($\rho=0.43$ on the Phase~I--III
labelling), and the two that move furthest are the two the original claim
leaned on.
Figure~\ref{fig:shap-decision}B makes the disagreement visible. The binary
\texttt{cur\_is\_weekend} takes $63.2\%$ of the split gain and $5.1\%$ of the
attribution: it is split on early and often, near the root, on a partition that
subsequent splits then refine, so it accumulates gain out of proportion to the
change it makes to any prediction. Conversely \texttt{onset\_hour} carries
$17.5\%$ of the gain and $43.1\%$ of the attribution.

\textbf{The corrected statement of the finding is therefore that four fifths,
not nineteen twentieths, of the decision forecaster's behaviour is calendar
--- and the conclusion drawn from it in Section~\ref{sec:crossmachine} is
unaffected.} A model that is $82\%$ plant timetable still transfers between any
two machines sharing a shift pattern and still transfers to nothing else,
which is exactly what the transfer matrices measure. What the attribution adds
is that the machine's own history is not negligible: the previous production
run contributes $13.3\%$ of the attribution against $2.3\%$ of the gain, so the
forecaster is not purely a clock.

The chance-constraint classifier keys on something different, and the
difference is what justifies fitting it separately rather than deriving it
from the point forecast. Elapsed idle time is its single largest input
($33.1\%$ of attribution, rank 1) against sixth and $4.3\%$ for the duration
regressor. The feasibility question --- \emph{will this episode last at least
$h_{\min}$ longer?} --- depends on how long the episode has already run in a
way the conditional mean does not.

\begin{figure}[t]
  \centering
  \includegraphics[width=\linewidth]{fig_task13a_decision_shap_phase4}
  \caption{Shapley attribution of the remaining-idle-duration forecaster on
  the minimum-dwell labelling. (A)~Per-prediction attributions on $7\,440$
  held-out decision epochs; the horizontal axis is symmetric-log because
  attributions span $\pm10^{5}$~s while the bulk lies inside $\pm10^{3}$~s.
  Colour is the feature's own value. (B)~Mean $|\phi|$ against split gain; the
  dashed line is where the two measures would agree.}
  \label{fig:shap-decision}
\end{figure}

\subsection{The window forecaster is an electrical model, and mostly a
persistence one}
\label{sec:explain:window}

The window forecaster and the duration forecaster are separate models serving
separate consumers, and their attributions have almost nothing in common. The
window model draws $40.8\%$ of its attribution from the raw electrical
channels, $23.9\%$ from the rolling statistics and $20.1\%$ from the derived
electrical features, against $11.6\%$ for everything calendar
(Figure~\ref{fig:shap-window}A). Set beside the $82\%$ of the previous
subsection, this is the quantitative form of the observation in
Section~\ref{sec:crossmachine} that the two layers depend on nearly disjoint
information --- and it explains why one transfers between machines and the
other does not.

Two results here are worth stating plainly because neither is convenient.

\textbf{Voltage is the single largest contributor}, at $20.0\%$ of the
attribution across its nine design columns and $11.8\%$ on the final observed
value alone, against $3.0\%$ of the split gain. Supply voltage is not a state
variable of the machine; it sags under load and recovers when the drive
unloads, so it is a legitimate but indirect indicator, and one that would not
survive a change of supply transformer or a different point of common
coupling. A deployment on another site should expect this contribution not to
transfer.

\textbf{Almost half the attribution sits on the final observed value.}
Grouping the design columns by what they summarise
(Figure~\ref{fig:shap-window}B) gives $43.5\%$ to the last sample of the
lookback window, $26.4\%$ to block means and $26.3\%$ to block standard
deviations. A model whose largest single input is the state the machine is in
right now is a refinement of ``nothing changes'', and it cannot be expected to
beat that reference by much. It does not: the measured margin over untrained
persistence is $1.45$~s of per-window MAE
(Section~\ref{sec:significance}), which the attribution here now explains
rather than merely reports.

\begin{figure}[t]
  \centering
  \includegraphics[width=\linewidth]{fig_task13a_window_shap}
  \caption{Shapley attribution of the window state forecaster, refitted at
  seed~0 under the protocol of Section~\ref{sec:forecasting} and verified to
  reproduce its stored metrics exactly. (A)~The twelve largest measured
  channels, summed over their nine design columns each; hollow bars are the
  split gain. (B)~The same attribution grouped by what each design column
  summarises.}
  \label{fig:shap-window}
\end{figure}

\subsection{The state layer separates STANDBY on power factor}
\label{sec:explain:state}

Methods~1 and 2 of the five-method proof rest on a physical assertion: an
energised but unloaded induction motor is identified by its power factor and
its magnetising-current ratio, not by its power magnitude. That assertion is
testable against the labels the pipeline actually produced, by fitting a
gradient-boosted surrogate on them and attributing it.

On the hard split --- STANDBY against $\{$WORKING, PEAK\_LOAD$\}$, restricted
to energised rows so that the trivial OFF boundary cannot dominate --- the
attribution ranks \texttt{current} first ($17.8\%$), \texttt{power\_factor}
second ($16.9\%$) and \texttt{active\_power} third ($12.3\%$). The four
physics-named channels carry $30.9\%$ of the attribution between them.

The result is clearest where the two measures disagree. By split gain,
\texttt{active\_power} takes $60.3\%$ and \texttt{power\_factor} $5.1\%$; by
Shapley attribution the order reverses. \textbf{Power magnitude is what the
trees split on; power factor is what changes the predictions.} That is the
physics argument confirmed by the data rather than asserted from motor theory
--- and it is confirmed on the labelling the paper now quotes, not only on the
one the argument was originally made against. Agreement with the corresponding
Phase~II ranking on the earlier labelling is Spearman $\rho=0.833$
($p=4\times10^{-7}$) with four of five top features shared, so the
minimum-dwell re-labelling does not move the explanation.

Figure~\ref{fig:shap-state}A shows the mechanism directly: for
\texttt{power\_factor} and \texttt{current}, the attributions pushing towards
STANDBY are the low feature values, which is the magnetising-current signature
stated in Section~\ref{sec:method:accept_reject}, read off the model.

The surrogate reproduces its targets to $99.74\%$ accuracy. \emph{That number
is not a validation result and should not be read as one}: the targets are the
state model's own labels, so the accuracy measures only that a tree can
re-express a mixture model's decision boundary. Only the ranking is a result.

\begin{figure}[t]
  \centering
  \includegraphics[width=\linewidth]{fig_task13b_state_shap_phase4}
  \caption{Attribution of a gradient-boosted surrogate fitted on the
  GMM--HMM labels, for the STANDBY-versus-productive split on energised rows.
  (A)~Per-row attributions; blue is a low feature value.
  (B)~Three importance measures with independent failure modes, each
  normalised to its own total.}
  \label{fig:shap-state}
\end{figure}

\subsection{The transition structure}
\label{sec:explain:transitions}

The temporal model is claimed to learn a physically admissible state graph.
Figure~\ref{fig:transitions} reports three matrices that are commonly
conflated: the HMM's estimated transition matrix, the frequencies realised by
Viterbi decoding, and the frequencies after the minimum-dwell constraint.
Transitions are never counted across an acquisition gap.

\textbf{The prediction the state definition makes is that a machine cannot
move between de-energised and loaded in one second.} On pelletizer~I the
learned probability of OFF\,$\rightarrow$\,PEAK\_LOAD is exactly zero, and the
decoded sequence contains \emph{zero} such transitions in $5.47$~million
readings. OFF\,$\rightarrow$\,WORKING occurs $126$ times
($6.1\times10^{-5}$ of the outgoing mass from OFF), which at one-second
sampling is the expected consequence of a machine crossing the
STANDBY band faster than the sampling interval rather than a
violation of the ordering.

\textbf{The minimum-dwell constraint acts on the transitions, not only on the
labels.} WORKING\,$\rightarrow$\,OFF falls from $124$ events to $25$, and the
expected dwell in STANDBY rises from $96$~s to $246$~s. The constraint is
therefore doing what it was introduced to do --- removing runs too short to be
physical --- rather than merely relabelling rows.

It also introduces a small artefact that should be stated: three
PEAK\_LOAD\,$\rightarrow$\,OFF transitions appear after the constraint that
the decoder never produced. Absorbing a sub-threshold intermediate run into
its neighbours can join two states that were not adjacent in the decoded
sequence. Three events in $5.47$~million is $5\times10^{-7}$ of the record and
changes no reported quantity, but it is a property of the operator rather than
of the machine.

Across the plant the structure holds where the state model does and fails where
it does not. The largest forbidden-transition probability among the four
motor-driven machines is $1.9\times10^{-4}$ (pelletizer~I's
WORKING\,$\rightarrow$\,OFF); on exhaust-fan~I it is $3.4\times10^{-3}$,
eighteen times larger --- the same machine whose state names
Section~\ref{sec:multimachine} showed to be mis-assigned, arriving here through
an independent measurement.

\begin{figure}[t]
  \centering
  \includegraphics[width=\linewidth]{fig_task13c_transition_matrix}
  \caption{Transition structure on pelletizer~I. (A)~The HMM's estimated
  matrix. (B)~Frequencies counted from the Viterbi-decoded sequence.
  (C)~The same after the minimum-dwell constraint. Colour is logarithmic;
  parenthesised numbers are event counts for the transitions the state
  definition rules out.}
  \label{fig:transitions}
\end{figure}

%  Figures are referenced by bare name; add to the preamble
%      \graphicspath{{../outputs/phase6/figures/}}
%  and use the PDF copies (both PDF and PNG are written).
%  Bibliography: paper/references.bib
%  Numbers below are generated by  python run_phase6.py  and stored in
%    outputs/phase6/task13a/<machine>/task13a_results.json
%    outputs/phase6/task13b/<machine>/task13b_results.json
%    outputs/phase6/task13c/task13c_results.json
%  Figures:
%    fig_task13a_decision_shap_phase4, fig_task13a_window_shap,
%    fig_task13b_state_shap_phase4, fig_task13c_transition_matrix,
%    fig_task13c_transition_graph
%
%  FORWARD REFERENCES.  This file cites four sections that do not yet
%  have a .tex of their own; they will resolve to ?? until those are
%  written, which is deliberate rather than an oversight.  The labels
%  expected are:
%    sec:forecasting    the Task 5 / Task 12A forecasting comparison
%    sec:significance   the Phase V significance testing
%    sec:multimachine   the Task 8 eight-machine analysis
%    sec:crossmachine   the Task 9 transfer study
%  Everything else resolves against method_gmm_labeller.tex,
%  related_work.tex and decision_layer.tex as they stand.
% =====================================================================

\section{What the Models Use}
\label{sec:explain}

The three stages of the framework are opaque in different ways. The state
layer assigns a label without saying which measured quantity carried the
decision; both forecasters report split-gain importances, which are known to
be biased toward high-cardinality splits and describe the fitted trees rather
than any individual prediction. This section replaces those importances with
Shapley attributions~\cite{lundberg2017shap}, which are additive, consistent,
and defined per prediction, and reports the transition structure the temporal
model learns.

One caveat governs everything below and is not repeated at each result. Every
target attributed here is derived from the pipeline's own GMM--HMM labels.
These attributions therefore explain what each model learned; they are not
evidence that the underlying state assignment is correct. Validation of the
state definition is the physics battery of
Section~\ref{sec:method:accept_reject} and the cross-machine evidence of
Section~\ref{sec:crossmachine}, not this section.

\subsection{Attribution method}
\label{sec:explain:method}

Every model attributed here is a gradient-boosted
tree~\cite{chen2016xgboost} --- the two forecasters, the chance-constraint
classifier, and the surrogate fitted to the state labels --- so exact tree SHAP
applies and no sampling approximation is needed. Attributions are
computed on held-out rows: $7\,440$ decision epochs for the duration
forecaster, $15\,000$ design rows for the window forecaster, and $20\,000$
energised rows for the state surrogate.

Shapley values satisfy $f(x)=\mathbb{E}[f]+\sum_j\phi_j(x)$, and we verify that
identity rather than assume it. The verification is performed on the model's
raw margin --- the quantity tree SHAP decomposes --- and on a \emph{relative}
scale: the duration forecaster's outputs are seconds and run to $10^{5}$,
while the tree traversal accumulates in single precision, so an absolute
tolerance written for probability outputs is not the right instrument. The
largest relative reconstruction error measured over the four forecaster
attributions is $2.0\times10^{-6}$, four orders of magnitude below the
smallest difference any claim below rests on; the state surrogate's
attributions satisfy the stricter absolute check unmodified.

\subsection{The decision layer's forecaster is a calendar model}
\label{sec:explain:decision}

Section~\ref{sec:crossmachine} reported that the remaining-duration forecaster
draws $94.5\%$ of its split gain from four calendar features, and inferred
from that a model carrying the plant timetable rather than machine physics.
The Shapley attribution confirms the conclusion and corrects the number.

Under the minimum-dwell labelling the seven calendar features carry $82.4\%$
of the mean $|\phi|$ against $96.7\%$ of the split gain; on the
Phase~I--III labelling the same quantities are $78.0\%$ and $96.2\%$. The
direction is the same on both labellings, and the gap between the two measures
is systematic rather than incidental: gain and Shapley agree only at
Spearman $\rho=0.72$ over the ten features ($\rho=0.43$ on the Phase~I--III
labelling), and the two that move furthest are the two the original claim
leaned on.
Figure~\ref{fig:shap-decision}B makes the disagreement visible. The binary
\texttt{cur\_is\_weekend} takes $63.2\%$ of the split gain and $5.1\%$ of the
attribution: it is split on early and often, near the root, on a partition that
subsequent splits then refine, so it accumulates gain out of proportion to the
change it makes to any prediction. Conversely \texttt{onset\_hour} carries
$17.5\%$ of the gain and $43.1\%$ of the attribution.

\textbf{The corrected statement of the finding is therefore that four fifths,
not nineteen twentieths, of the decision forecaster's behaviour is calendar
--- and the conclusion drawn from it in Section~\ref{sec:crossmachine} is
unaffected.} A model that is $82\%$ plant timetable still transfers between any
two machines sharing a shift pattern and still transfers to nothing else,
which is exactly what the transfer matrices measure. What the attribution adds
is that the machine's own history is not negligible: the previous production
run contributes $13.3\%$ of the attribution against $2.3\%$ of the gain, so the
forecaster is not purely a clock.

The chance-constraint classifier keys on something different, and the
difference is what justifies fitting it separately rather than deriving it
from the point forecast. Elapsed idle time is its single largest input
($33.1\%$ of attribution, rank 1) against sixth and $4.3\%$ for the duration
regressor. The feasibility question --- \emph{will this episode last at least
$h_{\min}$ longer?} --- depends on how long the episode has already run in a
way the conditional mean does not.

\begin{figure}[t]
  \centering
  \includegraphics[width=\linewidth]{fig_task13a_decision_shap_phase4}
  \caption{Shapley attribution of the remaining-idle-duration forecaster on
  the minimum-dwell labelling. (A)~Per-prediction attributions on $7\,440$
  held-out decision epochs; the horizontal axis is symmetric-log because
  attributions span $\pm10^{5}$~s while the bulk lies inside $\pm10^{3}$~s.
  Colour is the feature's own value. (B)~Mean $|\phi|$ against split gain; the
  dashed line is where the two measures would agree.}
  \label{fig:shap-decision}
\end{figure}

\subsection{The window forecaster is an electrical model, and mostly a
persistence one}
\label{sec:explain:window}

The window forecaster and the duration forecaster are separate models serving
separate consumers, and their attributions have almost nothing in common. The
window model draws $40.8\%$ of its attribution from the raw electrical
channels, $23.9\%$ from the rolling statistics and $20.1\%$ from the derived
electrical features, against $11.6\%$ for everything calendar
(Figure~\ref{fig:shap-window}A). Set beside the $82\%$ of the previous
subsection, this is the quantitative form of the observation in
Section~\ref{sec:crossmachine} that the two layers depend on nearly disjoint
information --- and it explains why one transfers between machines and the
other does not.

Two results here are worth stating plainly because neither is convenient.

\textbf{Voltage is the single largest contributor}, at $20.0\%$ of the
attribution across its nine design columns and $11.8\%$ on the final observed
value alone, against $3.0\%$ of the split gain. Supply voltage is not a state
variable of the machine; it sags under load and recovers when the drive
unloads, so it is a legitimate but indirect indicator, and one that would not
survive a change of supply transformer or a different point of common
coupling. A deployment on another site should expect this contribution not to
transfer.

\textbf{Almost half the attribution sits on the final observed value.}
Grouping the design columns by what they summarise
(Figure~\ref{fig:shap-window}B) gives $43.5\%$ to the last sample of the
lookback window, $26.4\%$ to block means and $26.3\%$ to block standard
deviations. A model whose largest single input is the state the machine is in
right now is a refinement of ``nothing changes'', and it cannot be expected to
beat that reference by much. It does not: the measured margin over untrained
persistence is $1.45$~s of per-window MAE
(Section~\ref{sec:significance}), which the attribution here now explains
rather than merely reports.

\begin{figure}[t]
  \centering
  \includegraphics[width=\linewidth]{fig_task13a_window_shap}
  \caption{Shapley attribution of the window state forecaster, refitted at
  seed~0 under the protocol of Section~\ref{sec:forecasting} and verified to
  reproduce its stored metrics exactly. (A)~The twelve largest measured
  channels, summed over their nine design columns each; hollow bars are the
  split gain. (B)~The same attribution grouped by what each design column
  summarises.}
  \label{fig:shap-window}
\end{figure}

\subsection{The state layer separates STANDBY on power factor}
\label{sec:explain:state}

Methods~1 and 2 of the five-method proof rest on a physical assertion: an
energised but unloaded induction motor is identified by its power factor and
its magnetising-current ratio, not by its power magnitude. That assertion is
testable against the labels the pipeline actually produced, by fitting a
gradient-boosted surrogate on them and attributing it.

On the hard split --- STANDBY against $\{$WORKING, PEAK\_LOAD$\}$, restricted
to energised rows so that the trivial OFF boundary cannot dominate --- the
attribution ranks \texttt{current} first ($17.8\%$), \texttt{power\_factor}
second ($16.9\%$) and \texttt{active\_power} third ($12.3\%$). The four
physics-named channels carry $30.9\%$ of the attribution between them.

The result is clearest where the two measures disagree. By split gain,
\texttt{active\_power} takes $60.3\%$ and \texttt{power\_factor} $5.1\%$; by
Shapley attribution the order reverses. \textbf{Power magnitude is what the
trees split on; power factor is what changes the predictions.} That is the
physics argument confirmed by the data rather than asserted from motor theory
--- and it is confirmed on the labelling the paper now quotes, not only on the
one the argument was originally made against. Agreement with the corresponding
Phase~II ranking on the earlier labelling is Spearman $\rho=0.833$
($p=4\times10^{-7}$) with four of five top features shared, so the
minimum-dwell re-labelling does not move the explanation.

Figure~\ref{fig:shap-state}A shows the mechanism directly: for
\texttt{power\_factor} and \texttt{current}, the attributions pushing towards
STANDBY are the low feature values, which is the magnetising-current signature
stated in Section~\ref{sec:method:accept_reject}, read off the model.

The surrogate reproduces its targets to $99.74\%$ accuracy. \emph{That number
is not a validation result and should not be read as one}: the targets are the
state model's own labels, so the accuracy measures only that a tree can
re-express a mixture model's decision boundary. Only the ranking is a result.

\begin{figure}[t]
  \centering
  \includegraphics[width=\linewidth]{fig_task13b_state_shap_phase4}
  \caption{Attribution of a gradient-boosted surrogate fitted on the
  GMM--HMM labels, for the STANDBY-versus-productive split on energised rows.
  (A)~Per-row attributions; blue is a low feature value.
  (B)~Three importance measures with independent failure modes, each
  normalised to its own total.}
  \label{fig:shap-state}
\end{figure}

\subsection{The transition structure}
\label{sec:explain:transitions}

The temporal model is claimed to learn a physically admissible state graph.
Figure~\ref{fig:transitions} reports three matrices that are commonly
conflated: the HMM's estimated transition matrix, the frequencies realised by
Viterbi decoding, and the frequencies after the minimum-dwell constraint.
Transitions are never counted across an acquisition gap.

\textbf{The prediction the state definition makes is that a machine cannot
move between de-energised and loaded in one second.} On pelletizer~I the
learned probability of OFF\,$\rightarrow$\,PEAK\_LOAD is exactly zero, and the
decoded sequence contains \emph{zero} such transitions in $5.47$~million
readings. OFF\,$\rightarrow$\,WORKING occurs $126$ times
($6.1\times10^{-5}$ of the outgoing mass from OFF), which at one-second
sampling is the expected consequence of a machine crossing the
STANDBY band faster than the sampling interval rather than a
violation of the ordering.

\textbf{The minimum-dwell constraint acts on the transitions, not only on the
labels.} WORKING\,$\rightarrow$\,OFF falls from $124$ events to $25$, and the
expected dwell in STANDBY rises from $96$~s to $246$~s. The constraint is
therefore doing what it was introduced to do --- removing runs too short to be
physical --- rather than merely relabelling rows.

It also introduces a small artefact that should be stated: three
PEAK\_LOAD\,$\rightarrow$\,OFF transitions appear after the constraint that
the decoder never produced. Absorbing a sub-threshold intermediate run into
its neighbours can join two states that were not adjacent in the decoded
sequence. Three events in $5.47$~million is $5\times10^{-7}$ of the record and
changes no reported quantity, but it is a property of the operator rather than
of the machine.

Across the plant the structure holds where the state model does and fails where
it does not. The largest forbidden-transition probability among the four
motor-driven machines is $1.9\times10^{-4}$ (pelletizer~I's
WORKING\,$\rightarrow$\,OFF); on exhaust-fan~I it is $3.4\times10^{-3}$,
eighteen times larger --- the same machine whose state names
Section~\ref{sec:multimachine} showed to be mis-assigned, arriving here through
an independent measurement.

\begin{figure}[t]
  \centering
  \includegraphics[width=\linewidth]{fig_task13c_transition_matrix}
  \caption{Transition structure on pelletizer~I. (A)~The HMM's estimated
  matrix. (B)~Frequencies counted from the Viterbi-decoded sequence.
  (C)~The same after the minimum-dwell constraint. Colour is logarithmic;
  parenthesised numbers are event counts for the transitions the state
  definition rules out.}
  \label{fig:transitions}
\end{figure}


% --- Component ablation
% =====================================================================
%  ablation.tex  --  PHASE VI, Task 13E
%
%  The component ablation: what each stage of the pipeline contributes,
%  measured in the metric that stage is responsible for.
%
%  Intended placement: Section 8 (\label{sec:ablation}), after the
%  explainability section.
%
%  Drop-in usage: % =====================================================================
%  ablation.tex  --  PHASE VI, Task 13E
%
%  The component ablation: what each stage of the pipeline contributes,
%  measured in the metric that stage is responsible for.
%
%  Intended placement: Section 8 (\label{sec:ablation}), after the
%  explainability section.
%
%  Drop-in usage: % =====================================================================
%  ablation.tex  --  PHASE VI, Task 13E
%
%  The component ablation: what each stage of the pipeline contributes,
%  measured in the metric that stage is responsible for.
%
%  Intended placement: Section 8 (\label{sec:ablation}), after the
%  explainability section.
%
%  Drop-in usage: \input{ablation}
%                 \input{../outputs/phase6/task13e/pelletizer-I/ablation_table}
%  Figures are referenced by bare name; add to the preamble
%      \graphicspath{{../outputs/phase6/figures/}}
%  The table needs \usepackage{booktabs}.
%  Numbers below are generated by  python run_phase6.py --only task13e
%  and stored in outputs/phase6/task13e/<machine>/.
%  The table itself is GENERATED -- do not edit ablation_table.tex by
%  hand; edit the script that writes it.
%
%  FORWARD REFERENCES: sec:forecasting, sec:significance (see the note
%  at the head of explainability.tex).
% =====================================================================

\section{Component Ablation}
\label{sec:ablation}

A framework assembled from five components invites one question above all
others: does each of them earn its place? Table~\ref{tab:ablation} answers it
by removing one component at a time and reporting what the removal costs.
Every cell is read from the stored output of the experiment that produced the
corresponding headline result, so the ablation and the claim it qualifies come
from the same fitted models and the same held-out data; nothing here is
refitted for the purpose.

\subsection{Why the table is blocked rather than cumulative}
\label{sec:ablation:structure}

The three stages are measured on three different units --- readings, windows
and factory days --- so a single column running the height of the table would
be meaningless. Two structural points follow, and both matter for reading it.

First, the state-layer metrics are held constant below the first block rather
than repeated: every forecast and decision configuration runs on the
minimum-dwell labelling, and printing its flicker rate nine times would assert
nine times what is asserted once.

Second, and less obviously, \textbf{the two forecasters in this framework are
different models and the table does not pretend otherwise}. The window
forecaster of Section~\ref{sec:forecasting} predicts the state at each step of
a 300\,s horizon and is scored on windows; the decision layer does not consume
it. What the decision layer consumes is the episode-duration forecaster of
Section~\ref{sec:decision:policies}, whose error is measured in seconds of
remaining idle time ($8\,954\pm156$\,s over ten seeds). Carrying the window
model's $F_1$ into the decision rows would attribute the economic result to a
model that plays no part in producing it.

\subsection{The state layer}
\label{sec:ablation:state}

Flicker --- the fraction of state runs shorter than the 10\,s physical
plausibility floor --- falls from $66.9\%$ under independent GMM assignment to
$44.5\%$ once Viterbi decoding imposes a temporal prior, and to $0.00\%$ once
an explicit minimum-dwell constraint is added. The median dwell rises from
3\,s to 13\,s to 92.5\,s across the same three configurations.

Two readings of this are worth separating. \textbf{The HMM contributes, but it
does not solve the problem}: a transition prior can only make short runs
unlikely, and at $10^{6}$ readings per record the likelihood term overwhelms
any prior that leaves the model able to change state at all. It takes an
explicit dwell floor to remove flicker outright, which is what
Section~\ref{sec:method:ghmm} argues on theoretical grounds and what this row
of the table measures.

\textbf{Neither component changes the physics compliance}, and the table says
so rather than hiding it: all three configurations pass 5/5, and the schedule
score moves by $0.06$ percentage points across the block. This is the expected
result and not a null one. The dwell constraint changes \emph{when} the machine
is said to change state, not \emph{which} cluster is standby; a table showing
it improving the physics score would indicate a leak between the two, not a
better model. For scale, the three clustering baselines outside the ablation
reach flicker rates of $62.5\%$ (DBSCAN), $78.2\%$ (K-Means) and $98.2\%$
(spectral clustering), the last of which also fails a physics check.

\subsection{The forecast layer}
\label{sec:ablation:forecast}

The block is anchored on an untrained reference --- the horizon continues the
last observed state --- because a comparison among trained models cannot show
that training helped.

The sequence model does not survive that comparison. The Seq2Seq LSTM reaches
STANDBY $F_1 = 0.561\pm0.064$ against persistence's $0.681$, and $20.90$\,s of
per-window MAE against $12.95$\,s: worse than not forecasting, by
$7.95$\,s, at $p_{\text{Holm}} = 8.6\times10^{-43}$
(Section~\ref{sec:significance}). The gradient-boosted tree is the one trained
forecaster that clears the reference, at $0.689\pm0.003$ and $11.50$\,s ---
$1.45$\,s better than persistence, $p_{\text{Holm}} = 7\times10^{-4}$, and
above the pre-registered 1\,s practical threshold.

The standard deviations in the block carry a second finding. The tree's
spread over 15 seeds is $0.003$ of $F_1$; the sequence model's over 5 seeds is
$0.064$, a factor of twenty-five. The swing between a good and a bad seed of
one neural architecture is larger than the entire spread between
architectures, which is why the forecasting comparison is reported as
mean\,$\pm$\,sd over seeds and not as a single-seed ranking.

\subsection{The decision layer}
\label{sec:ablation:decision}

Every decision row carries a 95\% day-block bootstrap interval over 18 factory
days, and the intervals are the point of the block. Two readings of them must
be kept apart. Each row \emph{on its own} is an interval against the status
quo, and every one of those spans zero except the oracle's: the proposed
policy's $+2.9$ USD/yr sits inside $[-55.9,+46.2]$, so no single row
demonstrates a saving. A \emph{paired contrast} between two rows is a different
and much sharper quantity, because every policy is scored on the same bootstrap
resample and the pairing removes the day-to-day variance that makes the
individual intervals wide. The contrasts quoted below are of the second kind,
and one of them survives.

Three components are separable, and their contributions differ by an order of
magnitude.

\textbf{Targeting the right functional of the predictive distribution is worth
more than anything else in the layer.} Replacing the conditional-mean
forecaster with a log-scale one --- the natural-looking choice, which estimates
the conditional median --- costs $31.5$ USD/yr, from $+2.9$ to $-28.6$. The
idle-duration distribution is heavy enough that its conditional mean exceeds
its conditional median by more than an order of magnitude at episode onset, and
the saving from de-energising is affine in the remaining duration, so the mean
is the sufficient statistic and a policy driven by the median waits far too
long to act. This is the single largest ablation effect in the table.

\textbf{The chance constraint is worth $17.8$ USD/yr}, from $-15.0$ without it
to $+2.9$ with it, and it is what moves the policy's point estimate across
zero. It also enforces the constraint it was introduced for. On the Phase~I--III
decision run, the only one for which both variants' violation counts are
recorded, the unconstrained policy leaves 19 minimum-off violations against the
proposed policy's 2; under the minimum-dwell labelling the proposed policy's
count is again 2, with no variation over ten seeds.

\textbf{The optimisation layer beats the fixed threshold it was built to
replace, and this is the one decision-layer contrast that is significant.} The
proposed policy returns $+68.9$ USD/yr more than the static break-even rule
($p = 0.001$, interval excluding zero, and above the 5 USD/yr economic
threshold). It does \emph{not} beat forecast-free ski-rental
($+2.8$, $p = 0.729$) and cannot be shown to beat the plant's own status quo
($+2.9$, $p = 0.828$).

\textbf{The conclusion the table supports is therefore narrow and should be
stated narrowly}: adaptive thresholds beat a fixed one; forecasting does not
demonstrably beat not forecasting on this machine and this sample. The oracle
row bounds what was available --- $+37.6$ USD/yr with an interval excluding
zero --- so the head-room is real and largely uncaptured, and the binding
constraint on any stronger claim is 18 factory days, not the estimator.

\begin{figure}[t]
  \centering
  \includegraphics[width=\linewidth]{fig_task13e_ablation}
  \caption{The ablation of Table~\ref{tab:ablation} in three panels, one per
  stage, because the three stages are measured in three units. (A)~Flicker
  under each temporal component, pooled over all states and binarised to
  STANDBY-versus-rest, which is the view the decision layer is sensitive to.
  (B)~Per-window duration error with the seed spread; the dashed line is the
  untrained reference. (C)~Annual saving with 95\% day-block bootstrap
  intervals; hollow markers mark intervals that include zero.}
  \label{fig:ablation}
\end{figure}

%                 \input{../outputs/phase6/task13e/pelletizer-I/ablation_table}
%  Figures are referenced by bare name; add to the preamble
%      \graphicspath{{../outputs/phase6/figures/}}
%  The table needs \usepackage{booktabs}.
%  Numbers below are generated by  python run_phase6.py --only task13e
%  and stored in outputs/phase6/task13e/<machine>/.
%  The table itself is GENERATED -- do not edit ablation_table.tex by
%  hand; edit the script that writes it.
%
%  FORWARD REFERENCES: sec:forecasting, sec:significance (see the note
%  at the head of explainability.tex).
% =====================================================================

\section{Component Ablation}
\label{sec:ablation}

A framework assembled from five components invites one question above all
others: does each of them earn its place? Table~\ref{tab:ablation} answers it
by removing one component at a time and reporting what the removal costs.
Every cell is read from the stored output of the experiment that produced the
corresponding headline result, so the ablation and the claim it qualifies come
from the same fitted models and the same held-out data; nothing here is
refitted for the purpose.

\subsection{Why the table is blocked rather than cumulative}
\label{sec:ablation:structure}

The three stages are measured on three different units --- readings, windows
and factory days --- so a single column running the height of the table would
be meaningless. Two structural points follow, and both matter for reading it.

First, the state-layer metrics are held constant below the first block rather
than repeated: every forecast and decision configuration runs on the
minimum-dwell labelling, and printing its flicker rate nine times would assert
nine times what is asserted once.

Second, and less obviously, \textbf{the two forecasters in this framework are
different models and the table does not pretend otherwise}. The window
forecaster of Section~\ref{sec:forecasting} predicts the state at each step of
a 300\,s horizon and is scored on windows; the decision layer does not consume
it. What the decision layer consumes is the episode-duration forecaster of
Section~\ref{sec:decision:policies}, whose error is measured in seconds of
remaining idle time ($8\,954\pm156$\,s over ten seeds). Carrying the window
model's $F_1$ into the decision rows would attribute the economic result to a
model that plays no part in producing it.

\subsection{The state layer}
\label{sec:ablation:state}

Flicker --- the fraction of state runs shorter than the 10\,s physical
plausibility floor --- falls from $66.9\%$ under independent GMM assignment to
$44.5\%$ once Viterbi decoding imposes a temporal prior, and to $0.00\%$ once
an explicit minimum-dwell constraint is added. The median dwell rises from
3\,s to 13\,s to 92.5\,s across the same three configurations.

Two readings of this are worth separating. \textbf{The HMM contributes, but it
does not solve the problem}: a transition prior can only make short runs
unlikely, and at $10^{6}$ readings per record the likelihood term overwhelms
any prior that leaves the model able to change state at all. It takes an
explicit dwell floor to remove flicker outright, which is what
Section~\ref{sec:method:ghmm} argues on theoretical grounds and what this row
of the table measures.

\textbf{Neither component changes the physics compliance}, and the table says
so rather than hiding it: all three configurations pass 5/5, and the schedule
score moves by $0.06$ percentage points across the block. This is the expected
result and not a null one. The dwell constraint changes \emph{when} the machine
is said to change state, not \emph{which} cluster is standby; a table showing
it improving the physics score would indicate a leak between the two, not a
better model. For scale, the three clustering baselines outside the ablation
reach flicker rates of $62.5\%$ (DBSCAN), $78.2\%$ (K-Means) and $98.2\%$
(spectral clustering), the last of which also fails a physics check.

\subsection{The forecast layer}
\label{sec:ablation:forecast}

The block is anchored on an untrained reference --- the horizon continues the
last observed state --- because a comparison among trained models cannot show
that training helped.

The sequence model does not survive that comparison. The Seq2Seq LSTM reaches
STANDBY $F_1 = 0.561\pm0.064$ against persistence's $0.681$, and $20.90$\,s of
per-window MAE against $12.95$\,s: worse than not forecasting, by
$7.95$\,s, at $p_{\text{Holm}} = 8.6\times10^{-43}$
(Section~\ref{sec:significance}). The gradient-boosted tree is the one trained
forecaster that clears the reference, at $0.689\pm0.003$ and $11.50$\,s ---
$1.45$\,s better than persistence, $p_{\text{Holm}} = 7\times10^{-4}$, and
above the pre-registered 1\,s practical threshold.

The standard deviations in the block carry a second finding. The tree's
spread over 15 seeds is $0.003$ of $F_1$; the sequence model's over 5 seeds is
$0.064$, a factor of twenty-five. The swing between a good and a bad seed of
one neural architecture is larger than the entire spread between
architectures, which is why the forecasting comparison is reported as
mean\,$\pm$\,sd over seeds and not as a single-seed ranking.

\subsection{The decision layer}
\label{sec:ablation:decision}

Every decision row carries a 95\% day-block bootstrap interval over 18 factory
days, and the intervals are the point of the block. Two readings of them must
be kept apart. Each row \emph{on its own} is an interval against the status
quo, and every one of those spans zero except the oracle's: the proposed
policy's $+2.9$ USD/yr sits inside $[-55.9,+46.2]$, so no single row
demonstrates a saving. A \emph{paired contrast} between two rows is a different
and much sharper quantity, because every policy is scored on the same bootstrap
resample and the pairing removes the day-to-day variance that makes the
individual intervals wide. The contrasts quoted below are of the second kind,
and one of them survives.

Three components are separable, and their contributions differ by an order of
magnitude.

\textbf{Targeting the right functional of the predictive distribution is worth
more than anything else in the layer.} Replacing the conditional-mean
forecaster with a log-scale one --- the natural-looking choice, which estimates
the conditional median --- costs $31.5$ USD/yr, from $+2.9$ to $-28.6$. The
idle-duration distribution is heavy enough that its conditional mean exceeds
its conditional median by more than an order of magnitude at episode onset, and
the saving from de-energising is affine in the remaining duration, so the mean
is the sufficient statistic and a policy driven by the median waits far too
long to act. This is the single largest ablation effect in the table.

\textbf{The chance constraint is worth $17.8$ USD/yr}, from $-15.0$ without it
to $+2.9$ with it, and it is what moves the policy's point estimate across
zero. It also enforces the constraint it was introduced for. On the Phase~I--III
decision run, the only one for which both variants' violation counts are
recorded, the unconstrained policy leaves 19 minimum-off violations against the
proposed policy's 2; under the minimum-dwell labelling the proposed policy's
count is again 2, with no variation over ten seeds.

\textbf{The optimisation layer beats the fixed threshold it was built to
replace, and this is the one decision-layer contrast that is significant.} The
proposed policy returns $+68.9$ USD/yr more than the static break-even rule
($p = 0.001$, interval excluding zero, and above the 5 USD/yr economic
threshold). It does \emph{not} beat forecast-free ski-rental
($+2.8$, $p = 0.729$) and cannot be shown to beat the plant's own status quo
($+2.9$, $p = 0.828$).

\textbf{The conclusion the table supports is therefore narrow and should be
stated narrowly}: adaptive thresholds beat a fixed one; forecasting does not
demonstrably beat not forecasting on this machine and this sample. The oracle
row bounds what was available --- $+37.6$ USD/yr with an interval excluding
zero --- so the head-room is real and largely uncaptured, and the binding
constraint on any stronger claim is 18 factory days, not the estimator.

\begin{figure}[t]
  \centering
  \includegraphics[width=\linewidth]{fig_task13e_ablation}
  \caption{The ablation of Table~\ref{tab:ablation} in three panels, one per
  stage, because the three stages are measured in three units. (A)~Flicker
  under each temporal component, pooled over all states and binarised to
  STANDBY-versus-rest, which is the view the decision layer is sensitive to.
  (B)~Per-window duration error with the seed spread; the dashed line is the
  untrained reference. (C)~Annual saving with 95\% day-block bootstrap
  intervals; hollow markers mark intervals that include zero.}
  \label{fig:ablation}
\end{figure}

%                 \input{../outputs/phase6/task13e/pelletizer-I/ablation_table}
%  Figures are referenced by bare name; add to the preamble
%      \graphicspath{{../outputs/phase6/figures/}}
%  The table needs \usepackage{booktabs}.
%  Numbers below are generated by  python run_phase6.py --only task13e
%  and stored in outputs/phase6/task13e/<machine>/.
%  The table itself is GENERATED -- do not edit ablation_table.tex by
%  hand; edit the script that writes it.
%
%  FORWARD REFERENCES: sec:forecasting, sec:significance (see the note
%  at the head of explainability.tex).
% =====================================================================

\section{Component Ablation}
\label{sec:ablation}

A framework assembled from five components invites one question above all
others: does each of them earn its place? Table~\ref{tab:ablation} answers it
by removing one component at a time and reporting what the removal costs.
Every cell is read from the stored output of the experiment that produced the
corresponding headline result, so the ablation and the claim it qualifies come
from the same fitted models and the same held-out data; nothing here is
refitted for the purpose.

\subsection{Why the table is blocked rather than cumulative}
\label{sec:ablation:structure}

The three stages are measured on three different units --- readings, windows
and factory days --- so a single column running the height of the table would
be meaningless. Two structural points follow, and both matter for reading it.

First, the state-layer metrics are held constant below the first block rather
than repeated: every forecast and decision configuration runs on the
minimum-dwell labelling, and printing its flicker rate nine times would assert
nine times what is asserted once.

Second, and less obviously, \textbf{the two forecasters in this framework are
different models and the table does not pretend otherwise}. The window
forecaster of Section~\ref{sec:forecasting} predicts the state at each step of
a 300\,s horizon and is scored on windows; the decision layer does not consume
it. What the decision layer consumes is the episode-duration forecaster of
Section~\ref{sec:decision:policies}, whose error is measured in seconds of
remaining idle time ($8\,954\pm156$\,s over ten seeds). Carrying the window
model's $F_1$ into the decision rows would attribute the economic result to a
model that plays no part in producing it.

\subsection{The state layer}
\label{sec:ablation:state}

Flicker --- the fraction of state runs shorter than the 10\,s physical
plausibility floor --- falls from $66.9\%$ under independent GMM assignment to
$44.5\%$ once Viterbi decoding imposes a temporal prior, and to $0.00\%$ once
an explicit minimum-dwell constraint is added. The median dwell rises from
3\,s to 13\,s to 92.5\,s across the same three configurations.

Two readings of this are worth separating. \textbf{The HMM contributes, but it
does not solve the problem}: a transition prior can only make short runs
unlikely, and at $10^{6}$ readings per record the likelihood term overwhelms
any prior that leaves the model able to change state at all. It takes an
explicit dwell floor to remove flicker outright, which is what
Section~\ref{sec:method:ghmm} argues on theoretical grounds and what this row
of the table measures.

\textbf{Neither component changes the physics compliance}, and the table says
so rather than hiding it: all three configurations pass 5/5, and the schedule
score moves by $0.06$ percentage points across the block. This is the expected
result and not a null one. The dwell constraint changes \emph{when} the machine
is said to change state, not \emph{which} cluster is standby; a table showing
it improving the physics score would indicate a leak between the two, not a
better model. For scale, the three clustering baselines outside the ablation
reach flicker rates of $62.5\%$ (DBSCAN), $78.2\%$ (K-Means) and $98.2\%$
(spectral clustering), the last of which also fails a physics check.

\subsection{The forecast layer}
\label{sec:ablation:forecast}

The block is anchored on an untrained reference --- the horizon continues the
last observed state --- because a comparison among trained models cannot show
that training helped.

The sequence model does not survive that comparison. The Seq2Seq LSTM reaches
STANDBY $F_1 = 0.561\pm0.064$ against persistence's $0.681$, and $20.90$\,s of
per-window MAE against $12.95$\,s: worse than not forecasting, by
$7.95$\,s, at $p_{\text{Holm}} = 8.6\times10^{-43}$
(Section~\ref{sec:significance}). The gradient-boosted tree is the one trained
forecaster that clears the reference, at $0.689\pm0.003$ and $11.50$\,s ---
$1.45$\,s better than persistence, $p_{\text{Holm}} = 7\times10^{-4}$, and
above the pre-registered 1\,s practical threshold.

The standard deviations in the block carry a second finding. The tree's
spread over 15 seeds is $0.003$ of $F_1$; the sequence model's over 5 seeds is
$0.064$, a factor of twenty-five. The swing between a good and a bad seed of
one neural architecture is larger than the entire spread between
architectures, which is why the forecasting comparison is reported as
mean\,$\pm$\,sd over seeds and not as a single-seed ranking.

\subsection{The decision layer}
\label{sec:ablation:decision}

Every decision row carries a 95\% day-block bootstrap interval over 18 factory
days, and the intervals are the point of the block. Two readings of them must
be kept apart. Each row \emph{on its own} is an interval against the status
quo, and every one of those spans zero except the oracle's: the proposed
policy's $+2.9$ USD/yr sits inside $[-55.9,+46.2]$, so no single row
demonstrates a saving. A \emph{paired contrast} between two rows is a different
and much sharper quantity, because every policy is scored on the same bootstrap
resample and the pairing removes the day-to-day variance that makes the
individual intervals wide. The contrasts quoted below are of the second kind,
and one of them survives.

Three components are separable, and their contributions differ by an order of
magnitude.

\textbf{Targeting the right functional of the predictive distribution is worth
more than anything else in the layer.} Replacing the conditional-mean
forecaster with a log-scale one --- the natural-looking choice, which estimates
the conditional median --- costs $31.5$ USD/yr, from $+2.9$ to $-28.6$. The
idle-duration distribution is heavy enough that its conditional mean exceeds
its conditional median by more than an order of magnitude at episode onset, and
the saving from de-energising is affine in the remaining duration, so the mean
is the sufficient statistic and a policy driven by the median waits far too
long to act. This is the single largest ablation effect in the table.

\textbf{The chance constraint is worth $17.8$ USD/yr}, from $-15.0$ without it
to $+2.9$ with it, and it is what moves the policy's point estimate across
zero. It also enforces the constraint it was introduced for. On the Phase~I--III
decision run, the only one for which both variants' violation counts are
recorded, the unconstrained policy leaves 19 minimum-off violations against the
proposed policy's 2; under the minimum-dwell labelling the proposed policy's
count is again 2, with no variation over ten seeds.

\textbf{The optimisation layer beats the fixed threshold it was built to
replace, and this is the one decision-layer contrast that is significant.} The
proposed policy returns $+68.9$ USD/yr more than the static break-even rule
($p = 0.001$, interval excluding zero, and above the 5 USD/yr economic
threshold). It does \emph{not} beat forecast-free ski-rental
($+2.8$, $p = 0.729$) and cannot be shown to beat the plant's own status quo
($+2.9$, $p = 0.828$).

\textbf{The conclusion the table supports is therefore narrow and should be
stated narrowly}: adaptive thresholds beat a fixed one; forecasting does not
demonstrably beat not forecasting on this machine and this sample. The oracle
row bounds what was available --- $+37.6$ USD/yr with an interval excluding
zero --- so the head-room is real and largely uncaptured, and the binding
constraint on any stronger claim is 18 factory days, not the estimator.

\begin{figure}[t]
  \centering
  \includegraphics[width=\linewidth]{fig_task13e_ablation}
  \caption{The ablation of Table~\ref{tab:ablation} in three panels, one per
  stage, because the three stages are measured in three units. (A)~Flicker
  under each temporal component, pooled over all states and binarised to
  STANDBY-versus-rest, which is the view the decision layer is sensitive to.
  (B)~Per-window duration error with the seed spread; the dashed line is the
  untrained reference. (C)~Annual saving with 95\% day-block bootstrap
  intervals; hollow markers mark intervals that include zero.}
  \label{fig:ablation}
\end{figure}

\input{../outputs/phase6/task13e/pelletizer-I/ablation_table}

% --- Limitations
% =====================================================================
%  limitations.tex  --  PHASE VII, Task 14D
%
%  A dedicated limitations section.  It exists because the phased
%  evaluation of this framework overturned several of its own earlier
%  claims, and the resulting boundaries are specific enough to be worth
%  stating in one place rather than distributing across the results.
%
%  Intended placement: Section 11 (\label{sec:limitations}), after the
%  ablation and before the conclusion.  Also carries
%  \label{sec:discussion}.
%
%  Drop-in usage: % =====================================================================
%  limitations.tex  --  PHASE VII, Task 14D
%
%  A dedicated limitations section.  It exists because the phased
%  evaluation of this framework overturned several of its own earlier
%  claims, and the resulting boundaries are specific enough to be worth
%  stating in one place rather than distributing across the results.
%
%  Intended placement: Section 11 (\label{sec:limitations}), after the
%  ablation and before the conclusion.  Also carries
%  \label{sec:discussion}.
%
%  Drop-in usage: % =====================================================================
%  limitations.tex  --  PHASE VII, Task 14D
%
%  A dedicated limitations section.  It exists because the phased
%  evaluation of this framework overturned several of its own earlier
%  claims, and the resulting boundaries are specific enough to be worth
%  stating in one place rather than distributing across the results.
%
%  Intended placement: Section 11 (\label{sec:limitations}), after the
%  ablation and before the conclusion.  Also carries
%  \label{sec:discussion}.
%
%  Drop-in usage: \input{limitations}
%  Bibliography: paper/references.bib
%
%  Every item below is traceable to a measured result; none is a
%  generic caveat.  Where a limitation is quantified, the number is the
%  one reported in the section named.
% =====================================================================

\section{Limitations}
\label{sec:limitations}
\label{sec:discussion}

The evaluation in this paper was run in phases, each testing the
preceding one, and several of the framework's original claims did not
survive that process. The boundaries this leaves are specific, and are
collected here rather than distributed through the results.

\subsection{The economic claim rests on eighteen factory days}
\label{sec:limitations:sample}

The decision layer is evaluated on $394$ held-out idle episodes spread
over \textbf{18 factory days}. The day-block bootstrap therefore has
eighteen units to resample, and no refinement of the estimator narrows
the resulting intervals: the proposed policy's $+2.9$\,USD/yr sits
inside $[-55.9,+46.2]$, and every policy's interval against the plant's
own behaviour spans zero except the oracle's
(Section~\ref{sec:significance:decision}).

\textbf{We therefore do not claim that the proposed policy saves money
against the status quo.} What the data support is narrower and is stated
as such throughout: adaptive thresholds beat the fixed threshold they
replace ($+68.9$\,USD/yr, $p=0.001$, paired on the same resample), and
the head-room is real but largely uncaptured. Whether forecasting beats
not forecasting on this machine is an open question that eighteen days
cannot answer, and the ten-seed spread of $\pm0.6$\,USD/yr against a
sampling spread of $\pm50$\,USD/yr confirms that the binding constraint
is the record, not the model.

A second consequence follows for anyone reproducing this work: the
scenario axis is not the uncertainty. The three scenarios differ in the
valuation of lost production, which is the one coefficient that could
not be measured and is therefore reported as an \emph{output} of the
break-even identity rather than assumed. A plant that supplies its own
figure collapses the three columns to one; it does not narrow the
interval.

\subsection{Two machines cannot be annualised, and the largest
per-day figure is one of them}
\label{sec:limitations:milling}

The milling machines' records cover $12.2$ of $43$ days, in $1\,082$ and
$580$ acquisition fragments. Twelve covered days support two one-week
bootstrap blocks, which is not enough to form a usable interval; the
intervals printed for those rows in Table~\ref{tab:multimachine} exist
only to make that visible.

This matters because their $5.5$\,h/day is the \emph{largest} per-day
standby figure among the four machines on which the state definition is
physically valid, and it is exactly the number a reader would want to
annualise. \textbf{It must not be annualised.} The same fragmentation
also breaks the decision layer on those machines --- forecast MAE
$18\,200$\,s against $9\,034$\,s on pelletizer~I, and a proposed policy
that loses over $1\,000$\,USD/yr
(Section~\ref{sec:multimachine:decision}) --- so the two limitations have
one cause.

\subsection{The window forecaster leans on a site-specific quantity}
\label{sec:limitations:voltage}

Supply voltage is the largest single channel contributor to the window
forecaster: $20.0\%$ of its Shapley attribution across nine design
columns, $11.8\%$ on the final observed value alone
(Section~\ref{sec:explain:window}). Voltage is not a state variable of
the machine. It sags under load and recovers when the drive unloads, so
it is a real but indirect indicator --- and one determined by this site's
point of common coupling, transformer impedance and the other loads
sharing the bus.

\textbf{A deployment on another site should not expect that contribution
to transfer}, and the zero-shot cross-site failure of
Section~\ref{sec:crossmachine:dataset} (Spearman $-0.245$, worse than
chance ordering) is consistent with it. Split gain ranks voltage tenth
of twenty-five channels and would have given no warning; the limitation
was found by attribution and would not otherwise have been stated at
all.

\subsection{The attributions are single-machine, and one step removed}
\label{sec:limitations:attribution}

Three separate qualifications apply to Section~\ref{sec:explain}.

\emph{One machine.} The transition structure is measured on all eight
machines, but the forecaster and state-layer attributions are computed
on pelletizer~I only. Whether power factor leads the STANDBY split on
pelletizer~II and the milling machines is a cheap experiment that has
not been run, and it is the obvious strengthening of the physics claim.

\emph{A surrogate is not the model.} The state layer is a GMM--HMM;
what is attributed is a gradient-boosted tree fitted to its output. The
rankings are meaningful and agree with two other measures, but the
mechanism is one step removed, and no attribution method for a
Viterbi-decoded mixture model is applied here.

\emph{Feature dependence is not modelled.} Tree SHAP with the
path-dependent perturbation splits credit between correlated features in
a way that depends on tree structure, and \texttt{active\_power},
\texttt{apparent\_power} and \texttt{current} are strongly correlated on
this record. The ordering \emph{among those three} should be treated as
less firm than the gap between them and the calendar or
rolling-statistic groups.

\subsection{What the ablation does not cover}
\label{sec:limitations:ablation}

Table~\ref{tab:ablation} is computed on one machine and, in its decision
block, on one scenario (S2). The other two scenarios are stored and give
the same ordering, but they are not in the table.

More importantly, \textbf{the strongest accept/reject result in this
paper is not a row of the ablation}. Conditioning the labelling on a
passing physics score collapses the STANDBY-hour spread from $\pm33.7\%$
to $\pm3.0\%$ (Section~\ref{sec:significance:seeds}) --- a larger effect
than any component the table does contain. It is absent because it is
measured over labelling seeds while every other row is measured over
held-out blocks, windows or days, and putting incommensurable units in
one table is precisely what Section~\ref{sec:ablation:structure} argues
against. A seed-resampled ablation of the battery, check by check, is
the experiment that would close this gap.

\subsection{Hyperparameters that are stated rather than learned}
\label{sec:limitations:hyper}

Four constants are set by argument and not by search, and each is a
place where a different plant would need to re-derive rather than
re-use.

The \emph{minimum dwell} is fixed at $10$\,s, equal to the flicker
threshold, so that it removes exactly what the metric defines as
implausible and nothing more. That avoids tuning it to its own metric
but does not make it learned; a hidden semi-Markov
model~\cite{yu2010hsmm} would estimate the duration distribution
instead. The \emph{minimum off time} is $600$\,s, with the saving flat
between $300$ and $1\,200$\,s. The \emph{chance-constraint confidence}
is a point value of $0.90$, stated before the sweep that later showed it
near-optimal; a full stochastic programme would price infeasibility
rather than constrain its probability. And $k=4$ is \emph{imposed}: the
auxiliary machines would be better described with two states and the
BIC scan says the driven machines want more than four
(Section~\ref{sec:multimachine:general}), so a per-machine $k$ selected
under the physics veto is a needed extension.

One further modelling gap belongs here. \textbf{Restart wear is priced
only through a per-day cardinality cap}, not as a cost per cycle with a
bearing-life model --- and the sensitivity sweep found that tightening
that cap to one restart per day \emph{more than doubles} the saving
($89$ against $39$\,USD/yr). Since the cap is applied chronologically it
cannot be selecting the best shutdowns, so a blunt restriction improving
the outcome means the policy's marginal shutdowns are loss-making: the
greedy online rule is not optimal under a binding cap, and a positive
expected-saving margin rather than a zero one is the obvious refinement.
It is reported rather than tuned away.

\subsection{Scope of the evidence}
\label{sec:limitations:scope}

\emph{One facility.} All eight machines share a plant, a shift calendar
and an acquisition system --- and the plant calendar is precisely what
the duration forecaster turns out to be using. Cross-facility transfer
of that component is untested on a comparable record, and the
single-channel evidence available suggests it would fail.

\emph{One comparable dataset, and it is single-channel.} The second
dataset can exercise the pipeline but not the physics battery: C3 and
C4, the two checks that do not depend on power magnitude, are undefined
without current and power-factor channels, and they are the two that
carry the discriminating power
(Section~\ref{sec:crossmachine:dataset}).
On a single-channel record the framework still produces states and still
consumes them, and the negative control shows it will report recoverable
standby on a device that cannot be shut down. A second \emph{six-channel}
industrial record would test the part of the framework that matters
most; none was available.

\emph{No contactor ground truth.} Nothing in this work observes the
machine's contactor directly, so the STANDBY/OFF boundary is inferred
throughout and its residual uncertainty is irreducible from these
measurements alone. Every state-time total is accordingly quoted to two
significant figures and conditional on the physics battery passing.

\subsection{Where the framework should not be deployed}
\label{sec:limitations:deploy}

Stated as criteria rather than caveats, the evidence in this paper rules
out four situations. \textbf{Non-motor loads}: the state definition is
specific to induction-motor drives, and on the fans and contactors there
is no energised-idle state to find. \textbf{Sub-watt standby power}: the
break-even point is then numerically undefined and every currency figure
is an artefact. \textbf{Heavily fragmented records}: the decision layer
inherits its forecaster's quality, and the milling machines show what
that costs. \textbf{Single-channel metering}: the validation battery
loses the two checks that carry its discriminating power. In the first
two cases the framework detects the condition itself and declines to
report a decision, which is the intended behaviour; in the last two it
does not, and an external check is required.

%  Bibliography: paper/references.bib
%
%  Every item below is traceable to a measured result; none is a
%  generic caveat.  Where a limitation is quantified, the number is the
%  one reported in the section named.
% =====================================================================

\section{Limitations}
\label{sec:limitations}
\label{sec:discussion}

The evaluation in this paper was run in phases, each testing the
preceding one, and several of the framework's original claims did not
survive that process. The boundaries this leaves are specific, and are
collected here rather than distributed through the results.

\subsection{The economic claim rests on eighteen factory days}
\label{sec:limitations:sample}

The decision layer is evaluated on $394$ held-out idle episodes spread
over \textbf{18 factory days}. The day-block bootstrap therefore has
eighteen units to resample, and no refinement of the estimator narrows
the resulting intervals: the proposed policy's $+2.9$\,USD/yr sits
inside $[-55.9,+46.2]$, and every policy's interval against the plant's
own behaviour spans zero except the oracle's
(Section~\ref{sec:significance:decision}).

\textbf{We therefore do not claim that the proposed policy saves money
against the status quo.} What the data support is narrower and is stated
as such throughout: adaptive thresholds beat the fixed threshold they
replace ($+68.9$\,USD/yr, $p=0.001$, paired on the same resample), and
the head-room is real but largely uncaptured. Whether forecasting beats
not forecasting on this machine is an open question that eighteen days
cannot answer, and the ten-seed spread of $\pm0.6$\,USD/yr against a
sampling spread of $\pm50$\,USD/yr confirms that the binding constraint
is the record, not the model.

A second consequence follows for anyone reproducing this work: the
scenario axis is not the uncertainty. The three scenarios differ in the
valuation of lost production, which is the one coefficient that could
not be measured and is therefore reported as an \emph{output} of the
break-even identity rather than assumed. A plant that supplies its own
figure collapses the three columns to one; it does not narrow the
interval.

\subsection{Two machines cannot be annualised, and the largest
per-day figure is one of them}
\label{sec:limitations:milling}

The milling machines' records cover $12.2$ of $43$ days, in $1\,082$ and
$580$ acquisition fragments. Twelve covered days support two one-week
bootstrap blocks, which is not enough to form a usable interval; the
intervals printed for those rows in Table~\ref{tab:multimachine} exist
only to make that visible.

This matters because their $5.5$\,h/day is the \emph{largest} per-day
standby figure among the four machines on which the state definition is
physically valid, and it is exactly the number a reader would want to
annualise. \textbf{It must not be annualised.} The same fragmentation
also breaks the decision layer on those machines --- forecast MAE
$18\,200$\,s against $9\,034$\,s on pelletizer~I, and a proposed policy
that loses over $1\,000$\,USD/yr
(Section~\ref{sec:multimachine:decision}) --- so the two limitations have
one cause.

\subsection{The window forecaster leans on a site-specific quantity}
\label{sec:limitations:voltage}

Supply voltage is the largest single channel contributor to the window
forecaster: $20.0\%$ of its Shapley attribution across nine design
columns, $11.8\%$ on the final observed value alone
(Section~\ref{sec:explain:window}). Voltage is not a state variable of
the machine. It sags under load and recovers when the drive unloads, so
it is a real but indirect indicator --- and one determined by this site's
point of common coupling, transformer impedance and the other loads
sharing the bus.

\textbf{A deployment on another site should not expect that contribution
to transfer}, and the zero-shot cross-site failure of
Section~\ref{sec:crossmachine:dataset} (Spearman $-0.245$, worse than
chance ordering) is consistent with it. Split gain ranks voltage tenth
of twenty-five channels and would have given no warning; the limitation
was found by attribution and would not otherwise have been stated at
all.

\subsection{The attributions are single-machine, and one step removed}
\label{sec:limitations:attribution}

Three separate qualifications apply to Section~\ref{sec:explain}.

\emph{One machine.} The transition structure is measured on all eight
machines, but the forecaster and state-layer attributions are computed
on pelletizer~I only. Whether power factor leads the STANDBY split on
pelletizer~II and the milling machines is a cheap experiment that has
not been run, and it is the obvious strengthening of the physics claim.

\emph{A surrogate is not the model.} The state layer is a GMM--HMM;
what is attributed is a gradient-boosted tree fitted to its output. The
rankings are meaningful and agree with two other measures, but the
mechanism is one step removed, and no attribution method for a
Viterbi-decoded mixture model is applied here.

\emph{Feature dependence is not modelled.} Tree SHAP with the
path-dependent perturbation splits credit between correlated features in
a way that depends on tree structure, and \texttt{active\_power},
\texttt{apparent\_power} and \texttt{current} are strongly correlated on
this record. The ordering \emph{among those three} should be treated as
less firm than the gap between them and the calendar or
rolling-statistic groups.

\subsection{What the ablation does not cover}
\label{sec:limitations:ablation}

Table~\ref{tab:ablation} is computed on one machine and, in its decision
block, on one scenario (S2). The other two scenarios are stored and give
the same ordering, but they are not in the table.

More importantly, \textbf{the strongest accept/reject result in this
paper is not a row of the ablation}. Conditioning the labelling on a
passing physics score collapses the STANDBY-hour spread from $\pm33.7\%$
to $\pm3.0\%$ (Section~\ref{sec:significance:seeds}) --- a larger effect
than any component the table does contain. It is absent because it is
measured over labelling seeds while every other row is measured over
held-out blocks, windows or days, and putting incommensurable units in
one table is precisely what Section~\ref{sec:ablation:structure} argues
against. A seed-resampled ablation of the battery, check by check, is
the experiment that would close this gap.

\subsection{Hyperparameters that are stated rather than learned}
\label{sec:limitations:hyper}

Four constants are set by argument and not by search, and each is a
place where a different plant would need to re-derive rather than
re-use.

The \emph{minimum dwell} is fixed at $10$\,s, equal to the flicker
threshold, so that it removes exactly what the metric defines as
implausible and nothing more. That avoids tuning it to its own metric
but does not make it learned; a hidden semi-Markov
model~\cite{yu2010hsmm} would estimate the duration distribution
instead. The \emph{minimum off time} is $600$\,s, with the saving flat
between $300$ and $1\,200$\,s. The \emph{chance-constraint confidence}
is a point value of $0.90$, stated before the sweep that later showed it
near-optimal; a full stochastic programme would price infeasibility
rather than constrain its probability. And $k=4$ is \emph{imposed}: the
auxiliary machines would be better described with two states and the
BIC scan says the driven machines want more than four
(Section~\ref{sec:multimachine:general}), so a per-machine $k$ selected
under the physics veto is a needed extension.

One further modelling gap belongs here. \textbf{Restart wear is priced
only through a per-day cardinality cap}, not as a cost per cycle with a
bearing-life model --- and the sensitivity sweep found that tightening
that cap to one restart per day \emph{more than doubles} the saving
($89$ against $39$\,USD/yr). Since the cap is applied chronologically it
cannot be selecting the best shutdowns, so a blunt restriction improving
the outcome means the policy's marginal shutdowns are loss-making: the
greedy online rule is not optimal under a binding cap, and a positive
expected-saving margin rather than a zero one is the obvious refinement.
It is reported rather than tuned away.

\subsection{Scope of the evidence}
\label{sec:limitations:scope}

\emph{One facility.} All eight machines share a plant, a shift calendar
and an acquisition system --- and the plant calendar is precisely what
the duration forecaster turns out to be using. Cross-facility transfer
of that component is untested on a comparable record, and the
single-channel evidence available suggests it would fail.

\emph{One comparable dataset, and it is single-channel.} The second
dataset can exercise the pipeline but not the physics battery: C3 and
C4, the two checks that do not depend on power magnitude, are undefined
without current and power-factor channels, and they are the two that
carry the discriminating power
(Section~\ref{sec:crossmachine:dataset}).
On a single-channel record the framework still produces states and still
consumes them, and the negative control shows it will report recoverable
standby on a device that cannot be shut down. A second \emph{six-channel}
industrial record would test the part of the framework that matters
most; none was available.

\emph{No contactor ground truth.} Nothing in this work observes the
machine's contactor directly, so the STANDBY/OFF boundary is inferred
throughout and its residual uncertainty is irreducible from these
measurements alone. Every state-time total is accordingly quoted to two
significant figures and conditional on the physics battery passing.

\subsection{Where the framework should not be deployed}
\label{sec:limitations:deploy}

Stated as criteria rather than caveats, the evidence in this paper rules
out four situations. \textbf{Non-motor loads}: the state definition is
specific to induction-motor drives, and on the fans and contactors there
is no energised-idle state to find. \textbf{Sub-watt standby power}: the
break-even point is then numerically undefined and every currency figure
is an artefact. \textbf{Heavily fragmented records}: the decision layer
inherits its forecaster's quality, and the milling machines show what
that costs. \textbf{Single-channel metering}: the validation battery
loses the two checks that carry its discriminating power. In the first
two cases the framework detects the condition itself and declines to
report a decision, which is the intended behaviour; in the last two it
does not, and an external check is required.

%  Bibliography: paper/references.bib
%
%  Every item below is traceable to a measured result; none is a
%  generic caveat.  Where a limitation is quantified, the number is the
%  one reported in the section named.
% =====================================================================

\section{Limitations}
\label{sec:limitations}
\label{sec:discussion}

The evaluation in this paper was run in phases, each testing the
preceding one, and several of the framework's original claims did not
survive that process. The boundaries this leaves are specific, and are
collected here rather than distributed through the results.

\subsection{The economic claim rests on eighteen factory days}
\label{sec:limitations:sample}

The decision layer is evaluated on $394$ held-out idle episodes spread
over \textbf{18 factory days}. The day-block bootstrap therefore has
eighteen units to resample, and no refinement of the estimator narrows
the resulting intervals: the proposed policy's $+2.9$\,USD/yr sits
inside $[-55.9,+46.2]$, and every policy's interval against the plant's
own behaviour spans zero except the oracle's
(Section~\ref{sec:significance:decision}).

\textbf{We therefore do not claim that the proposed policy saves money
against the status quo.} What the data support is narrower and is stated
as such throughout: adaptive thresholds beat the fixed threshold they
replace ($+68.9$\,USD/yr, $p=0.001$, paired on the same resample), and
the head-room is real but largely uncaptured. Whether forecasting beats
not forecasting on this machine is an open question that eighteen days
cannot answer, and the ten-seed spread of $\pm0.6$\,USD/yr against a
sampling spread of $\pm50$\,USD/yr confirms that the binding constraint
is the record, not the model.

A second consequence follows for anyone reproducing this work: the
scenario axis is not the uncertainty. The three scenarios differ in the
valuation of lost production, which is the one coefficient that could
not be measured and is therefore reported as an \emph{output} of the
break-even identity rather than assumed. A plant that supplies its own
figure collapses the three columns to one; it does not narrow the
interval.

\subsection{Two machines cannot be annualised, and the largest
per-day figure is one of them}
\label{sec:limitations:milling}

The milling machines' records cover $12.2$ of $43$ days, in $1\,082$ and
$580$ acquisition fragments. Twelve covered days support two one-week
bootstrap blocks, which is not enough to form a usable interval; the
intervals printed for those rows in Table~\ref{tab:multimachine} exist
only to make that visible.

This matters because their $5.5$\,h/day is the \emph{largest} per-day
standby figure among the four machines on which the state definition is
physically valid, and it is exactly the number a reader would want to
annualise. \textbf{It must not be annualised.} The same fragmentation
also breaks the decision layer on those machines --- forecast MAE
$18\,200$\,s against $9\,034$\,s on pelletizer~I, and a proposed policy
that loses over $1\,000$\,USD/yr
(Section~\ref{sec:multimachine:decision}) --- so the two limitations have
one cause.

\subsection{The window forecaster leans on a site-specific quantity}
\label{sec:limitations:voltage}

Supply voltage is the largest single channel contributor to the window
forecaster: $20.0\%$ of its Shapley attribution across nine design
columns, $11.8\%$ on the final observed value alone
(Section~\ref{sec:explain:window}). Voltage is not a state variable of
the machine. It sags under load and recovers when the drive unloads, so
it is a real but indirect indicator --- and one determined by this site's
point of common coupling, transformer impedance and the other loads
sharing the bus.

\textbf{A deployment on another site should not expect that contribution
to transfer}, and the zero-shot cross-site failure of
Section~\ref{sec:crossmachine:dataset} (Spearman $-0.245$, worse than
chance ordering) is consistent with it. Split gain ranks voltage tenth
of twenty-five channels and would have given no warning; the limitation
was found by attribution and would not otherwise have been stated at
all.

\subsection{The attributions are single-machine, and one step removed}
\label{sec:limitations:attribution}

Three separate qualifications apply to Section~\ref{sec:explain}.

\emph{One machine.} The transition structure is measured on all eight
machines, but the forecaster and state-layer attributions are computed
on pelletizer~I only. Whether power factor leads the STANDBY split on
pelletizer~II and the milling machines is a cheap experiment that has
not been run, and it is the obvious strengthening of the physics claim.

\emph{A surrogate is not the model.} The state layer is a GMM--HMM;
what is attributed is a gradient-boosted tree fitted to its output. The
rankings are meaningful and agree with two other measures, but the
mechanism is one step removed, and no attribution method for a
Viterbi-decoded mixture model is applied here.

\emph{Feature dependence is not modelled.} Tree SHAP with the
path-dependent perturbation splits credit between correlated features in
a way that depends on tree structure, and \texttt{active\_power},
\texttt{apparent\_power} and \texttt{current} are strongly correlated on
this record. The ordering \emph{among those three} should be treated as
less firm than the gap between them and the calendar or
rolling-statistic groups.

\subsection{What the ablation does not cover}
\label{sec:limitations:ablation}

Table~\ref{tab:ablation} is computed on one machine and, in its decision
block, on one scenario (S2). The other two scenarios are stored and give
the same ordering, but they are not in the table.

More importantly, \textbf{the strongest accept/reject result in this
paper is not a row of the ablation}. Conditioning the labelling on a
passing physics score collapses the STANDBY-hour spread from $\pm33.7\%$
to $\pm3.0\%$ (Section~\ref{sec:significance:seeds}) --- a larger effect
than any component the table does contain. It is absent because it is
measured over labelling seeds while every other row is measured over
held-out blocks, windows or days, and putting incommensurable units in
one table is precisely what Section~\ref{sec:ablation:structure} argues
against. A seed-resampled ablation of the battery, check by check, is
the experiment that would close this gap.

\subsection{Hyperparameters that are stated rather than learned}
\label{sec:limitations:hyper}

Four constants are set by argument and not by search, and each is a
place where a different plant would need to re-derive rather than
re-use.

The \emph{minimum dwell} is fixed at $10$\,s, equal to the flicker
threshold, so that it removes exactly what the metric defines as
implausible and nothing more. That avoids tuning it to its own metric
but does not make it learned; a hidden semi-Markov
model~\cite{yu2010hsmm} would estimate the duration distribution
instead. The \emph{minimum off time} is $600$\,s, with the saving flat
between $300$ and $1\,200$\,s. The \emph{chance-constraint confidence}
is a point value of $0.90$, stated before the sweep that later showed it
near-optimal; a full stochastic programme would price infeasibility
rather than constrain its probability. And $k=4$ is \emph{imposed}: the
auxiliary machines would be better described with two states and the
BIC scan says the driven machines want more than four
(Section~\ref{sec:multimachine:general}), so a per-machine $k$ selected
under the physics veto is a needed extension.

One further modelling gap belongs here. \textbf{Restart wear is priced
only through a per-day cardinality cap}, not as a cost per cycle with a
bearing-life model --- and the sensitivity sweep found that tightening
that cap to one restart per day \emph{more than doubles} the saving
($89$ against $39$\,USD/yr). Since the cap is applied chronologically it
cannot be selecting the best shutdowns, so a blunt restriction improving
the outcome means the policy's marginal shutdowns are loss-making: the
greedy online rule is not optimal under a binding cap, and a positive
expected-saving margin rather than a zero one is the obvious refinement.
It is reported rather than tuned away.

\subsection{Scope of the evidence}
\label{sec:limitations:scope}

\emph{One facility.} All eight machines share a plant, a shift calendar
and an acquisition system --- and the plant calendar is precisely what
the duration forecaster turns out to be using. Cross-facility transfer
of that component is untested on a comparable record, and the
single-channel evidence available suggests it would fail.

\emph{One comparable dataset, and it is single-channel.} The second
dataset can exercise the pipeline but not the physics battery: C3 and
C4, the two checks that do not depend on power magnitude, are undefined
without current and power-factor channels, and they are the two that
carry the discriminating power
(Section~\ref{sec:crossmachine:dataset}).
On a single-channel record the framework still produces states and still
consumes them, and the negative control shows it will report recoverable
standby on a device that cannot be shut down. A second \emph{six-channel}
industrial record would test the part of the framework that matters
most; none was available.

\emph{No contactor ground truth.} Nothing in this work observes the
machine's contactor directly, so the STANDBY/OFF boundary is inferred
throughout and its residual uncertainty is irreducible from these
measurements alone. Every state-time total is accordingly quoted to two
significant figures and conditional on the physics battery passing.

\subsection{Where the framework should not be deployed}
\label{sec:limitations:deploy}

Stated as criteria rather than caveats, the evidence in this paper rules
out four situations. \textbf{Non-motor loads}: the state definition is
specific to induction-motor drives, and on the fans and contactors there
is no energised-idle state to find. \textbf{Sub-watt standby power}: the
break-even point is then numerically undefined and every currency figure
is an artefact. \textbf{Heavily fragmented records}: the decision layer
inherits its forecaster's quality, and the milling machines show what
that costs. \textbf{Single-channel metering}: the validation battery
loses the two checks that carry its discriminating power. In the first
two cases the framework detects the condition itself and declines to
report a decision, which is the intended behaviour; in the last two it
does not, and an external check is required.


% --- Conclusion
% =====================================================================
%  conclusion.tex  --  PHASE VII, Task 14C
%
%  Section 12 (\label{sec:conclusion}).  Written after the limitations
%  section, so it does not have to hedge: the boundaries are already
%  stated and the conclusion can be direct about what was and was not
%  demonstrated.
%
%  Drop-in usage: % =====================================================================
%  conclusion.tex  --  PHASE VII, Task 14C
%
%  Section 12 (\label{sec:conclusion}).  Written after the limitations
%  section, so it does not have to hedge: the boundaries are already
%  stated and the conclusion can be direct about what was and was not
%  demonstrated.
%
%  Drop-in usage: % =====================================================================
%  conclusion.tex  --  PHASE VII, Task 14C
%
%  Section 12 (\label{sec:conclusion}).  Written after the limitations
%  section, so it does not have to hedge: the boundaries are already
%  stated and the conclusion can be direct about what was and was not
%  demonstrated.
%
%  Drop-in usage: \input{conclusion}
%  Bibliography: paper/references.bib
% =====================================================================

\section{Conclusion}
\label{sec:conclusion}

This paper set out to connect two capabilities that had matured
separately: inferring an industrial machine's operating state from its
electrical measurements, and deciding optimally when to de-energise it.
The connecting quantity is the duration of a non-productive interval,
which the first literature does not forecast and the second assumes is
given. We proposed a unified label-free framework that supplies it, and
evaluated every stage end to end on a public industrial benchmark of
eight machines.

Four results stand.

\textbf{Label-free state inference can be made trustworthy, but only if
the fit is gated.} Fitted at ten seeds on identical data, the labeller
places the standby cluster on a physically different state in three of
them, and the headline standby total moves by a third of its value.
Five circuit-level checks --- none of which uses the model's own
statistics --- separate those three from the other seven exactly, and
conditioning on them collapses the spread to $3\%$. An accept/reject
step is therefore part of the method, not an appendix to it, and this is
the paper's central methodological claim.

\textbf{Temporal coherence requires a duration model, not a stronger
prior.} A transition matrix can penalise a state change by at most
$\log(p_{\text{self}}/p_{\text{switch}})$ --- $9.23$ nats here --- while
the emission term routinely exceeds that by orders of magnitude, so at
over half of all timesteps the transition model cannot influence the
decode at all. Viterbi decoding takes flicker from $66.9\%$ to $44.5\%$;
an explicit minimum-dwell constraint takes it to zero, at a cost of
$0.19\%$ of rows and no measurable change to any physics or schedule
check, on all eight machines.

\textbf{Two of the framework's own components did not earn their
place, and are reported as such.} An untrained persistence reference
beats all five neural sequence forecasters, the originally proposed
sequence-to-sequence model by the largest margin of the seven; only a
gradient-boosted tree clears the reference, and by $1.45$\,s of
per-window MAE. On the economics, the proposed policy cannot be shown to
beat the plant's own behaviour on an 18-day held-out sample. What it
does beat, significantly and in every scenario, is the fixed-threshold
rule it was built to replace. The defensible decision-layer claim is
that \emph{adaptive thresholds beat a fixed one}, not that forecasting
beats not forecasting.

\textbf{The framework says where it does not apply.} Across the plant
the state definition is physically valid on the four motor-driven
machines and is rejected by the battery on the four auxiliary ones,
where no energised-idle state exists to find. Applied to a second site
it correctly characterises a CNC machining centre and --- on a
single-channel record, having lost the two checks that carry the
discriminating power --- fails to reject a solar inverter, a device that
cannot be shut down at all. A validation battery is only worth quoting
if it can reject something; this one rejects four machines out of eight
and, under reduced instrumentation, is shown to have a blind spot.

\medskip
\noindent
Set against the improvement the framework can extract, the honest
summary is that the facility's total attainable standby saving is under
$90$\,USD/yr, because the plant already de-energises for its long idle
gaps. \textbf{On this facility the framework's value is validated state
inference and decision quality, not recovered energy}, and that
conclusion belongs in the abstract rather than at the end of a
discussion. Whether a plant that does not already shut down manually,
or one with a larger energised-idle fraction, changes the arithmetic is
exactly the question the method now makes answerable from a meter
alone.

\subsection*{Future work}

Four extensions follow directly from the limitations of
Section~\ref{sec:limitations}. A \emph{longer record} is the only thing
that can settle the economic question: eighteen factory days is the
binding constraint on every currency interval here, and no estimator
narrows it. A \emph{second six-channel industrial dataset} would test
the physics battery rather than only the pipeline, which the available
single-channel record cannot do. A \emph{hidden semi-Markov
formulation}~\cite{yu2010hsmm} would learn the state-duration
distribution instead of imposing a constant, and a per-machine state
count selected under the physics veto would remove the $k=4$ assumption
that the eight-machine study shows to be wrong in both directions.
Finally, a \emph{seed-resampled ablation of the validation battery
itself}, check by check, would price the component this paper argues is
most important and is the one component its ablation table does not
contain.

%  Bibliography: paper/references.bib
% =====================================================================

\section{Conclusion}
\label{sec:conclusion}

This paper set out to connect two capabilities that had matured
separately: inferring an industrial machine's operating state from its
electrical measurements, and deciding optimally when to de-energise it.
The connecting quantity is the duration of a non-productive interval,
which the first literature does not forecast and the second assumes is
given. We proposed a unified label-free framework that supplies it, and
evaluated every stage end to end on a public industrial benchmark of
eight machines.

Four results stand.

\textbf{Label-free state inference can be made trustworthy, but only if
the fit is gated.} Fitted at ten seeds on identical data, the labeller
places the standby cluster on a physically different state in three of
them, and the headline standby total moves by a third of its value.
Five circuit-level checks --- none of which uses the model's own
statistics --- separate those three from the other seven exactly, and
conditioning on them collapses the spread to $3\%$. An accept/reject
step is therefore part of the method, not an appendix to it, and this is
the paper's central methodological claim.

\textbf{Temporal coherence requires a duration model, not a stronger
prior.} A transition matrix can penalise a state change by at most
$\log(p_{\text{self}}/p_{\text{switch}})$ --- $9.23$ nats here --- while
the emission term routinely exceeds that by orders of magnitude, so at
over half of all timesteps the transition model cannot influence the
decode at all. Viterbi decoding takes flicker from $66.9\%$ to $44.5\%$;
an explicit minimum-dwell constraint takes it to zero, at a cost of
$0.19\%$ of rows and no measurable change to any physics or schedule
check, on all eight machines.

\textbf{Two of the framework's own components did not earn their
place, and are reported as such.} An untrained persistence reference
beats all five neural sequence forecasters, the originally proposed
sequence-to-sequence model by the largest margin of the seven; only a
gradient-boosted tree clears the reference, and by $1.45$\,s of
per-window MAE. On the economics, the proposed policy cannot be shown to
beat the plant's own behaviour on an 18-day held-out sample. What it
does beat, significantly and in every scenario, is the fixed-threshold
rule it was built to replace. The defensible decision-layer claim is
that \emph{adaptive thresholds beat a fixed one}, not that forecasting
beats not forecasting.

\textbf{The framework says where it does not apply.} Across the plant
the state definition is physically valid on the four motor-driven
machines and is rejected by the battery on the four auxiliary ones,
where no energised-idle state exists to find. Applied to a second site
it correctly characterises a CNC machining centre and --- on a
single-channel record, having lost the two checks that carry the
discriminating power --- fails to reject a solar inverter, a device that
cannot be shut down at all. A validation battery is only worth quoting
if it can reject something; this one rejects four machines out of eight
and, under reduced instrumentation, is shown to have a blind spot.

\medskip
\noindent
Set against the improvement the framework can extract, the honest
summary is that the facility's total attainable standby saving is under
$90$\,USD/yr, because the plant already de-energises for its long idle
gaps. \textbf{On this facility the framework's value is validated state
inference and decision quality, not recovered energy}, and that
conclusion belongs in the abstract rather than at the end of a
discussion. Whether a plant that does not already shut down manually,
or one with a larger energised-idle fraction, changes the arithmetic is
exactly the question the method now makes answerable from a meter
alone.

\subsection*{Future work}

Four extensions follow directly from the limitations of
Section~\ref{sec:limitations}. A \emph{longer record} is the only thing
that can settle the economic question: eighteen factory days is the
binding constraint on every currency interval here, and no estimator
narrows it. A \emph{second six-channel industrial dataset} would test
the physics battery rather than only the pipeline, which the available
single-channel record cannot do. A \emph{hidden semi-Markov
formulation}~\cite{yu2010hsmm} would learn the state-duration
distribution instead of imposing a constant, and a per-machine state
count selected under the physics veto would remove the $k=4$ assumption
that the eight-machine study shows to be wrong in both directions.
Finally, a \emph{seed-resampled ablation of the validation battery
itself}, check by check, would price the component this paper argues is
most important and is the one component its ablation table does not
contain.

%  Bibliography: paper/references.bib
% =====================================================================

\section{Conclusion}
\label{sec:conclusion}

This paper set out to connect two capabilities that had matured
separately: inferring an industrial machine's operating state from its
electrical measurements, and deciding optimally when to de-energise it.
The connecting quantity is the duration of a non-productive interval,
which the first literature does not forecast and the second assumes is
given. We proposed a unified label-free framework that supplies it, and
evaluated every stage end to end on a public industrial benchmark of
eight machines.

Four results stand.

\textbf{Label-free state inference can be made trustworthy, but only if
the fit is gated.} Fitted at ten seeds on identical data, the labeller
places the standby cluster on a physically different state in three of
them, and the headline standby total moves by a third of its value.
Five circuit-level checks --- none of which uses the model's own
statistics --- separate those three from the other seven exactly, and
conditioning on them collapses the spread to $3\%$. An accept/reject
step is therefore part of the method, not an appendix to it, and this is
the paper's central methodological claim.

\textbf{Temporal coherence requires a duration model, not a stronger
prior.} A transition matrix can penalise a state change by at most
$\log(p_{\text{self}}/p_{\text{switch}})$ --- $9.23$ nats here --- while
the emission term routinely exceeds that by orders of magnitude, so at
over half of all timesteps the transition model cannot influence the
decode at all. Viterbi decoding takes flicker from $66.9\%$ to $44.5\%$;
an explicit minimum-dwell constraint takes it to zero, at a cost of
$0.19\%$ of rows and no measurable change to any physics or schedule
check, on all eight machines.

\textbf{Two of the framework's own components did not earn their
place, and are reported as such.} An untrained persistence reference
beats all five neural sequence forecasters, the originally proposed
sequence-to-sequence model by the largest margin of the seven; only a
gradient-boosted tree clears the reference, and by $1.45$\,s of
per-window MAE. On the economics, the proposed policy cannot be shown to
beat the plant's own behaviour on an 18-day held-out sample. What it
does beat, significantly and in every scenario, is the fixed-threshold
rule it was built to replace. The defensible decision-layer claim is
that \emph{adaptive thresholds beat a fixed one}, not that forecasting
beats not forecasting.

\textbf{The framework says where it does not apply.} Across the plant
the state definition is physically valid on the four motor-driven
machines and is rejected by the battery on the four auxiliary ones,
where no energised-idle state exists to find. Applied to a second site
it correctly characterises a CNC machining centre and --- on a
single-channel record, having lost the two checks that carry the
discriminating power --- fails to reject a solar inverter, a device that
cannot be shut down at all. A validation battery is only worth quoting
if it can reject something; this one rejects four machines out of eight
and, under reduced instrumentation, is shown to have a blind spot.

\medskip
\noindent
Set against the improvement the framework can extract, the honest
summary is that the facility's total attainable standby saving is under
$90$\,USD/yr, because the plant already de-energises for its long idle
gaps. \textbf{On this facility the framework's value is validated state
inference and decision quality, not recovered energy}, and that
conclusion belongs in the abstract rather than at the end of a
discussion. Whether a plant that does not already shut down manually,
or one with a larger energised-idle fraction, changes the arithmetic is
exactly the question the method now makes answerable from a meter
alone.

\subsection*{Future work}

Four extensions follow directly from the limitations of
Section~\ref{sec:limitations}. A \emph{longer record} is the only thing
that can settle the economic question: eighteen factory days is the
binding constraint on every currency interval here, and no estimator
narrows it. A \emph{second six-channel industrial dataset} would test
the physics battery rather than only the pipeline, which the available
single-channel record cannot do. A \emph{hidden semi-Markov
formulation}~\cite{yu2010hsmm} would learn the state-duration
distribution instead of imposing a constant, and a per-machine state
count selected under the physics veto would remove the $k=4$ assumption
that the eight-machine study shows to be wrong in both directions.
Finally, a \emph{seed-resampled ablation of the validation battery
itself}, check by check, would price the component this paper argues is
most important and is the one component its ablation table does not
contain.


% ---------------------------------------------------------------------
%  BACK MATTER.  Placed before the bibliography, which is where every
%  target venue puts it.  The data-availability statement is factual as
%  written except for the FILL_IN items: the code repository and the
%  IMDELD DOI are checked, the second-site records are not published and
%  their provenance has to be stated by the author.
% ---------------------------------------------------------------------
\section*{Acknowledgements}

FILL\_IN: funding source and grant number, or the sentence ``This
research received no specific grant from any funding agency in the
public, commercial, or not-for-profit sectors.'' if unfunded.

\section*{Data Availability}

The IMDELD dataset analysed in Sections~\ref{sec:states}%
--\ref{sec:crossmachine} is publicly available from IEEE DataPort at
\url{https://doi.org/10.21227/cg5v-dk02}~\cite{imdeld2018dataset}. All
experiment code, the stored per-phase outputs from which every table and
figure in this paper is generated, and the assembly script that builds
this document are available at
\url{https://github.com/Ananthvarshan/spark-energy-prediction}. The two
single-channel second-site records used in
Section~\ref{sec:crossmachine:dataset} are FILL\_IN: state whether these
are available on request, under what licence, and from whom.

\section*{Declaration of Competing Interest}

FILL\_IN: ``The authors declare that they have no known competing
financial interests or personal relationships that could have appeared
to influence the work reported in this paper.'' -- or the disclosure.

% `plain` is chosen because it ships with every TeX distribution and
% therefore always builds.  For submission, swap in the venue's style --
% \bibliographystyle{IEEEtran} or \bibliographystyle{elsarticle-num} --
% which is the only change the bibliography needs.
\bibliographystyle{plain}
\bibliography{references}

\end{document}
